% Options for packages loaded elsewhere
\PassOptionsToPackage{unicode}{hyperref}
\PassOptionsToPackage{hyphens}{url}
\PassOptionsToPackage{dvipsnames,svgnames,x11names}{xcolor}
\documentclass[
  english,
  11pt,
  a4paper,
]{article}
\usepackage{xcolor}
\usepackage[margin=20mm]{geometry}
\usepackage{amsmath,amssymb}
\setcounter{secnumdepth}{-\maxdimen} % remove section numbering
\usepackage{iftex}
\ifPDFTeX
  \usepackage[T1]{fontenc}
  \usepackage[utf8]{inputenc}
  \usepackage{textcomp} % provide euro and other symbols
\else % if luatex or xetex
  \usepackage{unicode-math} % this also loads fontspec
  \defaultfontfeatures{Scale=MatchLowercase}
  \defaultfontfeatures[\rmfamily]{Ligatures=TeX,Scale=1}
\fi
\usepackage{lmodern}
\ifPDFTeX\else
  % xetex/luatex font selection
\fi
% Use upquote if available, for straight quotes in verbatim environments
\IfFileExists{upquote.sty}{\usepackage{upquote}}{}
\IfFileExists{microtype.sty}{% use microtype if available
  \usepackage[]{microtype}
  \UseMicrotypeSet[protrusion]{basicmath} % disable protrusion for tt fonts
}{}
\makeatletter
\@ifundefined{KOMAClassName}{% if non-KOMA class
  \IfFileExists{parskip.sty}{%
    \usepackage{parskip}
  }{% else
    \setlength{\parindent}{0pt}
    \setlength{\parskip}{6pt plus 2pt minus 1pt}}
}{% if KOMA class
  \KOMAoptions{parskip=half}}
\makeatother
\ifLuaTeX
\usepackage[bidi=basic,shorthands=off]{babel}
\else
\usepackage[bidi=default,shorthands=off]{babel}
\fi
\ifLuaTeX
  \usepackage{selnolig} % disable illegal ligatures
\fi
\setlength{\emergencystretch}{3em} % prevent overfull lines
\providecommand{\tightlist}{%
  \setlength{\itemsep}{0pt}\setlength{\parskip}{0pt}}
\usepackage{amsmath,amssymb}
\usepackage{mathrsfs}
\usepackage{newunicodechar}
\newunicodechar{₂}{\textsubscript{2}}
\newunicodechar{∎}{\ensuremath{\blacksquare}}
\newunicodechar{□}{\ensuremath{\square}}
\allowdisplaybreaks
\setlength{\emergencystretch}{3em}
\sloppy
\relpenalty=100
\binoppenalty=100
\usepackage{fvextra}
\usepackage{needspace}
\setmainfont{LibertinusSerif-Regular.otf}[BoldFont=LibertinusSerif-Bold.otf,ItalicFont=LibertinusSerif-Italic.otf,BoldItalicFont=LibertinusSerif-BoldItalic.otf]
\setmonofont{DejaVuSansMono.ttf}
\makeatletter\renewcommand{\@pnumwidth}{3em}\renewcommand{\@tocrmarg}{4em}\makeatother
\RecustomVerbatimEnvironment{verbatim}{Verbatim}{breaklines=true,breakanywhere=true,fontsize=\small}
\AtBeginDocument{\pdfstringdefDisableCommands{\def\varepsilon{epsilon}}}
\usepackage{bookmark}
\IfFileExists{xurl.sty}{\usepackage{xurl}}{} % add URL line breaks if available
\urlstyle{same}
\hypersetup{
  pdftitle={Derived categories of sheaves},
  pdfauthor={The Stacks project authors and credited programme contributors},
  pdflang={en},
  colorlinks=true,
  linkcolor={Maroon},
  filecolor={Maroon},
  citecolor={Blue},
  urlcolor={Blue},
  pdfcreator={LaTeX via pandoc}}

\title{Derived categories of sheaves}
\usepackage{etoolbox}
\makeatletter
\providecommand{\subtitle}[1]{% add subtitle to \maketitle
  \apptocmd{\@title}{\par {\large #1 \par}}{}{}
}
\makeatother
\subtitle{10 lessons and 1 common reading · Published eleven-unit course
edition; writing-AI self-check; independent AI review, human review and
formal verification not claimed}
\author{The Stacks project authors and credited programme contributors}
\date{3 October 2026}

\begin{document}
\maketitle

{
\setcounter{tocdepth}{2}
\tableofcontents
}
\section*{About this edition}\label{about-edition}
\addcontentsline{toc}{section}{About this edition}

This downloadable edition contains the 10 lessons and 1 common reading
in the selected reader edition of \emph{Derived categories of sheaves}.
It preserves the included lesson text, mathematics, references,
exercises and supplied solutions.

Source status: Published eleven-unit course edition; writing-AI
self-check; independent AI review, human review and formal verification
not claimed. Includes all eleven units in the published reading order,
full source/authorship notice and GNU Free Documentation License. Linked
preceding programme material, including Methods of Algebra Volume 2, is
not reproduced in this volume.

Original writing provenance: The Stacks project authors (adapted
statements and proofs); GPT-6.1 Sol (OpenAI), Ultra (course exposition
and self-check); GPT-6 Astra, Ultra and GPT-5.6 Sol, Ultra (credited
source contributions); Claude Opus 5.5 (Anthropic, effort unrecorded);
one K-flat contribution model is unverified. Exact roles are retained in
the source notice.. Source-preserving format conversion and packaging:
OpenAI Codex --- GPT-6 Astra, Ultra effort (source-preserving format
conversion and layout checks). The source-specific author, contributor
and review statements remain authoritative. Format validation adds no
claim of human editing, independent mathematical review or course
completion. This is a source-preserving edition, not an everyday-English
rewrite.

\href{https://kokunoyumeto.github.io/open-math-courses/courses/derived-categories-and-sheaf-operations/}{Read
the course online} ·
\href{https://kokunoyumeto.github.io/program-matematika-indonesia/en/}{Return
to the mathematics programme}.

This combined course is distributed under the GNU Free Documentation
License, version 1.2 or any later version, with no Invariant Sections,
Front-Cover Texts or Back-Cover Texts. The unchanged source/authorship
notice and full licence follow the lessons. Original CC0 contributions
retain their separate dedication; Stacks-derived material is not
relicensed as CC0. This portable edition changes formatting and
navigation only.

The PDF and EPUB share the same ordered content selection. Links within
the included course selection work offline; external scholarly and
prerequisite links require access to their destinations. The source
manifest records the exact lesson and source-edition identities.

\section{Sheaves of modules on a ringed
space}\label{sheaves-of-modules-on-a-ringed-space}

\emph{Written and edited by GPT-6.1 Sol (OpenAI), in Codex at Ultra,
October 2026. Mathematically self-checked by the writing AI. Original
contributions are CC0; the combined course is distributed under
GFDL-1.2-or-later. Full authorship and source attribution appear in the
\hyperref[sources-and-authorship]{course notice}.}

The first task in derived sheaf theory is to decide what an exact
sequence means. A map can lift every germ and still fail to lift a
global section. A cokernel must therefore be constructed locally,
whereas a kernel can be constructed directly on sections. Once this
distinction is clear, the basic operations on sheaves acquire precise
exactness properties. These properties will determine which operations
need resolutions in later lessons.

The prerequisites are the elementary definitions of a topological space,
a ring, a left module, a functor, a limit and an adjunction. Section 1
recalls presheaves and proves the filtered-colimit, module-category,
sheafification and stalk facts used below. In particular,
\hyperref[sheaves-of-modules-on-a-ringed-space--sheafification-and-the-stalk-test]{Lemma
1.3} supplies the germ construction, its universal property and the
isomorphism test without assuming an abelian category of sheaves.
Definitions and the germ construction are also freely readable in the
Stacks project authors'
\href{https://github.com/KokunoYumeto/unofficial-stacks-project-ai-drafts/blob/565b10e987aba5969b21145a0833f42d69f96790/sheaves.tex}{\emph{Sheaves
on Spaces}}, in the AI Integrated Stacks Project edition.

After this lesson, \emph{Injective modules and bounded-below derived
functors} will construct resolutions using skyscrapers. \emph{Flat
modules and K-flat resolutions} will use extensions by zero to construct
flat resolutions. The distinction between inverse image and module
pullback is the starting point for \emph{Derived pullback and
pushforward}.

The main construction source is the Stacks project authors'
\emph{Sheaves of Modules}, in the
\href{https://github.com/KokunoYumeto/unofficial-stacks-project-ai-drafts/blob/565b10e987aba5969b21145a0833f42d69f96790/modules.tex}{AI
Integrated Stacks Project edition}. The kernels, cokernels, colimits and
sheaf operations developed below supply the foundations for the
resolution lessons. Every proof prerequisite for this lesson is supplied
below. Source attribution and the licence for adapted passages appear in
the \hyperref[sources-and-authorship]{course notice}.

\subsection{1. Coefficients, germs and finite
relations}\label{sheaves-of-modules-on-a-ringed-space--1-coefficients-germs-and-finite-relations}

Let \(X\) be any topological space and let \(\mathcal O\) be a sheaf of
associative rings with identity. Commutativity is not assumed. A sheaf
of left \(\mathcal O\)-modules assigns to each open \(V\) an
\(\mathcal O(V)\)-module, with restriction maps compatible with
restriction of scalars. Write \(\operatorname{Mod}(\mathcal O)\) for
their category. The zero ring and the empty space are allowed. All
indexing sets and diagrams lie in one fixed universe.

A stalk \(F_x\) consists of germs of sections near \(x\). To define its
action by \(\mathcal O_x\), represent the scalar and section on a common
neighbourhood and multiply there. Changing representatives does not
change the resulting germ. The ring identities and module identities
follow on such a common neighbourhood. In particular the action is
meaningful for noncommutative rings, with the order of multiplication
preserved.

\textbf{Lemma 1.1 (filtered colimits of modules).} For a fixed ring
\(A\), filtered colimits in \(\operatorname{Mod}(A)\) preserve short
exact sequences.

\textbf{Proof.} A category is filtered if every finite collection of
objects and arrows has a compatible map to a common object.
Equivalently, it is nonempty, any pair of objects maps to a common
object, and any two parallel arrows become equal after a further arrow.
Elements of a filtered colimit can be represented at one object. A
finite list of elements can be moved to one common object. Equality of
two representatives can be checked after maps to one further object: the
relations defining a colimit involve finitely many arrows, and
filteredness moves all of them to a common object and equalizes parallel
arrows. Thus a representative is zero precisely when it becomes zero
after some further arrow.

Given a filtered diagram of exact sequences
\(0\to A_i\to B_i\to C_i\to0\), suppose an element of \(\varinjlim A_i\)
maps to zero in \(\varinjlim B_i\). Represent it at \(i\), and move to
an object where its image in \(B\) is zero. Injectivity of \(A\to B\) at
that object makes the representative zero there. This proves injectivity
on colimits. If a representative in \(B_i\) maps to zero in the colimit
of \(C_i\), move to an object where its image in \(C\) is zero and use
exactness there to lift it from \(A\). Finally every representative in
\(C_i\) lifts from \(B_i\). These give exactness at the middle and
surjectivity at the right. The constructions respect addition and the
\(A\)-action. \(\square\)

In a neighbourhood diagram used to construct a stalk, the rings
\(\mathcal O(V)\) vary. We apply Lemma 1.1 to the underlying abelian
groups, which all have the fixed coefficient ring \(\mathbb Z\). The
induced equalities of kernels and cokernels respect the
\(\mathcal O_x\)-action, since every scalar germ and each finite list of
section germs can be represented on a common neighbourhood. Thus there
is no assumption that the neighbourhood diagram has one fixed section
ring.

The word \emph{filtered} cannot be removed. For a pushout diagram
written \(L\leftarrow M\to R\), consider the following componentwise
short exact sequence of diagrams:

\[(\mathbb Z\leftarrow0\to0)
\longrightarrow(\mathbb Z\xleftarrow{1}\mathbb Z\to0)
\longrightarrow(0\leftarrow\mathbb Z\to0).\]

The first map is the identity on the left object and the zero inclusion
on the other objects; the second is the identity on the middle object
and the quotient elsewhere. The maps commute with both arrows. The three
pushouts are respectively \(\mathbb Z,0,0\). The first resulting map is
not injective. This example already occurs on a one-point space.

\textbf{Lemma 1.2 (the elementary module category).} Left modules over
any ring form an abelian category. Kernels are submodules, cokernels are
quotient modules, and a sequence is exact if its ordinary image equals
its ordinary kernel. Small limits exist as compatible tuples in
products; small colimits exist as a direct sum modulo the relations from
the arrows.

\textbf{Proof.} The kernel of \(u:M\to N\) is \(\{m:u(m)=0\}\); its
universal property follows because a map killed by \(u\) lands there.
The cokernel is \(N/u(M)\); a map vanishing on \(u(M)\) factors uniquely
through that quotient. The map \(M/\ker u\to u(M)\),
\([m]\mapsto u(m)\), is linear, onto and one-to-one. Thus the canonical
coimage-to-image map is an isomorphism. Morphisms add, composition is
bilinear, and finite direct sums are simultaneously products and
coproducts. These facts establish the abelian-category axioms and the
stated exactness criterion. For a diagram, compatible tuples in the
product have the universal property of its limit. Quotienting the direct
sum by all elements \(m_i-u(m_i)\) associated with arrows \(i\to j\)
gives the colimit\textquotesingle s universal property. These
constructions are sets in the chosen universe. \(\square\)

Sheafification will be denoted \(P^+\).

\subsubsection{Sheafification and the stalk
test}\label{sheaves-of-modules-on-a-ringed-space--sheafification-and-the-stalk-test}

\textbf{Lemma 1.3 (sheafification and the stalk test).} Let \(X\) be any
topological space.

\begin{enumerate}
\tightlist
\item
  Every presheaf of sets \(P\) has a sheaf \(P^+\) and a map
  \(\eta:P\to P^+\) that induces a bijection on every stalk. Every
  section of \(P^+\) is locally the image of a section of \(P\). For
  each sheaf \(H\), composition with \(\eta\) gives a bijection
  \(\operatorname{Hom}(P^+,H)\cong\operatorname{Hom}(P,H)\).
\item
  If \(P\) is separated, \(\eta_U\) is injective for each open \(U\).
\item
  For a presheaf of rings, the same construction preserves addition,
  multiplication and the identity. For a presheaf of left modules over a
  sheaf of associative unital rings \(\mathcal O\), it produces a sheaf
  of left \(\mathcal O\)-modules with the same universal property for
  linear maps. Commutativity is unnecessary.
\item
  A map of sheaves of sets whose stalk maps are all bijective is an
  isomorphism. This also holds for sheaves of rings and of modules.
\end{enumerate}

\textbf{Proof.} Recall the elementary definitions. A presheaf gives
sections on opens and restriction maps compatible with composition. A
sheaf glues every family of sections agreeing on overlaps to a unique
section on the union. A separated presheaf has the uniqueness part of
this condition. Empty covers are included, so a sheaf of sets has
exactly one section on the empty open.

For an open \(U\), define \(P^+(U)\) to be the families
\((a_x)_{x\in U}\), with \(a_x\in P_x\), having the following property:
each \(x\) has a neighbourhood \(V\subset U\) and \(p\in P(V)\) such
that \(a_y=p_y\) for every \(y\in V\). Restrictions discard coordinates,
and \(\eta_U(p)=(p_x)_{x\in U}\). Compatible families on a cover glue
coordinate by coordinate. At a point, any local representative from a
covering member represents the glued family on that same neighbourhood.
This proves existence and uniqueness of gluing in \(P^+\), including its
one empty family on the empty open.

A germ of a section of \(P^+\) at \(x\) has a local representative
\(\eta(p)\); hence \(\eta_x\) is surjective. If two germs from \(P_x\)
have the same image, their images as families agree on a neighbourhood
of \(x\). Their coordinates at \(x\) are therefore equal in \(P_x\),
proving injectivity. If \(P\) is separated and \(\eta_U(p)=\eta_U(q)\),
equality of germs gives a cover of \(U\) on which \(p\) and \(q\) agree.
Separatedness then gives \(p=q\), proving (2).

Sections of any sheaf \(H\) with the same germs agree: germ equality
gives equality on a neighbourhood of each point, and the uniqueness part
of the sheaf condition gives equality on the union. Given \(b:P\to H\),
choose local representatives \(p_i\) of a family in \(P^+(U)\). On an
overlap, \(b(p_i)\) and \(b(p_j)\) have the same germs, so they agree.
Glue them in \(H\). Changing representatives produces sections with the
same germs and hence the same glued section. The construction commutes
with restrictions and extends \(b\). Any extension must agree with it
locally on all representatives, hence globally by uniqueness of gluing.
This proves the universal property in (1).

For (3), form the algebraic operations coordinatewise on germs. To check
local representability of a sum or product, restrict the finitely many
representatives to one common neighbourhood and perform the
corresponding operation there. Zero and the identity are locally
represented by their section values. All ring identities therefore hold
on germs and on their families, with the given order of multiplication.
The extension of a ring map preserves these operations, as is checked on
those local representatives. For a module presheaf, define
\((c\cdot a)_x=c_xa_x\) for \(c\in\mathcal O(U)\). Near a point,
represent \(a\) by \(p\) and use \(c|_V p\); thus the action preserves
local representability. The module identities and restriction
compatibility follow from the section identities. The universal
extension of a linear map is linear because that identity can be checked
on local representatives and then glued. On the empty open the ring and
module have their unique zero element. None of these arguments uses an
abelian category of sheaves.

Finally let \(u:F\to G\) have bijective stalk maps. Injectivity on
stalks implies injectivity on sections, by the germ-equality argument
above. Given \(g\in G(U)\), lift its germ at each point, represent the
lift on a neighbourhood, and shrink until its image equals the
restriction of \(g\). The local lifts agree on overlaps by sectionwise
injectivity. They glue to an \(f\in F(U)\) with \(u(f)=g\). Thus every
\(u_U\) is bijective. Its inverse commutes with restrictions, since
restricting a lift is the unique lift on the smaller open. For rings and
modules the inverses preserve the operations because the original
bijections do. This proves (4). \(\square\)

The germ construction follows the Stacks project authors'
\href{https://github.com/KokunoYumeto/unofficial-stacks-project-ai-drafts/blob/565b10e987aba5969b21145a0833f42d69f96790/sheaves.tex\#L1522}{\emph{Sheaves
on Spaces}, sheafification}, in the AI Integrated Stacks Project
edition. The proof above includes the stalk-isomorphism argument and the
associative-ring module operations in full. Adapted material retains
GFDL-1.2-or-later; the expansion and checks are by GPT-6.1 Sol (OpenAI),
at Ultra.

\subsection{2. The abelian category of module
sheaves}\label{sheaves-of-modules-on-a-ringed-space--2-the-abelian-category-of-module-sheaves}

\textbf{Theorem 2.1.} Let \(\mathcal O\) be any sheaf of rings on \(X\).

\begin{enumerate}
\tightlist
\item
  \(\operatorname{Mod}(\mathcal O)\) is additive. A morphism
  \(\varphi:F\to G\) has a kernel, computed on sections:
  \((\operatorname{Ker}\varphi)(U)=\ker\varphi_U\). It has a cokernel:
  the sheafification of \(U\mapsto\operatorname{coker}\varphi_U\).
\item
  For each \(x\), the stalk functor \(F\mapsto F_x\), from
  \(\operatorname{Mod}(\mathcal O)\) to
  \(\operatorname{Mod}(\mathcal O_x)\), is additive, and
  \((\operatorname{Ker}\varphi)_x=\ker\varphi_x\),
  \((\operatorname{Coker}\varphi)_x=\operatorname{coker}\varphi_x\).
\item
  \(\operatorname{Mod}(\mathcal O)\) is abelian, and each stalk functor
  is exact.
\item
  A sequence \(F\to G\to H\) with zero composite is exact at \(G\) if
  and only if \(F_x\to G_x\to H_x\) is exact at \(G_x\) for every \(x\).
  A morphism \(\varphi\) is a monomorphism (epimorphism, isomorphism) if
  and only if every \(\varphi_x\) is one.
\end{enumerate}

\textbf{Proof.} (1) Morphisms add on sections, and composition is
bilinear. The zero sheaf is a zero object. Finite direct sums are formed
on sections: \(U\mapsto F(U)\oplus G(U)\) is a sheaf, because sections
glue coordinatewise. Let \(\varphi:F\to G\). The presheaf
\(U\mapsto\ker\varphi_U\) is a sheaf. Indeed, let a section \(s\) of
\(F\) lie in the kernel on each set of an open cover. Then
\(\varphi(s)\) vanishes on that cover, so \(\varphi(s)=0\), since \(G\)
is a sheaf. Sections of the kernel that agree on overlaps glue in \(F\),
and by this remark the glued section lies in the kernel. Each
\(\varphi_U\) is \(\mathcal O(U)\)-linear, so the kernel is a
subpresheaf of left modules. A morphism \(E\to F\) whose composite with
\(\varphi\) is zero lands in it on every open set. So it is the kernel.
For the cokernel, let \(P(U)=\operatorname{coker}\varphi_U\), a presheaf
of left \(\mathcal O\)-modules. By
\hyperref[sheaves-of-modules-on-a-ringed-space--sheafification-and-the-stalk-test]{Lemma
1.3}, morphisms \(P^+\to H\) into a sheaf \(H\) correspond to morphisms
of presheaves \(P\to H\). These are the morphisms \(G\to H\) that vanish
after \(\varphi\). So \(P^+\) is the cokernel.

(2) The stalk \(F_x=\varinjlim_{U\ni x}F(U)\) is a left module over
\(\mathcal O_x=\varinjlim_{U\ni x}\mathcal O(U)\), and the functor is
additive. The open neighbourhoods of \(x\) form a filtered system, and
filtered colimits are exact (Lemma 1.1). Hence
\((\operatorname{Ker}\varphi)_x=\varinjlim_U\ker\varphi_U=\ker\varphi_x\).
By
\hyperref[sheaves-of-modules-on-a-ringed-space--sheafification-and-the-stalk-test]{Lemma
1.3}, \(\eta\) is bijective on stalks, so \((P^+)_x=P_x\); and by (Lemma
1.1) again,
\(P_x=\varinjlim_U\operatorname{coker}\varphi_U=\operatorname{coker}\varphi_x\).

(3) The coimage of \(\varphi\) is the cokernel of
\(\operatorname{Ker}\varphi\to F\). The image is the kernel of
\(G\to\operatorname{Coker}\varphi\). By (2), at each \(x\) the canonical
map \(\operatorname{Coim}\varphi\to\operatorname{Im}\varphi\) becomes
the canonical map
\(\operatorname{Coim}\varphi_x\to\operatorname{Im}\varphi_x\) of
\(\mathcal O_x\)-modules. That map is an isomorphism, because modules
over a ring form an abelian category. By
\hyperref[sheaves-of-modules-on-a-ringed-space--sheafification-and-the-stalk-test]{Lemma
1.3}(4) the map of sheaves is an isomorphism. Its inverse is
\(\mathcal O\)-linear, as the map is. So
\(\operatorname{Mod}(\mathcal O)\) is abelian. The stalk functor
preserves kernels and cokernels by (2), so it is exact.

(4) Exactness at \(G\) means that
\(\operatorname{Im}(F\to G)\to\operatorname{Ker}(G\to H)\) is an
isomorphism. Both sides are built from kernels and cokernels. By (2),
the map at \(x\) is the same map for the stalk sequence. If every stalk
sequence is exact,
\hyperref[sheaves-of-modules-on-a-ringed-space--sheafification-and-the-stalk-test]{Lemma
1.3}(4) gives exactness. Conversely, exact functors preserve exactness.
A sheaf whose stalks all vanish is zero, since each section then
vanishes on an open cover. Hence \(\varphi\) is a monomorphism iff
\(\operatorname{Ker}\varphi=0\) iff every \(\ker\varphi_x=0\). The same
holds with cokernels. An isomorphism is a morphism that is both.
\(\square\)

The proof never uses commutativity of \(\mathcal O\). In particular,
Theorem 2.1 holds for the constant sheaf \(k_X\) of any unital ring
\(k\), commutative or not.

\subsection{3. Limits, colimits, stalks, and sections of
sums}\label{sheaves-of-modules-on-a-ringed-space--3-limits-colimits-stalks-and-sections-of-sums}

\textbf{Theorem 3.1.} Let \(\mathcal O\) be any sheaf of rings on \(X\).

\begin{enumerate}
\tightlist
\item
  Every small diagram \((F_i)\) in \(\operatorname{Mod}(\mathcal O)\)
  has a limit, computed on sections. So \(\Gamma(U,-)\) commutes with
  limits.
\item
  Every small diagram has a colimit: the sheafification of the presheaf
  colimit. Stalks commute with colimits:
  \((\varinjlim_iF_i)_x=\varinjlim_i(F_i)_x\).
\item
  Filtered colimits are exact.
\item
  Finite direct sums are computed on sections.
\item
  If \(U\) is a quasi-compact open set, the canonical map
  \(\bigoplus_i\Gamma(U,F_i)\to\Gamma(U,\bigoplus_iF_i)\) is bijective
  for every family. Without quasi-compactness it can fail to be onto
  (Exercise 5).
\item
  Products exist and are computed on sections, but they need not be
  exact (Exercise 6).
\end{enumerate}

\textbf{Proof.} (1) Let \(L(U)=\varprojlim_iF_i(U)\), with
coordinatewise restriction and module structure. It is a sheaf. A local
section on a cover is a compatible family on each piece. Glue each
coordinate in \(F_i\). The glued coordinates again form a compatible
family, because the transition maps are morphisms of sheaves and
sections of a sheaf are determined locally. A morphism \(E\to L\) is the
same as a compatible family of morphisms \(E\to F_i\), open set by open
set. So \(L\) is the limit, and
\(\Gamma(U,L)=\varprojlim_i\Gamma(U,F_i)\).

(2) Let \(P(U)=\varinjlim_iF_i(U)\). By
\hyperref[sheaves-of-modules-on-a-ringed-space--sheafification-and-the-stalk-test]{Lemma
1.3}, morphisms \(P^+\to H\) into a sheaf correspond to compatible
families of morphisms \(F_i\to H\). So \(P^+\) is the colimit. At \(x\),
\[(P^+)_x=P_x=\varinjlim_{U\ni x}\varinjlim_iF_i(U)=\varinjlim_i\varinjlim_{U\ni x}F_i(U)=\varinjlim_i(F_i)_x ,\]
since colimits commute with colimits.

(3) Let \(0\to A_i\to B_i\to C_i\to0\) be exact, naturally in \(i\) over
a filtered category. At each \(x\), by (2), the stalk sequence of the
colimits is the filtered colimit of the exact sequences
\(0\to(A_i)_x\to(B_i)_x\to(C_i)_x\to0\). It is exact by (Lemma 1.1).
Since exactness is checked on stalks
(\hyperref[sheaves-of-modules-on-a-ringed-space--2-the-abelian-category-of-module-sheaves]{Theorem
2.1}(4)), the colimit sequence is exact.

(4) This was shown in the proof of
\hyperref[sheaves-of-modules-on-a-ringed-space--2-the-abelian-category-of-module-sheaves]{Theorem
2.1}(1).

(5) Let \(S=\bigoplus_iF_i\), the sheafification of
\(P(V)=\bigoplus_iF_i(V)\). The germ at \(x\) of \(p=(p_i)\in P(V)\) is
\((p_{i,x})_i\in\bigoplus_i(F_i)_x\). If all germs of \(p\) vanish, each
\(p_i\) has zero germs and is zero. So \(P\) is separated, and
\(\theta:P(V)\to S(V)\) is injective for every \(V\). For \(V=U\) this
is the injectivity of the canonical map. For surjectivity, let
\(t\in S(U)\). By
\hyperref[sheaves-of-modules-on-a-ringed-space--sheafification-and-the-stalk-test]{Lemma
1.3}, every point of \(U\) has an open neighbourhood \(V\subset U\) and
some \(p\in P(V)\) with \(\theta(p)=t|_V\). As \(U\) is quasi-compact,
finitely many such pairs \((V_1,p_1),\ldots,(V_m,p_m)\) cover \(U\).
Each \(p_k\) has finitely many nonzero coordinates. Let \(J\) be the
finite set of indices that occur, and \(S_J=\bigoplus_{i\in J}F_i\). By
(4), \(S_J\) is a sheaf. It is also a subpresheaf of \(P\) that contains
all \(p_k\). On \(V_k\cap V_l\), the restrictions of \(p_k\) and \(p_l\)
have the same image under the injective map \(\theta\), so they are
equal. Glue them to \(s\in S_J(U)=\bigoplus_{i\in J}F_i(U)\). Then
\(\theta(s)\) and \(t\) agree on each \(V_k\), so \(\theta(s)=t\).

(6) Products are limits, so (1) applies. Exercise 6 proves that a
product of epimorphisms need not be an epimorphism. \(\square\)

\subsection{4. Exact inverse
image}\label{sheaves-of-modules-on-a-ringed-space--4-exact-inverse-image}

\textbf{Theorem 4.1.} Let \(f:X\to Y\) be continuous and
\(\mathcal O_Y\) a sheaf of rings on \(Y\).

\begin{enumerate}
\tightlist
\item
  On \(X\), \(f^{-1}\mathcal O_Y\) is a sheaf of rings, with stalks
  \((f^{-1}\mathcal O_Y)_x=\mathcal O_{Y,f(x)}\). For
  \(G\in\operatorname{Mod}(\mathcal O_Y)\), \(f^{-1}G\) is a sheaf of
  left \(f^{-1}\mathcal O_Y\)-modules, and the canonical map
  \(G_{f(x)}\to(f^{-1}G)_x\) is an isomorphism of
  \(\mathcal O_{Y,f(x)}\)-modules.
\item
  \(f^{-1}:\operatorname{Mod}(\mathcal O_Y)\to\operatorname{Mod}(f^{-1}\mathcal O_Y)\)
  is exact and commutes with colimits.
\item
  For a unital ring \(k\), \(f^{-1}k_Y=k_X\) as sheaves of rings. So
  \(f^{-1}:\operatorname{Mod}(k_Y)\to\operatorname{Mod}(k_X)\) is exact,
  and \((f^{-1}G)_x=G_{f(x)}\) as \(k\)-modules. This holds for every
  unital ring \(k\), commutative or not.
\item
  \(f^{-1}\) is left adjoint to
  \(f_*:\operatorname{Mod}(f^{-1}\mathcal O_Y)\to\operatorname{Mod}(\mathcal O_Y)\).
\end{enumerate}

All four assertions are proved below. \emph{References:}
{[}Kashiwara--Schapira 1990, Chapter II{]}; {[}Stacks, Tag 01AJ{]} for
exactness on sheaves of abelian groups.

\textbf{Proof of (1)--(3).} Recall that \(f^{-1}G\) is the
sheafification of the presheaf
\(V\mapsto\varinjlim_{W\supset f(V)}G(W)\), where \(W\) runs over the
open sets containing \(f(V)\). We use its description by germs. For
\(V\subset X\) open, \(f^{-1}G(V)\) is the set of maps \(s\) that give
each \(x\in V\) an element \(s(x)\in G_{f(x)}\), such that each point of
\(V\) has an open neighbourhood \(V'\subset V\), an open
\(W\supset f(V')\) and a section \(g\in G(W)\) with \(s(x')=g_{f(x')}\)
for all \(x'\in V'\). The condition is local, so this is a sheaf of
sets. Take \(G=\mathcal O_Y\) to get \(f^{-1}\mathcal O_Y\). Sum,
product and the action are defined pointwise, through the ring
\(\mathcal O_{Y,f(x)}\) and its module \(G_{f(x)}\). They preserve the
local condition: if near a point \(s\) is represented by \(g\in G(W)\)
and \(a\) by \(h\in\mathcal O_Y(W')\), then near that point \(as\) is
represented by \(h|_{W\cap W'}\,g|_{W\cap W'}\). This gives the ring and
module structures. Evaluation at \(x\) is a ring map, and a module map
over it.

Evaluation \((f^{-1}G)_x\to G_{f(x)}\) is onto: \(g\in G(W)\) with
\(W\ni f(x)\) defines a section over \(f^{-1}(W)\). It is one-to-one: if
\(s(x)=0\), represent \(s\) near \(x\) by \(g\in G(W)\). Then
\(g_{f(x)}=0\), so \(g\) vanishes on an open \(W_0\ni f(x)\), and \(s\)
vanishes on \(V'\cap f^{-1}(W_0)\). The canonical map
\(G_{f(x)}\to(f^{-1}G)_x\) sends the germ of \(g\) to the germ of the
section defined by \(g\). It is inverse to evaluation. This proves (1).

(2) Exactness can be checked on stalks
(\hyperref[sheaves-of-modules-on-a-ringed-space--2-the-abelian-category-of-module-sheaves]{Theorem
2.1}(4)). By (1), the stalk of \(f^{-1}G\) at \(x\) is \(G_{f(x)}\), so
on stalks \(f^{-1}\) acts as the identity, and \(f^{-1}\) is exact. For
a colimit, the canonical map
\(\varinjlim_if^{-1}G_i\to f^{-1}\varinjlim_iG_i\) is, at \(x\), the map
\(\varinjlim_i(G_i)_{f(x)}\to(\varinjlim_iG_i)_{f(x)}\). This map is
bijective, because stalks commute with colimits
(\hyperref[sheaves-of-modules-on-a-ringed-space--3-limits-colimits-stalks-and-sections-of-sums]{Theorem
3.1}(2)). So the canonical map is an isomorphism by
\hyperref[sheaves-of-modules-on-a-ringed-space--sheafification-and-the-stalk-test]{Lemma
1.3}(4).

(3) A germ of a section of \(k_Y\) is the germ of a locally constant
function. So the sections of \(f^{-1}k_Y\) over \(V\) are the maps
\(V\to k\) that are constant near each point, that is, the sections of
\(k_X\). The ring operations agree pointwise. The argument never uses
commutativity of \(k\). \(\square\)

\textbf{Proof of (4).} Let \(G\) be an \(\mathcal O_Y\)-module and \(F\)
an \(f^{-1}\mathcal O_Y\)-module. A map \(\alpha:f^{-1}G\to F\) sends a
section \(g\in G(W)\) to the image under \(\alpha\) of its pullback
section on \(f^{-1}(W)\). This defines
\[\beta_W:G(W)\longrightarrow F(f^{-1}(W))=(f_*F)(W).\] The maps commute
with restriction and are \(\mathcal O_Y(W)\)-linear through the
canonical map \(\mathcal O_Y(W)\to(f^{-1}\mathcal O_Y)(f^{-1}(W))\).

Conversely, suppose \(\beta:G\to f_*F\) is \(\mathcal O_Y\)-linear.
Locally a section \(s\) of \(f^{-1}G\) on \(V\) is represented by a
section \(g\in G(W)\), with \(f(V)\subset W\). Send it to
\(\beta_W(g)|_V\). If two representatives have the same germ at \(x\),
their germs in \(G_{f(x)}\) agree, so they agree on some smaller
neighbourhood \(W'\) of \(f(x)\). Their proposed images therefore agree
near \(x\). This proves independence of representatives; the proposed
local images glue uniquely. Multiplication is checked after locally
representing the scalar by a section of \(\mathcal O_Y\), so the glued
map is \(f^{-1}\mathcal O_Y\)-linear. The two constructions undo one
another on local representatives, and hence everywhere. They commute
with maps of \(G\) and \(F\); this is the adjunction. \(\square\)

\subsubsection{Module pullback and extension of
scalars}\label{sheaves-of-modules-on-a-ringed-space--module-pullback-and-extension-of-scalars}

The inverse image changes the space. A morphism of ringed spaces also
changes the coefficient ring. Let

\[f:(X,\mathcal O_X)\longrightarrow(Y,\mathcal O_Y)\]

consist of a continuous map and a specified homomorphism
\(f^{-1}\mathcal O_Y\to\mathcal O_X\). Write \(A=f^{-1}\mathcal O_Y\)
and \(B=\mathcal O_X\). Then \(B\) is a \((B,A)\)-bimodule: its right
\(A\)-action is multiplication by the image of \(A\). For an
\(A\)-module \(N\), the sheaf \(B\otimes_A N\) is the sheafification of
\(V\mapsto B(V)\otimes_{A(V)}N(V)\). It is a sheaf of left
\(B\)-modules.

\textbf{Lemma 4.3 (tensor and stalks).} There is a natural isomorphism

\[(B\otimes_A N)_x\cong B_x\otimes_{A_x}N_x.\]

\textbf{Proof.} A tensor on either side is a finite sum of elementary
tensors. Choose representatives for its finitely many factors on one
neighbourhood of \(x\); this constructs the inverse of the canonical map
from the left side to the right. The tensor relations are finite sums of
additivity and balancing relations \(ba\otimes n=b\otimes an\). Any
finite collection of germs witnessing those relations can again be
represented on a common neighbourhood. Thus two representing expressions
give the same germ exactly when the corresponding tensors agree on the
right. Sheafification preserves stalks, so this verifies both
injectivity and surjectivity for the displayed sheaf tensor product.
\(\square\)

\textbf{Theorem 4.4 (pullback adjunction and exactness).} Define

\[f^*G=\mathcal O_X\otimes_{f^{-1}\mathcal O_Y}f^{-1}G.\]

For sheaves of left modules over arbitrary unital coefficient rings,
\(f^*\) is left adjoint to \(f_*\). It is right exact; \(f_*\) is left
exact. At a point,

\[(f^*G)_x=\mathcal O_{X,x}\otimes_{\mathcal O_{Y,f(x)}}G_{f(x)}.\]

\textbf{Proof.} First, extension of scalars is left adjoint to
forgetting the \(B\)-action. Given a \(B\)-linear map
\(t:B\otimes_A N\to F\), define \(a:N\to F\) by
\(n\mapsto t(1\otimes n)\). It is \(A\)-linear by the balancing
relation. Conversely an \(A\)-linear map \(a\) defines
\(b\otimes n\mapsto ba(n)\) on elementary tensors. This is balanced,
hence gives a map on the presheaf tensor product and uniquely on its
sheafification. The two formulas undo one another, commute with
restriction, and are natural in both modules. Combining them with
Theorem 4.1(4) gives

\[\operatorname{Hom}_{\mathcal O_X}(f^*G,F)
\cong\operatorname{Hom}_{f^{-1}\mathcal O_Y}(f^{-1}G,F)
\cong\operatorname{Hom}_{\mathcal O_Y}(G,f_*F).\]

Here the middle term regards \(F\) as an \(f^{-1}\mathcal O_Y\)-module
through the specified ring map. Thus the adjunction has the required
coefficient structures. The stalk formula follows from Lemma 4.3 and
Theorem 4.1(1).

For completeness, tensor over a ring is right exact. If
\(N\to N'\to N''\to0\) is exact, present \(N''\) as the quotient of
\(N'\) by the image of \(N\). A balanced map on \(B\times N''\) is
exactly a balanced map on \(B\times N'\) vanishing on
\(B\times\operatorname{im}N\). The defining universal property of tensor
therefore identifies \(B\otimes N''\) with the cokernel of
\(B\otimes N\to B\otimes N'\). This also proves surjectivity of the map
to that cokernel. Apply this fact to the stalk formula and use Theorem
2.1 to obtain right exactness of \(f^*\).

Finally \((f_*F)(V)=F(f^{-1}V)\). Kernels of module-sheaf maps are
computed on sections, so \(f_*\) preserves kernels. It preserves finite
products as well, again by their sectionwise description. These are the
left-exactness assertions. No commutativity was used. \(\square\)

The tensor of two \emph{left} \(\mathcal O\)-modules used in later
lessons requires commutative \(\mathcal O\). The tensor in Theorem 4.4
has a specified right action on its first factor and therefore does not
need that hypothesis.

\textbf{Example 4.5 (an exact inverse image with a nonexact pullback).}
Take both spaces to be a point and use the coefficient map
\(\mathbb Z\to\mathbb Z/2\). Inverse image on abelian groups is the
identity. Module pullback is \(\mathbb Z/2\otimes_{\mathbb Z}-\). The
injection \(\mathbb Z\xrightarrow{2}\mathbb Z\) becomes the zero map
\(\mathbb Z/2\to\mathbb Z/2\), so pullback is not left exact. If the
specified structure map \(f^{-1}\mathcal O_Y\to\mathcal O_X\) is an
isomorphism, the formula instead identifies \(f^*\) with \(f^{-1}\), and
pullback is exact.

\subsection{5. Extension by zero and
skyscrapers}\label{sheaves-of-modules-on-a-ringed-space--5-extension-by-zero-and-skyscrapers}

\textbf{Lemma 5.1 (extension by zero from an open set).} Let
\(j:U\hookrightarrow X\) be open and \(B\) a sheaf of left
\(\mathcal O|_U\)-modules. For \(V\subset X\) open put
\[j_!B(V)=\{\,s\in B(U\cap V):\ \operatorname{supp}s\text{ is closed in }V\,\}.\]

\begin{enumerate}
\tightlist
\item
  \(j_!B\) is a sheaf of left \(\mathcal O\)-modules, with
  \((j_!B)_x=B_x\) for \(x\in U\) and \((j_!B)_x=0\) for \(x\notin U\).
  So \(j_!\) is exact.
\item
  \(\operatorname{Hom}_{\mathcal O}(j_!B,G)\cong\operatorname{Hom}_{\mathcal O|_U}(B,G|_U)\),
  naturally in \(B\) and \(G\).
\item
  \(\operatorname{Hom}_{\mathcal O}(j_!\mathcal O_U,G)\cong G(U)\) by
  \(\varphi\mapsto\varphi(1)\). For opens \(V\subset W\), the morphism
  \(j_{V!}\mathcal O_V\to j_{W!}\mathcal O_W\) that corresponds to
  \(1\in\mathcal O(V)=(j_{W!}\mathcal O_W)(V)\) is a monomorphism, and
  under these identifications it induces restriction \(G(W)\to G(V)\).
\end{enumerate}

\textbf{Proof.} (1) Supports of sections are closed in their domains.
Closedness of a support in \(V\) is a local condition on \(V\), so
sections of \(B\) glue within \(j_!B\), and \(j_!B\) is a sheaf. Module
operations do not enlarge supports. If \(x\in U\), the opens inside
\(U\) are cofinal among neighbourhoods of \(x\), and there \(j_!B=B\);
so \((j_!B)_x=B_x\). If \(x\notin U\) and \(s\in j_!B(V)\) with
\(x\in V\), then \(x\notin\operatorname{supp}s\), and this support is
closed in \(V\); so \(s\) vanishes near \(x\) and the stalk is \(0\).
Exactness is checked on stalks
(\hyperref[sheaves-of-modules-on-a-ringed-space--2-the-abelian-category-of-module-sheaves]{Theorem
2.1}(4)), and on each stalk \(j_!\) acts as the identity or as zero; so
\(j_!\) is exact.

(2) A morphism \(j_!B\to G\) restricts over \(U\) to a morphism
\(B\to G|_U\). Conversely, let \(\alpha:B\to G|_U\) and
\(s\in j_!B(V)\). The section \(\alpha(s)\) on \(U\cap V\) and the
section \(0\) on \(V\setminus\operatorname{supp}s\) agree on the
overlap. The two open sets cover \(V\), because
\(\operatorname{supp}s\subset U\cap V\). Glue them to a section of \(G\)
over \(V\). This defines the inverse bijection; linearity and naturality
are checked on the two pieces.

(3) By (2), a morphism \(j_!\mathcal O_U\to G\) is a left-linear map
\(\mathcal O_U\to G|_U\), and such a map is \(a\mapsto as\) with \(s\)
the image of \(1\in\mathcal O(U)\). For \(V\subset W\),
\((j_{W!}\mathcal O_W)(V)=\mathcal O(V)\). At \(x\in V\) the morphism is
the identity of \(\mathcal O_x\); at other points its source stalk is
\(0\). So it is injective on every stalk, hence a monomorphism by
\hyperref[sheaves-of-modules-on-a-ringed-space--2-the-abelian-category-of-module-sheaves]{Theorem
2.1}(4). Composing with it sends the morphism with
\(\varphi(1)=s\in G(W)\) to the one with value \(s|_V\) at \(1\).
\(\square\)

\textbf{Lemma 5.2 (skyscrapers).} Let \(x\in X\) and let \(J\) be a left
\(\mathcal O_x\)-module. Put \(x_*J(V)=J\) if \(x\in V\) and \(0\)
otherwise, with identity restrictions between opens containing \(x\),
and let \(\mathcal O(V)\) act through \(\mathcal O(V)\to\mathcal O_x\).
Then \(x_*J\) is a sheaf of left \(\mathcal O\)-modules, and
\[\operatorname{Hom}_{\mathcal O}(F,x_*J)\cong\operatorname{Hom}_{\mathcal O_x}(F_x,J),
\tag{5.2}\] naturally in \(F\) and \(J\).

\textbf{Proof.} Let \(x\in V\) and let \((V_i)\) cover \(V\). Some
\(V_i\) contains \(x\). Compatible sections agree on all \(V_i\) that
contain \(x\), since their pairwise overlaps contain \(x\); this common
value is the glued section, and it is unique. If \(x\notin V\), all
sections are \(0\). A morphism \(\varphi:F\to x_*J\) has components
\(\varphi_V\), for \(x\in V\), that are compatible with restriction. So
they define an \(\mathcal O_x\)-linear map \(F_x\to J\). Conversely,
given \(g:F_x\to J\), put \(\varphi_V(s)=g(s_x)\) if \(x\in V\) and
\(0\) otherwise. These constructions are inverse. \(\square\)

\subsection{6. Three local and global
computations}\label{sheaves-of-modules-on-a-ringed-space--6-three-local-and-global-computations}

\textbf{Example 6.1 (the two-point space).} Let \(S=\{o,c\}\), with
opens \(\varnothing,\{o\},S\). Fix any ring \(A\), and give \(S\) the
constant coefficient sheaf \(A_S\). A module sheaf is exactly a
homomorphism \(r:M\to N\) of left \(A\)-modules: its sections on \(S\)
are \(M\), on \(\{o\}\) are \(N\), and the only nontrivial restriction
is \(r\). Every covering of \(S\) contains \(S\), and the sheaf
condition imposes no further restriction. The stalk at \(c\) is \(M\),
since its only neighbourhood is \(S\); the stalk at \(o\) is \(N\).
Kernels, cokernels and exactness are therefore computed componentwise.

For \(j:\{o\}\to S\), restriction sends \((M\to N)\) to \(N\), extension
by zero sends \(Q\) to \((0\to Q)\), and direct image sends \(Q\) to
\((Q\xrightarrow{1}Q)\). A section of \(j_!Q\) on \(S\) has support
contained in \(\{o\}\); that subset is not closed in \(S\), so only the
zero section is allowed. For the closed point, \(c_*Q=(Q\to0)\). For the
open point, the skyscraper is \(o_*Q=(Q\xrightarrow{1}Q)\). Thus a
skyscraper at a point need not have zero stalks at all other points. No
closed-point assumption appeared in Lemma 5.2.

As a concrete cokernel, map \((A\xrightarrow{1}A)\) to \((A\to0)\) by
the identity at \(c\) and zero at \(o\). Its kernel is \((0\to A)\),
which is precisely \(j_!A\). This gives the short exact sequence

\[0\to j_!A\to o_*A\to c_*A\to0.\]

It records the open and closed pieces of this space without any
separation hypothesis. In the lesson on supports the same space will
make a localization triangle completely explicit.

\textbf{Example 6.2 (a surjection with no global lift).} On
\(X=\mathbb C^{\times}\), let \(\mathcal C\) be the sheaf of continuous
complex-valued functions, written additively, and let
\(\mathcal C^{\times}\) be the sheaf of continuous nowhere-zero
functions, written multiplicatively. They are sheaves of abelian groups,
hence left \(\mathbb Z_X\)-modules. They are not being regarded as
modules over the sheaf of complex-valued functions. Exponentiation gives

\[0\longrightarrow(2\pi i\mathbb Z)_X
\longrightarrow\mathcal C\xrightarrow{\exp}\mathcal C^{\times}
\longrightarrow0.\]

The kernel consists of continuous functions with values in the discrete
set \(2\pi i\mathbb Z\), so it is the indicated locally constant sheaf.
Every nonzero complex number has a neighbourhood on which a continuous
logarithm is defined: choose a sufficiently small disk avoiding zero and
a continuous branch of the argument there. For a nowhere-zero function,
continuity places a neighbourhood of each point in such a disk.
Composition gives a local logarithm, proving stalkwise surjectivity.

The global section \(z\mapsto z\) has no continuous logarithm on \(X\).
Suppose one existed, and restrict it to \(z=e^{2\pi it}\),
\(0\le t\le1\). The function

\[t\longmapsto\frac{g(e^{2\pi it})-2\pi it}{2\pi i}\]

takes values in \(\mathbb Z\) and is continuous, so it is constant. This
forces \(g(e^{2\pi i})-g(1)=2\pi i\), although the two arguments are the
same point. The contradiction proves failure of global surjectivity. The
cokernel sheaf is zero even though the cokernel on global sections is
nonzero.

\textbf{Example 6.3 (extension by zero at an endpoint).} Let
\(X=[0,1]\), \(U=(0,1)\), and let \(\mathcal O_X\) be the sheaf of
continuous real-valued functions. A global section of
\(j_!\mathcal O_U\) is a continuous function on \(U\) whose support is
closed in \(X\) and contained in \(U\). Its support is compact and
misses both endpoints. Consequently the section vanishes on a
neighbourhood of each endpoint. Conversely any such function has the
required support and extends by zero to a continuous function on \(X\).
The constant function one on \(U\) is a section of \(j_*\mathcal O_U\),
but not of \(j_!\mathcal O_U\) on \(X\). At either endpoint
\(j_!\mathcal O_U\) has zero stalk; at an interior point it has stalk
\(\mathcal O_x\). These stalks, rather than the global section module,
give the later flatness calculation.

\subsection{7. Exercises with complete
solutions}\label{sheaves-of-modules-on-a-ringed-space--7-exercises-with-complete-solutions}

The first two exercises test the exactness criteria. The next two
compare constructions. The last two concern finiteness and topology.

\textbf{Exercise 1 (easy: an isomorphism criterion).} Show directly that
a morphism of module sheaves with isomorphisms on all stalks has an
inverse sheaf morphism. Explain why the inverse is module-linear.

\textbf{Solution.} If two sections have the same image, their germs have
the same image at every point. Stalkwise injectivity makes their germs
equal, hence the sections equal. Given a section \(t\) in the target,
lift each of its germs, choose local representatives, and shrink until
the image of each representative is \(t\) on that neighbourhood. On
overlaps the representatives have the same image, so the just-proved
injectivity makes them agree. They glue to a unique section in the
source. These inverses commute with restriction by uniqueness. The
inverse sends \(a t\) to \(a s\), because the forward map sends \(a s\)
to \(a t\), again with a unique preimage. It preserves sums for the same
reason. Thus it is a linear inverse sheaf morphism.

\textbf{Exercise 2 (medium: filtered and unfiltered colimits).} Prove
exactness of filtered colimits of module sheaves. Then verify every
arrow in the pushout counterexample in Section 1 and compute its three
pushouts.

\textbf{Solution.} Stalks of a colimit are the colimits of the stalks by
Theorem 3.1(2). For a diagram of short exact sequences, the stalk
diagrams are exact by Theorem 2.1. Lemma 1.1 gives exactness after
taking their filtered colimits. Theorem 2.1 then gives exactness of the
sheaf sequence. In the counterexample, at the left object the sequence
is \(0\to\mathbb Z\xrightarrow{1}\mathbb Z\to0\to0\), at the middle
object it is \(0\to0\to\mathbb Z\xrightarrow{1}\mathbb Z\to0\), and at
the right all terms vanish. The square for the arrow middle-to-left in
the first morphism is \(0\to\mathbb Z\) followed by the identity on the
left, and on the other side the zero inclusion followed by the identity
middle-to-left; both composites from the zero group are zero. In the
second morphism, both composites from the middle \(\mathbb Z\) to the
last left object zero are zero. The middle-to-right squares are zero
throughout. A module pushout of \(L\leftarrow M\to R\) is
\((L\oplus R)/\{(a(m),-b(m)):m\in M\}\), as its universal property
shows. This gives \(\mathbb Z\), \(0\) and \(0\). Thus the pushout
functor does not preserve the displayed injection of diagrams.

\textbf{Exercise 3 (medium: extension by zero).} Given a map
\(B\to G|_U\), construct its adjoint \(j_!B\to G\) on every open set.
Show this construction gives the identity in both adjunction
compositions, and deduce exactness of \(j_!\).

\textbf{Solution.} For \(s\in j_!B(V)\), glue its image on \(V\cap U\)
with zero on \(V\setminus\operatorname{supp}s\). These sets cover \(V\),
and the images agree on the intersection because \(s\) is zero there.
Restriction commutes with gluing, as do sums and scalar multiplication.
Restricting the resulting map to \(U\) recovers the original map, since
\(j_!B|_U=B\). Starting with a map \(j_!B\to G\), the reconstructed
section has the same restriction on \(V\cap U\) and is zero off the
support; these are also the restrictions of the original image.
Uniqueness of gluing recovers it. At a point of \(U\), the stalk functor
of \(j_!\) is the identity; elsewhere it is zero. Both preserve exact
sequences, so stalkwise exactness proves exactness of \(j_!\).

\textbf{Exercise 4 (medium: scalar extension).} For a homomorphism of
possibly noncommutative rings \(A\to B\), write the unit and counit of
\(B\otimes_A-\dashv\operatorname{Res}_A^B\). Check both triangle
identities. Apply the formulas to the pullback adjunction in Theorem
4.4.

\textbf{Solution.} The unit is \(M\to\operatorname{Res}(B\otimes_A M)\),
\(m\mapsto1\otimes m\). It is \(A\)-linear because
\(1\otimes am=a\otimes m\), with \(a\) acting through \(A\to B\). The
counit is \(B\otimes_A\operatorname{Res}N\to N\),
\(b\otimes n\mapsto bn\), which is balanced and \(B\)-linear. The first
composite sends \(b\otimes m\) to \(b\otimes(1\otimes m)\) and then to
\(b\otimes m\). The second sends \(n\) to \(1\otimes n\) and then to
\(n\). These verify the triangle identities. For sheaves the formulas
hold on local tensor representatives and pass through sheafification.
Compose with the inverse-image/direct-image adjunction: the pullback
unit sends a section \(g\) over \(V\subset Y\) to \(1\otimes f^{-1}g\)
over \(f^{-1}V\). The pullback counit sends a local tensor
\(b\otimes s\), where \(s\) is locally represented by a section of
\(f_*F\), to \(b\) times its restriction to the relevant open subset of
\(X\). These are precisely the formulas in the proof of Theorem 4.4.

\textbf{Exercise 5 (hard: sections of an infinite direct sum).} Let
\(X=\mathbb N\) be discrete, let \(A\ne0\) be a ring, and let \(F_n\) be
the skyscraper \(A\) at \(n\). Compute \(\Gamma(X,\bigoplus_nF_n)\) and
compare it with \(\bigoplus_n\Gamma(X,F_n)\).

\textbf{Solution.} A sheaf on a discrete space has sections on an open
subset \(V\) equal to the product of its stalks at the points of \(V\).
Indeed, the singleton opens cover \(V\), their intersections are empty,
and the sheaf condition says that all choices of singleton sections glue
uniquely. The direct sum sheaf has stalk \(A\) at each \(n\), since at
that stalk only \(F_n\) is nonzero. Its sections on \(X\) are therefore
\(\prod_nA\). In contrast \(\Gamma(X,F_n)=A\), so the sum of the global
sections is \(\bigoplus_nA\). The canonical map is the inclusion of
finitely supported sequences into all sequences, and is not onto: the
sequence with every entry one is absent. The open \(X\) is not
quasi-compact, as the singleton cover has no finite subcover. This does
not contradict Theorem 3.1(5).

\textbf{Exercise 6 (hard: products need not preserve epimorphisms).} Let
\(X=\mathbb R\), let \(k\) be a field, and put \(U_n=(-1/n,1/n)\). For
each \(n\), extend the identity germ at zero to a morphism
\(j_{n!}k_{U_n}\to0_*k\). Prove it is an epimorphism and show that the
product of these maps is not an epimorphism.

\textbf{Solution.} By Lemma 5.2 a map to \(0_*k\) is determined by the
map of stalks at zero. Choose the identity \(k\to k\). At zero it is
onto; at every other point the target stalk is zero because zero is
closed in \(\mathbb R\). Thus Theorem 2.1 makes it an epimorphism.
Products are computed on sections. A germ at zero of a section of
\(\prod_nj_{n!}k_{U_n}\) is represented on one common neighbourhood
\(V\) of zero. Choose a connected interval \(W\) containing zero inside
\(V\). Each factor restricts to a locally constant function on
\(W\cap U_n\), whose support is closed in \(W\). For every sufficiently
large \(n\), the interval \(U_n\) is contained in \(W\). A nonzero
constant on \(U_n\) would have support \(U_n\), which is not closed in
\(W\); hence that factor is zero. The images of the germ at zero thus
have only finitely many nonzero coordinates. The target product sheaf
has stalk \(\prod_nk\) at zero: on every neighbourhood of zero its
sections are that product and all restrictions are identities. The germ
with every coordinate one is not in the image. The product map is not
onto on the stalk, hence is not an epimorphism. This argument uses one
common neighbourhood for all factors; commuting a stalk with an infinite
product would lose exactly that condition.

\subsection{References}\label{sheaves-of-modules-on-a-ringed-space--references}

\begin{itemize}
\tightlist
\item
  \textbf{{[}Stacks{]}} The Stacks project authors, \emph{The Stacks
  project}, chapter \emph{Sheaves of Modules},
  \href{https://github.com/KokunoYumeto/unofficial-stacks-project-ai-drafts/blob/565b10e987aba5969b21145a0833f42d69f96790/modules.tex\#L60}{kernels
  and cokernels},
  \href{https://github.com/KokunoYumeto/unofficial-stacks-project-ai-drafts/blob/565b10e987aba5969b21145a0833f42d69f96790/modules.tex\#L149}{the
  abelian category}, and
  \href{https://github.com/KokunoYumeto/unofficial-stacks-project-ai-drafts/blob/565b10e987aba5969b21145a0833f42d69f96790/modules.tex\#L207}{limits
  and colimits}, in the AI Integrated Stacks Project edition.
  Stacks-derived material retains GFDL-1.2-or-later; the edition is an
  unofficial AI-integrated derivative.
\item
  \textbf{{[}Kashiwara--Schapira 1990{]}} Masaki Kashiwara and Pierre
  Schapira, \emph{Sheaves on Manifolds}, Grundlehren der mathematischen
  Wissenschaften 292, Springer, 1990. For module sheaves and sheaf
  operations.
\item
  \textbf{{[}Kashiwara--Schapira 2006{]}} Masaki Kashiwara and Pierre
  Schapira, \emph{Categories and Sheaves}, Grundlehren der
  mathematischen Wissenschaften 332, Springer, 2006. Scholarly reference
  for categorical constructions. The statements used in this lesson have
  complete proofs above.
\end{itemize}

\section{Complexes, cones and
localization}\label{complexes-cones-and-localization}

\emph{Written and edited by GPT-6.1 Sol (OpenAI), in Codex at Ultra,
October 2026. Mathematically self-checked by the writing AI. Original
contributions are CC0; the combined course is distributed under
GFDL-1.2-or-later. Full authorship and source attribution appear in the
\hyperref[sources-and-authorship]{course notice}.}

An exact sequence contains more information than its three objects: it
specifies how a quotient attaches to a subobject. A complex contains
still more attachment information, in each degree. Mapping cones retain
that information when a map is replaced by a triangle. Localization then
lets us compare complexes that have the same cohomology without
requiring a chain homotopy equivalence between them.

This common reading precedes the resolution lessons. Its algebra
prerequisite is Wen-Wei Li's \emph{Methods of Algebra, Volume 2}:
\href{https://kokunoyumeto.github.io/methods-of-algebra-volume-2-en/\#unit-chapter2-unit-021}{``Definition
of an Abelian Category''}, and
\href{https://kokunoyumeto.github.io/methods-of-algebra-volume-2-en/\#unit-chapter2-unit-023}{``Some
Diagram Lemmas''}, especially the exactness criterion by epimorphic
lifting, the Snake Lemma and the Five Lemma.
\href{https://kokunoyumeto.github.io/methods-of-algebra-volume-2-en/\#unit-chapter3-unit-037}{``Complexes
on an Abelian Category''} proves the long exact cohomology sequence
using that Snake Lemma. We recall its construction here and supply the
cone and localization arguments in full.

The cone and triangulation arguments follow the Stacks project authors'
\emph{Derived Categories}, in the
\href{https://github.com/KokunoYumeto/unofficial-stacks-project-ai-drafts/blob/565b10e987aba5969b21145a0833f42d69f96790/derived.tex}{AI
Integrated Stacks Project edition}. We develop localization by roofs and
common refinements, then prove that cone triangles descend to the
derived category. Source attribution and the licence for adapted
passages appear in the \hyperref[sources-and-authorship]{course notice}.

\subsection{1. Maps of complexes and
homotopies}\label{complexes-cones-and-localization--1-maps-of-complexes-and-homotopies}

Let \(\mathcal A\) be an abelian category. A \textbf{cochain complex}
\(K\) consists of objects \(K^n\) and differentials
\(d_K^n:K^n\to K^{n+1}\) satisfying \(d_K^{n+1}d_K^n=0\). A chain map
\(f:K\to L\) satisfies \(d_Lf=fd_K\). Its cohomology is

\[H^n(K)=\ker(d_K^n)/\operatorname{im}(d_K^{n-1}).\]

We set \(K[r]^n=K^{n+r}\), with differential \((-1)^r d_K\). A homotopy
from \(g\) to \(f\) is a family \(s^n:K^n\to L^{n-1}\) with
\(f-g=d_Ls+sd_K\). These homotopies form an equivalence relation: zero,
negation and addition give reflexivity, symmetry and transitivity.
Precomposing or postcomposing a homotopy with a chain map gives a
homotopy of composites. Thus homotopy classes form an additive category
\(K(\mathcal A)\), with the same objects as the category of complexes.
Its finite sums and zero object are those of complexes. A
\textbf{homotopy equivalence} is an isomorphism in this category. A
complex is \textbf{contractible} when its identity is null-homotopic.

The group-valued Hom complex, as distinct from an internal sheaf Hom, is

\[\operatorname{Hom}^{r}(K,L)=\prod_{n\in\mathbb Z}
\operatorname{Hom}_{\mathcal A}(K^n,L^{n+r}),\qquad
D(h)=d_Lh-(-1)^rhd_K.
\tag{1.1}\]

\textbf{Lemma 1.1.} Formula (1.1) defines a complex of abelian groups,
and

\[H^r\operatorname{Hom}^{\bullet}(K,L)
=\operatorname{Hom}_{K(\mathcal A)}(K,L[r]).
\tag{1.2}\]

\textbf{Proof.} Applying the differential twice cancels the two terms
involving \(d_Lhd_K\); the other terms vanish because \(d_L^2=d_K^2=0\).
A degree-\(r\) family is a cocycle precisely when \(d_Lh=(-1)^rhd_K\),
which is the chain-map equation into \(L[r]\). If \(h=D(v)\) with \(v\)
of degree \(r-1\), then

\[h=d_Lv+(-1)^rvd_K
=d_{L[r]}\bigl((-1)^rv\bigr)+\bigl((-1)^rv\bigr)d_K.\]

So boundaries are exactly the null-homotopic maps into \(L[r]\).
Conversely a homotopy \(s\) gives the boundary of \((-1)^rs\). This
proves (1.2). Composition of homogeneous families obeys

\[D(gh)=D(g)h+(-1)^{\deg g}gD(h):\]

expanding the right side cancels its two middle terms and leaves
\(d_Mgh-(-1)^{\deg g+\deg h}ghd_K\). In particular composition descends
to cohomology and hence to the homotopy category. \(\square\)

\textbf{Lemma 1.2 (cohomology of a short exact sequence).} A degreewise
short exact sequence \(0\to A\xrightarrow{a}B\xrightarrow{b}C\to0\) has
a natural long exact cohomology sequence. Homotopic maps induce the same
cohomology map.

\textbf{Proof.} For the second assertion, on \(\ker d_K^n\) the summand
\(sd_K\) vanishes and \(d_Ls\) has image in
\(\operatorname{im}d_L^{n-1}\). It therefore induces zero on cohomology.

For the first assertion, apply the Snake Lemma to the exact rows

\[\operatorname{coker}d_A^{n-2}\longrightarrow
\operatorname{coker}d_B^{n-2}\longrightarrow
\operatorname{coker}d_C^{n-2}\longrightarrow0,\]
\[0\longrightarrow\ker d_A^n\longrightarrow
\ker d_B^n\longrightarrow\ker d_C^n,\]

whose vertical maps are induced by \(d^{n-1}\). The kernels of these
vertical maps are \(H^{n-1}\); their cokernels are \(H^n\). The first
row is right exact by the cokernel part of the Snake Lemma applied to
the degreewise exact diagram with vertical maps \(d^{n-2}\); the second
is left exact by its kernel part for \(d^n\). Thus the new Snake Lemma
sequence joins \(H^{n-1}(C)\) to \(H^n(A)\) and gives exactly the six
consecutive terms of the cohomology sequence. Varying \(n\) joins the
sequences. Naturality follows from the naturality of the Snake Lemma
construction by kernels, cokernels and pullbacks.

In modules, its boundary has the familiar precise description: lift a
cocycle \(c\) to \(b\), write \(d_Bb=a(a')\), and send \([c]\) to
\([a']\). A changed lift adds a boundary, and a changed cocycle by a
boundary also changes \(a'\) by a boundary. In an arbitrary abelian
category these lifts are performed after pullback along an epimorphism,
as in the exactness criterion in the prerequisite; the kernel and
cokernel description just given supplies the resulting canonical
morphism. \(\square\)

A \textbf{quasi-isomorphism} is a map inducing an isomorphism on every
cohomology object. A complex is \textbf{acyclic} when all its cohomology
objects vanish. A homotopy equivalence is a quasi-isomorphism by Lemma
1.2, but an acyclic complex need not be contractible.

\subsection{2. Cones and triangles from split
sequences}\label{complexes-cones-and-localization--2-cones-and-triangles-from-split-sequences}

For \(f:K\to L\), define

\[C(f)^n=L^n\oplus K^{n+1},\qquad
d(l,k)=(d_Ll+fk,-d_Kk).
\tag{2.1}\]

Squaring this differential gives zero because \(f\) is a chain map. Let
\(i:L\to C(f)\) be inclusion and \(p:C(f)\to K[1]\) projection. Our
\textbf{cone triangle} is

\[K\xrightarrow fL\xrightarrow iC(f)\xrightarrow{-p}K[1].
\tag{2.2}\]

The minus sign is intentional. With this convention, a triangle attached
to a short exact sequence has the Snake Lemma\textquotesingle s positive
boundary.

\textbf{Lemma 2.1 (comparison of cones).} If \(bf-f'a=d h+h d\), the
maps

\[c^n=\begin{pmatrix}b^n&h^{n+1}\\0&a^{n+1}\end{pmatrix}:
C(f)^n\longrightarrow C(f')^n\]

give a morphism of cone triangles in \(K(\mathcal A)\).

\textbf{Proof.} The diagonal entries of \(d c=c d\) are the chain-map
equations for \(a,b\). Equality of the upper-right entries is
\(d h+f'a=bf-hd\), exactly the assumed homotopy equation. The lower-left
entries are zero. Also \(c i=i'b\) and \(p'c=a[1]p\) hold already as
chain maps. This proves all triangle squares, including the one with
negative projections. The choice of \(h\) need not be unique; the third
triangle map is not canonically determined by a square in the homotopy
category. \(\square\)

\textbf{Lemma 2.2 (split sequences).} Suppose \(0\to A\to B\to C\to0\)
is split in each degree. Write

\[B^n=A^n\oplus C^n,\qquad
d_B=\begin{pmatrix}d_A&\delta\\0&d_C\end{pmatrix}.
\tag{2.3}\]

Then \(\delta:C\to A[1]\) is a chain map, and the triangle
\(A\to B\to C\xrightarrow\delta A[1]\) is isomorphic to the cone
triangle of \(A\to B\).

\textbf{Proof.} The upper-right entry of \(d_B^2=0\) is
\(d_A\delta+\delta d_C=0\), the asserted chain-map equation. In degree
\(n\), the cone of inclusion is \(A^n\oplus C^n\oplus A^{n+1}\). Define

\[r(a,c,a')=c,\qquad t(c)=(0,c,-\delta c).\]

Both are chain maps by (2.3), and \(rt=1\). The homotopy
\(H(a,c,a')=(0,0,a)\), landing in degree \(n-1\), satisfies

\[(dH+Hd)(a,c,a')=(a,0,a'+\delta c)=(1-tr)(a,c,a').\]

Thus \(r,t\) are inverse in \(K\). Inclusion of \(B\) is carried to its
quotient map into \(C\), and \((-p)t=\delta\). These give the claimed
triangle isomorphism. Formula (2.3) also shows that \(H^n(\delta)\) is
exactly the boundary constructed in Lemma 1.2. \(\square\)

\textbf{Lemma 2.3 (replacement by a split injection).} Every map
\(f:K\to L\) factors as a degreewise split injection
\(K\to\widetilde L\) followed by a homotopy equivalence
\(\widetilde L\to L\). A finite string of maps can be replaced by
degreewise split injections, with a commuting diagram of homotopy
equivalences to the original string.

\textbf{Proof.} Put \(\widetilde L=L\oplus C(1_K)\), with degree-\(n\)
coordinates \((l,k,k')\), differential \((d_Ll,d_Kk+k',-d_Kk')\), and
injection \(k\mapsto(fk,k,0)\). Projection \(\pi\) to \(L\) has
chain-map section \(s(l)=(l,0,0)\). The homotopy \(H(l,k,k')=(0,0,k)\)
gives

\[(dH+Hd)(l,k,k')=(0,k,-d_Kk)+(0,0,d_Kk+k')=(0,k,k')=1-s\pi.\]

The injection is split by projection onto \(K^n\). For a string, start
with its first object, perform this construction on the first map, then
on the composite from the new second object to the original third
object, and continue. Every resulting square commutes as a square of
chain maps. Shifts and these replacements preserve boundedness above,
below, or in both directions. \(\square\)

\textbf{Lemma 2.4 (representable exactness and cone uniqueness).} For a
cone triangle, applying \(\operatorname{Hom}_K(T,-)\) or
\(\operatorname{Hom}_K(-,T)\) gives a long exact sequence after shifts.
A morphism of cone triangles which is an isomorphism on two objects is
an isomorphism on the third.

\textbf{Proof.} There is a chain isomorphism

\[\operatorname{Hom}^{\bullet}(T,C(f))
=C\bigl(\operatorname{Hom}^{\bullet}(T,f)\bigr).\]

It follows by distributing products over the two summands in (2.1); its
differential has exactly the matrix (2.1). For the contravariant version
put \(U=\operatorname{Hom}^{\bullet}(L,T)\),
\(V=\operatorname{Hom}^{\bullet}(K,T)\), and \(u=f^*:U\to V\). A
degree-\(r\) map from the cone has coordinates
\((a,b)\in U^r\oplus V^{r-1}\) and differential

\[(a,b)\longmapsto\bigl(Da,Db+(-1)^{r+1}u(a)\bigr).\]

The map \((a,b)\mapsto((-1)^{r-1}b,a)\) identifies this complex with
\(C(u)[-1]\), whose differential is \((v,a)\mapsto(-Dv-u(a),Da)\). Lemma
1.2 for these group complexes and Lemma 1.1 give both exact sequences.
Signs of individual arrows do not change exactness.

For a morphism of triangles, apply either representable functor and take
five consecutive terms with the possibly unknown map in the middle. The
four surrounding maps are isomorphisms when the other two object maps
are isomorphisms. The Five Lemma makes the middle map an isomorphism for
every \(T\). If \(c:Z\to Z'\) induces bijections
\(\operatorname{Hom}(T,Z)\to\operatorname{Hom}(T,Z')\) for every \(T\),
surjectivity for \(T=Z'\) gives a right inverse \(v\), and injectivity
for \(T=Z\) gives \(vc=1\). Thus \(c\) is an isomorphism. Rotating the
exact sequences proves the assertion for any two object maps. In
particular cones of a fixed morphism are unique up to an isomorphism
fixing its first two objects. Lemma 2.1 gives existence of that
comparison, and this argument gives invertibility. \(\square\)

\subsection{3. The triangulated structure of the homotopy
category}\label{complexes-cones-and-localization--3-the-triangulated-structure-of-the-homotopy-category}

A triangle means \(X\xrightarrow fY\xrightarrow gZ\xrightarrow hX[1]\).
It is \textbf{distinguished} in \(K(\mathcal A)\) when it is isomorphic
to (2.2). A triangulated category is an additive category with an
additive shift equivalence and distinguished triangles satisfying four
requirements: identity and completion triangles exist and are closed
under isomorphism; rotation \((f,g,h)\mapsto(g,h,-f[1])\) preserves
distinguished triangles in both directions; a commuting square between
their first maps extends to a morphism of triangles; and composable maps
have compatible cone triangles forming an octahedron. We give the exact
octahedron maps in the proof, rather than using a picture as an axiom.

\textbf{Theorem 3.1.} These definitions make \(K(\mathcal A)\) a
triangulated category. The full subcategories of bounded-below,
bounded-above and bounded complexes are triangulated.

\textbf{Proof.} The isomorphism and completion requirements follow from
the definition. The cone \(C(1_X)\) is contractible by Lemma 2.3, so its
cone triangle is isomorphic to \(X\xrightarrow1X\to0\to X[1]\).

For rotation, \(0\to L\to C(f)\to K[1]\to0\) is split in each degree,
with off-diagonal map \(f[1]\). By Lemma 2.2 its triangle \((i,p,f[1])\)
is distinguished. Changing the sign on its third object identifies it
with \((i,-p,-f[1])\), the rotation of (2.2). Conversely, suppose a
rotated triangle is represented by a split sequence
\(0\to A\to B\to C\to0\) with boundary \(\delta\). The inverse rotation
is \(C[-1]\xrightarrow{-\delta[-1]}A\to B\to C\). Formula (2.3)
identifies \(B\) with \(C(-\delta[-1])\) by the map
\((a,c)\mapsto(a,-c)\); its projection into \(C\) becomes the negative
cone projection. Hence this inverse rotation is distinguished. Every
cone triangle has a split-sequence representative: use Lemma 2.3 on its
first map, Lemma 2.1 and Lemma 2.4 to compare cones, and Lemma 2.2 on
the new split injection. Thus the converse applies to all distinguished
triangles.

The square-extension requirement is precisely Lemma 2.1 after replacing
the rows by cone triangles.

For the octahedron take composable maps
\(A\xrightarrow\alpha B\xrightarrow\beta C\). Lemma 2.3 replaces this
string by degreewise split injections. Lemmas 2.1, 2.2 and 2.4 identify
their three cone triangles with the triangles of the three quotient
sequences. Choose degreewise decompositions \(B=A\oplus Q_1\) and
\(C=A\oplus Q_1\oplus Q_3\). Their differentials have the forms

\[d_B=\begin{pmatrix}d_A&u\\0&d_1\end{pmatrix},\qquad
d_C=\begin{pmatrix}d_A&u&v\\0&d_1&w\\0&0&d_3\end{pmatrix}.
\tag{3.1}\]

Set \(Q_2=C/A=Q_1\oplus Q_3\), with differential
\(\left(\begin{smallmatrix}d_1&w\\0&d_3\end{smallmatrix}\right)\). The
three boundary maps are \(u\), \((u,v)\), and \((v,w)\). The split
sequence \(0\to Q_1\to Q_2\to Q_3\to0\) has boundary \(w\), so its
triangle is distinguished. Its first two maps are inclusion
\(a:Q_1\to Q_2\) and projection \(b:Q_2\to Q_3\). The triples
\((1_A,\beta,a)\) and \((\alpha,1_C,b)\) are morphisms of the three
original triangles. All their squares commute strictly except possibly
the second boundary square. The difference

\[(v,w)b-\alpha[1](u,v)=\begin{pmatrix}-u&0\\0&w\end{pmatrix}\]

is the homotopy boundary for \(H:Q_2^n\to B[1]^{n-1}=B^n\),
\(H(q_1,q_3)=(0,q_1)\): substituting (3.1) into \(d_{B[1]}H+Hd_{Q_2}\)
gives \((-u q_1,w q_3)\). Finally \(w=p_1[1](v,w)\), where
\(p_1:B\to Q_1\). Thus

\[Q_1\xrightarrow aQ_2\xrightarrow bQ_3
\xrightarrow{p_1[1]\delta_3}Q_1[1]\]

is the required fourth triangle, with exactly the required comparison
maps. Transport this diagram through the homotopy equivalences of the
string. For any initially chosen distinguished triangles of the three
maps, Lemma 2.4 compares them by isomorphisms fixing their first two
objects. Transport through those too. This proves the full octahedron
requirement, not only a special choice of cones.

All complexes, shifts, cones and replacements used here preserve each
stated boundedness condition. The same proof therefore works in the
three bounded homotopy categories. \(\square\)

\textbf{Corollary 3.2.} An additive functor
\(\Phi:\mathcal A\to\mathcal B\) induces a triangulated functor on
homotopy categories. A quasi-isomorphism \(f\) is characterized by
acyclicity of \(C(f)\).

\textbf{Proof.} Applying \(\Phi\) degreewise preserves differentials,
homotopies, shifts, finite sums and every entry of (2.1), hence cone
triangles. For the second assertion apply Lemma 1.2 to
\(0\to L\to C(f)\to K[1]\to0\); its connecting map is \(H(f[1])\).
Exactness proves that the middle complex is acyclic exactly when each
\(H^n(f)\) is invertible. \(\square\)

\subsection{4. Fractions at
quasi-isomorphisms}\label{complexes-cones-and-localization--4-fractions-at-quasi-isomorphisms}

We work in a larger universe if necessary: objects and arrows from the
chosen small-module universe form sets in that larger universe. This
makes the following constructions legitimate before proving that the
derived Hom sets are small in the original universe. The resolution
lessons will establish that stronger smallness conclusion for module
sheaves.

Let \(S\) be the quasi-isomorphisms in \(K(\mathcal A)\). It contains
identities, is closed under composition, and has the two-out-of-three
property, because these assertions hold for isomorphisms after every
\(H^n\).

\textbf{Lemma 4.1 (fraction squares and cancellation).} Given
\(f:M\to Y\) and \(t:N\to Y\) with \(t\in S\), there are \(r:P\to M\) in
\(S\) and \(a:P\to N\) with \(fr=ta\) in \(K\). If \(tf=tg\) for
\(t:Y\to Z\) in \(S\), there is \(r:P\to X\) in \(S\) with \(fr=gr\).
The dual assertions, with all arrows reversed, also hold.

\textbf{Proof.} For the square take \(P=C((f,-t):M\oplus N\to Y)[-1]\).
Its coordinates in degree \(n\) are \((y,m,n')\), with differential
\((-d_Yy-fm+tn',d_Mm,d_Nn')\). Let \(r,a\) be projections to \(m,n'\).
Their required equality holds up to the homotopy \((y,m,n')\mapsto-y\).
The kernel of \(r\) is \(C(-t)[-1]\), which is acyclic by Corollary 3.2.
Since \(r\) is degreewise split onto, Lemma 1.2 makes \(r\) a
quasi-isomorphism. This proves the square assertion.

For cancellation put \(v=f-g\). Choose a homotopy \(tv=d_Zh+hd_X\). The
map \((-h,v):X\to C(t)[-1]\) is a chain map: its first
component\textquotesingle s chain-map equation is \(d_Zh-tv=-hd_X\),
exactly the chosen homotopy equation. Its second component projects to
\(v\). Write this map as \(e:X\to E\), where \(E=C(t)[-1]\) is acyclic.
Projection \(r:C(e)[-1]\to X\) is a quasi-isomorphism, since its kernel
\(E[-1]\) is acyclic. Moreover \(er\) is null-homotopic by projection
onto the first cone coordinate with a minus sign, exactly as in the
square computation. Therefore \(vr=0\).

For the dual square take the cone of \((t,-f):M\to N\oplus Y\); the
inclusion of \(Y\) has acyclic quotient \(C(t)\), and the two composites
agree up to the cone\textquotesingle s homotopy. For the dual
cancellation, representable exactness in Lemma 2.4 factors a map \(v\)
satisfying \(vt=0\) through the acyclic cone \(C(t)\). If this factor is
\(e:C(t)\to Z\), the inclusion \(Z\to C(e)\) is a quasi-isomorphism and
kills \(e\) up to homotopy. It consequently kills \(v\). These are the
asserted reversed constructions. \(\square\)

A \textbf{roof} from \(X\) to \(Y\) is
\(X\xleftarrow{s}M\xrightarrow{f}Y\), with \(s\in S\). Declare two roofs
equivalent if they have a common refinement: maps \(u:P\to M\),
\(v:P\to M'\) such that \(su=s'v\in S\) and \(fu=f'v\). Two-out-of-three
then implies \(u,v\in S\).

\textbf{Theorem 4.2 (localization).} Roof classes form a category
\(D(\mathcal A)\). The map \(Q:K(\mathcal A)\to D(\mathcal A)\),
\(f\mapsto(1,f)\), inverts quasi-isomorphisms and is their localization.
Its morphism sets are abelian groups, and \(Q\) is additive. Moreover
\(Qf=Qg\) if and only if \(fr=gr\) for some quasi-isomorphism \(r\) into
their common source.

\textbf{Proof.} Equivalence is reflexive and symmetric. For
transitivity, two refinements through a middle roof give two
quasi-isomorphisms into its middle object. Lemma 4.1 gives a common
square; two-out-of-three makes both new arrows quasi-isomorphisms.
Composing the refinements gives the third required refinement. Thus the
relation is transitive.

Compose roofs \((s,f):X\to Y\) and \((t,g):Y\to Z\) by using Lemma 4.1
to choose \(r:P\to M\) in \(S\) and \(a:P\to N\) with \(fr=ta\), and set
the composite to \((sr,ga)\). Two choices \((r,a)\), \((r',a')\) have a
common refinement of \(r,r'\). On it, \(ta=ta'\); cancellation in Lemma
4.1 refines once more to make \(a=a'\). Their resulting roofs are
therefore equivalent. Replacing the first roof by a refinement is
handled by a common square between its refinement arrow and \(r\);
replacing the second by a refinement is handled by a square between
\(a\) and that refinement arrow. Each produces precisely a refinement of
the original composite. Since any two equivalent roofs admit such
refinements, composition is well defined.

For associativity take three roofs \((s,f),(t,g),(u,h)\). First form a
square \(g r=u b\), with \(r:P\to N\) in \(S\). Next form
\(f r'=t r c\), with \(r':P'\to M\) in \(S\). For the other order of
composition, choose \((r',rc)\) as the square for the first two roofs.
Their numerator is \(grc=ubc\), so choose the identity on \(P'\) for the
last square. Both parenthesizations give \((sr',hbc)\). Independence of
choices, already proved, gives associativity for every choice. The roof
\((1,1)\) is an identity by choosing identity squares.

For \(s:M\to X\) in \(S\), the roof \((s,1_M)\) is inverse to \(Q(s)\);
either composite refines the identity roof. A functor \(F\) inverting
\(S\) must send \((s,f)\) to \(F(f)F(s)^{-1}\). Common refinements leave
this expression unchanged, and the composition square proves that it
preserves composition. It defines the unique factorization through
\(Q\). Natural transformations also descend: a naturality square for
\(s\) remains natural after inverting \(F(s)\) and \(G(s)\), and every
roof is a composite of such an inverse and an ordinary arrow.

For \(Qf=Qg\), a common refinement of their identity roofs is exactly an
\(r\in S\) with \(fr=gr\); the converse is the same refinement. This is
the stated equality criterion.

To add roofs, first give them a common denominator \(s\) by Lemma 4.1
and set \((s,f)+(s,g)=(s,f+g)\). If two choices or representatives are
made, refine their denominators to a common one and use cancellation to
equalize each pair of numerators. Adding the equal numerators proves
independence. Negation and zero are \((s,-f)\) and \((s,0)\); the group
laws follow after common refinement. Composition is bilinear. For
example, for a fixed second denominator \(t\), construct squares for two
numerators \(f,g\), then refine their source arrows together. The
equalities \(fr=ta\), \(gr=tb\) imply \((f+g)r=t(a+b)\), so the
composite distributes over their sum. Distributivity in the other
variable follows directly after giving its two roofs a common
denominator. Thus \(Q\) is additive. Its original finite-sum inclusions
and projections retain their equations, including \(i_1p_1+i_2p_2=1\),
and so remain biproducts. The zero object\textquotesingle s identity
remains zero, making it a zero object. This proves additivity of the
localized category. \(\square\)

\subsection{5. Triangles after
localization}\label{complexes-cones-and-localization--5-triangles-after-localization}

\textbf{Theorem 5.1.} A triangle in \(D(\mathcal A)\) is distinguished
when it is isomorphic to the image of a cone triangle. With this
definition \(D(\mathcal A)\) is triangulated, \(Q\) is triangulated, and
each \(H^n\) descends to a functor that takes distinguished triangles to
long exact cohomology sequences.

\textbf{Proof.} Shift acts on roofs by shifting both arrows; its inverse
is the negative shift. Identity triangles descend. A roof
\((s,f):X\to Y\) is, through the isomorphism \(Q(s):M\to X\), the first
map of the image of the cone triangle of \(f\). Thus every localized map
has a completion. Rotation in both directions follows from Theorem 3.1.

We give the full square-lifting argument. Suppose \(bf=f'a\) in \(D\),
for chain-homotopy classes \(f:X\to Y\), \(f':X'\to Y'\), and localized
maps \(a:X\to X'\), \(b:Y\to Y'\). Write \(a\) as
\(X\xleftarrow{s}U\xrightarrow{a_0}X'\) and \(b\) as
\(Y\xleftarrow{t}V\xrightarrow{b_0}Y'\). Form a square \(tv=fsu\), with
\(u:W\to U\) a quasi-isomorphism. In \(D\), the maps \(b_0v\) and
\(f'a_0u\) are equal. The equality criterion of Theorem 4.2 gives a
quasi-isomorphism \(k:W'\to W\) equalizing them in \(K\). We now have a
map of first arrows \(W'\xrightarrow{vk}V\) to \(X\xrightarrow fY\),
whose vertical maps \(suk,t\) are quasi-isomorphisms, and a map to
\(X'\xrightarrow{f'}Y'\), whose vertical maps are \(a_0uk,b_0\). Lemma
2.1 extends each square to cone triangles. The first cone comparison is
a quasi-isomorphism by their cohomology sequences and the Five Lemma.
Invert it in \(D\) and compose with the second comparison. This gives
precisely the missing third map for the original localized square. It
proves the square-extension axiom.

We also need cone uniqueness to handle initially chosen octahedra. For
fixed \(T\), the roofs construction writes \(\operatorname{Hom}_D(T,-)\)
as the filtered colimit of \(\operatorname{Hom}_K(M,-)\) over
denominators \(M\to T\), ordered by refinements. This index category is
filtered: common squares give common refinements, and cancellation
equalizes parallel refinements. Filtered colimits of abelian groups are
exact by Lemma 1.1 of \emph{Sheaves of modules on a ringed space}. Apply
that result in the larger universe when required. Lemma 2.4 therefore
gives long exact representable sequences in \(D\). The dual fraction
description, furnished by the dual statements in Lemma 4.1 and the same
proof as Theorem 4.2, gives the contravariant sequences as well. The
Five Lemma and the inverse construction in Lemma 2.4 show that two
isomorphisms in a morphism of localized triangles imply the third.
Combined with square extension, this gives uniqueness of completion
triangles up to an isomorphism fixing their first two objects.

Finally take two composable localized maps represented by
\(X\xleftarrow{s}U\xrightarrow aY\) and
\(Y\xleftarrow{t}V\xrightarrow bZ\). Form \(th=ar\), with \(r:P\to U\) a
quasi-isomorphism. The string \(P\xrightarrow hV\xrightarrow bZ\) maps
to the original string through the isomorphisms \(Q(sr),Q(t),1_Z\).
Theorem 3.1 supplies its octahedron, whose image supplies an octahedron
for the original maps. Cone uniqueness just proved compares its three
chosen triangles with any three initially specified triangles, fixing
their first two objects. Transport gives the full localized octahedron
axiom.

Every \(H^n\) inverts quasi-isomorphisms and therefore descends by
Theorem 4.2. Its long exact sequence on an image triangle is Lemma 1.2
applied to the cone, with the signs in (2.2); an isomorphic triangle has
the transported sequence. This finishes the proof. \(\square\)

\textbf{Proposition 5.2 (short exact sequences and boundedness).} Every
short exact sequence of complexes gives a distinguished triangle in
\(D(\mathcal A)\). A complex whose cohomology vanishes below \(a\) is
quasi-isomorphic to a complex whose terms vanish below \(a\).

\textbf{Proof.} For \(0\to A\xrightarrow fB\xrightarrow gC\to0\), the
map \(r:C(f)\to C\), \((b,a')\mapsto g(b)\), is a chain map. Its kernel
is \(C(1_A)\), under the identification \(\ker g=A\), and is
contractible. Lemma 1.2 makes \(r\) a quasi-isomorphism. Thus the cone
triangle identifies with \(A\to B\to C\to A[1]\) in \(D\), with boundary
\((-p)r^{-1}\). When the original sequence splits degreewise, Lemma 2.2
identifies this boundary with its positive Snake Lemma boundary. In
general the same agreement follows from the lift formula in Lemma 1.2,
or from its naturality applied to \(0\to A\to B\to C\to0\) and the cone
construction.

For the last assertion use the good truncation \(\tau_{\ge a}K\): it is
zero below \(a\), equals \(K^a/\operatorname{im}d_K^{a-1}\) in degree
\(a\), and equals \(K^n\) above \(a\). The outgoing differential factors
through this quotient because \(d^2=0\). The natural map from \(K\) is
an isomorphism on cohomology at and above \(a\), by kernels and images,
and both cohomologies vanish below \(a\). It is the required
quasi-isomorphism. The dual good truncation uses \(\ker d_K^b\) at the
top degree \(b\). \(\square\)

We write \(D^+(\mathcal A)\) for the full subcategory whose cohomology
vanishes sufficiently far below zero, and similarly \(D^-\) and \(D^b\).
They are triangulated: shifts preserve the relevant bounds, and the long
exact sequence bounds the cohomology of a cone whenever its first two
objects have those bounds. Proposition 5.2 provides bounded
representatives where needed. It is cohomological boundedness, rather
than the terms of a particular representative, that defines these
subcategories.

\subsection{6. Exercises with complete
solutions}\label{complexes-cones-and-localization--6-exercises-with-complete-solutions}

\textbf{Exercise 1 (easy).} For a module \(M\), show that the complex
\(M\xrightarrow1M\) in degrees \(n,n+1\) is contractible. Explain why
this proves that adding it to any complex does not change its object in
\(K\).

\textbf{Solution.} Take the only nonzero homotopy component to be the
identity from degree \(n+1\) to degree \(n\). Then \(dH+Hd\) is the
identity in both nonzero degrees. Projection from the direct sum to the
original complex has its inclusion as inverse up to the homotopy that is
zero on the original summand and the indicated contraction on the new
summand. Thus it is an isomorphism in \(K\).

\textbf{Exercise 2 (medium).} The complex
\(\mathbb Z\xrightarrow2\mathbb Z\to\mathbb Z/2\) in degrees \(-1,0,1\),
with the quotient as second differential, is acyclic. Prove that it is
not contractible.

\textbf{Solution.} Multiplication by two is injective, its image is the
quotient map\textquotesingle s kernel, and that quotient map is onto.
Thus all cohomology groups vanish. A contraction at degree \(1\) would
give a homomorphism \(s:\mathbb Z/2\to\mathbb Z\) such that the quotient
composed with \(s\) is the identity. Every such homomorphism is zero,
since \(\mathbb Z\) has no element of order two. No contraction exists.
Localization makes this complex zero, because its map from the zero
complex is a quasi-isomorphism, whereas the homotopy category retains it
as a nonzero object.

\textbf{Exercise 3 (medium).} A map \(f:M[0]\to N[0]\) of modules has a
cone supported in degrees \(-1,0\). Compute its cohomology and the signs
in its cone triangle.

\textbf{Solution.} The cone is \(M\xrightarrow fN\), in those degrees.
Its cohomology is \(\ker f\) in degree \(-1\), \(\operatorname{coker}f\)
in degree \(0\), and zero otherwise. The triangle is
\(M[0]\xrightarrow fN[0]\to C(f)\xrightarrow{-p}M[1]\), so the
degree-\(-1\) component of the last map is \(-1_M\). This sign is
consistent with the split-sequence calculation of Lemma 2.2; replacing
it silently by a positive projection would change the adopted triangle
convention.

\textbf{Exercise 4 (hard).} Let \(I\) be a complex such that
\(\operatorname{Hom}^{\bullet}(A,I)\) is acyclic for every acyclic
\(A\). Without using roofs, prove that
\(\operatorname{Hom}_K(L,I)\to\operatorname{Hom}_D(L,I)\) is bijective.

\textbf{Solution.} If \(s:M\to N\) is a quasi-isomorphism, its cone is
acyclic. The contravariant Hom-cone identity in Lemma 2.4 implies that
\(s^*:\operatorname{Hom}^{\bullet}(N,I)\to\operatorname{Hom}^{\bullet}(M,I)\)
is a quasi-isomorphism. The induced degree-zero cohomology map is
therefore a bijection. By Lemma 1.1 it is the map on homotopy classes
into the target complex. Thus \(H(L)=\operatorname{Hom}_K(L,I)\)
descends to a contravariant functor on \(D\). For \(v:L\to I\) in \(D\),
send \(v\) to \(H(v)(1_I)\). On homotopy classes this is plainly inverse
to the canonical map in one direction. The other composite is a natural
endomorphism of \(\operatorname{Hom}_D(-,I)\) which fixes \(1_I\).
Naturality holds on ordinary arrows, on their inverted
quasi-isomorphisms by inverting the naturality squares, and hence on
every localized arrow. Applied to \(v\), it forces that composite to fix
\(v\). The two maps are inverse. This will be the comparison mechanism
for injective resolutions.

\section{Injective modules, flasque sheaves and bounded-below derived
functors}\label{injective-modules-and-bounded-below-derived-functors}

\emph{Written and edited by GPT-6.1 Sol (OpenAI), in Codex at Ultra,
October 2026. Mathematically self-checked by the writing AI. Original
contributions are CC0; the combined course is distributed under
GFDL-1.2-or-later. Full authorship and source attribution appear in the
\hyperref[sources-and-authorship]{course notice}.}

A global lifting problem can be replaced by an exact complex of sheaves
for which global lifting is easy. Injective modules make extensions of
maps possible; flasque sheaves make extensions of sections possible.
This lesson relates the two ideas and constructs the bounded-below
resolutions used to define right derived functors.

The prerequisite is
\hyperref[sheaves-of-modules-on-a-ringed-space]{Sheaves of modules on a
ringed space}, especially Theorem 2.1, Theorem 3.1 and Lemmas 5.1--5.2,
together with \hyperref[complexes-cones-and-localization]{Complexes,
cones and localization}, Lemmas 1.1--1.2, 2.4, Theorems 3.1, 4.2, 5.1
and Proposition 5.2. We work with left modules over an arbitrary sheaf
\(\mathcal O\) of associative unital rings. No separation condition is
imposed on the space. Commutativity is unnecessary throughout this
lesson.

The construction sources are the Stacks project authors'
\emph{Injectives}, \emph{Derived Categories} and \emph{Cohomology of
Sheaves}, in the
\href{https://github.com/KokunoYumeto/unofficial-stacks-project-ai-drafts/tree/565b10e987aba5969b21145a0833f42d69f96790}{AI
Integrated Stacks Project edition}. The flasque extension argument
corresponds to Tag 09SY. Spaltenstein's \emph{Resolutions of unbounded
complexes} provides historical background. Source attribution and the
licence for adapted passages appear in the
\hyperref[sources-and-authorship]{course notice}.

\subsection{1. Extending maps and extending
sections}\label{injective-modules-and-bounded-below-derived-functors--1-extending-maps-and-extending-sections}

An object \(I\) of an abelian category is \textbf{injective} if each map
from a subobject \(M\subset N\) into \(I\) extends to \(N\).
Equivalently, the contravariant functor \(\operatorname{Hom}(-,I)\) is
exact: surjectivity at \(\operatorname{Hom}(M,I)\) is precisely the
extension condition, while exactness at the earlier two terms follows
from kernels and cokernels. A sheaf \(F\) is \textbf{flasque}, or
\textbf{flabby}, if every restriction \(F(V)\to F(U)\), for opens
\(U\subset V\), is onto. It suffices that every \(F(X)\to F(U)\) be
onto, since this map factors through \(F(V)\).

\textbf{Lemma 1.1 (injectives are flasque and restrict to injectives).}
An injective module sheaf is flasque. Its restriction to any open subset
is injective in the module category of that open subset.

\textbf{Proof.} For opens \(U\subset V\), Lemma 5.1 of the preceding
lesson constructs a monomorphism
\(j_{U!}\mathcal O_U\to j_{V!}\mathcal O_V\). Under the isomorphisms
\(\operatorname{Hom}(j_{W!}\mathcal O_W,I)=I(W)\), precomposition is
restriction \(I(V)\to I(U)\). Injectivity therefore makes this
restriction onto. For the second claim, extension by zero is exact and
left adjoint to restriction. Consequently

\[\operatorname{Hom}_{\mathcal O_U}(-,I|_U)
\cong\operatorname{Hom}_{\mathcal O}(j_{U!}(-),I)\]

is a composite of exact functors. This is the required injectivity.
\(\square\)

\textbf{Lemma 1.2 (sections lift across a flasque kernel).} Suppose

\[0\longrightarrow F\longrightarrow G\longrightarrow H\longrightarrow0\]

is an exact sequence of module sheaves and \(F\) is flasque. For every
open \(U\), the sequence \(0\to F(U)\to G(U)\to H(U)\to0\) is exact. If
\(G\) is also flasque, then \(H\) is flasque.

\textbf{Proof.} Kernels are computed on sections, so only surjectivity
on sections requires proof. Fix \(s\in H(U)\). Consider pairs \((V,t)\)
where \(V\subset U\) is open and \(t\in G(V)\) maps to \(s|_V\). Order
the pairs by extension: \((V,t)\le(V',t')\) if \(V\subset V'\) and
\(t'|_V=t\). This is a set, because opens and sections are sets. It is
nonempty, containing the zero section on the empty open. A chain of
pairs has an upper bound. Take the union of its opens and glue its
compatible sections there. Any two chain members are comparable, which
gives compatibility on their intersection. The sheaf condition gives the
glued section and its uniqueness. Zorn\textquotesingle s lemma gives a
maximal pair \((V,t)\).

Choose \(x\in U\). The epimorphism of sheaves \(G\to H\) is onto on
stalks. Lift the germ of \(s\) at \(x\), represent the lift on an open
neighbourhood, and shrink until its image equals \(s\) there. We obtain
an open \(W\subset U\) containing \(x\) and a section \(t'\in G(W)\)
mapping to \(s|_W\). On \(W\cap V\), the difference \(t'-t\) has zero
image in \(H\). Since \(F\) is the kernel sheaf, it is the image of a
unique section \(r'\in F(W\cap V)\). Flasqueness extends \(r'\) to
\(r\in F(W)\). Replace \(t'\) by \(t'-r\), where \(r\) is mapped into
\(G(W)\). Its image in \(H\) is unchanged, and it now agrees with \(t\)
on \(W\cap V\). Glue the two sections to a lift of \(s\) on \(W\cup V\).
Maximality forces \(W\cup V=V\); thus \(x\in V\). Since \(x\) was
arbitrary, \(V=U\). This proves surjectivity.

If \(G\) is flasque, take opens \(U\subset V\) and \(s\in H(U)\). Lift
it to \(G(U)\) by the first part. Extend that lift to \(G(V)\), then map
it to \(H(V)\). The result restricts to \(s\). Thus \(H\) is flasque.
The proof uses neither commutativity of \(\mathcal O\) nor separation of
\(X\). \(\square\)

The first assertion needs only the kernel to be flasque. This is useful
even when the middle sheaf is not flasque. In particular, an exact
sequence beginning with a flasque sheaf has an exact section sequence on
every open; an arbitrary exact sheaf sequence need not, as the
exponential example of the preceding lesson shows.

\textbf{Proposition 1.3 (flasque resolutions compute zero higher
cohomology of a flasque sheaf).} Suppose module sheaves have enough
injectives. Resolve a flasque sheaf \(F\) by injectives,

\[0\longrightarrow F\longrightarrow I^0\longrightarrow I^1\longrightarrow\cdots.\]

For every open \(U\), the augmented complex of sections is exact. For
every continuous map with a compatible coefficient-ring morphism
\(f:X\to Y\), the augmented complex
\(0\to f_*F\to f_*I^0\to f_*I^1\to\cdots\) is exact. Consequently \(F\)
is acyclic for \(\Gamma(U,-)\) and \(f_*\), once these right derived
functors are constructed below.

\textbf{Proof.} Set \(F^0=F\) and
\(F^{n+1}=\operatorname{coker}(F^n\to I^n)\). The resolution gives short
exact sequences \(0\to F^n\to I^n\to F^{n+1}\to0\). By Lemma 1.1,
\(I^n\) is flasque. Starting from \(F^0\), Lemma 1.2 inductively makes
each \(F^{n+1}\) flasque and makes each of these short exact sequences
exact on \(U\). The image of \(I^{n-1}(U)\to I^n(U)\) is therefore
\(F^n(U)\), which is also the kernel of the next map. This proves
exactness of the section complex.

For an open \(V\subset Y\), apply the same argument to \(U=f^{-1}V\).
These section modules are exactly \((f_*I^n)(V)\). Thus the augmented
pushforward complex is exact even on every open\textquotesingle s
sections, and in particular is exact as a complex of sheaves. Applying
the later definition \(R^q\Phi(F)=H^q(\Phi(I^\bullet))\) gives zero for
\(q>0\) for either functor \(\Phi\). There is no need for a description
of higher direct images to prove this consequence, so it does not depend
on a later local-description theorem. \(\square\)

The existence assumption in Proposition 1.3 will be discharged by the
explicit functorial embedding construction in the next section. Its
statement separates an elementary lifting argument from the choice of a
resolution. It is not an appeal to an unproved acyclicity theorem.

\subsection{2. Functorial injective
embeddings}\label{injective-modules-and-bounded-below-derived-functors--2-functorial-injective-embeddings}

We now discharge the existence assumption of Proposition 1.3. First we
prove an extension test for modules over a ring. Then we construct a
functorial embedding using character modules. Choosing injective hulls
independently would not supply this functoriality.

\textbf{Lemma 2.1 (Baer\textquotesingle s criterion).} A left module
\(E\) over a ring \(A\) is injective if every \(A\)-linear map from a
left ideal of \(A\) to \(E\) extends to \(A\). \textbf{Proof.}
Injectivity immediately gives the stated extension property. Conversely,
suppose that property holds. Let \(N\subset M\) be left \(A\)-modules
and \(f:N\to E\) linear. Order its extensions \((N',f')\), with
\(N\subset N'\subset M\), by restriction. A chain has an upper bound:
take the union of its submodules and the compatible union of its maps.
Zorn\textquotesingle s lemma gives a maximal extension \((N',f')\). If
\(x\in M\setminus N'\), put \(J=\{a\in A:ax\in N'\}\). This is a left
ideal, and \(a\mapsto f'(ax)\) is a linear map \(J\to E\). By assumption
it extends to \(b:A\to E\). Define
\[f''(n+ax)=f'(n)+b(a)\quad(n\in N',\ a\in A).\] If \(n+ax=n'+a'x\),
then \((a-a')x=n'-n\), so \(a-a'\in J\) and \(b(a-a')=f'(n'-n)\). Thus
the formula is well defined. It is linear and extends \(f'\) to the
strictly larger submodule \(N'+Ax\), contradicting maximality. Hence
\(N'=M\), and every map from a submodule extends. This is injectivity.
Only the usual axiom of choice, in its Zorn form, is used; no existence
theorem for injectives is assumed. \(\square\)

\textbf{Lemma 2.2 (functorial injective embeddings of modules).} Let
\(A\) be a ring. Put
\(E_A=\operatorname{Hom}_{\mathbb Z}(A,\mathbb Q/\mathbb Z)\), a left
\(A\)-module by \((a\varphi)(b)=\varphi(ba)\). For a left \(A\)-module
\(M\) put
\[I_A(M)=\prod_{\varphi\in\operatorname{Hom}_A(M,E_A)}E_A,\qquad \iota_M(m)=(\varphi(m))_\varphi .
\tag{2.1}\] Then \(I_A(M)\) is injective and \(\iota_M\) is injective.
For \(u:M\to N\), the formula
\(I_A(u)\bigl((y_\varphi)_\varphi\bigr)=(y_{\chi\circ u})_{\chi}\), with
\(\chi\) running over \(\operatorname{Hom}_A(N,E_A)\), makes \(I_A\) a
functor and \(\iota\) a natural transformation.

\textbf{Proof.} \(E_A\) is a left module:
\(((aa')\varphi)(b)=\varphi(baa')=(a(a'\varphi))(b)\). The maps
\[\Phi(\varphi)=\bigl(m\mapsto\varphi(m)(1)\bigr),\qquad \Psi(\psi)=\bigl(m\mapsto(b\mapsto\psi(bm))\bigr)\]
are inverse bijections between \(\operatorname{Hom}_A(M,E_A)\) and
\(\operatorname{Hom}_{\mathbb Z}(M,\mathbb Q/\mathbb Z)\), natural in
\(M\). Here \(\Psi(\psi)\) is \(A\)-linear, since
\(\Psi(\psi)(am)(b)=\psi(bam)=(a\,\Psi(\psi)(m))(b)\). The group
\(\mathbb Q/\mathbb Z\) is an injective \(\mathbb Z\)-module. To see
this, use Baer\textquotesingle s criterion (Lemma 2.1): the ideals of
\(\mathbb Z\) are the \(n\mathbb Z\). For \(n=0\) there is nothing to
extend. For \(n\geq1\), a map \(n\mathbb Z\to\mathbb Q/\mathbb Z\)
sending \(n\) to \(q\) extends to \(\mathbb Z\) by sending \(1\) to any
\(q'\) with \(nq'=q\), which exists since \(\mathbb Q/\mathbb Z\) is
divisible. Exact sequences of \(A\)-modules are exact as groups. So
\(\operatorname{Hom}_A(-,E_A)\cong\operatorname{Hom}_{\mathbb Z}(-,\mathbb Q/\mathbb Z)\)
is exact, and \(E_A\) is injective. A product of injectives is
injective, because
\(\operatorname{Hom}(-,\prod E)=\prod\operatorname{Hom}(-,E)\) and
products of exact sequences of groups are exact: injectivity and kernels
are coordinatewise, and a preimage of each coordinate of a target
element can be chosen independently. This last assertion uses choice for
a set-indexed family, and concerns groups, not exactness of products of
sheaves. So \(I_A(M)\) is injective. If \(m\neq0\), define \(\psi\) on
the subgroup \(\mathbb Zm\) by \(\psi(m)=1/n+\mathbb Z\) if \(m\) has
finite order \(n\geq2\), and \(\psi(m)=1/2+\mathbb Z\) if \(m\) has
infinite order. Extend \(\psi\) to \(M\), as \(\mathbb Q/\mathbb Z\) is
injective. Then \(\varphi=\Psi(\psi)\) has
\(\varphi(m)(1)=\psi(m)\neq0\), so \(\iota_M(m)\neq0\). Finally
\(I_A(u)(\iota_M(m))=(\chi(u(m)))_\chi=\iota_N(u(m))\), and the
displayed formula gives \(I_A(vu)=I_A(v)I_A(u)\) and
\(I_A(\mathrm{id})=\mathrm{id}\). \(\square\)

\textbf{Theorem 2.3 (functorial injective embeddings).} Let
\(\mathcal O\) be any sheaf of rings on \(X\). Put
\[E(F)=\prod_{x\in X}x_*\bigl(I_{\mathcal O_x}(F_x)\bigr),\] where
\(I_{\mathcal O_x}\) is the functor (2.1) for the ring \(\mathcal O_x\),
and let \(\iota_F:F\to E(F)\) have \(x\)-component the morphism that
corresponds under (5.2) of the preceding lesson to
\(\iota_{F_x}:F_x\to I_{\mathcal O_x}(F_x)\). Then \(E(F)\) is
injective, \(\iota_F\) is a monomorphism, and \(E\) is a functor with
\(\iota\) natural. So \(\operatorname{Mod}(\mathcal O)\) has enough
injectives, with functorial injective embeddings.

\emph{Source:} the Stacks project,
\href{https://github.com/KokunoYumeto/unofficial-stacks-project-ai-drafts/blob/565b10e987aba5969b21145a0833f42d69f96790/injectives.tex\#L498}{Tag
01DI in its AI Integrated Stacks Project edition}, gives the
commutative-ring construction. The full character-module proof here also
checks the associative unital-ring generality; proof contributions by
GPT-6 Astra at Ultra are edited by GPT-6.1 Sol at Ultra.

\textbf{Proof.} Products exist and are computed on sections
(\hyperref[sheaves-of-modules-on-a-ringed-space--3-limits-colimits-stalks-and-sections-of-sums]{Theorem
3.1 of the preceding lesson}(1)), and a product of injectives is
injective, as in the proof of Lemma 2.2. Each factor is injective: for
an injective stalk module J, the adjunction of Lemma 5.2 of the
preceding lesson identifies Hom(-, x\_*J) with Hom over the stalk ring
from the exact stalk functor into J. This is an exact contravariant
functor. A product of injective sheaves is injective because its Hom
functor is the product of their exact group-valued Hom functors;
coordinatewise preimages prove surjectivity. Thus no assertion that
products of sheaves are exact is needed. Explicitly, \(\iota_F\) sends
\(s\in F(V)\) to \((\iota_{F_x}(s_x))_{x\in V}\). If this is zero, each
\(s_x=0\), as \(\iota_{F_x}\) is injective, so \(s=0\). Kernels are
computed on sections
(\hyperref[sheaves-of-modules-on-a-ringed-space--2-the-abelian-category-of-module-sheaves]{Theorem
2.1 of the preceding lesson}(1)), so \(\iota_F\) is a monomorphism. For
\(u:F\to G\) put \(E(u)=\prod_xx_*(I_{\mathcal O_x}(u_x))\). Naturality
of \(\iota\) holds factor by factor, by Lemma 2.2. \(\square\)

The construction gives an injective embedding, not an injective hull.
For example, over the field with two elements, the character module of
the ring has two elements. For the one-dimensional module, there are two
linear maps into that character module, including the zero map. The
indicated product is two-dimensional, and evaluation embeds the original
one-dimensional module into it. A different one-dimensional subspace has
zero intersection with the embedded image. Thus the embedding is not
essential, while a hull would require essentiality. The functorial
construction is useful precisely without a claim of minimality.

The embedding functor need not be additive in the morphism variable. Its
product coordinates are indexed by all maps into a fixed injective
cogenerator, and a module map changes coordinates by precomposition.
Identities, composition and naturality of the embedding hold by the
explicit formula. The later passage to derived functors uses comparison
maps unique up to chain homotopy, not additivity of this embedding
functor.

\subsection{3. Complexes that admit comparison
maps}\label{injective-modules-and-bounded-below-derived-functors--3-complexes-that-admit-comparison-maps}

A complex \(I\) is \textbf{K-injective} if every chain map \(A\to I\)
from an acyclic complex is null-homotopic. Equivalently,
\(\operatorname{Hom}^{\bullet}(A,I)\) is acyclic for every acyclic
\(A\): by Lemma 1.1 of the common reading, its degree-\(r\) cohomology
is \(\operatorname{Hom}_K(A,I[r])\), and shifting the acyclic source
gives the equivalence. This definition concerns a complex, not just its
individual terms.

\textbf{Lemma 3.1 (closure properties).} Shifts, cones of maps between
K-injectives, homotopy-equivalent complexes and degreewise products of
K-injective complexes are K-injective, whenever the indicated products
exist.

\textbf{Proof.} Shifting the source in (1.2) of the common reading
proves the shift assertion. For a map \(I\to J\), the covariant Hom-cone
identity of its Lemma 2.4 identifies
\(\operatorname{Hom}^{\bullet}(A,C(I\to J))\) with the cone of a map
between two acyclic group complexes. It is therefore acyclic. A homotopy
equivalence induces inverse maps up to homotopy after applying Hom, so
it preserves the condition. Finally,

\[\operatorname{Hom}^{\bullet}\left(A,\prod_\lambda I_\lambda\right)
\cong\prod_\lambda\operatorname{Hom}^{\bullet}(A,I_\lambda).\]

Products of exact sequences of groups are exact by coordinatewise
lifting and choice, as proved in Lemma 2.2. The product on the right is
consequently acyclic. This uses exactness of products of groups, and
does not presume that products of module sheaves are exact. \(\square\)

\textbf{Lemma 3.2 (bounded-below complexes of injectives are
K-injective).} Let \(I\) be a complex of injectives with \(I^n=0\) for
\(n<a\), and let \(A\) be acyclic. Then every chain map \(f:A\to I\) is
null-homotopic.

\textbf{Proof.} We build maps \(s^n:A^n\to I^{n-1}\) with
\[f^n=d\,s^n+s^{n+1}d\qquad\text{for all }n .
\tag{3.1}\] Put \(s^n=0\) for \(n\leq a\). Then (3.1) holds for \(n<a\),
where both sides are \(0\). Suppose \(s^k\) is built for \(k\leq n\) and
(3.1) holds for all indices below \(n\). Put
\(g=f^n-d\,s^n:A^n\to I^n\). Using (3.1) at \(n-1\), and \(d^2=0\),
\[g\,d=f^nd-d\,s^nd=d\,f^{n-1}-d\bigl(f^{n-1}-d\,s^{n-1}\bigr)=0 .\] So
\(g\) vanishes on the image of \(d:A^{n-1}\to A^n\). As \(A\) is
acyclic, that image is the kernel of \(d:A^n\to A^{n+1}\). Thus \(g\)
factors uniquely through the epimorphism from \(A^n\) onto
\(\operatorname{im}d\), giving a map from that subobject of \(A^{n+1}\)
to \(I^n\). Since \(I^n\) is injective, the factored map extends to
\(s^{n+1}:A^{n+1}\to I^n\). Then \(s^{n+1}d=g\), which is (3.1) at
\(n\). \(\square\)

The induction needs a lower bound. An acyclic K-injective complex is
contractible, since its own identity is one of the maps to which the
defining vanishing condition applies. A termwise injective complex
without a lower bound need not have this property; the singular-ring
example below gives a counterexample.

\textbf{Theorem 3.3 (maps into a K-injective complex).} For a
K-injective \(I\), a quasi-isomorphism \(s:M\to N\) induces a bijection
\[s^*:\operatorname{Hom}_{K(\mathcal A)}(N,I)
\longrightarrow\operatorname{Hom}_{K(\mathcal A)}(M,I).\] Consequently
the natural map from either homotopy-category Hom group to the
corresponding derived-category Hom group with target \(I\) is bijective.

\textbf{Proof.} The complex \(\operatorname{Hom}^\bullet(A,I)\) is
acyclic when \(A\) is acyclic: its degree-\(r\) cohomology is
\(\operatorname{Hom}_{K(\mathcal A)}(A,I[r])\), which vanishes because a
shift of \(A\) is again acyclic. Put
\(U=\operatorname{Hom}^\bullet(N,I)\),
\(V=\operatorname{Hom}^\bullet(M,I)\), and \(u=s^*:U\to V\). In degree
\(r\), write a map out of \(\operatorname{Cone}(s)\) as
\((a,b)\in U^r\oplus V^{r-1}\). Its differential is
\[(a,b)\longmapsto (d_Ua,\ d_Vb+(-1)^{r+1}u(a)).\] Thus
\((a,b)\mapsto((-1)^{r-1}b,a)\) is a chain isomorphism onto
\(\operatorname{Cone}(u)[-1]\): its differential is
\((v,a)\mapsto(-d_Vv-u(a),d_Ua)\). This explicitly identifies the Hom
complex out of \(\operatorname{Cone}(s)\) with the shifted cone of
\[\operatorname{Hom}^\bullet(N,I)\longrightarrow
\operatorname{Hom}^\bullet(M,I).\] The cone of \(s\) is acyclic. Thus
this map of Hom complexes is a quasi-isomorphism; its degree-zero
cohomology gives the displayed bijection.

Here is a direct use of the localization universal property that avoids
assuming a choice of roofs. The contravariant functor
\(H(L)=\operatorname{Hom}_K(L,I)\) sends every quasi-isomorphism to an
isomorphism by the first paragraph, so it factors through
\(D(\mathcal A)^{\mathrm{op}}\). For \(v:L\to I\) in \(D\), apply this
factored functor to \(v\) and to the element \(1_I\); this gives an
element of \(\operatorname{Hom}_K(L,I)\). The resulting map is inverse
to the natural map into \(\operatorname{Hom}_D(L,I)\): one composite
fixes every chain-homotopy class, and the other is a natural
endomorphism of the representable functor \(\operatorname{Hom}_D(-,I)\)
fixing \(1_I\). Naturality for inverted quasi-isomorphisms follows by
inverting their naturality squares. Every derived morphism is generated
by ordinary morphisms and these inverses, so the second composite fixes
all morphisms too. \(\square\)

\textbf{Corollary 3.4 (the full characterization).} For a complex \(I\),
the following are equivalent: it is K-injective;
\(\operatorname{Hom}_K(A,I)=0\) for every acyclic \(A\); every
quasi-isomorphism into the source induces a bijection on homotopy
classes of maps to \(I\); the natural map
\(\operatorname{Hom}_K(L,I)\to\operatorname{Hom}_D(L,I)\) is bijective
for every \(L\).

\textbf{Proof.} The first two conditions are the definition. Theorem 3.3
proves the third and fourth from them. Conversely, the third condition
applied to \(0\to A\) makes \(\operatorname{Hom}_K(A,I)=0\). The fourth
does too, since any acyclic \(A\) becomes a zero object in \(D\) by the
quasi-isomorphism \(0\to A\). Thus all conditions are equivalent. A
quasi-isomorphism between K-injectives is a homotopy equivalence: the
third condition supplies an inverse on one side, and its injectivity
forces the other composite to be the identity. \(\square\)

\subsection{4. Right derived functors on bounded-below
complexes}\label{injective-modules-and-bounded-below-derived-functors--4-right-derived-functors-on-bounded-below-complexes}

For an additive functor \(\Phi\), applying it degreewise gives a functor
on homotopy categories, but it need not preserve quasi-isomorphisms. We
seek \(R\Phi:D^+(\mathcal A)\to D^+(\mathcal B)\) with a comparison on
\(K^+(\mathcal A)\). Its \textbf{right-derived universal property} is
this: every natural comparison \(\Phi\to T\circ Q\), with \(T\) a
functor on \(D^+(\mathcal A)\), factors uniquely as
\(\Phi\to R\Phi\circ Q\to T\circ Q\). This makes \(R\Phi\) the initial
comparison out of \(\Phi\).

\textbf{Theorem 4.1.} Let \(\mathcal A\) and \(\mathcal B\) be abelian,
and assume that every object of \(\mathcal A\) embeds in an injective
object. Let \(\Phi:\mathcal A\to\mathcal B\) be additive.

\begin{enumerate}
\tightlist
\item
  Every bounded-below complex \(K\) has a qis \(K\to I\) into a
  bounded-below complex of injectives. If \(K^n=0\) for \(n<a\), it can
  be chosen injective in each degree, with \(I^n=0\) for \(n<a\).
\item
  Such an \(I\) is K-injective (Lemma 3.2). Hence
  \(\operatorname{Hom}_{K(\mathcal A)}(L,I)=\operatorname{Hom}_{D(\mathcal A)}(L,I)\)
  for every \(L\). A morphism into \(I\) extends along a qis, uniquely
  up to homotopy.
\item
  The rule \(R\Phi(K)=\Phi(I)\) gives a triangulated functor
  \(D^+(\mathcal A)\to D^+(\mathcal B)\), and this functor is the right
  derived functor \(R\Phi\). If \(\Phi\) is left exact, \(R^0\Phi=\Phi\)
  and \((R^i\Phi)_i\) is a universal \(\delta\)-functor.
\item
  (Leray) If \(\Phi\) is left exact and \(J\) is a bounded-below complex
  with \(R^i\Phi(J^n)=0\) for all \(i>0\) and all \(n\), then
  \(\Phi(J)\to R\Phi(J)\) is an isomorphism. Left exactness cannot be
  dropped. For a general additive \(\Phi\), the terms must also satisfy
  \(\Phi(J^n)\cong R^0\Phi(J^n)\), which holds by (3) when \(\Phi\) is
  left exact. For \(\Phi=-\otimes_{\mathbb Z}\mathbb Z/2\) on
  \(\operatorname{Mod}(\mathbb Z)\) and \(J=\mathbb Z\) in degree \(0\),
  the injective resolution \(\mathbb Q\to\mathbb Q/\mathbb Z\) gives
  \(R\Phi(\mathbb Z)=0\), because \(\mathbb Q\) and
  \(\mathbb Q/\mathbb Z\) are divisible, while
  \(\Phi(\mathbb Z)=\mathbb Z/2\).
\item
  For \(\mathcal A=\operatorname{Mod}(\mathcal O)\), with \(\mathcal O\)
  any sheaf of rings, (1)--(4) hold by Theorem 2.3. This gives \(Rf_*\),
  \(R\Gamma(U,-)\), and \(R\operatorname{Hom}_{\mathcal O}(G,-)\) on
  \(D^+\). The support functor is treated in \emph{Sections with support
  and the localization triangle}.
\end{enumerate}

We have proved the comparison assertion in Section 3. We now construct
resolutions and establish the derived-functor properties. The
localization, triangulated structure and cohomology exact sequences used
in these proofs are Theorems 3.1, 4.2, 5.1 and Lemma 1.2 of
\emph{Complexes, cones and localization}. The full proofs of
\href{https://github.com/KokunoYumeto/unofficial-stacks-project-ai-drafts/blob/565b10e987aba5969b21145a0833f42d69f96790/derived.tex\#L6630}{injective
resolutions},
\href{https://github.com/KokunoYumeto/unofficial-stacks-project-ai-drafts/blob/565b10e987aba5969b21145a0833f42d69f96790/derived.tex\#L9987}{K-injective
comparison} and
\href{https://github.com/KokunoYumeto/unofficial-stacks-project-ai-drafts/blob/565b10e987aba5969b21145a0833f42d69f96790/derived.tex\#L6075}{acyclic
terms} in AI Integrated Stacks Project are proof sources. The explicit
construction and expanded checks are by GPT-6 Astra at Ultra, checked
and edited by GPT-6.1 Sol at Ultra, under GFDL-1.2-or-later.

\subsubsection{Constructing the resolution one degree at a
time}\label{injective-modules-and-bounded-below-derived-functors--constructing-the-resolution-one-degree-at-a-time}

\textbf{Proof of Theorem 4.1(1).} The elementary operation is to replace
one term of a complex by an injective object. Suppose \(C\) is a complex
and choose a monomorphism \(u:C^n\to E\), with \(E\) injective. Form the
pushout \[P=(E\oplus C^{n+1})/\operatorname{im}(u,-d_C^n).\] Replace
\(C^n\) by \(E\) and \(C^{n+1}\) by \(P\). The incoming differential is
\(u d_C^{n-1}\), the new differential \(E\to P\) sends \(e\) to
\([e,0]\), and the outgoing differential sends \([e,c]\) to
\(d_C^{n+1}c\). These maps square to zero: \([u(c),0]=[0,d_Cc]\), and
\(d_C^2=0\). This gives a complex \(C'\) and a chain map \(C\to C'\).

The map \(C^{n+1}\to P\) is a monomorphism. Indeed, the pullback of the
monomorphism \((u,-d_C):C^n\to E\oplus C^{n+1}\) along the
second-summand inclusion is zero, because its projection to \(E\) has
kernel \(\ker u=0\). Thus that summand has zero intersection with the
subobject being quotiented. This is an argument by kernels and pullbacks
in an arbitrary abelian category. The map of complexes is termwise
monic. Quotienting in degree \(n+1\) also by the image of \(C^{n+1}\)
leaves \(E/u(C^n)\). The induced differential from degree \(n\) is the
identity of that quotient. Consequently the quotient complex is
\[E/u(C^n)\ \xrightarrow{\ 1\ }\ E/u(C^n)\] in degrees \(n,n+1\), and
zero in other degrees. It is contractible, so the cohomology exact
sequence proves that \(C\to C'\) is a quasi-isomorphism.

Start at \(n=a\), then perform the operation at \(a+1,a+2,\ldots\). Once
degree \(n\) has been replaced, later operations leave that object
unchanged; once degree \(n+1\) has been replaced, the differential out
of degree \(n\) is unchanged as well. Define \(I\) using these eventual
objects and differentials. This uses no infinite colimit in
\(\mathcal A\): each degree is defined at a finite stage. All its terms
are injective, it is zero below \(a\), and the map \(K\to I\) is
termwise monic. For fixed \(r\), the three terms and two differentials
computing \(H^r\) have stabilized after finitely many steps. Every step
was a quasi-isomorphism, so \(H^r(K)\to H^r(I)\) is an isomorphism. This
proves (1). \(\square\)

\subsubsection{Comparison, triangulation and
universality}\label{injective-modules-and-bounded-below-derived-functors--comparison-triangulation-and-universality}

\textbf{Proof of Theorem 4.1(2)--(3).} Choose a resolution
\(q_K:K\to I_K\). Use part (1). A chain map \(f:K\to L\) has a unique
homotopy-class extension \(\bar f:I_K\to I_L\) satisfying
\(\bar f q_K=q_L f\) in \(K(\mathcal A)\), by Theorem 3.3. Uniqueness
proves compatibility with composition and identities. If \(f\) is a
quasi-isomorphism, \(\bar f\) is one too. A quasi-isomorphism between
K-injectives is a homotopy equivalence: apply Theorem 3.3 first to
obtain a left inverse and again to show the other composite is the
identity.

An additive \(\Phi\) preserves chain homotopies, shifts and mapping
cones, since these involve only finite direct sums, differentials and
sums of maps. Therefore \(\Phi(\bar f)\) is invertible in
\(K(\mathcal B)\) whenever \(\bar f\) is a homotopy equivalence. The
assignment consequently factors through the localization of
\(K^+(\mathcal A)\). This localization identifies with the full
\(D^+(\mathcal A)\): Proposition 5.2 of the common reading gives its
bounded-below representatives, while Theorem 3.3 applied both in \(K^+\)
and in \(K\) computes all maps between such representatives by the same
maps into their bounded-below injective resolutions. Thus inclusion is
fully faithful as well as essentially surjective. All images lie in
\(D^+(\mathcal B)\). Different choices of \(I_K\) give a unique natural
comparison isomorphism: extend the identity between resolutions and
apply the preceding argument. These comparisons satisfy the cocycle
identity by uniqueness.

The functor is additive because the unique extension of \(f+g\) is the
sum of the unique extensions of \(f\) and \(g\). It is triangulated as
follows. A triangle in \(D^+\) is an image of a cone triangle up to
isomorphism, by Theorem 5.1 of the common reading. Replace its first two
objects by bounded-below injective resolutions. Theorem 3.3 extends its
first map, and Lemma 2.1 of the common reading compares its cone with
the cone of the extended map. The comparison is a quasi-isomorphism by
the cohomology exact sequences. Shifts and cones of K-injectives are
K-injective by Lemma 3.1, and those resolutions\textquotesingle{} cones
remain bounded below with injective terms, since finite sums of
injectives are injective. Thus the triangle can be computed by this
triangle of resolutions. Applying \(\Phi\) preserves its cone and shift,
proving the claim.

We verify the stated universal property. Applying \(\Phi\) to \(q_K\)
gives a natural comparison \(\eta_K:\Phi(K)\to R\Phi(K)\) in
\(D(\mathcal B)\). If \(T:D^+(\mathcal A)\to D^+(\mathcal B)\) is
another functor with a natural map \(\theta_K:\Phi(K)\to T(K)\) on
\(K^+(\mathcal A)\), the only possible induced map \(R\Phi(K)\to T(K)\)
is \[\Phi(I_K)\xrightarrow{\theta_{I_K}}T(I_K)
\xrightarrow{T(q_K)^{-1}}T(K).\] It is natural by the comparison
equation for \(\bar f\), it composes with \(\eta\) to \(\theta\), and it
is forced by that equation at \(q_K\). To see uniqueness explicitly, at
a K-injective resolution we may use the identity as its resolution, so
the comparison there is the identity. Naturality at \(q_K\) then forces
the displayed formula. The result is independent of that choice by the
unique resolution comparison. This is the right-derived-functor
property.

Define \(R^i\Phi(M)=H^i(R\Phi(M[0]))\). If \(\Phi\) is left exact, the
initial exact sequence \(0\to M\to I^0\to I^1\) gives
\(R^0\Phi(M)=\Phi(M)\). The resolution has no negative terms, so
\(R^i\Phi=0\) for \(i<0\). Short exact sequences give distinguished
triangles by Proposition 5.2 of the common reading. Their images under
\(R\Phi\) give natural long exact cohomology sequences. Thus the
nonnegative functors \(R^i\Phi\), with these connecting maps, form a
cohomological \(\delta\)-functor. Universality means that a natural
transformation in degree zero into any other such \(\delta\)-functor
extends uniquely to every degree, commuting with all connecting maps.

Here are the details of universality of that \(\delta\)-functor. Each
monomorphism \(M\to E\) into an injective kills \(R^i\Phi(M)\) after
mapping to \(R^i\Phi(E)\) for \(i>0\), because \(E[0]\) is already a
resolution. Let \(T^i\) be any other cohomological \(\delta\)-functor
and choose a natural map \(R^0\Phi\to T^0\). In
\(0\to M\to E\to C\to0\), exactness says that
\[R^0\Phi(C)\longrightarrow R^1\Phi(M)
\quad\text{is onto},\qquad
R^{i-1}\Phi(C)\xrightarrow{\ \sim\ }R^i\Phi(M)\quad(i>1).\] For \(i=1\),
compose the chosen degree-zero map at \(C\) with the boundary into
\(T^1(M)\). This vanishes on the image of \(R^0\Phi(E)\), by naturality
and exactness of \(T\), so it descends uniquely. Inductively, compose
the degree-\((i-1)\) map at \(C\) with the boundary into \(T^i(M)\) and
the inverse of the displayed isomorphism. Given a map \(M\to M'\),
extend its composite into an injective containing \(M'\) across
\(M\to E\); this gives a morphism of the two short exact sequences.
Induction and naturality of boundary maps prove naturality of the
constructed degree-\(i\) map. Two embeddings of the same \(M\) are
compared with their diagonal embedding into the finite sum of the two
injectives. The same argument proves independence of the embedding.
Compatibility with arbitrary connecting maps follows by embedding the
middle term of a short exact sequence into an injective and using the
induced diagram of cokernels; the defining construction is precisely the
connecting-map square in that diagram. Finally the displayed surjection
and isomorphisms force every degree of any extension. Thus the extension
exists uniquely as a morphism of \(\delta\)-functors. This proves
universality. \(\square\)

\subsubsection{Why acyclic terms suffice in the bounded-below
case}\label{injective-modules-and-bounded-below-derived-functors--why-acyclic-terms-suffice-in-the-bounded-below-case}

In the universality argument above, compatibility with an arbitrary
boundary can be checked without any further resolution theorem. For
\(0\to M\to N\to P\to0\), embed \(N\) in an injective \(E\) and put
\(C=E/M\). The identity on \(M\), the embedding \(N\to E\), and the
induced map \(P\to C\) give a morphism to \(0\to M\to E\to C\to0\).
Naturality of each \(\delta\)-functor gives
\[\delta_{M,N,P}^i=\delta_{M,E,C}^i\circ R^i\Phi(P\to C),\] and the same
formula for \(T\). The constructed maps commute with the second boundary
by their definition and with \(P\to C\) by their established naturality.
They therefore commute with the first boundary as well.

\textbf{Proof of Theorem 4.1(4)--(5).} Call an object \(A\) right
\(\Phi\)-acyclic when the comparison \(\Phi(A)[0]\to R\Phi(A[0])\) is an
isomorphism. For left exact \(\Phi\), part (3) shows that this is
equivalent to \(R^i\Phi(A)=0\) for \(i>0\). For a general additive
functor the degree-zero comparison must additionally be an isomorphism,
as the statement requires.

First let \(J\) have only finitely many nonzero terms, all right
\(\Phi\)-acyclic. Remove its bottom term using the degreewise split
sequence \[0\longrightarrow \sigma_{\ge a+1}J\longrightarrow J
\longrightarrow J^a[-a]\longrightarrow0,\] where \(\sigma\) denotes
brutal truncation, with the indicated boundary differential set to zero.
Applying \(\Phi\) still gives a degreewise split exact sequence. The
associated cone triangles and the natural comparison with \(R\Phi\)
show, by induction on the number of terms and the cohomology exact
sequences, that \(\Phi(J)\to R\Phi(J)\) is an isomorphism.

For a bounded-below \(J\), fix an integer \(r\) and use
\[0\longrightarrow \sigma_{\ge r+2}J\longrightarrow J
\longrightarrow \sigma_{\le r+1}J\longrightarrow0.\] The last complex is
bounded. The first, and also its image under \(\Phi\), vanish in degrees
below \(r+2\). By (1) it has an injective resolution vanishing below
\(r+2\), so its derived image also has no cohomology below \(r+2\). The
two long exact sequences therefore identify degree-\(r\) cohomology of
both middle terms with that of the bounded last terms. The already
proved bounded case gives the desired isomorphism in degree \(r\). This
holds for every \(r\), proving (4). Part (5) now follows by applying the
construction with the injective embeddings of Theorem 2.3. No
countability, commutativity, spectral sequence, or exactness of infinite
products of sheaves was used. \(\square\)

\subsection{5. Examples: ordinary, flasque and unbounded
resolutions}\label{injective-modules-and-bounded-below-derived-functors--5-examples-ordinary-flasque-and-unbounded-resolutions}

On a one-point space with coefficient ring \(R\), sheaves are simply
\(R\)-modules. Theorem 2.3 reduces to the character-module product of
Lemma 2.2. It is a functorial injective embedding. It is not generally a
hull, as the example over the field with two elements in Section 2
demonstrates. For \(R=\mathbb Z\), an ordinary injective resolution of
\(\mathbb Z\) is

\[0\longrightarrow\mathbb Z\longrightarrow\mathbb Q
\longrightarrow\mathbb Q/\mathbb Z\longrightarrow0.\]

Both nonzero resolution terms are divisible, and Baer\textquotesingle s
criterion makes them injective. Thus, for
\(\Phi=\operatorname{Hom}_{\mathbb Z}(\mathbb Z/n,-)\), the degree-zero
term after applying \(\Phi\) is zero and the degree-one term is the
subgroup of \(\mathbb Q/\mathbb Z\) killed by \(n\). It is cyclic of
order \(n\), by the map \(1\mapsto 1/n+\mathbb Z\). Hence
\(R^1\Phi(\mathbb Z)=\mathbb Z/n\) and all other nonnegative derived
values vanish. The right derived functor detects the extension
information absent from the original Hom group.

For a different additive functor,
\(\Phi=-\otimes_{\mathbb Z}\mathbb Z/2\), the same resolution has zero
image: a divisible group has zero quotient by twice itself. Thus its
right derived value at \(\mathbb Z[0]\) is zero, although
\(\Phi(\mathbb Z)=\mathbb Z/2\). This is precisely why higher vanishing
alone does not define right acyclicity for a general additive functor.
Later, left resolutions will give the derived tensor functor appropriate
to tensor\textquotesingle s right exactness.

The constant sheaf \(\underline{\mathbb Z}\) on \([0,1]\) is not
flasque. On the disconnected open \((0,1/3)\cup(2/3,1)\), the section
with value zero on the first component and one on the second cannot
extend to a locally constant function on the connected interval. There
is nevertheless a direct flasque resolution. Put

\[C^0(F)=\prod_{x\in X}x_*F_x.\]

Its sections on \(U\) are \(\prod_{x\in U}F_x\). Restriction forgets
coordinates, so it is onto: fill missing coordinates with zero. These
sections form a sheaf because compatible families glue coordinate by
coordinate. The germ map \(F\to C^0(F)\) is monic by the stalk
criterion. Let \(F^1\) be its cokernel, form \(C^0(F^1)\), and continue
with successive cokernels. The resulting augmented complex

\[0\to F\to C^0(F)\to C^0(F^1)\to C^0(F^2)\to\cdots\]

is exact by its construction, and all its resolution terms are flasque.
It is the Godement flasque resolution. Proposition 1.3 and Theorem
4.1(4) show that its section or pushforward complex computes the
corresponding derived functor. Its terms need not be injective;
acyclicity suffices because the resolution is bounded below.

An unbounded K-injective example goes in the other direction. On a
point, put

\[I=\prod_{m\ge0}(\mathbb Q/\mathbb Z)[m].\]

Each factor is bounded below and injective, hence K-injective by Lemma
3.2. Lemma 3.1 makes the product K-injective. In degree \(n\le0\),
exactly the factor with \(m=-n\) is nonzero, so
\(I^n=\mathbb Q/\mathbb Z\); in positive degrees it is zero. Its
differentials are zero. Thus it is not bounded below. In fact this
product represents the product in the derived category too: Theorem 3.3
and the Hom-product identity of Lemma 3.1 give

\[\operatorname{Hom}_D(K,I)
\cong\prod_{m\ge0}\operatorname{Hom}_D(K,(\mathbb Q/\mathbb Z)[m])\]

for every \(K\).

Termwise injectivity alone does not give such a conclusion. Let
\(R=k[\epsilon]/(\epsilon^2)\) for a field \(k\). Its only proper
nonzero ideal is \((\epsilon)\): an element with nonzero constant term
is a unit, while a nonzero scalar multiple of \(\epsilon\) generates
that ideal. An \(R\)-linear map \((\epsilon)\to R\) sends \(\epsilon\)
to \(c\epsilon\), because its image is killed by \(\epsilon\).
Multiplication by \(c\) extends it to \(R\). Baer\textquotesingle s
criterion therefore proves that \(R\) is injective over itself. The
bi-infinite complex

\[\cdots\xrightarrow\epsilon R\xrightarrow\epsilon R
\xrightarrow\epsilon R\xrightarrow\epsilon\cdots\]

is termwise injective and acyclic, since
\(\ker(\epsilon)=\epsilon R=\operatorname{im}(\epsilon)\). A contraction
would require \(1=\epsilon s^n+s^{n+1}\epsilon\) in every degree. Each
\(R\)-endomorphism is multiplication by an element of \(R\), so the
right side has image in \(\epsilon R\) and cannot be the identity. Its
identity is a non-null-homotopic map from an acyclic complex into
itself. It is not K-injective. The failure comes from the missing
starting degree in Lemma 3.2\textquotesingle s induction.

\subsection{6. Exercises with complete
solutions}\label{injective-modules-and-bounded-below-derived-functors--6-exercises-with-complete-solutions}

\textbf{Exercise 1 (easy).} At an arbitrary point \(x\), prove directly
that a skyscraper \(x_*J\) of an injective \(\mathcal O_x\)-module is
injective. Do not assume that \(x\) is a closed point.

\textbf{Solution.} For a monomorphism \(F\to G\), exactness of stalks
gives \(F_x\to G_x\) monic. A map \(F\to x_*J\) corresponds, under the
stalk-skyscraper adjunction in Lemma 5.2 of the preceding lesson, to
\(F_x\to J\). Extend it to \(G_x\to J\) by injectivity and transport
back through that adjunction. This extends the original sheaf map. No
assertion about the closure of \(x\) occurs in the proof; the
sheaf\textquotesingle s values and the adjunction already apply to every
point.

\textbf{Exercise 2 (medium).} Give a second proof that injective module
sheaves are flasque, using their embedding in Theorem 2.3 and a
splitting, instead of extension by zero.

\textbf{Solution.} An injective \(I\) embeds into \(E(I)\). Extend
\(1_I\) across this monomorphism to a retraction \(r:E(I)\to I\). Each
\(E(I)(U)\) is a product of stalk injectives indexed by \(x\in U\); its
restrictions are onto by filling missing coordinates with zero. For
\(U\subset V\), include a section \(s\in I(U)\) in \(E(I)(U)\), extend
it to \(E(I)(V)\), and apply \(r(V)\). Naturality of the retraction
makes its restriction equal \(s\). This proves flasqueness.

\textbf{Exercise 3 (medium).} Use dimension shifting to prove that a
flasque \(F\) has \(H^q(U,F)=0\) for \(q>0\), where
\(H^q(U,F)=R^q\Gamma(U,F)\).

\textbf{Solution.} Embed \(F\) into an injective \(I\) and put
\(C=I/F\). By Lemmas 1.1 and 1.2, \(C\) is flasque and \(I(U)\to C(U)\)
is onto. Since \(I\) is already its own injective resolution, all its
positive derived section functors vanish. The long exact sequence
therefore makes \(H^1(U,F)\) the cokernel of \(I(U)\to C(U)\), hence
zero, and gives \(H^q(U,F)\cong H^{q-1}(U,C)\) for \(q>1\). Induct on
\(q\), applying the same first-step proof to each flasque quotient. This
proves vanishing on every open, without an assumption of compactness or
separation.

\textbf{Exercise 4 (hard).} Prove all implications in Corollary 3.4,
including injectivity and surjectivity of the map to derived Hom, using
the roofs of the common reading.

\textbf{Solution.} Assume every map from an acyclic complex to \(I\) is
null-homotopic. Shifts have the same property, so both outer groups in
the contravariant cone exact sequence for a quasi-isomorphism
\(s:M\to N\) vanish. Thus
\(s^*:\operatorname{Hom}_K(N,I)\to\operatorname{Hom}_K(M,I)\) is
bijective. A derived map \(L\to I\) has a roof
\(L\xleftarrow{s}M\xrightarrow{f}I\). The unique homotopy class
\(g:L\to I\) with \(gs=f\) maps to this roof, proving surjectivity. If
two homotopy classes have the same derived image, the equality criterion
in Theorem 4.2 of the common reading equalizes them after some
quasi-isomorphism into \(L\). Injectivity of its induced \(s^*\)
equalizes the two classes before localization, proving injectivity.
Conversely, bijectivity to derived Hom implies vanishing from acyclic
sources because those sources are zero in \(D\). The source-comparison
condition implies the same vanishing by applying it to \(0\to A\) for
acyclic \(A\). Vanishing is exactly K-injectivity, completing the
equivalences.

\textbf{Exercise 5 (hard).} Suppose \(\Phi\) is left exact, \(J\) is
bounded below, and every term \(J^n\) is \(\Phi\)-acyclic. Explain
precisely why one cannot prove the comparison in degree \(r\) by
truncating only at degree \(r\). Give the correct truncation and prove
the comparison.

\textbf{Solution.} A quotient brutal truncation ending at degree \(r\)
deletes \(d^r\), which is needed to determine \(\ker d^r\) and hence
\(H^r\). Moreover the removed tail starts at degree \(r+1\), whose
derived cohomology in degree \(r+1\) can occur as the next term of the
long exact sequence. Instead use the split sequence with tail
\(\sigma_{\ge r+2}J\) and quotient \(\sigma_{\le r+1}J\). Both the
tail\textquotesingle s ordinary image and its derived image have zero
cohomology in degrees \(r,r+1\), because the bounded-below resolution
preserves its lower term bound. The long exact sequences identify the
degree-\(r\) middle cohomologies with those of the bounded quotient.
Induction on the quotient\textquotesingle s finitely many terms, using
the split bottom-term sequence, proves its comparison. Thus the
degree-\(r\) comparison for \(J\) is an isomorphism. The two-degree
buffer is the reason this argument works without a spectral sequence.

The next lesson proves the existence of resolutions for all complexes in
a Grothendieck abelian category. The unbounded existence proof supplies
a replacement for the lower-bound induction; the comparison theorem
already established here continues to apply unchanged.

\section{K-injective resolutions in Grothendieck abelian
categories}\label{k-injective-resolutions-in-grothendieck-categories}

\emph{Written and edited by GPT-6.1 Sol (OpenAI), in Codex at Ultra,
October 2026. Mathematically self-checked by the writing AI. Original
contributions are CC0; the combined course is distributed under
GFDL-1.2-or-later. Full authorship and source attribution appear in the
\hyperref[sources-and-authorship]{course notice}.}

A bounded-below resolution has a first degree. For an unbounded complex,
starting farther and farther to the left does not by itself produce a
compatible resolution. The solution developed here is to enforce the
desired extension and homotopy properties on a set of small tests, then
take a sufficiently long increasing union. The size estimates ensure
that every test into the union appeared at an earlier stage where the
construction dealt with it.

Read \hyperref[sheaves-of-modules-on-a-ringed-space]{Sheaves of modules
on a ringed space}, Theorems 2.1 and 3.1 and Lemmas 5.1--5.2;
\hyperref[injective-modules-and-bounded-below-derived-functors]{Injective
modules and bounded-below derived functors}, Section 3 and Theorem 4.1;
and \hyperref[complexes-cones-and-localization]{Complexes, cones and
localization}, Theorems 3.1, 4.2 and 5.1. Those readings supply abelian
categories, exact filtered colimits of module sheaves, Hom complexes,
cones, localization, and the comparison theorem for K-injectives. We use
sets and choice in a fixed universe. A larger universe temporarily
accommodates localization; this lesson proves that its Hom sets lie in
the original universe.

The construction follows the Stacks project authors' \emph{Injectives},
sections on Grothendieck categories and K-injective complexes, and
\emph{Derived Categories}, in the
\href{https://github.com/KokunoYumeto/unofficial-stacks-project-ai-drafts/tree/565b10e987aba5969b21145a0833f42d69f96790}{AI
Integrated Stacks Project edition}. The chapter credits Christian Serpé
and the relation to Nicolas Spaltenstein's \emph{Resolutions of
unbounded complexes}. The size bounds and transfinite construction are
developed below before their derived-functor applications. Source
attribution and the licence for adapted passages appear in the
\hyperref[sources-and-authorship]{course notice}.

\subsection{1. The Grothendieck
hypotheses}\label{k-injective-resolutions-in-grothendieck-categories--1-the-grothendieck-hypotheses}

An abelian category has \textbf{AB3} if it has all small coproducts,
\textbf{AB4} if those coproducts are exact, and \textbf{AB5} if it has
coproducts and exact filtered colimits. The starred conditions refer
respectively to products, exact products, and exact cofiltered limits. A
\textbf{generator} \(U\) means that \(\operatorname{Hom}(U,-)\) is
faithful: distinct morphisms become distinct after precomposition with
some map from \(U\). A locally small abelian category satisfying AB5 and
having a generator is a \textbf{Grothendieck abelian category}.
``Locally small'' means that each Hom collection is a set in the chosen
universe.

AB3 gives every small colimit. For a diagram \((A_i)\), take the
cokernel of the map from the sum over its arrows to the sum over its
objects, whose component for \(i\to j\) is the inclusion at \(i\) minus
the diagram map followed by the inclusion at \(j\). Maps out of this
cokernel are exactly compatible families of maps out of the diagram,
which proves its universal property. AB5 implies AB4: a coproduct of
short exact sequences is the filtered colimit of the exact finite
subsums. Exact filtered colimits also commute with cohomology, since
they preserve the kernels and cokernels defining it. These facts will be
used at successor and limit stages.

In an abelian category with coproducts the generator condition has two
useful equivalent forms. Every proper subobject \(N\subsetneq M\) is
escaped by some map \(U\to M\); and every object is a quotient of a
coproduct of copies of \(U\). Faithfulness implies the first form by
applying it to the nonzero quotient map \(M\to M/N\). That form implies
the second: sum all maps from \(U\) into \(M\); their image cannot be
proper. The second implies faithfulness: if a morphism is nonzero, an
epimorphism from a coproduct of copies of \(U\) cannot kill it on every
summand. We use all three forms without assuming that \(U\) is
projective.

\textbf{Proposition 1.1 (a generator).} Let \(\mathcal O\) be any sheaf
of rings on \(X\), and \(G=\bigoplus_{U}j_{U!}\mathcal O_U\), the sum
over all open sets \(U\). Then:

\begin{itemize}
\tightlist
\item
  (a) for every subsheaf \(N\subsetneq M\) there is a morphism
  \(G\to M\) that does not factor through \(N\);
\item
  (b) \(\operatorname{Hom}_{\mathcal O}(G,-)\) is faithful: a nonzero
  morphism \(M\to M'\) stays nonzero after composing with some
  \(G\to M\);
\item
  (c) every \(M\) receives an epimorphism from a direct sum of copies of
  \(G\).
\end{itemize}

Together with Theorem 3.1 of the first lesson, which gives all colimits
and exact filtered colimits, this makes
\(\operatorname{Mod}(\mathcal O)\) a Grothendieck abelian category as
defined above.

\emph{Reference:}
\href{https://github.com/KokunoYumeto/unofficial-stacks-project-ai-drafts/blob/565b10e987aba5969b21145a0833f42d69f96790/cohomology.tex\#L6992}{AI
Integrated Stacks Project, unbounded complexes} states this generator
for commutative \(\mathcal O\).

\textbf{Proof.} A section \(s\in M(U)\) gives
\(e_s:j_{U!}\mathcal O_U\to M\) with \(e_s(1)=s\) (Lemma 5.1 of the
first lesson). Composing with the projection \(G\to j_{U!}\mathcal O_U\)
gives \(\tilde e_s:G\to M\). (a) A subobject of \(M\) in
\(\operatorname{Mod}(\mathcal O)\) is a subsheaf, since kernels are
computed on sections. If \(N\neq M\), there are \(U\) and
\(s\in M(U)\setminus N(U)\). If \(\tilde e_s\) factored through \(N\),
then \(s=e_s(1)\) would lie in \(N(U)\). (b) If \(\varphi:M\to M'\) is
nonzero, then \(\varphi(s)\neq0\) for some section \(s\in M(U)\). So
\(\varphi\circ\tilde e_s\neq0\): on the summand \(j_{U!}\mathcal O_U\)
it sends \(1\) to \(\varphi(s)\). (c) The sum of all \(\tilde e_s\),
over all pairs \((U,s)\), is a morphism \(\bigoplus_{(U,s)}G\to M\)
whose image contains every section. So it is onto on stalks, hence an
epimorphism. \(\square\)

The three forms of generator are equivalent by the argument preceding
the proposition.

\subsection{2. Sizes and functorial injective
embeddings}\label{k-injective-resolutions-in-grothendieck-categories--2-sizes-and-functorial-injective-embeddings}

In the bounded-below case there was a first degree from which to
construct a homotopy. An unbounded complex has no such starting point.
Instead we will arrange that every map from an acyclic complex is
null-homotopic. There are too many acyclic complexes to attach them all
at once. The first part of the proof reduces the tests to a set of small
complexes; the second kills those tests and makes the final terms
injective.

For the construction let \(\mathcal A\) be any locally small
Grothendieck abelian category: it has small coproducts, exact filtered
colimits, and a generator \(U\). Here ``generator'' means that
\(\operatorname{Hom}(U,-)\) detects distinct morphisms. Proposition 1.1
and Theorem 3.1 of the first lesson give these properties for our module
sheaves. Sets, well-orderings and the axiom of choice are understood in
the fixed universe. The argument uses elementary abelian-category
kernels, images, quotients and pushouts, but not the construction of the
derived category.

We use two consequences of the filtered-colimit axiom. Coproducts are
exact: a coproduct is the filtered colimit of its finite subsums, and
finite direct sums are exact. Also filtered colimits commute with
cohomology of complexes: exactness preserves kernels and cokernels in
the formula for cohomology. Thus a filtered union of acyclic complexes
is acyclic, and a filtered composite of monic quasi-isomorphisms is a
monic quasi-isomorphism.

\subsubsection{Size bounds and maps into increasing
unions}\label{k-injective-resolutions-in-grothendieck-categories--size-bounds-and-maps-into-increasing-unions}

\textbf{Lemma 2.1 (a set of subobjects and bounded representatives).}
For every object \(B\), its subobjects form a set. Write \(s(B)\) for
the cardinality of this set. A subobject or quotient of \(B\) has size
at most \(s(B)\). For each infinite cardinal \(\mu\), there is a set of
representatives of the isomorphism classes of objects with
\(s(B)\leq\mu\).

\emph{Proof.} Associate to a subobject \(V\subset B\) the subset of
\(\operatorname{Hom}(U,B)\) consisting of maps that factor through
\(V\). This subset determines \(V\): the images of all maps \(U\to V\)
fill \(V\). Indeed, if their sum had a proper image \(W\subsetneq V\),
the nonzero quotient map \(V\to V/W\) would, by the generator property,
have nonzero composite with some map \(U\to V\), a contradiction. Hence
subobjects inject into a power set of a Hom set.

Subobjects of a subobject remain subobjects of \(B\). Subobjects of a
quotient correspond, by inverse image, to subobjects of \(B\) containing
its kernel. This proves the size inequalities.

For each proper subobject \(V\subsetneq B\), choose a map \(a_V:U\to B\)
not factoring through \(V\). Such a map exists by applying the generator
property to \(B\to B/V\). The sum of these maps is onto: a proper image
would be one of the subobjects the chosen maps were required to escape.
There are at most \(s(B)\) summands. Therefore, if \(s(B)\leq\mu\),
\(B\) is a quotient of \(U^{(\mu)}\), after adding zero summands if
necessary. Its isomorphism class is represented by the quotient by one
of the set of subobjects of \(U^{(\mu)}\). This proves the last
assertion. \(\square\)

\textbf{Lemma 2.2 (smallness for monic systems).} Suppose that
\((B_\beta)_{\beta<\lambda}\) is an increasing ordinal-indexed system of
subobjects with union \(B\), and \(\operatorname{cf}(\lambda)>s(A)\).
Every map \(A\to B\) factors through some \(B_\beta\).

\emph{Proof.} The inverse images \(A_\beta=A\times_B B_\beta\) form an
increasing family of subobjects of \(A\), with union \(A\). To see the
last equality without using elements, take filtered colimits of the
kernel diagrams for
\[A\oplus B_\beta\longrightarrow B,\qquad (a,b)\longmapsto f(a)-b.\]
Exactness identifies the resulting kernel with \(A\times_B B=A\). At
most \(s(A)\) distinct subobjects occur among the \(A_\beta\). Choose an
index for each distinct one. The cofinality assumption bounds all those
indices by a single \(\gamma<\lambda\), so \(A_\gamma\) contains every
\(A_\beta\), and is \(A\). The pullback property gives the required
factorization. \(\square\)

Recall that the cofinality of a limit ordinal is the least cardinality
of an unbounded subset of it. We will always choose a limit ordinal with
cofinality larger than a specified infinite cardinal \(\mu\); one may
take the initial ordinal of \(\mu^+\). Indeed, a family of at most
\(\mu\) ordinals below \(\mu^+\) has union of cardinality at most
\(\mu\), hence is bounded below \(\mu^+\). This also explains why
countably many factorization stages below such an ordinal have a common
upper bound.

\subsubsection{Enough injectives, with a functorial
choice}\label{k-injective-resolutions-in-grothendieck-categories--enough-injectives-with-a-functorial-choice}

\textbf{Lemma 2.3 (the generator extension test).} An object \(E\) is
injective if and only if every map \(V\to E\), where \(V\subset U\),
extends to \(U\).

\emph{Proof.} Necessity is the definition of injectivity. Conversely,
take a monomorphism \(A\subset B\) and a map \(A\to E\). Order its
extensions to subobjects of \(B\) by restriction. This is a set by Lemma
2.1 and local smallness. The union of a chain, with the compatible map
to \(E\), is an upper bound, so Zorn\textquotesingle s lemma gives a
maximal extension \(a:A'\to E\).

If \(A'\neq B\), the nonzero quotient \(B\to B/A'\) has nonzero
composite with some \(\psi:U\to B\). Put \(V=\psi^{-1}(A')\). The map
\(a\psi:V\to E\) extends to \(b:U\to E\) by assumption. It vanishes on
\(\ker\psi\), since its restriction there is \(a\psi=0\), so it descends
to \(\operatorname{im}\psi\). This descended map and \(a\) agree on
\(A'\cap\operatorname{im}\psi\). They therefore give a map
\(A'+\operatorname{im}\psi\to E\), strictly extending \(a\), a
contradiction. Thus \(A'=B\). \(\square\)

\textbf{Theorem 2.4 (functorial injective embeddings).} Every object
\(M\) has a natural monomorphism into an injective object \(E(M)\).

\textbf{Proof.} For an object \(M\), form the pushout
\[\begin{array}{ccc}
\displaystyle\bigoplus_{(V,a)}V&\longrightarrow&M\\
\big\downarrow&&\big\downarrow\\
\displaystyle\bigoplus_{(V,a)}U&\longrightarrow&\mathcal E(M),
\end{array}
\qquad
(V,a):\ V\subset U,\quad a:V\to M .
\tag{2.1}\] This is a set-indexed coproduct by Lemma 2.1. The left arrow
is monic because coproducts are exact; its pushout \(M\to\mathcal E(M)\)
is monic. Every selected map \(V\to M\) extends to
\(U\to\mathcal E(M)\). A map \(M\to M'\) takes the summand labelled
\((V,a)\) to the summand labelled by its composite, so \(\mathcal E\)
and the embedding are functorial.

Iterate \(\mathcal E\), taking unions at limit stages. Choose a fixed
\(\lambda_0\) with \(\operatorname{cf}(\lambda_0)>s(U)\), and call the
union \(E(M)\). Given \(V\subset U\) and a map \(V\to E(M)\), Lemma 2.2
makes it factor through a stage before \(\lambda_0\), since
\(s(V)\leq s(U)\). At the next stage (2.1) extends it to \(U\). Lemma
2.3 now proves that \(E(M)\) is injective. We have constructed
functorial monomorphisms \(M\to E(M)\). The functor \(E\) need not be
additive; no subsequent argument assumes that it is. \(\square\)

\subsection{3. Testing on a set of acyclic
complexes}\label{k-injective-resolutions-in-grothendieck-categories--3-testing-on-a-set-of-acyclic-complexes}

\textbf{Lemma 3.1 (bounded small lifts).} There is an infinite cardinal
\(\kappa\), depending only on \(\mathcal A\) and \(U\), such that every
nonzero acyclic complex contains a nonzero bounded-above acyclic
subcomplex all of whose terms have size at most \(\kappa\).

\emph{Proof.} For an infinite cardinal \(\mu\), set
\[c(\mu)=\max\{\mu,\aleph_0,s(U^{(\mu)})\}.\] This function is
nondecreasing: an inclusion of index sets gives a split monomorphism
between the corresponding coproducts.

First consider an epimorphism \(p:B\to C\), with \(s(C)\leq\mu\). For
each proper subobject \(C'\subset C\), the map \(B\to C/C'\) is nonzero.
Choose \(b_{C'}:U\to B\) whose composite is nonzero. The image \(B'\) of
their sum maps onto \(C\), because its image escapes every proper
\(C'\). It is a quotient of at most \(\mu\) copies of \(U\), and hence
\(s(B')\leq c(\mu)\). Notice that this argument does not assert that
maps from \(U\) lift through epimorphisms: \(U\) was not assumed
projective.

Put \(\kappa_0=\max\{s(U),\aleph_0\}\), \(\kappa_{r+1}=c(\kappa_r)\),
and \(\kappa=\sup_{r<\omega}\kappa_r\). Let \(A^\bullet\) be acyclic and
let \(\psi:U\to A^n\). Set \[N^n=\operatorname{im}\psi,\qquad
N^{n+1}=d(N^n),\qquad N^j=0\quad(j>n+1).\] The resulting complex is
already exact at and above degree \(n+1\). Both nonzero terms have size
at most \(\kappa_0\). If \(N^j\) has been constructed for \(j\geq k\),
use the epimorphism \[(d^{k-1})^{-1}\!\left(\ker(N^k\to N^{k+1})\right)
\longrightarrow \ker(N^k\to N^{k+1}).\] It is onto because \(A^\bullet\)
is acyclic. The small-lift construction gives a subobject \(N^{k-1}\)
mapping onto that kernel. If \(s(N^k)\leq\kappa_r\), it can be chosen
with \(s(N^{k-1})\leq\kappa_{r+1}\). Descending through all integers
constructs a bounded-above acyclic subcomplex containing
\(\operatorname{im}\psi\), with every term of size at most \(\kappa\).
Only finite iterates of \(c\) are used at each degree; we do not assume
\(c(\kappa)=\kappa\).

If \(A^\bullet\neq0\), choose a degree with \(A^n\neq0\) and a nonzero
\(\psi:U\to A^n\). The constructed subcomplex is nonzero. \(\square\)

Choose a set \(\mathscr S\) of representatives of all bounded-above
acyclic complexes with term sizes at most \(\kappa\). Such a set exists:
Lemma 2.1 supplies a set of choices for each term, and local smallness
supplies a set of differentials between each pair of choices; there are
only countably many degrees. Restrict that set of complexes to the ones
satisfying \(d^2=0\), acyclicity and boundedness above.

\textbf{Lemma 3.2 (testing K-injectivity on a set).} Suppose that every
term of \(I^\bullet\) is injective and every chain map
\(S^\bullet\to I^\bullet\), for \(S^\bullet\in\mathscr S\), is
null-homotopic. Then \(I^\bullet\) is K-injective.

\emph{Proof.} Take an acyclic \(A^\bullet\). Starting with \(F_0=0\),
build an increasing continuous filtration by acyclic subcomplexes. If
\(F_\beta\neq A\), the quotient \(A/F_\beta\) is nonzero and acyclic.
Choose in it a nonzero small acyclic subcomplex from Lemma 3.1, and let
\(F_{\beta+1}\) be its inverse image. At a limit stage take the union.
Successor exactness follows from the short exact sequence of complexes;
limit exactness follows from exact filtered colimits. The construction
exhausts \(A\): the subcomplexes of \(A\) form a set, since they are
among the sequences of subobjects of its terms. A strictly increasing
chain longer than that set is impossible. Every nonzero successor
quotient \(F_{\beta+1}/F_\beta\) is isomorphic to an element of
\(\mathscr S\).

Let \(f:A\to I\) be a chain map. We construct null-homotopies
\(h_\beta\) of \(f|_{F_\beta}\) that agree on earlier stages. At a
successor, extend each component \(h_\beta^n:F_\beta^n\to I^{n-1}\) to a
map \(\widetilde h^n:F_{\beta+1}^n\to I^{n-1}\), using injectivity. The
map \[g=f|_{F_{\beta+1}}-d_I\widetilde h-\widetilde h\,d_F\] is a chain
map, as \(d_I^2=d_F^2=0\). It vanishes on \(F_\beta\), and so descends
to a chain map \(\bar g:F_{\beta+1}/F_\beta\to I\). By the test
assumption \(\bar g=d_I k+k d\) for a degree-minus-one map \(k\). With
\(\pi:F_{\beta+1}\to F_{\beta+1}/F_\beta\), put
\(h_{\beta+1}=\widetilde h+k\pi\). This extends \(h_\beta\) and
satisfies \(f=d_Ih_{\beta+1}+h_{\beta+1}d_F\).

At a limit stage the compatible components define maps out of the union,
and the same identity holds because it holds on every earlier stage. At
exhaustion they give a null-homotopy of \(f\). Thus every map from every
acyclic complex to \(I\) is null-homotopic, exactly the definition of
K-injectivity. This proof does not pass cohomology through an inverse
limit. \(\square\)

\subsection{4. Constructing the unbounded
resolution}\label{k-injective-resolutions-in-grothendieck-categories--4-constructing-the-unbounded-resolution}

\textbf{Theorem 4.1.} In a locally small Grothendieck abelian category,
every complex \(M\) has a functorial, termwise monic quasi-isomorphism
\(M\to I(M)\), where each term of \(I(M)\) is injective and \(I(M)\) is
K-injective.

\textbf{Proof.}

There are two operations, with different purposes.

For \(S\in\mathscr S\), let \(C(S)=\operatorname{Cone}(1_S)\), with
\[C(S)^n=S^n\oplus S^{n+1},\qquad
d(a,b)=(da+b,-db).\] The map \(S\to C(S)\), \(a\mapsto(a,0)\), is monic
and null-homotopic: \(a\mapsto(0,a)\) is a homotopy for it. The same
formula contracts \(C(S)\). Its quotient by \(S\) is \(S[1]\), which is
acyclic. Define \(P(M)\) by pushing out all these inclusions along all
chain maps \(S\to M\): \[\begin{array}{ccc}
\displaystyle\bigoplus_{(S,w)}S&\xrightarrow{\ \sum w\ }&M\\
\big\downarrow&&\big\downarrow\\
\displaystyle\bigoplus_{(S,w)}C(S)&\longrightarrow&P(M).
\end{array}
\tag{4.1}\] Then \(M\to P(M)\) is termwise monic. Its quotient is the
coproduct of the acyclic complexes \(S[1]\), hence is acyclic. It is
therefore a quasi-isomorphism. Every chosen \(w\) becomes null-homotopic
in \(P(M)\), using the corresponding cone. Indexing by all maps makes
\(P\) and its embedding functorial.

The second operation uses the embeddings \(i_A:A\to E(A)\) constructed
above. Set \[J(M)^n=E(M^n)\oplus E(M^{n+1}),\qquad
d_J(a,b)=(b,0),\qquad
u^n=(i_{M^n},\,i_{M^{n+1}}d_M).
\tag{4.2}\] The equation \(d_Ju=ud_M\) follows from \(d_M^2=0\), and
\(u\) is monic because its first component is monic. Let
\(Q=J(M)/u(M)\), with quotient map \(v:J(M)\to Q\), and define
\[N(M)=\operatorname{Cone}(v)[-1],\qquad
N(M)^n=Q^{n-1}\oplus J(M)^n,\] \[d_N(q,j)=(-d_Qq-v(j),d_Jj),\qquad
M\longrightarrow N(M),\quad m\longmapsto(0,u(m)).
\tag{4.3}\] These formulas are identities of component morphisms in the
abelian category; they do not require objects to have underlying sets of
elements. They show directly that the last map is a monic chain map. Its
degree-\(n\) map factors through the injective subobject
\(0\oplus J(M)^n\) of \(N(M)^n\). A finite direct sum of injectives is
injective, since a map into it extends component by component.

The quotient \(N(M)/M\) has terms \(Q^{n-1}\oplus Q^n\), differential
\[D(q,a)=(-dq-a,da),\] and contraction \(h(q,a)=(0,-q)\): substitution
gives \(Dh+hD=1\). Thus \(M\to N(M)\) is a quasi-isomorphism. Naturality
of \(i_A\) gives functoriality of every map in (4.2)--(4.3), even though
\(E\) need not be additive.

The roles of the two operations are worth separating: \[\begin{gathered}
T_\beta\xrightarrow{\text{kill test maps}}P(T_\beta),\\
P(T_\beta)\xrightarrow{\text{insert injectives}}T_{\beta+1}.
\end{gathered}
\tag{4.4}\] Both arrows are monic quasi-isomorphisms. Neither
intermediate complex is asserted to be K-injective.

\subsubsection{The final union and its two
checks}\label{k-injective-resolutions-in-grothendieck-categories--the-final-union-and-its-two-checks}

Choose once and for all a limit ordinal \(\lambda\) of cofinality
greater than \(\kappa\). Define \[T_0(M)=M,\qquad
T_{\beta+1}(M)=N(P(T_\beta(M))),\qquad
T_\gamma(M)=\varinjlim_{\beta<\gamma}T_\beta(M)\] at limit stages, and
put \(I(M)=T_\lambda(M)\). All transitions are monic quasi-isomorphisms,
including at limits by exactness of filtered colimits. Hence
\(M\to I(M)\) is monic and a quasi-isomorphism. The choices of
\(U,\mathscr S,\lambda_0,\lambda\) are fixed independently of \(M\), so
the construction is functorial in \(M\).

First, every term \(I(M)^n\) is injective. Given \(V\subset U\) and a
map \(V\to I(M)^n\), smallness factors it through \(T_\beta(M)^n\) for
some \(\beta<\lambda\). Its composite into \(T_{\beta+1}(M)^n\) factors
through the injective object \(J(P(T_\beta(M)))^n\) by (4.3). Extend
\(V\to J(P(T_\beta(M)))^n\) to \(U\), and compose into \(I(M)^n\). Lemma
2.3 proves injectivity. This is a proof about the final terms, not an
assertion that filtered colimits of injectives are automatically
injective.

Second, let \(S\in\mathscr S\) and \(w:S\to I(M)\) be a chain map. Every
component factors through a stage by Lemma 2.2, since
\(s(S^n)\leq\kappa\). There are countably many components. As
\(\operatorname{cf}(\lambda)>\kappa\geq\aleph_0\), all components factor
through a single \(T_\beta(M)\). They already commute with the
differentials there: their commutators vanish after the monomorphism
\(T_\beta(M)^{n+1}\to I(M)^{n+1}\), and hence vanish before it. Thus
\(w\) factors as a chain map. Operation \(P\) makes that map
null-homotopic at the next stage, so it is null-homotopic in \(I(M)\).
Lemma 3.2 proves K-injectivity.

Apply this construction to
\(\mathcal A=\operatorname{Mod}(\mathcal O)\). It proves the resolution
assertion of Theorem 4.1, including termwise injectivity, naturality and
noncommutative coefficients. The unbounded derived-functor consequences
are proved in Section 5 below. The derived category was constructed in
the common reading; the existence proof here only uses complexes and
homotopies. \(\square\)

\subsection{5. Derived functors, adjoints and
products}\label{k-injective-resolutions-in-grothendieck-categories--5-derived-functors-adjoints-and-products}

Let \(Q:K(\mathcal A)\to D(\mathcal A)\) denote localization. By Theorem
3.3 of the preceding lesson, maps into a K-injective complex satisfy

\[\operatorname{Hom}_{K}(K,I)=\operatorname{Hom}_{D}(QK,QI).
\tag{5.1}\]

That theorem also proves that a quasi-isomorphism between K-injectives
is a homotopy equivalence. We now use it without boundedness
restrictions. Although the construction of \(I(M)\) is functorial on
chain maps, its functor need not be additive. The additive functor on
the derived category is obtained from (5.1), rather than by claiming
additivity of the transfinite construction.

\textbf{Theorem 5.1 (right derived functors on all complexes).} Suppose
every complex of an abelian category \(\mathcal A\) has a
quasi-isomorphism \(q_K:K\to I_K\) into a K-injective complex. For every
exact functor of triangulated categories
\(F:K(\mathcal A)\to\mathcal T\), there is an exact right derived
functor

\[RF:D(\mathcal A)\longrightarrow\mathcal T,
\qquad RF(K)=F(I_K).\]

Its comparison \(\eta:F\to RF\circ Q\) is initial among natural
comparisons from \(F\) to a functor on \(D(\mathcal A)\). In particular,
for any additive \(\Phi:\mathcal A\to\mathcal B\), with \(\mathcal B\)
abelian, \(R\Phi:D(\mathcal A)\to D(\mathcal B)\) is computed by
applying \(\Phi\) degreewise to a K-injective resolution.

\textbf{Proof.} The full subcategory
\(K_{\mathrm{inj}}\subset K(\mathcal A)\) of K-injectives is closed
under shifts and cones by Lemma 3.1 of the preceding lesson. Equation
(5.1) shows that \(Q\) restricted to this subcategory is fully faithful,
and the resolution hypothesis makes it essentially surjective. Define an
inverse \(R\) explicitly: send \(K\) to \(I_K\), and a derived morphism
\(v:QK\to QL\) to the unique homotopy class \(Rv:I_K\to I_L\) satisfying

\[Q(Rv)=Q(q_L)\,v\,Q(q_K)^{-1}.
\tag{5.2}\]

Existence and uniqueness follow from (5.1). The same formula proves
preservation of identities, composition and addition. For an ordinary
homotopy class \(f:K\to L\), injectivity of the map in (5.1) gives
\((Rf)q_K=q_Lf\) already in \(K(\mathcal A)\).

The inverse \(R\) is exact. Its shift comparison is the unique
isomorphism \(I_{K[1]}\cong I_K[1]\) inducing the corresponding
isomorphism in \(D\). A morphism between K-injectives has a K-injective
cone. For any distinguished triangle in \(D\), compare its first arrow
with the unique arrow between the chosen K-injectives. The cone triangle
of that arrow maps under \(Q\) to a distinguished triangle with the same
first arrow. The triangle morphism axiom and the long exact
representable sequences, proved in the common reading, supply an
isomorphism on the third objects compatible with all arrows. Full
faithfulness lifts this isomorphism and its equations to
\(K_{\mathrm{inj}}\). Thus \(R\) transports distinguished triangles to
distinguished triangles, with the stated shift comparison. Set
\(RF=F\circ R\) and \(\eta_K=F(q_K)\); the equality just proved gives
naturality.

To verify universality, take any functor
\(T:D(\mathcal A)\to\mathcal T\) and natural \(\theta:F\to TQ\). Define

\[\alpha_K=T(Qq_K)^{-1}\theta_{I_K}:
F(I_K)\longrightarrow T(QK).
\tag{5.3}\]

Naturality of \(\theta\) for the representative \(Rv\), together with
(5.2), proves naturality of \(\alpha\) for every derived morphism \(v\).
Naturality for \(q_K\) gives \(\alpha_K\eta_K=\theta_K\). This
factorization is unique: at a K-injective object \(I\), its resolution
map \(q_I\) is a homotopy equivalence, so \(F(q_I)\) is invertible and
forces \(\alpha_I\). Naturality for the isomorphism \(Qq_K\) then forces
\(\alpha_K\). Consequently the displayed construction has exactly the
right-derived universal property. It also proves independence of the
chosen resolutions and comparisons.

Finally, an additive \(\Phi\) preserves chain homotopies, shifts, and
finite direct sums, and carries the defining matrix of a cone to that of
the cone of the image map. Composing its functor on homotopy categories
with \(K(\mathcal B)\to D(\mathcal B)\) is therefore exact. Apply the
result to this \(F\). On bounded-below cohomology it agrees with the
preceding lesson: its bounded-below injective resolution is K-injective,
and (5.1) uniquely compares it with any unbounded resolution.
\(\square\)

For a Grothendieck category, (5.1) also settles the size of derived
morphisms. \(\operatorname{Hom}_{K}(K,I_L)\) is the quotient of the set
of chain maps by the set-valued homotopy relation; chain maps and
homotopies are subsets of countable products of original-universe Hom
sets. Thus \(D(\mathcal A)\) is locally small in the original universe,
including when \(\mathcal A\) is a large category of sheaves.

\textbf{Proposition 5.2 (exact left adjoints).} Let
\(u:\mathcal A\to\mathcal B\) be exact and left adjoint to an additive
\(v:\mathcal B\to\mathcal A\). Then \(v\) preserves K-injectives.

\textbf{Proof.} The adjunction on objects identifies chain maps and
their homotopies componentwise, hence

\[\operatorname{Hom}_{K(\mathcal A)}(A,vI)
=\operatorname{Hom}_{K(\mathcal B)}(uA,I).\]

If \(A\) is acyclic, exactness of \(u\) makes \(uA\) acyclic and
K-injectivity makes the right side zero. This is the definition of
K-injectivity of \(vI\). \(\square\)

In particular, restriction \(j^{-1}\) to an open subset preserves
K-injectives because its left adjoint \(j_!\) is exact. For a continuous
\(f:X\to Y\) with \(\mathcal O_X=f^{-1}\mathcal O_Y\), the exact left
adjoint \(f^{-1}\) shows that \(f_*\) preserves K-injectives. For an
arbitrary morphism of ringed spaces, \(Rf_*\) still exists by Theorems
4.1 and 5.1 applied to module sheaves; preservation by the underived
\(f_*\) requires an exact left adjoint, for example flat module
pullback. Section 6 gives a failure when pullback is not exact.

\textbf{Proposition 5.3 (existence of products in \(\mathcal A\)).} A
locally small abelian category with a generator and small coproducts has
all small products. In particular a Grothendieck category satisfies
\(\mathrm{AB3}^{*}\).

\textbf{Proof.} Fix a family \((A_i)_{i\in J}\). For any cone
\((f_i:T\to A_i)\), let \(N\subset T\) be the sum of all subobjects on
which every \(f_i\) vanishes. This is a set of subobjects by Lemma 2.1,
so its sum exists as the image of their coproduct. Each \(f_i\) vanishes
on the sum. The induced family on \(B=T/N\) is \textbf{jointly monic},
meaning that a map killed by all its components is zero. Indeed, if a
subobject of \(B\) is killed by every component, its inverse image in
\(T\) would enlarge \(N\); hence it is zero. Apply this to the image of
a map killed by the family. This constructs a jointly monic quotient of
every cone, without using products or intersections to define it.

For a jointly monic cone on \(B\), the map of sets

\[\operatorname{Hom}(U,B)\longrightarrow
\prod_{i\in J}\operatorname{Hom}(U,A_i)\]

is injective: the difference of two maps with equal images is killed by
every component. Let \(\nu\) be the cardinality of the product on the
right and set \(\mu=\max(2^\nu,\aleph_0)\). The subobject-to-subset
argument in Lemma 2.1 gives \(s(B)\le\mu\). That lemma supplies a set of
representatives of all possible \(B\), and local smallness supplies a
set of all their cones to \((A_i)\). Take the coproduct \(C\) of all
these jointly monic cones, with the induced maps \(C\to A_i\), and
divide by their common vanishing subobject as in the first paragraph.
Call the result \(P\).

Every cone on \(T\) factors through its jointly monic quotient \(B\),
through a representative included as a summand of \(C\), and then
through \(P\). Thus a map \(T\to P\) inducing the given cone exists. It
is unique because the cone on \(P\) is jointly monic. This is the
universal property of \(\prod_i A_i\). The empty product is zero and the
construction still works: its only jointly monic cone has zero source.
\(\square\)

The proposition proves existence, not exactness, of products. Products
of sheaves already existed in the first lesson, but their epimorphisms
can fail to be onto on stalks.

\textbf{Proposition 5.4 (products and coproducts in the derived
category).} A Grothendieck category has all small derived products and
coproducts. A coproduct is represented by the termwise sum of any
representatives. A product is represented by the termwise product of
K-injective resolutions.

\textbf{Proof.} Put \(P=\prod_i I_i\) for K-injectives \(I_i\).
Componentwise projections identify the complex of abelian groups
\(\operatorname{Hom}^\bullet(A,P)\) with
\(\prod_i\operatorname{Hom}^\bullet(A,I_i)\). Products of abelian groups
are exact: kernels are coordinatewise and coordinatewise lifts prove
surjectivity using choice. Therefore if \(A\) is acyclic, the product
Hom complex is acyclic. Thus \(P\) is K-injective. For any \(K\),
exactness of those group products and (5.1) give

\[\operatorname{Hom}_{D}(K,P)
=\prod_i\operatorname{Hom}_{D}(K,I_i),\]

which proves the product property. The proof assumes exact products in
abelian groups, not in \(\mathcal A\).

For the sum \(S=\bigoplus_i K_i\), and any K-injective \(I\), the Hom
complex identifies with \(\prod_i\operatorname{Hom}^\bullet(K_i,I)\).
The same exactness gives
\(\operatorname{Hom}_K(S,I)=\prod_i\operatorname{Hom}_K(K_i,I)\).
Equation (5.1) and a K-injective resolution of every target prove the
coproduct property in \(D\). AB4 ensures that changing representatives
by quasi-isomorphisms changes their sum by a quasi-isomorphism.
\(\square\)

\subsection{6. Examples that isolate the
hypotheses}\label{k-injective-resolutions-in-grothendieck-categories--6-examples-that-isolate-the-hypotheses}

\textbf{Unbounded termwise injectivity.} Over
\(R=k[\epsilon]/(\epsilon^2)\), the complex with \(R\) in every integer
degree and differential multiplication by \(\epsilon\) is acyclic,
because \(\ker(\epsilon)=\epsilon R=\operatorname{im}(\epsilon)\). Its
terms are injective: the only ideals are \(0,\epsilon R,R\), and a map
\(\epsilon R\to R\) sends \(\epsilon\) to an element \(\epsilon a\), so
it extends by sending \(1\) to \(a\); the Baer criterion was proved in
the preceding lesson. Its identity is not null-homotopic, because each
expression \(\epsilon h^n+h^{n+1}\epsilon\) has image in \(\epsilon R\)
and cannot equal \(1_R\). Taking the acyclic complex itself as a source
proves that it is not K-injective. Theorem 4.1 must change this complex,
although its terms are already injective.

\textbf{Pushforward can lose K-injectivity.} Take one-point ringed
spaces with ring map \(\mathbb Z\to\mathbb F_2\). The module pushforward
is restriction of scalars. The complex \(\mathbb F_2[0]\) is K-injective
over \(\mathbb F_2\): a linear map from a subspace extends after
choosing a basis, so every vector space is injective, and the
bounded-below lemma applies. Consider the acyclic complex of abelian
groups

\[A^\bullet=(\mathbb Z\xrightarrow{2}\mathbb Z
\xrightarrow{\mathrm{mod}\,2}\mathbb F_2),
\qquad\text{in degrees }0,1,2.\]

The map \(A\to\mathbb F_2[0]\) that is reduction modulo two in degree
zero and zero elsewhere is a chain map. A null-homotopy would imply
\(\mathrm{mod}\,2=h^1\circ2\), impossible for a homomorphism
\(h^1:\mathbb Z\to\mathbb F_2\). Thus the pushed-forward complex is not
K-injective. Its left adjoint \(\mathbb F_2\otimes_{\mathbb Z}-\) is not
exact, since it sends the injection \(2:\mathbb Z\to\mathbb Z\) to zero.
This matches the hypothesis of Proposition 5.2 exactly.

\textbf{An unbounded derived product.} The injective group
\(\mathbb Q/\mathbb Z\) placed in degree \(-m\) is
\((\mathbb Q/\mathbb Z)[m]\). The product of these complexes for
\(m\ge0\) has one copy of \(\mathbb Q/\mathbb Z\) in every nonpositive
degree, zero elsewhere, and zero differential. Proposition 5.4 makes it
K-injective and the product in \(D(\mathbb Z)\). It has nonzero
cohomology in arbitrarily negative degrees, so it is not bounded below.
This also shows why resolving only bounded-below targets would miss
objects of the unbounded category.

\subsection{7. Exercises with checked
solutions}\label{k-injective-resolutions-in-grothendieck-categories--7-exercises-with-checked-solutions}

\textbf{Exercise 1 (easy: a concrete generator).} For a discrete finite
space \(X\) with constant ring \(R\), identify the generator \(G\) of
Proposition 1.1 and prove directly that it detects morphisms. Does the
proposition require \(R\) to be commutative?

\textbf{Solution.} A sheaf of modules is a tuple \((M_x)_{x\in X}\), and
\(j_{U!}\mathcal O_U\) has stalk \(R\) at points of \(U\) and zero
elsewhere. Thus \(G_x=\bigoplus_{U\ni x}R\). If a tuple of maps has a
nonzero component at \(x\), choose an element there with nonzero image.
The summand belonging to the singleton \(\{x\}\) gives a map into the
source tuple detecting that element. For a left module, \(R\to M_x\),
\(a\mapsto am\), is left linear for every associative unital \(R\). No
commutativity is used. The empty-space module category has only its zero
object, and its zero generator is faithful there.

\textbf{Exercise 2 (medium: \(\mathrm{AB3}^{*}\) without
\(\mathrm{AB4}^{*}\)).} Give the size bound that makes the product
construction in Proposition 5.3 a set-sized construction. Explain why it
proves no exactness statement about products of sheaves.

\textbf{Solution.} With \(H_i=\operatorname{Hom}(U,A_i)\), every jointly
monic cone on \(B\) embeds \(\operatorname{Hom}(U,B)\) into the set
\(\prod_i H_i\). A subobject of \(B\) is determined by the subset of
maps from \(U\) that factor through it. Hence
\(s(B)\le2^{|\prod_i H_i|}\), and Lemma 2.1 provides a set of
representatives. For each representative, the cones form a subset of
\(\prod_i\operatorname{Hom}(B,A_i)\), also a set. Their union is a set,
so their coproduct is permitted by AB3. This argument concerns a
universal property. It supplies no common open neighbourhood on which
infinitely many stalkwise lifts exist. The product-of-epimorphisms
counterexample in Section 6 of the first lesson therefore remains
compatible with AB3*.

\textbf{Exercise 3 (medium: the unbounded direct image).} For an
arbitrary morphism of ringed spaces, construct \(Rf_*\), explain its
agreement with bounded-below \(Rf_*\), and say when \(f_*I\) can itself
be used as a K-injective complex on the target.

\textbf{Solution.} The module categories are Grothendieck by Proposition
1.1 and the first lesson. Choose a K-injective resolution \(K\to I_K\)
by Theorem 4.1. Define \(Rf_*K\) as \(f_*I_K\) in the target derived
category, with morphisms obtained by applying \(f_*\) to the unique
comparison homotopy classes. Additivity of \(f_*\) makes it preserve
homotopies and cones, so Theorem 5.1 proves existence, exactness and
uniqueness. A bounded-below injective resolution is K-injective and is
uniquely homotopy equivalent to any chosen K-injective resolution of the
same object, giving agreement. If the left adjoint \(f^*\) is exact,
Proposition 5.2 proves that \(f_*I_K\) is K-injective on the target. The
one-point example in Section 6 proves that this final claim cannot be
made for arbitrary ring maps.

\textbf{Exercise 4 (hard: what a transfinite step does for groups).}
Work in \(\operatorname{Mod}(\mathbb Z)\), with generator
\(U=\mathbb Z\). For \(M=\mathbb Z[0]\), describe an explicit nonzero
small acyclic complex that occurs among the tests, and write the
contraction that kills any map from it in \(P(M)\). Then check the
contraction of \(N(P(M))/P(M)\). Explain why a countable sequence of
steps has not been shown sufficient.

\textbf{Solution.} The complex
\(S=(\mathbb Z\xrightarrow{2}\mathbb Z\to\mathbb Z/2)\) in degrees
\(0,1,2\) is bounded above and acyclic, with terms having at most
countably many subobjects. After replacing it by its representative it
belongs to \(\mathscr S\). Every selected chain map \(w:S\to M\) has an
attached copy of \(C(S)\). If \(e:C(S)\to P(M)\) is its pushout map,
then the degree-minus-one map \(a\mapsto e(0,a)\) satisfies
\(dh+hd=jw\), because \(d(0,a)=(a,-da)\). For this particular \(M\),
\(w^0:\mathbb Z\to\mathbb Z\) can be multiplication by any integer and
the other components vanish; the same formula kills all these maps.

Put \(L=P(M)\), \(Q=J(L)/L\). In \(N(L)/L\), write a term as
\((q,a)\in Q^{n-1}\oplus Q^n\). Its differential is
\(D(q,a)=(-dq-a,da)\) and \(h(q,a)=(0,-q)\). Direct substitution gives
\(Dh(q,a)=(q,-dq)\) and \(hD(q,a)=(0,dq+a)\), whose sum is \((q,a)\).
Thus this quotient is contractible. Finally, a test has countably many
components. Their first factorization stages could be unbounded in
\(\omega\); no single finite stage would then contain the chain map. The
chosen cofinality greater than \(\kappa\ge\aleph_0\) ensures a common
stage and makes the next-step killing argument valid.

\textbf{Exercise 5 (hard: compatible homotopies at a limit).} In Lemma
3.2, explain why choosing an unrelated null-homotopy at each filtration
stage is insufficient. Derive the successor correction that makes the
chosen homotopies agree.

\textbf{Solution.} A map out of a union is induced by compatible maps on
its subobjects. Unrelated null-homotopies need not restrict to one
another, so they define no homotopy on the union. Suppose \(h_\beta\)
has been fixed. Termwise injectivity extends it to an arbitrary graded
map \(\widetilde h\) on \(F_{\beta+1}\). The error
\(g=f-d\widetilde h-\widetilde h d\) is a chain map and restricts to
zero on \(F_\beta\). It factors through
\(\pi:F_{\beta+1}\to F_{\beta+1}/F_\beta\). This quotient is a small
acyclic test, so its error map is \(dk+kd\). Define
\(h_{\beta+1}=\widetilde h+k\pi\). Then \(dh_{\beta+1}+h_{\beta+1}d=f\),
and its restriction is the prescribed \(h_\beta\), because \(\pi\)
vanishes there. At a limit, the universal property of the colimit
assembles these compatible components and their equations. This uses
exact filtered colimits to keep the filtration acyclic, but no exact
inverse-limit assertion.

The next lesson constructs K-flat resolutions by a different method.
This theorem is not needed for their existence; the two resolution tools
address different variables of derived operations.

Sources and licensing: the Stacks project authors, \emph{Injectives},
\href{https://github.com/KokunoYumeto/unofficial-stacks-project-ai-drafts/blob/565b10e987aba5969b21145a0833f42d69f96790/injectives.tex}{the
complete pinned chapter}, especially the proofs corresponding to Tags
0E8N and 079C--079P; \emph{Derived Categories},
\href{https://github.com/KokunoYumeto/unofficial-stacks-project-ai-drafts/blob/565b10e987aba5969b21145a0833f42d69f96790/derived.tex\#L9967}{the
pinned K-injective section}, for comparison, derived functors, products
and exact adjoints. The elementary large-cofinality argument in Section
2 is also the proof of
\href{https://github.com/KokunoYumeto/unofficial-stacks-project-ai-drafts/blob/565b10e987aba5969b21145a0833f42d69f96790/sets.tex\#L217}{Stacks,
Tag 05N3, in the pinned AI Integrated edition}. These adapted passages
retain GFDL-1.2-or-later. N. Spaltenstein, \emph{Resolutions of
unbounded complexes}, Compositio Mathematica 65 (1988), 121--154, is
\href{https://www.numdam.org/item/CM_1988__65_2_121_0/}{available from
Numdam}. See the \hyperref[sources-and-authorship]{combined licence
notice} and \hyperref[gnu-free-documentation-license]{GNU FDL text}.

\section{Flat modules and K-flat
resolutions}\label{flat-modules-and-k-flat-resolutions}

\emph{Written and edited by GPT-6.1 Sol (OpenAI), in Codex at Ultra,
October 2026. Mathematically self-checked by the writing AI. Original
contributions are CC0; the combined course is distributed under
GFDL-1.2-or-later. Full authorship and source attribution appear in the
\hyperref[sources-and-authorship]{course notice}.}

Tensoring a resolution should preserve the cohomology information it was
meant to carry. Flat modules solve this problem for a single tensor
factor in one degree. Unbounded complexes require a stronger condition,
K-flatness. The two conditions are independent for unbounded complexes:
a contractible complex can have nonflat terms, and an acyclic complex of
free terms can fail to be K-flat.

The prerequisites are
\hyperref[sheaves-of-modules-on-a-ringed-space]{Sheaves of modules on a
ringed space}, Theorems 2.1, 3.1 and 4.4 and Lemmas 5.1--5.2, and
\hyperref[complexes-cones-and-localization]{Complexes, cones and
localization}, Lemmas 1.1--1.2, 2.4 and Proposition 5.2. We construct
K-flat resolutions independently of the K-injective existence theorem.
This avoids a dependency cycle when derived tensor products and Hom are
developed.

Throughout this lesson \(\mathcal O\) is a sheaf of \textbf{commutative}
unital rings on an arbitrary topological space \(X\). Tensor
totalizations use direct sums, including for unbounded complexes. We
impose no separation, compactness, Noetherian or finite
homological-dimension assumption.

The construction sources are the Stacks project authors' \emph{Modules
on Ringed Spaces}, ``Flat modules'', and \emph{Cohomology of Sheaves},
``Flat resolutions'', in the
\href{https://github.com/KokunoYumeto/unofficial-stacks-project-ai-drafts/tree/565b10e987aba5969b21145a0833f42d69f96790}{AI
Integrated Stacks Project edition}. We develop flat-module calculations,
cycle attachments and compatible bounded resolutions. Spaltenstein's
\emph{Resolutions of unbounded complexes} provides historical
background. Source attribution and the licence for adapted passages
appear in the \hyperref[sources-and-authorship]{course notice}.

\subsection{1. Flat sheaves and exact tensor
sequences}\label{flat-modules-and-k-flat-resolutions--1-flat-sheaves-and-exact-tensor-sequences}

A sheaf \(M\) is \textbf{flat} if \(-\otimes_{\mathcal O}M\) is exact.
The tensor sheaf is the sheafification of the presheaf
\(U\mapsto F(U)\otimes_{\mathcal O(U)}M(U)\). Its balanced universal
property follows by gluing the presheaf balanced maps, using structured
sheafification from the first lesson. Tensor is right exact: a balanced
map on a quotient is precisely a balanced map on the original module
vanishing on the relations defining that quotient. Thus flatness amounts
to preservation of injections.

\subsubsection{Flat building
blocks}\label{flat-modules-and-k-flat-resolutions--flat-building-blocks}

\textbf{Lemma 1.1 (local covers).} Every module sheaf is a quotient of a
direct sum of sheaves \(B_U=j_{U!}\mathcal O_U\).

\textbf{Proof.} For an open \(U\subset X\), Lemma 5.1 of the first
lesson gives \[\operatorname{Hom}_{\mathcal O}(B_U,M)=M(U).
\tag{1.1}\] For every module sheaf \(M\), form the set-indexed sum
\[L(M)=\bigoplus_{U\subset X\ {\rm open}}\ \bigoplus_{s\in M(U)}B_U .
\tag{1.2}\] The map \(\lambda_M:L(M)\to M\) sends the generator indexed
by \(s\) to \(s\), using (1.1). Every germ is represented by such a
local section, so \(\lambda_M\) is an epimorphism. \(\square\)

This uses local sections on all opens. A sheaf need not be generated by
its global sections.

\textbf{Lemma 1.2 (flatness and stalks).} A module sheaf \(M\) is flat
if and only if every \(M_x\) is flat over \(\mathcal O_x\). In
particular \(B_U\), \(L(M)\), and direct sums of flat module sheaves are
flat.

\textbf{Proof.} The tensor sheaf is the sheafification of the
sectionwise tensor presheaf. Its stalk is
\[(F\otimes_{\mathcal O}M)_x=F_x\otimes_{\mathcal O_x}M_x .
\tag{1.3}\] Here is the relevant finite-relation argument. A tensor is a
finite sum of elementary tensors. Represent all its factors on one
neighbourhood of \(x\). An equality between tensors is generated by
finitely many additive and balancing relations. Represent their scalar
and module entries on a common smaller neighbourhood; each equality of
germs in this finite certificate holds after a further shrinking. This
proves both surjectivity and injectivity of (1.3), even though the
section rings vary with the neighbourhood.

If every \(M_x\) is flat, (1.3) and Theorem 2.1 of the first lesson show
that tensoring with \(M\) preserves exact sequences. Conversely, an
injection \(A\to B\) of \(\mathcal O_x\)-modules gives an injection
\(x_*A\to x_*B\) of module sheaves: these maps are injective on
sections. If \(M\) is flat, tensoring preserves this injection, and its
stalk at \(x\) is \(A\otimes M_x\to B\otimes M_x\). Tensor products are
right exact, by their presentation with generators and relations, so
this proves that \(M_x\) is flat.

The stalks of \(B_U\) are \(\mathcal O_x\) or zero (Lemma 5.1 of the
first lesson); both are flat. Stalks commute with direct sums. A direct
sum of flat modules is flat because its tensor functor is the direct sum
of their exact tensor functors, and direct sums of exact module
sequences remain exact coordinatewise. This proves the final assertions.
\(\square\)

\textbf{Lemma 1.3 (flat quotients and flat exact sequences).} If
\(0\to M'\to M\to M''\to0\) is exact and \(M''\) is flat, then tensoring
with any sheaf \(N\) preserves that exact sequence. Consequently flat
sheaves are closed under extensions and kernels of epimorphisms between
flat sheaves. Filtered colimits of flat sheaves are flat, and a direct
summand of a flat sheaf is flat.

\textbf{Proof.} We first prove the flat-quotient assertion over a ring
\(R\). Resolve a test module \(N\) by a complex \(F^\bullet\) of free
modules in degrees at most zero. Choose a free module on the elements of
\(N\), map it onto \(N\), choose a free module onto that kernel, and
repeat. This construction gives \(H^0(F)=N\) and zero negative
cohomology. Flatness of \(M''\) makes \(F\otimes_R M''\) exact in
negative degrees, by tensoring the short exact sequences of cycles. Free
terms make

\[0\to F\otimes M'\to F\otimes M\to F\otimes M''\to0\]

termwise exact, because a free tensor is a direct sum of the original
exact sequence. The long exact sequence of cohomology gives an injection
\(H^0(F\otimes M')\to H^0(F\otimes M)\), since the preceding
\(H^{-1}(F\otimes M'')\) is zero. Right exactness identifies these
degree-zero groups with \(N\otimes M'\) and \(N\otimes M\). This proves
the assertion without presupposing a Tor theorem. Apply it at each
stalk, using Lemma 1.2, to obtain the sheaf assertion.

For an injection \(N\to N_1\), compare the two exact tensor rows
supplied by this assertion. If the outside sheaves \(M',M''\) are flat,
their vertical maps are monic and so is the middle one: an element in
its kernel maps to zero in the right column and then is forced to be
zero by the left column. This is the same kernel argument on every
stalk. If \(M,M''\) are flat, the injection
\(N\otimes M'\to N\otimes M\), followed by the monomorphism into
\(N_1\otimes M\), shows that the vertical map for \(M'\) is monic as
well.

Tensor commutes with colimits by its balanced universal property.
Exactness of filtered colimits, proved in the first lesson, therefore
proves flatness of a filtered colimit of flats. Finally
\(N\otimes(M_1\oplus M_2)=(N\otimes M_1)\oplus(N\otimes M_2)\); if this
map preserves each injection, so does each summand. \(\square\)

The flat-quotient assertion is often phrased as vanishing of first Tor
against a flat module. The proof here establishes the exact sequence
needed before Tor is defined in the next lesson.

\subsection{2. K-flat complexes and their closure
properties}\label{flat-modules-and-k-flat-resolutions--2-k-flat-complexes-and-their-closure-properties}

For cochain complexes define

\[(A\otimes P)^n=\bigoplus_{p+q=n}A^p\otimes P^q,
\qquad d(a\otimes p)=d_Aa\otimes p+(-1)^{|a|}a\otimes d_Pp.
\tag{2.1}\]

Two applications of \(d\) have opposite mixed terms, while the remaining
terms contain \(d_A^2\) and \(d_P^2\); hence \(d^2=0\). A complex \(P\)
is \textbf{K-flat} if \(A\otimes P\) is acyclic for every acyclic \(A\).
The graded flip \(a\otimes p\mapsto(-1)^{|a||p|}p\otimes a\) is an
isomorphism of complexes: on the two differential terms the sign
exponents differ by exactly the degree increased in that term. It
interchanges the tensor variables.

\textbf{Lemma 2.1 (tensor respects homotopies and triangles).} For any
fixed complex \(A\), tensoring with \(A\) is an exact functor on the
homotopy category. A K-flat \(P\) therefore carries quasi-isomorphisms
in the other variable to quasi-isomorphisms.

\textbf{Proof.} If \(f-g=dh+hd\), then
\(H(a\otimes p)=(-1)^{|a|}a\otimes h(p)\) gives
\(1_A\otimes(f-g)=dH+Hd\); expanding the mixed terms cancels them. The
shift comparison \(A\otimes P[1]\to(A\otimes P)[1]\) multiplies a
component in \(A^p\) by \((-1)^p\). For the cone convention of the
common reading, the same factor on the shifted second summand gives

\[A\otimes\operatorname{Cone}(f)
\cong\operatorname{Cone}(1_A\otimes f).
\tag{2.2}\]

Indeed, the off-diagonal term on the left is \((-1)^pf\); the factor
\((-1)^p\) makes it the off-diagonal term on the right. On the second
diagonal block, the resulting differential is the negative of the total
differential on \(A\otimes P\). The comparison respects the cone
inclusion and the negative projection used in the distinguished
triangle. Thus shifts and cone triangles are preserved. The graded flip
gives the assertion in the other variable. A quasi-isomorphism has
acyclic cone. If \(P\) is K-flat, its tensor with that cone is acyclic,
and (2.2) proves the last assertion. \(\square\)

\textbf{Lemma 2.2 (stalks and pullbacks).} A complex \(P\) of module
sheaves is K-flat if and only if every stalk \(P_x\) is K-flat over
\(\mathcal O_x\). Pullback along any morphism of commutative ringed
spaces preserves flat sheaves and K-flat complexes.

\textbf{Proof.} Stalks commute with the tensor sheaf by Lemma 1.2 and
with direct sums by the first lesson. Hence
\((A\otimes P)_x=A_x\otimes P_x\). If the stalks of \(P\) are K-flat, an
acyclic \(A\) has acyclic tensor on all stalks, and the tensor sheaf
complex is acyclic. Conversely, for an acyclic complex \(C\) over
\(\mathcal O_x\), the complex of skyscrapers \(x_*C\) is acyclic: the
skyscraper functor is exact on sections, as established in Lemma 5.2 of
the first lesson. Its stalk at \(x\) is \(C\). Applying K-flatness of
\(P\) to this complex and taking that stalk proves K-flatness of
\(P_x\).

For a ring map \(R\to S\) and K-flat \(P\) over \(R\), any acyclic
complex \(C\) of \(S\)-modules remains acyclic after restriction of
scalars, and

\[C\otimes_S(S\otimes_R P)=C\otimes_R P\]

is acyclic. The isomorphism sends \(c\otimes(s\otimes p)\) to
\(sc\otimes p\); its inverse sends \(c\otimes p\) to
\(c\otimes(1\otimes p)\), and both commute with the total differentials.
This proves base change of K-flatness. The identical one-degree argument
proves base change of flat modules. The stalk formula for pullback from
the first lesson, with \(R=\mathcal O_{Y,f(x)}\) and
\(S=\mathcal O_{X,x}\), proves the sheaf statements. No flatness of
\(S\) over \(R\) is needed. \(\square\)

\textbf{Lemma 2.3 (operations preserving K-flatness).} A flat sheaf in a
single degree is K-flat. Direct sums and filtered colimits of K-flat
complexes are K-flat. If
\[0\longrightarrow P\longrightarrow Q\longrightarrow R\longrightarrow0
\tag{2.3}\] is a degreewise split exact sequence of complexes and
\(P,R\) are K-flat, then \(Q\) is K-flat. Every bounded-above complex of
flat sheaves is K-flat.

\textbf{Proof.} If \(A\) is acyclic and \(M\) is flat, tensor the exact
sequences \(0\to Z^nA\to A^n\to Z^{n+1}A\to0\) with \(M\). They identify
the kernels and images in \(A\otimes M\), proving its acyclicity.
Shifting a complex does not change this conclusion.

Tensor totalization commutes with direct sums and filtered colimits.
These operations are exact for module sheaves by Theorem 3.1 of the
first lesson (an arbitrary sum is the filtered colimit of its finite
partial sums). They therefore commute with cohomology: kernels and
images in an exact filtered colimit are the colimits of the
corresponding kernels and images. Apply this to \(A\otimes P_i\) to
obtain the assertions about sums and colimits.

After tensoring (2.3) with any complex \(A\), the sequence remains
degreewise split exact, including after totalization. Its long exact
cohomology sequence proves the extension assertion. The splitting is
only a splitting of the underlying graded sheaves; it need not respect
the differentials.

A bounded complex of flat sheaves is K-flat by induction on its length,
using its bottom-degree term as the quotient and the remaining tail as
the subcomplex in (2.3). If \(P\) is merely bounded above, its brutal
lower truncations \(P^{\geq m}\) are bounded complexes and are
subcomplexes of \(P\). As \(m\) decreases their union is \(P\). The
filtered-colimit assertion proves the last statement. \(\square\)

\textbf{Lemma 2.4 (tensor products and two out of three).} The tensor
product of two K-flat complexes is K-flat. Shifts, direct summands, and
homotopy equivalent complexes preserve K-flatness. Two out of three
objects in a distinguished triangle of the homotopy category being
K-flat forces the third to be K-flat. The same two-out-of-three
statement holds in a short exact sequence of complexes whose quotient
has flat terms.

\textbf{Proof.} Direct-sum reindexing identifies
\(A\otimes(P\otimes Q)\) with \((A\otimes P)\otimes Q\). On a term of
degrees \(a,p,q\), both differentials have signs
\(1,(-1)^a,(-1)^{a+p}\). If \(A\) is acyclic, successively applying
K-flatness of \(P\) and \(Q\) gives acyclicity. Shifts and homotopy
equivalences follow from Lemma 2.1. A direct summand of an acyclic
tensor complex is acyclic, giving the summand assertion. Tensoring a
distinguished triangle gives a distinguished triangle by that lemma, and
its long exact cohomology sequence gives two out of three. For a short
exact sequence with flat quotient terms, Lemma 1.3 makes its tensor with
each term of any complex exact. Direct sums preserve exactness, so
totalization gives a short exact sequence of tensor complexes. Its long
exact sequence supplies the last claim. \(\square\)

Only \textbf{filtered} colimits occur in Lemma 2.3. An arbitrary
cokernel of K-flat complexes need not be K-flat; Section 5 supplies the
smallest counterexample.

\subsection{3. A surjective resolution by attaching
cycles}\label{flat-modules-and-k-flat-resolutions--3-a-surjective-resolution-by-attaching-cycles}

\textbf{Theorem 3.1 (K-flat resolutions).} Every complex \(G\) of
\(\mathcal O\)-modules has a quasi-isomorphism \(P\to G\), surjective in
each degree, where \(P\) is K-flat and every \(P^n\) is flat.

\textbf{Proof.} For a sheaf \(M\), let \(D_n(M)\) have \(M\) in degrees
\(n,n+1\), identity differential between them, and zero elsewhere. A map
\(M\to G^n\) extends uniquely to a chain map \(D_n(M)\to G\), with upper
component its composite with \(d_G\). Start with \[\begin{gathered}
P_0=\bigoplus_{n\in\mathbb Z}D_n(L(G^n)),\\
\epsilon_0:P_0\longrightarrow G.
\end{gathered}
\tag{3.1}\] Use \(\lambda_{G^n}\) on the lower term of the \(n\)-th
summand. The map is onto in every degree. Each disk is contractible.
Tensoring a contraction \(h\) on the second factor with a complex \(A\)
gives the contraction \(a\otimes p\mapsto(-1)^{|a|}a\otimes h(p)\). Thus
\(P_0\) is K-flat; its terms are flat by Lemma 1.2.

Suppose \(P_s\to G\) has been built, and put \(K_s=\ker(\epsilon_s)\), a
subcomplex. For every integer \(q\), open \(U\), and section
\[z\in Z^q(K_s)(U),\] add a copy of \(B_U\) in degree \(q-1\). Write its
generator locally as \(e_z\), and prescribe
\[d e_z=z,\qquad \epsilon_{s+1}(e_z)=0.
\tag{3.2}\] These are morphisms of sheaves by (1.1). The two conditions
needed for them to define a complex over \(G\) hold: \(dz=0\) and
\(\epsilon_s(z)=0\). Explicitly, if \(V_s^n\) is the sum of the newly
added \(B_U\)\textquotesingle s in degree \(n\), the underlying graded
sheaf is \(P_s\oplus V_s\), with differential
\[d(p,v)=(d_{P_s}p+\alpha_s(v),0),\] where
\(\alpha_s:V_s^n\to P_s^{n+1}\) sends each generator to its chosen
cycle. Since \(d_{P_s}\alpha_s=0\), its square is zero.

The inclusion \(P_s\to P_{s+1}\) is degreewise split, and the quotient
is \(V_s\) with zero differential. It is a direct sum of shifts of the
flat sheaves \(B_U\), hence K-flat. Lemma 2.3 proves by induction that
every \(P_s\) is K-flat. Its terms remain direct sums of the flat
building blocks.

Set \(P=\varinjlim_{s\geq0}P_s\). The maps to \(G\) agree, and the map
\(\epsilon:P\to G\) is termwise onto already from \(P_0\). Its terms are
flat, and \(P\) is K-flat by Lemma 2.3. Exactness of filtered colimits
gives \[\ker(\epsilon)=\varinjlim_s K_s,\qquad
H^q(\ker\epsilon)=\varinjlim_s H^q(K_s).
\tag{3.3}\] Each transition \(H^q(K_s)\to H^q(K_{s+1})\) is zero.
Indeed, a stalk class is locally represented by a cycle section \(z\);
the next stage supplies \(e_z\) in the kernel with \(de_z=z\). Thus
every term in the colimit on the right of (3.3) maps to zero at the next
stage, and the colimit vanishes. The kernel of \(\epsilon\) is acyclic.
The long exact sequence for \(0\to\ker\epsilon\to P\to G\to0\) proves
that \(\epsilon\) is a quasi-isomorphism. \(\square\)

The stages can create new cycles; the next stage kills them. A germ
involves a finite stage, so a countable sequence of these stages
suffices. Nothing in the construction asks sections to lift through an
arbitrary sheaf epimorphism: generators are added for sections that
actually exist on each open.

\textbf{Lemma 3.2 (making another resolution termwise surjective).} If
\(Q\to G\) is a quasi-isomorphism with \(Q\) K-flat, then
\(Q\oplus P_0\to G\), the sum of that map and (3.1), is a
termwise-surjective K-flat resolution. If \(Q\) also has flat terms, so
does this enlargement.

\textbf{Proof.} The new map is onto because \(\epsilon_0\) is. The
complex \(P_0\) is contractible, and the inclusion \(Q\to Q\oplus P_0\)
is a quasi-isomorphism commuting with the maps to \(G\). The new map is
therefore a quasi-isomorphism. The direct sum is K-flat by Lemma 2.3;
this uses no termwise-flatness assumption on \(Q\). If its terms are
flat, Lemma 1.2 gives flatness of the enlarged terms as well.
\(\square\)

\textbf{Remark 3.3 (what the construction does not prove).} Extensions
by zero of structure sheaves need not be projective objects in the
category of module sheaves. The construction gives a resolution suitable
for \textbf{tensor products}, not a K-projective resolution for
computing arbitrary Hom functors. Neither an inverse limit of injective
truncations nor exactness of products of sheaves was used.

\subsection{4. Compatible bounded resolutions and
factorization}\label{flat-modules-and-k-flat-resolutions--4-compatible-bounded-resolutions-and-factorization}

The cycle construction proves the unbounded existence theorem directly.
We also prove the stronger compatible bounded-above form used in the
Stacks treatment, with every term a direct sum of the flat building
blocks \(B_U\).

\textbf{Lemma 4.1 (bounded-above resolutions).} If \(C\) vanishes in
degrees greater than \(b\), there is a termwise-surjective
quasi-isomorphism \(E\to C\), with \(E\) vanishing above \(b\) and its
terms direct sums of \(B_U\).

\textbf{Proof.} At degree \(n\), replace the current term \(C^n\) by its
surjective cover \(u:L(C^n)\to C^n\). Replace degree \(n-1\) by the
pullback \(C^{n-1}\times_{C^n}L(C^n)\), with differential into degree
\(n\) its second projection. The next differential is \(d_Cu\); the
preceding differential sends \(c\) to \((d_Cc,0)\). These prescriptions
have square zero because the old differential had square zero. They give
a termwise epimorphic chain map to the old complex; the pullback
projection is epic because it is a pullback of \(u\). Its kernel is
\(\ker u\) in degrees \(n-1,n\), with identity differential, so the
replacement is a quasi-isomorphism.

Perform these replacements for \(n=b,b-1,b-2,\ldots\). A fixed term
changes at most twice and is finally a chosen cover \(L(-)\). Define
\(E\), its differential and its map to the original complex by these
stabilized components. The finite compositions into any fixed original
term are epic. For cohomology in degree \(r\), the three terms and two
differentials in degrees \(r-1,r,r+1\) stabilize after finitely many
replacements. At that point the finite composition is a
quasi-isomorphism, so the stabilized map is an isomorphism on \(H^r\).
This proves the result without assuming exact inverse limits.
\(\square\)

\textbf{Lemma 4.2 (an explicit extension over a cone).} Given
\(a:P\to H\) and a termwise-surjective quasi-isomorphism
\(u:Q\to C(a)\), write \(u(q)=(u_Hq,u_Pq)\), where
\(u_P:Q^n\to P^{n+1}\). Define

\[\begin{gathered}
N^n=P^n\oplus Q^n,\\
d_N(p,q)=(d_Pp-u_Pq,d_Qq),\\
f(p,q)=a(p)+u_Hq.
\end{gathered}
\tag{4.1}\]

Then \(P\to N\) is a degreewise split inclusion, its quotient is \(Q\),
\(f:N\to H\) is a termwise-surjective quasi-isomorphism, and \(f|_P=a\).

\textbf{Proof.} The chain-map equation for \(u\) gives
\(d_Pu_P+u_Pd_Q=0\) and \(d_Hu_H+a u_P=u_Hd_Q\). These identities give
\(d_N^2=0\) and \(d_Hf=f d_N\) by substitution. Surjectivity of \(u\)
followed by the graded projection \(C(a)^n\to H^n\) proves surjectivity
of \(u_H\), and hence of \(f\).

Compare the triangle from \(0\to P\to N\to Q\to0\) with the cone
triangle of \(a\). The maps on \(P,N,Q\) are respectively \(1,f,u\). The
first square commutes strictly. In the second square, \(u\pi-i f\) is
the homotopy boundary of \(h(p,q)=(0,-p)\) into \(C(a)^{n-1}\), because

\[d_Ch+h d_N=(-a(p),u_Pq)=u\pi-i f.\]

The boundary of the split extension is \(-u_P:Q\to P[1]\), since that is
its off-diagonal differential, and it equals the negative cone
projection composed with \(u\). Thus all squares commute on cohomology.
The outside vertical maps are an identity and a quasi-isomorphism; the
long exact cohomology sequences and the Five Lemma prove that \(f\) is a
quasi-isomorphism. \(\square\)

\textbf{Proposition 4.3 (compatible bounded-above system).} For any
complex \(G\), there are termwise-surjective quasi-isomorphisms
\(E_n\to\tau_{\le n}G\), for \(n=1,2,\ldots\), with \(E_n\) bounded
above and all its terms sums of \(B_U\). They form a commutative direct
system whose transitions are termwise split injections with quotients
termwise sums of \(B_U\). The induced map \(\varinjlim E_n\to G\) is a
termwise-surjective quasi-isomorphism from a K-flat complex of flat
terms.

\textbf{Proof.} Recall that \(\tau_{\le n}G\) agrees with \(G\) below
\(n\), has \(\ker d_G^n\) at \(n\), and vanishes above \(n\). Its
inclusions into the next truncation are chain maps, and their union is
\(G\). Lemma 4.1 supplies \(E_1\). Suppose \(E_n\to\tau_{\le n}G\) has
been constructed. Compose it into \(H=\tau_{\le n+1}G\), and take the
cone \(C(a)\). This cone is bounded above. Resolve it by a
termwise-surjective \(Q\to C(a)\) as in Lemma 4.1 and apply Lemma 4.2.
Set \(E_{n+1}=N\) and its map to \(H\) equal to \(f\). Its terms are
sums of the desired blocks, it is bounded above, and the displayed
formula makes the transition square strictly commute. Its graded
quotient by \(E_n\) is \(Q\), with the required terms.

Exact filtered colimits commute with cohomology, so the union map is a
quasi-isomorphism. It is onto in each degree because the truncations
exhaust \(G\). Each \(E_n\) is K-flat by Lemma 2.3, so their filtered
colimit is K-flat. Its terms are flat by Lemma 1.3, or directly from the
chosen graded splittings as sums of the flat blocks. \(\square\)

\textbf{Corollary 4.4 (factoring a map from a K-flat complex).} If
\(a:P\to H\) and \(P\) is K-flat, then \(a\) factors through a K-flat
\(N\to H\) that is a termwise-surjective quasi-isomorphism. If \(P\) has
flat terms, \(N\) can have flat terms too.

\textbf{Proof.} Resolve \(C(a)\) by Theorem 3.1 and apply Lemma 4.2. The
split exact sequence \(0\to P\to N\to Q\to0\) makes \(N\) K-flat by
Lemma 2.3. Its terms \(P^n\oplus Q^n\) are flat when the terms of \(P\)
are. The factorization is strict in this construction; the corresponding
homotopy-category statement therefore follows. \(\square\)

\textbf{Lemma 4.5 (comparison after tensoring).} A quasi-isomorphism
between K-flat complexes remains a quasi-isomorphism after tensoring
with any complex. An acyclic K-flat complex consequently has acyclic
tensor with every complex.

\textbf{Proof.} For a complex \(A\), choose a K-flat resolution
\(F\to A\) by Theorem 3.1. If \(P\to Q\) is a quasi-isomorphism between
K-flats, in the tensor square with rows \(F\otimes P\to F\otimes Q\) and
\(A\otimes P\to A\otimes Q\), the vertical maps are quasi-isomorphisms
by K-flatness of \(P,Q\), and the top map is one by K-flatness of \(F\).
Two out of three forces the bottom map to be one. Apply this with
\(P\to0\) for the last assertion. \(\square\)

\subsection{5. Examples and failed
shortcuts}\label{flat-modules-and-k-flat-resolutions--5-examples-and-failed-shortcuts}

\textbf{K-flat with nonflat terms.} The disk \(D_0(\mathbb Z/2)\) is
contractible. The contraction used for the disks in Theorem 3.1 makes
its tensor with every complex contractible, so it is K-flat. Its nonzero
terms are not flat over \(\mathbb Z\): tensoring the injection
\(2:\mathbb Z\to\mathbb Z\) gives zero on the nonzero group
\(\mathbb Z/2\).

\textbf{Free terms without K-flatness.} Over
\(R=k[\epsilon]/(\epsilon^2)\), let \(P\) have \(R\) in every integer
degree and differential \(\epsilon\). It is acyclic since both kernel
and image are \(\epsilon R\), and its terms are free. Let \(F\) be its
left half in degrees at most zero, with augmentation
\(F\to k=R/(\epsilon)\). Then \(F\) resolves \(k[0]\): the negative
degrees are exact and the degree-zero cokernel is \(k\). It is K-flat by
the bounded-above lemma. Its tensor with the acyclic \(P\) is therefore
acyclic. But \(k\otimes P\) has one copy of \(k\) in every degree and
zero differential. If \(P\) were K-flat, it would carry \(F\to k[0]\) to
a quasi-isomorphism, which is impossible. More explicitly, the acyclic
test \(C(F\to k[0])\) has tensor with \(P\) of cohomology \(k\) in every
degree, by the cone exact sequence. This displays the test that fails.

\textbf{Why a general colimit is different.} The cokernel of
\(2:\mathbb Z[0]\to\mathbb Z[0]\) is \((\mathbb Z/2)[0]\). Both source
and target are K-flat, but the cokernel is not: tensoring the acyclic
complex \(\mathbb Z\xrightarrow{2}\mathbb Z\to\mathbb Z/2\) in degrees
\(0,1,2\) gives a complex with zero first differential and nonzero
degree-zero cohomology. Thus the word ``filtered'' in the colimit
theorem cannot be dropped.

\textbf{Local generators need not be projective.} On \([0,1]\) with the
constant ring sheaf \(\mathbb Z\), let \(j:(0,1)\hookrightarrow[0,1]\).
The sequence
\(0\to j_!\mathbb Z\to\mathbb Z_{[0,1]}\to i_*\mathbb Z_{\{0,1\}}\to0\)
is exact on stalks. Its global-section map in the last two terms is the
diagonal \(\mathbb Z\to\mathbb Z^2\), so it is not onto. Therefore
\(\operatorname{Hom}(B_X,-)=\Gamma(X,-)\) is not exact and the flat
building block \(B_X=\mathcal O\) is not projective. The resolution is
suited to tensor products; it supplies no K-projectivity theorem for
Hom.

\subsection{6. Exercises with checked
solutions}\label{flat-modules-and-k-flat-resolutions--6-exercises-with-checked-solutions}

\textbf{Exercise 1 (easy: summands).} Prove that a direct summand of a
flat sheaf is flat, and that a direct summand of a K-flat complex is
K-flat.

\textbf{Solution.} Tensor distributes over the finite direct sum
defining the splitting. For an injection of test sheaves, injectivity of
the direct sum of the two tensor maps forces injectivity of each
component. Right exactness gives flatness of each sheaf summand. For
complexes, cohomology of a finite direct sum is the direct sum of
cohomologies. If tensoring with an acyclic test has acyclic sum, each
tensor summand has zero cohomology in every degree. That is K-flatness.

\textbf{Exercise 2 (medium: the converse stalk test).} Why can the stalk
K-flatness theorem be proved using a skyscraper complex even when the
point is not closed?

\textbf{Solution.} The skyscraper functor in the first lesson is defined
on every point: sections are the given stalk module on opens containing
the point and zero otherwise. An exact sequence of these modules gives
an exact section sequence on every open, hence an exact sheaf sequence.
Its stalk at the chosen point is the original module, because every
neighbourhood there contributes that module with identity restrictions.
No assertion about stalks at other points is needed. For an acyclic
stalk complex \(C\), use \(x_*C\) as the sheaf test, then take its
tensor stalk \(C\otimes P_x\). Thus the proof applies without a
closed-point assumption.

\textbf{Exercise 3 (medium: the bound direction).} Prove that a
bounded-above complex of flat sheaves is K-flat and explain why the same
proof does not establish the assertion for a bounded-below complex of
flat sheaves.

\textbf{Solution.} A bounded flat complex is built from its individual
terms by graded split exact sequences: remove its lowest-degree term as
quotient, leaving the higher-degree tail as subcomplex. Lemma 2.3 gives
induction on length. If \(P\) is bounded above, the brutal lower
truncations \(P^{\ge m}\) are bounded subcomplexes, because the
differential raises degree. Their inclusions as \(m\) decreases form a
filtered system with union \(P\), so filtered closure proves the
assertion. Brutal upper truncations do not form subcomplexes when the
outgoing differential at their top degree is nonzero. Thus the proposed
bounded-below analogue lacks the system of inclusions used in this
proof; no such analogue is inferred.

\textbf{Exercise 4 (hard: a two-term Sierpiński example).} Let
\(X=\{\eta,s\}\) have opens \(\varnothing,\{\eta\},X\) and constant
structure ring \(\mathbb Z\). Set \(B=B_X\), \(E=B_{\{\eta\}}\), and let
\(i:E\to B\) be the canonical inclusion. Let \(M\) have section diagram
\(M(X)=\mathbb Z/2\to M(\{\eta\})=0\), and \(T=j_{\eta!}(\mathbb Z/2)\).
Construct a bounded-above resolution by sums of \(B,E\) of the two-term
complex \(G=(M\xrightarrow{0}T)\) in degrees \(0,1\), and the compatible
system of Proposition 4.3 for this \(G\).

\textbf{Solution.} On this space a module sheaf is exactly its
restriction map from \(X\) to \(\{\eta\}\); the only other sheaf
condition is on the empty set. Here \(B\) has diagram
\(\mathbb Z\xrightarrow{1}\mathbb Z\), and \(E\) has \(0\to\mathbb Z\).
Define

\[P^{-2}=E,\quad P^{-1}=E\oplus B,\quad
P^0=B\oplus E,\quad P^1=E,\]

with zero other terms and differentials

\[t\longmapsto(-2t,i(t)),\qquad
(a,b)\longmapsto(i(a)+2b,0),\qquad
(b,c)\longmapsto2c.\]

Both adjacent composites vanish. The augmentation reduces the \(B\)
coordinate in degree zero onto \(M\) and reduces the degree-one \(E\)
onto \(T\); it is epic on every stalk. At \(s\), the \(E\) stalks vanish
and the complex is \(\mathbb Z\xrightarrow{2}\mathbb Z\) in degrees
\(-1,0\), resolving \(M_s\). At \(\eta\), the first three components
contributed by \(M\) are \(\mathbb Z\to\mathbb Z^2\to\mathbb Z\), with
maps \((-2,1)\) and \((1,2)\). The first map is injective, its image is
the kernel of the second, and the second is onto. The remaining \(E\)
coordinates give \(\mathbb Z\xrightarrow{2}\mathbb Z\) in degrees
\(0,1\), resolving \(T_\eta\) in degree one. Thus the augmentation is a
quasi-isomorphism. Each term is a sum of the required flat blocks, and
bounded-above flatness proves K-flatness. Since \(\tau_{\le n}G=G\) for
all \(n\ge1\), take every \(E_n=P\) and every transition the identity.
These are split injections with zero quotient, which is the empty sum of
flat blocks, and meet all properties of Proposition 4.3.

\textbf{Exercise 5 (hard: the extension signs).} In Lemma 4.2, verify
both the chain-map equation for \(f\) and the homotopy in the second
triangle square. What breaks if only the sign of \(u_P\) in \(d_N\) is
reversed?

\textbf{Solution.} The first component of the equation \(d_Cu=u d_Q\) is
\(d_Hu_H+a u_P=u_Hd_Q\). Therefore

\[d_Hf(p,q)=a d_Pp+u_Hd_Qq-a u_Pq=f d_N(p,q).\]

For \(h(p,q)=(0,-p)\), the two homotopy terms are \((-a p,d_Pp)\) and
\((0,-d_Pp+u_Pq)\), whose sum is \((-a p,u_Pq)=u\pi-i f\). Replacing
just the off-diagonal sign gives \(f d_N=a d_Pp+a u_Pq+u_Hd_Qq\), which
differs from \(d_Hf\) by \(2a u_Pq\) in general. It also reverses the
extension boundary without reversing the cone projection. The coherent
negative signs in the displayed construction are essential.

The next lesson uses these resolutions to define the derived tensor
product and Tor sheaves, proving choice independence from Lemma 4.5 and
the localization universal property.

Sources and licensing: Stacks project authors,
\href{https://github.com/KokunoYumeto/unofficial-stacks-project-ai-drafts/blob/565b10e987aba5969b21145a0833f42d69f96790/modules.tex\#L2669}{\emph{Modules
on Ringed Spaces}, flat modules};
\href{https://github.com/KokunoYumeto/unofficial-stacks-project-ai-drafts/blob/565b10e987aba5969b21145a0833f42d69f96790/cohomology.tex\#L6239}{\emph{Cohomology
of Sheaves}, flat resolutions}; and
\href{https://github.com/KokunoYumeto/unofficial-stacks-project-ai-drafts/blob/565b10e987aba5969b21145a0833f42d69f96790/derived.tex\#L9561}{\emph{Derived
Categories}, compatible direct systems}. The complete used proofs were
read and their short module/cone arguments expanded here. The cycle
construction is an expository alternative to the
source\textquotesingle s compatible bounded resolutions, not a new
existence theorem. N. Spaltenstein, \emph{Resolutions of unbounded
complexes}, Compositio Mathematica 65 (1988), 121--154, is
\href{https://www.numdam.org/item/CM_1988__65_2_121_0/}{available from
Numdam}. See the \hyperref[sources-and-authorship]{combined licence
notice} and \hyperref[gnu-free-documentation-license]{GNU FDL text}.

\section{The derived tensor product and Tor
sheaves}\label{derived-tensor-products-and-tor-sheaves}

\emph{Written and edited by GPT-6.1 Sol (OpenAI), in Codex at Ultra,
October 2026. Mathematically self-checked by the writing AI. Original
contributions are CC0; the combined course is distributed under
GFDL-1.2-or-later. Full authorship and source attribution appear in the
\hyperref[sources-and-authorship]{course notice}.}

Ordinary tensor products preserve quotients, but can lose an injection.
For example, tensoring \(\mathbb Z\xrightarrow{2}\mathbb Z\) with
\(\mathbb Z/2\) turns an injection into the zero map. A resolution keeps
the lost information in negative cohomological degrees. The derived
tensor product packages this information without depending on a chosen
resolution, and Tor sheaves name its cohomology for two module sheaves.

The prerequisite is \hyperref[flat-modules-and-k-flat-resolutions]{Flat
modules and K-flat resolutions}, Lemmas 1.1--1.3 and 2.1--2.4, Theorem
3.1, Lemmas 3.2 and 4.1, and Lemma 4.5. We also use its earlier
prerequisite, \hyperref[complexes-cones-and-localization]{Complexes,
cones and localization}, Theorems 4.2 and 5.1 and Proposition 5.2. In
particular, roofs and their common refinements have already been proved,
rather than assumed.

Throughout, \(X\) is an arbitrary topological space and \(\mathcal O\)
is a sheaf of \textbf{commutative unital rings}. Complexes are
cohomologically indexed and can be unbounded in both directions. Their
tensor product means the direct-sum totalization

\[\begin{aligned}
(A\otimes B)^n&=\bigoplus_{p+q=n}A^p\otimes_{\mathcal O}B^q,\\
d(a\otimes b)&=d_Aa\otimes b\\
&\quad+(-1)^p a\otimes d_Bb .
\end{aligned}
\tag{0.1}\]

We write \(K(\mathcal O)\) for the homotopy category, \(D(\mathcal O)\)
for its localization at quasi-isomorphisms, and \(Q\) for localization.
Cone triangles have the convention \((f,i,-p)\) of the common reading.
No compactness, separation, finite dimension or Noetherian assumption
enters the construction.

The construction follows the Stacks project authors' \emph{Cohomology of
Sheaves}, ``Flat resolutions'', especially Tags 06YG, 06YH, 08BP and
08BQ, in the
\href{https://github.com/KokunoYumeto/unofficial-stacks-project-ai-drafts/blob/565b10e987aba5969b21145a0833f42d69f96790/cohomology.tex}{AI
Integrated Stacks Project edition}. We prove the comparison and
coherence maps before defining Tor sheaves. Source attribution and the
licence for adapted passages appear in the
\hyperref[sources-and-authorship]{course notice}.

\subsection{1. The two tensor
comparisons}\label{derived-tensor-products-and-tor-sheaves--1-the-two-tensor-comparisons}

A K-flat complex \(P\) has \(A\otimes P\) acyclic whenever \(A\) is
acyclic. The two variables in a tensor comparison play different roles.

\textbf{Lemma 1.1 (tensor invariance).} If \(P\) is K-flat and
\(u:A\to B\) is a quasi-isomorphism, then \(u\otimes1_P\) is a
quasi-isomorphism. If \(v:P\to P'\) is a quasi-isomorphism
\textbf{between K-flat complexes}, then \(1_G\otimes v\) is a
quasi-isomorphism for every complex \(G\).

\textbf{Proof.} Lemma 2.1 of the K-flat lesson identifies the tensor of
a cone with the cone of the tensored map. When the cone is in the second
variable, the shifted summand acquires the factor \((-1)^{|a|}\); the
graded flip gives the other variable. The cone of \(u\) is acyclic, so
its tensor with \(P\) is acyclic. This proves the first assertion.

For the second, choose a K-flat resolution \(R\to G\), using Theorem 3.1
of that lesson. Consider the commuting square

\[\begin{array}{ccc}
R\otimes P&\longrightarrow&R\otimes P'\\
\downarrow&&\downarrow\\
G\otimes P&\longrightarrow&G\otimes P'.
\end{array}
\tag{1.1}\]

Both vertical maps are quasi-isomorphisms by the first assertion, using
respectively \(P\) and \(P'\). The top map is one because \(R\) is
K-flat, after interchanging the tensor factors. On every cohomology
sheaf, three arrows are isomorphisms and the square commutes; the fourth
is therefore an isomorphism. This proves the second assertion for
arbitrary, possibly unbounded and nonflat, \(G\). In particular an
acyclic K-flat complex tensors acyclicly with every complex, by applying
this assertion to its map to zero. \(\square\)

The second assertion would be false for an arbitrary quasi-isomorphism
in place of \(v\). For example
\((\mathbb Z\xrightarrow{2}\mathbb Z)\to\mathbb Z/2[0]\), in degrees
\(-1,0\), becomes a map from a zero-differential two-term complex to a
one-term complex after tensoring with \(\mathbb Z/2\). It loses an
isomorphism in degree \(-1\). Its target is not K-flat.

\textbf{Theorem 1.2 (one-variable definition and independence).} For a
K-flat resolution \(\epsilon:P\to F\), the rule \(G\mapsto G\otimes P\)
defines an exact functor

\[-\otimes_{\mathcal O}^{\mathbf L}F:D(\mathcal O)\longrightarrow D(\mathcal O).
\tag{1.2}\]

The functors from different resolutions have canonical natural
comparison isomorphisms satisfying the identity and cocycle laws. If
\(R\to G\) is K-flat, both arrows in

\[G\otimes P\ \longleftarrow\ R\otimes P\ \longrightarrow\ R\otimes F
\tag{1.3}\]

are quasi-isomorphisms. Thus either variable, or both variables, may be
resolved.

\textbf{Proof.} Tensor respects homotopies, shifts and cone triangles by
Lemma 2.1 of the K-flat lesson. Lemma 1.1 shows that \(Q(-\otimes P)\)
inverts quasi-isomorphisms, so the localization universal property,
Theorem 4.2 of the common reading, gives (1.2). Theorem 5.1 there
identifies distinguished triangles with images of cone triangles; their
tensor images are distinguished, so the functor is exact.

For the comparison, first let \(\epsilon_i:P_i\to F\), \(i=1,2\), be
termwise-surjective K-flat resolutions. Form the ordinary fibre product
of complexes

\[W=P_1\times_F P_2.
\tag{1.4}\]

Its projection onto \(P_1\) is termwise surjective with kernel
\(\ker\epsilon_2\); similarly for the other projection. Each kernel is
acyclic, by the short exact cohomology sequence. Hence both projections
are quasi-isomorphisms. Choose a termwise-surjective K-flat resolution
\(T\to W\). Its maps \(r_i:T\to P_i\) are quasi-isomorphisms between
K-flats. Lemma 1.1 gives a natural comparison \(c_{12}(G)\) from
\(Q(G\otimes P_1)\) to \(Q(G\otimes P_2)\), given by

\[c_{12}(G)=Q(1\otimes r_2)Q(1\otimes r_1)^{-1}.
\tag{1.5}\]

Two choices \(T,T'\to W\) have a common refinement: resolve
\(T\times_W T'\) by a termwise-surjective K-flat complex. The
projections are again surjective quasi-isomorphisms. All maps to the
\(P_i\) commute strictly, so after tensoring and inverting the
refinement arrows the two formulas (1.5) agree.

For three resolutions, resolve \(P_1\times_F P_2\times_F P_3\). Its
projections to the pairwise products are surjective quasi-isomorphisms,
with acyclic kernels. The previous refinement argument allows every
pairwise comparison to be computed through this one model. Cancelling
the middle invertible arrow proves the cocycle law.

For the identity law put \(W=P\times_F P\). Its diagonal
\(\delta:P\to W\) is a quasi-isomorphism, because either projection is
one and its composite with \(\delta\) is the identity. Lemma 3.2 of the
K-flat lesson enlarges \(\delta\) to a termwise-surjective K-flat
resolution \(P\oplus D\to W\), with \(D\) contractible. The two maps
from this resolution to \(P\) agree after precomposition with
\(P\to P\oplus D\). Tensoring that inclusion with \(G\) gives an
invertible derived map by Lemma 1.1. The two tensor maps therefore agree
in the derived category, proving the identity law. The maps are natural
in \(G\), first for chain maps and homotopies and then for roofs by
localization.

If a chosen resolution is not termwise surjective, Lemma 3.2 of the
K-flat lesson enlarges it by a contractible disk complex. Its inclusion
into the enlargement is a quasi-isomorphism between K-flats over \(F\).
Lemma 1.1 supplies the comparison to the enlargement. Here is
compatibility for two enlargements \(E_1,E_2\) of \(P\). The inclusions
give \(b:P\to W=E_1\times_F E_2\); it is a quasi-isomorphism because
\(W\to E_1\) and \(P\to E_1\) are. Resolve \(W\) by a surjective
\(T\to W\). The projection \(T\times_W P\to P\) is a surjective
quasi-isomorphism, with kernel \(\ker(T\to W)\). Its other projection to
\(T\) is a quasi-isomorphism by two out of three after composing into
\(W\). Resolve this fibre product by a K-flat complex. Its commuting
maps to \(P,T,E_1,E_2\), all quasi-isomorphisms between K-flats, show
after tensoring that the comparison of the enlargements intertwines
their inclusions from \(P\). This proves independence of the
enlargement.

Finally the left arrow of (1.3) is a quasi-isomorphism because \(P\) is
K-flat; the right one is a quasi-isomorphism because \(R\) is K-flat.
This proves balancing. \(\square\)

\subsection{2. Morphisms, associativity and
symmetry}\label{derived-tensor-products-and-tor-sheaves--2-morphisms-associativity-and-symmetry}

Choosing resolutions of objects does not by itself define the action on
morphisms. We now justify that action and its compatibility with every
comparison above.

\textbf{Lemma 2.1 (localizing the K-flat models).} Let
\(K_{\rm fl}(\mathcal O)\) be the full subcategory of \(K(\mathcal O)\)
on K-flat complexes, and let \(S_{\rm fl}\) consist of its
quasi-isomorphisms. The inclusion induces an equivalence

\[K_{\rm fl}(\mathcal O)[S_{\rm fl}^{-1}]
\ \simeq\ D(\mathcal O).
\tag{2.1}\]

\textbf{Proof.} Lemma 2.4 of the K-flat lesson shows that this full
subcategory contains zero and is closed under shifts, finite sums and
cones. All fraction squares and cancellation constructions in Lemma 4.1
of the common reading use these operations. When their starting objects
are K-flat, their new cone or shifted cone is therefore K-flat.
Consequently the roof construction and equivalence by common refinements
in Theorem 4.2 there apply within \(K_{\rm fl}\).

Every object of \(D(\mathcal O)\) is isomorphic to a K-flat resolution,
proving essential surjectivity. A roof \(P\leftarrow Z\to R\) in the
ambient homotopy category, with \(P,R\) K-flat, can be refined by a
K-flat resolution \(T\to Z\). The composite \(T\to P\) is a
quasi-isomorphism. Thus it has a roof entirely in \(K_{\rm fl}\),
proving fullness. If two such roofs are equal in \(D(\mathcal O)\),
Theorem 4.2 supplies a common refinement, possibly with a nonflat middle
complex \(Z\). Resolve that \(Z\) by a K-flat \(T\). The composites are
a common refinement within \(K_{\rm fl}\), proving faithfulness. These
operations use homotopy classes of maps; strict equality of arbitrary
resolution lifts is not assumed. \(\square\)

\textbf{Theorem 2.2 (the tensor bifunctor).} There is a bifunctor

\[\otimes_{\mathcal O}^{\mathbf L}:
D(\mathcal O)\times D(\mathcal O)\longrightarrow D(\mathcal O)
\tag{2.2}\]

exact in each variable, whose value is represented by \(P\otimes R\) for
K-flat models of its two arguments. It has canonical associativity,
symmetry and unit isomorphisms, and commutes with restriction to any
open subset.

\textbf{Proof.} On \(K_{\rm fl}\times K_{\rm fl}\) use ordinary total
tensor. It respects homotopies and sends a quasi-isomorphism in either
variable to one by Lemma 1.1. Localize successively in each variable. To
check that the two resulting morphism actions commute, first note that
\(f\otimes1\) and \(1\otimes g\) commute for chain maps of degree zero.
They therefore commute with inverse denominators too: invert either side
of the commuting square. A general roof is a numerator followed by an
inverse denominator. This proves the commuting rule for all pairs of
roofs and defines a bifunctor on the two localizations. Lemma 2.1
transfers it to \(D(\mathcal O)\times D(\mathcal O)\).

Cones of maps between K-flats are K-flat. Cone squares and tensor shift
comparisons thus prove exactness on their localization. More explicitly,
any roof between K-flats can be written \(Qf\,(Qs)^{-1}\) with the
middle object K-flat; its cone triangle is isomorphic to the image of
the cone triangle of \(f\). Tensor preserves the latter, so it preserves
the former. Theorem 1.2 identifies this bifunctor with the one-variable
construction, including its resolution comparisons.

For three K-flat models \(P,R,T\), the ordinary association map is

\[\begin{gathered}
(P\otimes R)\otimes T\longrightarrow P\otimes(R\otimes T),\\
(p\otimes r)\otimes t\longmapsto p\otimes(r\otimes t).
\end{gathered}
\tag{2.3}\]

Its inverse removes the parentheses. Both totalizations have the same
direct sum indexed by \(i+j+k=n\). Their differentials on an elementary
tensor have the three terms with signs \(1,(-1)^i,(-1)^{i+j}\), so (2.3)
is a chain isomorphism. The tensor product of two K-flats is K-flat by
Lemma 2.4 of the previous lesson. Thus the left and right sides compute
the two iterated derived products. Naturality for chain maps extends to
roofs and inverse denominators as above, proving a natural derived
associativity isomorphism.

For \(|p|=i\) and \(|r|=j\), symmetry is represented by

\[c_{P,R}(p\otimes r)=(-1)^{ij}r\otimes p.
\tag{2.4}\]

For the \(d_Pp\) term the two routes have sign exponents \((i+1)j\) and
\(ij+j\); for the \(d_Rr\) term they have exponents \(i+i(j+1)\) and
\(ij\), equal modulo two. This checks the chain-map equation. Its square
has sign \((-1)^{2ij}=1\), giving its inverse. The unit is
\(\mathcal O[0]\), which is K-flat; the usual multiplication
\(P\otimes\mathcal O\to P\) is a chain isomorphism.

These identifications are coherent. For four factors every path around
the associativity pentagon sends a tensor to the same tensor with the
final parentheses. The unit triangle does the same after multiplying by
the degree-zero scalar. In the symmetry hexagon, moving a degree-\(i\)
tensor past degrees \(j,k\) gives either \(i(j+k)\) or \(ij+ik\); the
signs agree. These are equalities of chain maps, hence of derived maps.
Naturality and (2.1) imply that replacing any model preserves these
equalities.

Finally, restriction \(j^*\) to an open \(U\) is exact and commutes with
tensor and direct sums, by the stalk and sheafification formulas of the
first lesson. It preserves K-flatness by Lemma 2.2 of the K-flat lesson.
Restricting resolutions \(P\to F\), \(R\to G\) therefore gives
resolutions on \(U\), and the ordinary tensor identification gives

\[(F\otimes_{\mathcal O}^{\mathbf L}G)|_U
\cong F|_U\otimes_{\mathcal O_U}^{\mathbf L}G|_U.
\tag{2.5}\]

On local elementary tensors this is the identity, so it commutes with
the association, symmetry and unit maps and with all roof comparisons.
\(\square\)

The tensor sign in (2.4) concerns the degrees of the tensors. A morphism
in the homotopy or derived category has degree zero, so no additional
sign is inserted when composing the two morphism actions. Shifts carry
the signed identifications already checked in Lemma 2.1 of the K-flat
lesson.

\subsection{3. Tor sheaves and the flatness
criterion}\label{derived-tensor-products-and-tor-sheaves--3-tor-sheaves-and-the-flatness-criterion}

For module sheaves \(F,G\), viewed in degree zero, and integers
\(p\ge0\), define

\[\operatorname{Tor}^{\mathcal O}_p(F,G)
=H^{-p}(F\otimes_{\mathcal O}^{\mathbf L}G).
\tag{3.1}\]

They are sheaves of \(\mathcal O\)-modules and are functorial in both
arguments. The cohomological minus sign is essential: ordinary
projective or flat resolutions have their resolving terms in nonpositive
degrees.

\textbf{Proposition 3.1 (degree zero and stalks).} There is a natural
identification \(\operatorname{Tor}_0(F,G)=F\otimes G\). The derived
tensor of two sheaves has no positive cohomology. Moreover

\[\operatorname{Tor}^{\mathcal O}_p(F,G)_x
\cong\operatorname{Tor}^{\mathcal O_x}_p(F_x,G_x).
\tag{3.2}\]

\textbf{Proof.} Lemma 4.1 of the K-flat lesson supplies a resolution
\(P\to F[0]\) by flat terms, with \(P^n=0\) for \(n>0\). It is K-flat by
Lemma 2.3 there. Its degree-zero cohomology is \(F\), so
\(\operatorname{coker}(d_P^{-1})=F\). Tensoring this presentation with
\(G\) is right exact and gives

\[H^0(P\otimes G)
=\operatorname{coker}(d_P^{-1}\otimes1_G)=F\otimes G.\]

The tensor has no positive terms, hence no positive cohomology. The
augmentation induces the displayed degree-zero identification;
naturality follows either by the roof construction or by the uniqueness
of the degree-zero quotient map.

Stalks are exact and commute with tensor and direct sums, by Theorems
2.1 and 3.1 of the first lesson and Lemma 1.2 of the K-flat lesson.
Hence \((P\otimes G)_x=P_x\otimes_{\mathcal O_x}G_x\) and taking stalks
commutes with cohomology. Lemma 2.2 of that lesson says \(P_x\) is
K-flat, and its map to \(F_x\) is a quasi-isomorphism. It therefore
computes the tensor over the stalk ring, proving (3.2). \(\square\)

\textbf{Proposition 3.2 (long exact Tor sequence).} A short exact
sequence \(0\to F_1\to F_2\to F_3\to0\) gives a natural long exact
sequence

\[\begin{gathered}
\cdots\longrightarrow\operatorname{Tor}_p(F_1,G)\\
\to\operatorname{Tor}_p(F_2,G)
\to\operatorname{Tor}_p(F_3,G)\\
\longrightarrow\operatorname{Tor}_{p-1}(F_1,G)
\longrightarrow\cdots\\
\longrightarrow\operatorname{Tor}_1(F_3,G)
\longrightarrow F_1\otimes G\\
\longrightarrow F_2\otimes G
\longrightarrow F_3\otimes G\longrightarrow0.
\end{gathered}
\tag{3.3}\]

There is the analogous sequence in the second variable.

\textbf{Proof.} Proposition 5.2 of the common reading turns the short
exact sequence into a distinguished triangle. The exact tensor functor
of Theorem 2.2 gives a distinguished triangle after tensoring with
\(G\). Its long exact cohomology sequence exists by Theorem 5.1 of that
reading. Replace \(H^{-p}\) by (3.1) and \(H^0\) by Proposition 3.1; the
subsequent positive cohomology is zero. The connecting map is the usual
connecting cohomology map with the cone sign convention fixed above. The
construction is natural for maps of short exact sequences. Symmetry
gives the second-variable assertion. \(\square\)

\textbf{Theorem 3.3 (Tor detects flatness).} For a module sheaf \(F\),
the following are equivalent:

\begin{enumerate}
\tightlist
\item
  \(F\) is flat.
\item
  \(\operatorname{Tor}_1(F,G)=0\) for every module sheaf \(G\).
\item
  \(\operatorname{Tor}_p(F,G)=0\) for every \(p>0\) and every module
  sheaf \(G\).
\end{enumerate}

\textbf{Proof.} If \(F\) is flat, \(F[0]\) is K-flat by Lemma 2.3 of the
preceding lesson. The one-variable construction therefore computes
\(F\otimes^{\mathbf L}G\) as the ordinary sheaf \(F\otimes G\) in degree
zero. This proves \(1\Rightarrow3\Rightarrow2\).

For \(2\Rightarrow1\), take any injection \(G\to H\), with cokernel
\(C\). The second-variable version of (3.3) contains

\[\operatorname{Tor}_1(F,C)\longrightarrow F\otimes G
\longrightarrow F\otimes H.\]

The left sheaf is zero, so the right arrow is injective. Tensor is
already right exact, so it is exact and \(F\) is flat. This proves all
implications without a finiteness assumption. \(\square\)

The quantifier ``every \(G\)'' cannot be replaced by the single test
\(G=F\). The continuous-function example below makes this failure
explicit.

\subsubsection{Quotients by
ideals}\label{derived-tensor-products-and-tor-sheaves--quotients-by-ideals}

\textbf{Lemma 3.4.} For ideals \(I,J\) in a commutative ring \(R\),

\[\begin{gathered}
\operatorname{Tor}_0^R(R/I,R/J)\\
{}=R/(I+J),\\
\operatorname{Tor}_1^R(R/I,R/J)\\
{}=(I\cap J)/(IJ).
\end{gathered}
\tag{3.4}\]

For \(p\ge2\), \(\operatorname{Tor}_p^R(R/I,R/J)\cong
\operatorname{Tor}_{p-1}^R(I,R/J)\). The same assertions hold for ideal
sheaves, with intersections, sums and products taken as sheaf ideals.

\textbf{Proof.} Apply (3.3) to \(0\to I\to R\to R/I\to0\), with second
factor \(R/J\). Flatness of \(R\) annihilates all its positive Tor. The
resulting first segment is

\[\begin{gathered}
0\longrightarrow\operatorname{Tor}_1(R/I,R/J)
\longrightarrow I\otimes_R R/J
\\
\longrightarrow R/J
\longrightarrow (R/I)\otimes_R R/J\longrightarrow0.
\end{gathered}\]

The presentation of tensor by balancing relations identifies
\(I\otimes_R R/J\) with \(I/IJ\): the map sends \(i\otimes(r+J)\) to
\(ri+IJ\), and its inverse sends \(i+IJ\) to \(i\otimes1\). The kernel
of \(I/IJ\to R/J\) is \((I\cap J)/IJ\), and the cokernel is \(R/(I+J)\).
The remaining exact segments give the dimension-shifting assertion. For
sheaf ideals perform these calculations on every stalk; stalks preserve
the relevant kernels and images, including ideal products, so they give
the stated sheaf identities. \(\square\)

Formula (3.4) computes only the first two groups without extra
information on \(I\). The higher groups require a resolution or the
displayed dimension shift. Exercise 1 supplies a complete all-degree
computation for two integer ideals.

\subsection{4. Two geometric
calculations}\label{derived-tensor-products-and-tor-sheaves--4-two-geometric-calculations}

\subsubsection{Skyscrapers on a smooth one-dimensional local
model}\label{derived-tensor-products-and-tor-sheaves--skyscrapers-on-a-smooth-one-dimensional-local-model}

Take \(X=\mathbb R\) with \(\mathcal O\) the sheaf of real analytic
functions, and let \(k_a=a_*\mathbb R\), with the action given by
evaluation at \(a\). Multiplication by \(t-a\) gives an exact sequence

\[0\longrightarrow\mathcal O
\xrightarrow{\ t-a\ }\mathcal O
\longrightarrow k_a\longrightarrow0.
\tag{4.1}\]

Here is a stalk proof. At \(b\ne a\), the germ \(t-a\) is invertible. At
\(a\), a germ of an analytic function has a convergent power series in
\(t-a\). Vanishing at \(a\) means its constant coefficient is zero, so
dividing by \(t-a\) shifts that series and gives another convergent
series on a smaller interval. Conversely every such multiple vanishes at
\(a\). Multiplication is injective because a continuous function with
\((t-a)f(t)=0\) vanishes away from \(a\) and hence also at \(a\).
Evaluation is surjective on germs, by constants. The first
lesson\textquotesingle s stalk criterion proves (4.1).

The two \(\mathcal O\) terms give a bounded flat, hence K-flat,
resolution in degrees \(-1,0\). For the same point \(a\), tensoring it
with \(k_a\) makes its differential zero. Therefore

\[\begin{gathered}
\operatorname{Tor}_0(k_a,k_a)=k_a,\\
\operatorname{Tor}_1(k_a,k_a)=k_a,\\
\operatorname{Tor}_p(k_a,k_a)=0\quad(p\ge2).
\end{gathered}
\tag{4.2}\]

For \(a\ne b\), tensoring (4.1) with \(k_b\) makes the differential
multiplication by the nonzero real number \(b-a\), an isomorphism. All
Tor sheaves are then zero. The nonzero degree-one sheaf at a coincident
point records the kernel lost by ordinary tensor; its location is
visible directly from the stalk computation.

\subsubsection{Continuous functions and two axes in the
plane}\label{derived-tensor-products-and-tor-sheaves--continuous-functions-and-two-axes-in-the-plane}

Now let \(X=\mathbb R^2\) and \(\mathcal O\) be the sheaf of real-valued
continuous functions. Set \(Z=\{y=0\}\), \(Z'=\{x=0\}\), and write
\(I_Z,I_{Z'}\) for the sheaf ideals of functions vanishing on these
subsets. Define

\[\mathcal O_Z=\mathcal O/I_Z,\qquad
\mathcal O_{Z'}=\mathcal O/I_{Z'}.\]

These are the continuous-function sheaves along the axes, extended to
\(X\): a continuous function on an axis extends locally by projection
onto that axis, proving surjectivity of the restriction map on germs.
Its kernel is precisely the indicated vanishing ideal.

For these ideals

\[I_Z\cap I_{Z'}=I_ZI_{Z'}.
\tag{4.3}\]

Indeed, for a real continuous function \(f\) vanishing on both axes, put
\(h=\sqrt{|f|}\) and define \(g=f/\sqrt{|f|}\) where \(f\ne0\), with
\(g=0\) on its zero set. Both functions are continuous: at the zero set
their absolute values are \(\sqrt{|f|}\), which tends to zero. Both
vanish on each axis and \(f=gh\). This proves inclusion into the product
on every germ; the reverse inclusion follows from vanishing. Lemma 3.4
gives

\[\operatorname{Tor}_1(\mathcal O_Z,\mathcal O_{Z'})=0.
\tag{4.4}\]

The ordinary tensor is the real skyscraper at the origin. Outside the
origin at least one quotient has zero stalk. At the origin the sum of
the two ideals is the maximal ideal of germs vanishing there. To prove
this last assertion, decompose such a germ \(f\) as

\[f=\frac{x^2f}{x^2+y^2}+\frac{y^2f}{x^2+y^2},
\tag{4.5}\]

with both terms defined to be zero at the origin. Their absolute values
are bounded by \(|f|\), so they are continuous there. The first vanishes
on \(Z'\), and the second on \(Z\). Thus
\(\mathcal O/(I_Z+I_{Z'})=0_*\mathbb R\), including its evaluation
action. Equations (4.4) and (4.5) compute degree one and degree zero for
the required closed-subset example.

This behaviour of first Tor does not imply flatness. Let
\(R=\mathcal O_0\), \(\mathfrak m\) its evaluation kernel, and
\(k=R/\mathfrak m=\mathbb R\). The same square-root factorization gives
\(\mathfrak m=\mathfrak m^2\), so \(\operatorname{Tor}_1^R(k,k)=0\). But
multiplication by the coordinate \(x\) is injective on \(R\): a
continuous germ annihilated by \(x\) vanishes on the dense complement of
the vertical axis and therefore everywhere on a neighbourhood. The
two-term free resolution of \(R/(x)\) tensors with \(k\) to a
zero-differential complex, giving

\[\operatorname{Tor}_1^R(k,R/(x))=k\ne0.
\tag{4.6}\]

Thus \(k\) is nonflat by Theorem 3.3. Notice that \(I_{Z',0}\ne(x)\):
the continuous germ \(\sqrt{|x|}\) vanishes on the vertical axis, but
cannot be divided by \(x\) continuously at the origin. In the analytic
curve calculation division by its parameter was possible. The ideals in
these examples must therefore be calculated in their actual function
rings.

\subsubsection{Bounded inputs with an unbounded
answer}\label{derived-tensor-products-and-tor-sheaves--bounded-inputs-with-an-unbounded-answer}

On a one-point space put \(R=k[\varepsilon]/(\varepsilon^2)\) and view
\(k=R/(\varepsilon)\) in degree zero. The resolution

\[\cdots\xrightarrow{\varepsilon}R
\xrightarrow{\varepsilon}R
\xrightarrow{\varepsilon}R\longrightarrow k
\tag{4.7}\]

has its last \(R\) in degree zero. In each negative degree the kernel
and image of multiplication by \(\varepsilon\) are the same ideal
\((\varepsilon)\), and the degree-zero cokernel is \(k\). It is a
bounded-above free resolution, hence K-flat. Tensoring with \(k\) kills
every differential, so \(\operatorname{Tor}_p^R(k,k)=k\) for every
\(p\ge0\). Even two sheaves in degree zero can therefore have derived
tensor with infinitely many negative cohomology sheaves. The
perfect-complex lesson later proves the finite bounds furnished by a
finite locally free model.

\subsection{5. Exercises with checked
solutions}\label{derived-tensor-products-and-tor-sheaves--5-exercises-with-checked-solutions}

\textbf{Exercise 1 (easy: ideals on a one-point space).} Let \(R\) be a
commutative ring and \(I,J\) ideals. Derive formulas for Tor in degrees
zero and one. Then compute every
\(\operatorname{Tor}_p^{\mathbb Z}(\mathbb Z/m,\mathbb Z/n)\), for
positive integers \(m,n\ge2\). Explain why the first formulas alone do
not determine all higher groups for arbitrary ideals.

\textbf{Solution.} Tensor \(0\to I\to R\to R/I\to0\) with \(R/J\) in the
derived sense. The vanishing of positive Tor against \(R\) embeds first
Tor as the kernel of \(I/IJ\to R/J\). This kernel is \((I\cap J)/IJ\),
and the cokernel is \(R/(I+J)\). The next segments identify higher Tor
with \(\operatorname{Tor}_{p-1}(I,R/J)\), which need not vanish.

For \(\mathbb Z/m\), use the injective multiplication map in the free
resolution \(\mathbb Z\xrightarrow{m}\mathbb Z\), in degrees \(-1,0\).
Tensor with \(\mathbb Z/n\). Put \(d=\gcd(m,n)\). Its cokernel is
\(\mathbb Z/d\), and its kernel consists of the multiples of \(n/d\)
modulo \(n\), also a cyclic group of order \(d\). Thus

\[\operatorname{Tor}^{\mathbb Z}_p(\mathbb Z/m,\mathbb Z/n)
=
\begin{cases}
\mathbb Z/d,&p=0,1,\\
0,&p\ge2.
\end{cases}\]

The first-degree ideal formula agrees:
\(I\cap J=\operatorname{lcm}(m,n)\mathbb Z\) and \(IJ=mn\mathbb Z\). In
contrast, for the dual-number ring in (4.7), the ideals
\(I=J=(\varepsilon)\) have all positive Tor equal to \(k\), by that
explicit free resolution. Thus the higher-degree answer needs
information about the actual ideal or resolution.

\textbf{Exercise 2 (medium: associativity with changed models).} Let
\(P,R,T\) resolve \(F,G,H\). Prove the derived associativity
isomorphism, including independence when a different K-flat model
\(V\to P\otimes R\) is used for the intermediate object. Check the
symmetry hexagon on tensors of degrees \(1,1,-1\).

\textbf{Solution.} The product \(P\otimes R\) is K-flat, so
\((P\otimes R)\otimes T\) computes the left-associated product without
another resolution. The right-associated product is computed by
\(P\otimes(R\otimes T)\). Their summands are indexed by the same
triples, and their differential terms have signs
\(1,(-1)^i,(-1)^{i+j}\). Reassociation is therefore a chain isomorphism
with no extra sign.

The map \(V\to P\otimes R\) is a quasi-isomorphism between K-flats.
Lemma 1.1 makes \(V\otimes T\to(P\otimes R)\otimes T\) a
quasi-isomorphism. Its composite with reassociation gives the comparison
from this alternative intermediate model. For two such choices, resolve
their termwise-surjective enlargements\textquotesingle{} fibre product;
the resulting comparisons agree by Theorem 1.2. Changing the individual
\(P,R,T\) is handled by the same comparison and the naturality of the
ordinary association map. Consequently the isomorphism belongs to the
derived objects, rather than to the displayed choice alone.

For the stated degrees, moving the degree-one factor across the combined
block has sign \((-1)^{1(1-1)}=1\). Moving it successively past the two
odd-degree factors gives \((-1)^{1\cdot1}(-1)^{1\cdot(-1)}=(-1)(-1)=1\).
This checks that instance of the hexagon; the calculation
\(i(j+k)=ij+ik\) proves every instance.

\textbf{Exercise 3 (medium: which Tor tests suffice?).} Prove Theorem
3.3 from the derived construction. Show also that it is enough to test
\(\operatorname{Tor}_1(F,x_*N)=0\) for every point \(x\) and every
\(\mathcal O_x\)-module \(N\). Explain the role of ``every'' using the
continuous germ ring at the origin.

\textbf{Solution.} A flat \(F\) in degree zero is K-flat, so its
ordinary tensor with \(G[0]\) already computes the derived tensor and
has no negative cohomology. Thus flatness implies all positive Tor
vanish, which implies first Tor vanishes. Conversely, the exact sequence
for an injection \(G\to H\) has
\(\operatorname{Tor}_1(F,\operatorname{coker}(G\to H))\) immediately
before \(F\otimes G\to F\otimes H\). Vanishing of this first Tor makes
that map injective. Right exactness then proves flatness.

For the point tests, take the stalk at \(x\). Proposition 3.1 and
\((x_*N)_x=N\), from Lemma 5.2 of the first lesson, give
\(\operatorname{Tor}_1^{\mathcal O_x}(F_x,N)=0\). For any injection of
stalk-ring modules apply the same Tor exact sequence, proving that
\(F_x\) is flat. The stalk criterion, Lemma 1.2 of the K-flat lesson,
proves \(F\) flat.

In the continuous germ ring \(R\), the residue \(k\) has
\(\operatorname{Tor}_1(k,k)=\mathfrak m/\mathfrak m^2=0\), but
\(\operatorname{Tor}_1(k,R/(x))=k\). This particular additional test
detects the failure. Testing against the residue itself therefore does
not replace testing against all modules.

\textbf{Exercise 4 (hard: only one resolution for an unbounded input).}
Over \(\mathbb Z\), let \(F=\mathbb Z/2[0]\), and let \(G\) have
\(\mathbb Z/2\) in every integer degree with zero differentials. Compute
\(G\otimes^{\mathbf L}F\) by resolving only \(F\). Prove in general that
the resulting one-resolution answer agrees with resolving both
variables. Compare it with the ordinary tensor.

\textbf{Solution.} Use \(P=(\mathbb Z\xrightarrow{2}\mathbb Z)\) in
degrees \(-1,0\), with its augmentation to \(F\). This is a bounded
free, hence K-flat, resolution. In total degree \(n\), \(G\otimes P\)
has two summands, \(G^n\otimes P^0\) and \(G^{n+1}\otimes P^{-1}\). Both
are \(\mathbb Z/2\). The \(G\) differential is zero, and the
contribution of \(2\) on \(P\) is zero after tensoring; the total
differential is therefore zero. Hence

\[\begin{gathered}
H^n(G\otimes^{\mathbf L}F)
=(\mathbb Z/2)\oplus(\mathbb Z/2),\\
\text{for every }n\in\mathbb Z.
\end{gathered}\]

No resolution of \(G\) was needed for this computation. To verify
agreement in general, choose any K-flat \(R\to G\). The map
\(R\otimes P\to G\otimes P\) is a quasi-isomorphism by K-flatness of
\(P\), and \(R\otimes P\to R\otimes F\) is a quasi-isomorphism by
K-flatness of \(R\). This gives (1.3) with no boundedness assumption. If
\(P\) is replaced by another K-flat resolution, the common-refinement
comparison in Theorem 1.2 makes the resulting isomorphisms agree
naturally.

The ordinary tensor \(G\otimes F\) has only one \(\mathbb Z/2\) in each
degree. The augmentation map \(G\otimes P\to G\otimes F\) is the
projection to the \(P^0\) summand, so it is not a quasi-isomorphism. The
target \(F\) is not K-flat; it cannot serve as the target K-flat model
in the second assertion of Lemma 1.1.

\textbf{Exercise 5 (medium: where can tensor cohomology occur?).} For
complexes \(A,B\), put
\(\Sigma(A)=\{x:H^n(A)_x\ne0\text{ for some }n\}\). Prove

\[\Sigma(A\otimes^{\mathbf L}B)\subset\Sigma(A)\cap\Sigma(B).\]

Use this to compute the derived tensor of sheaves supported at two
distinct points of a Hausdorff space.

\textbf{Solution.} Choose K-flat resolutions \(P\to A\), \(R\to B\). At
a point outside \(\Sigma(A)\), the stalk \(A_x\) is acyclic, so \(P_x\)
is an acyclic K-flat complex. Lemma 1.1 says that its tensor with every
complex, in particular \(R_x\), is acyclic. Since stalks commute with
total tensor and cohomology, all stalk cohomology of \(P\otimes R\)
vanishes there. The same argument applies outside \(\Sigma(B)\), proving
the inclusion.

For a sheaf supported at one point the only possibly nonzero cohomology
stalk is at that point. If the two points are distinct the two sets are
disjoint, so every cohomology stalk of the derived tensor is zero.
Exactness detection by stalks makes the product zero in the derived
category. On a Hausdorff space ordinary skyscrapers have precisely this
stalk support: a different point has a neighbourhood avoiding the chosen
point. This recovers the distinct-point part of the curve example
without its parameter resolution.

The next lesson applies K-flat models to pullback and K-injective models
to pushforward, and proves their derived adjunction. The tensor
bifunctor and its comparisons established here will make the
pullback\textquotesingle s compatibility with tensor natural.

Sources and licensing: the tensor-invariance and construction arguments
of Lemma 1.1 and Theorem 1.2 are taken from the Stacks project, tags
06YA, 06YG and 06YH, in its AI Integrated Stacks Project edition (GFDL),
checked and edited by GPT-6.1 Sol (OpenAI), at Ultra, with the full
fibre-product comparison included. The Tor definition and flatness proof
are taken from the Stacks project, tags 08BP and 08BQ, in that edition
(GFDL), checked and edited by the same writing AI; the explicit K-flat
computation is given above. The pinned
\href{https://github.com/KokunoYumeto/unofficial-stacks-project-ai-drafts/blob/565b10e987aba5969b21145a0833f42d69f96790/cohomology.tex\#L6562}{complete
source passage} and its preceding comparison lemmas were read. The
localization, coherence, geometric calculations and exercises are
expanded here with their actual conditions. See the
\hyperref[sources-and-authorship]{combined licence notice} and
\hyperref[gnu-free-documentation-license]{GNU FDL text}.

\section{Derived pullback and
pushforward}\label{derived-pullback-and-pushforward}

\emph{Written and edited by GPT-6.1 Sol (OpenAI), in Codex at Ultra,
October 2026. Mathematically self-checked by the writing AI. Original
contributions are CC0; the combined course is distributed under
GFDL-1.2-or-later. Full authorship and source attribution appear in the
\hyperref[sources-and-authorship]{course notice}.}

Pullback changes both the space and the coefficient ring. Pushforward
collects sections over inverse images of opens. Their ordinary
adjunction survives passage to derived categories, but proving this
requires two different resolutions. A K-flat model controls the tensor
in pullback, and a K-injective model controls pushforward. The proof
must also work when pushforward of that K-injective model is not itself
K-injective.

The prerequisites are
\hyperref[derived-tensor-products-and-tor-sheaves]{The derived tensor
product and Tor sheaves}, Lemma 1.1, Theorem 1.2, Lemma 2.1 and Theorem
2.2, and
\hyperref[k-injective-resolutions-in-grothendieck-categories]{K-injective
resolutions in Grothendieck categories}, Theorems 4.1 and 5.1 and
Proposition 5.2. Their earlier prerequisites supply the chain-level
adjunction (first lesson, Theorem 4.4), the open-extension adjunction
(Lemma 5.1 there), maps into K-injectives (bounded-derived lesson,
Theorem 3.3), and roof cancellation (common reading, Theorem 4.2). For
the spectral-sequence exercise we use the lower-bound preservation in
Theorem 4.1 of the bounded-derived lesson.

Unless explicitly stated otherwise, all structure sheaves are
commutative and unital, and all derived categories are unbounded. Let
\(f:(X,\mathcal O_X)\to(Y,\mathcal O_Y)\) be a morphism of ringed
spaces. Write

\[f^*G=\mathcal O_X\otimes_{f^{-1}\mathcal O_Y}f^{-1}G.
\tag{0.1}\]

We retain cohomological indexing, direct-sum tensor totalization and the
common reading\textquotesingle s negative cone-projection convention.
These results require no separation, Noetherian, compactness or finite
homological-dimension hypothesis.

The construction sources are the Stacks project authors'
\emph{Cohomology of Sheaves}, ``Derived pullback'', ``Cohomology of
unbounded complexes'' and ``Some properties of K-injective complexes'',
in the
\href{https://github.com/KokunoYumeto/unofficial-stacks-project-ai-drafts/blob/565b10e987aba5969b21145a0833f42d69f96790/cohomology.tex}{AI
Integrated Stacks Project edition}. The proofs use resolutions, roofs
and adjunction to establish the comparison maps with their actual
coefficient hypotheses. Source attribution and the licence for adapted
passages appear in the \hyperref[sources-and-authorship]{course notice}.

\subsection{1. Derived pullback on K-flat
models}\label{derived-pullback-and-pushforward--1-derived-pullback-on-k-flat-models}

\textbf{Theorem 1.1.} Choosing a K-flat resolution \(P\to G\) and
setting

\[Lf^*G=f^*P
\tag{1.1}\]

defines a well-defined exact functor
\(Lf^*:D(\mathcal O_Y)\to D(\mathcal O_X)\). It carries K-flat models to
K-flat models. For composable \(f:X\to Y\), \(g:Y\to Z\), there are
canonical natural isomorphisms

\[Lf^*Lg^*G\cong L(gf)^*G.
\tag{1.2}\]

\textbf{Proof.} Pullback preserves K-flatness by Lemma 2.2 of the K-flat
lesson. It also preserves quasi-isomorphisms \textbf{between} K-flats.
Indeed, at \(x\) their pullback map is

\[\mathcal O_{X,x}\otimes_{\mathcal O_{Y,f(x)}}P_{f(x)}
\longrightarrow
\mathcal O_{X,x}\otimes_{\mathcal O_{Y,f(x)}}P'_{f(x)}.\]

The stalk models are K-flat; Lemma 1.1 of the tensor lesson applies to
their quasi-isomorphism and the possibly nonflat module
\(\mathcal O_{X,x}\). Thus this map is a quasi-isomorphism on every
stalk. Pullback is additive, so it preserves homotopies, shifts and
cones. Lemma 2.1 of the tensor lesson identifies the localization of
K-flat models with the full derived category. Localizing pullback on
those models gives the claimed exact functor and canonical comparisons
for all choices and roofs.

Choose \(P\) K-flat on \(Z\). Then \(g^*P\) is K-flat on \(Y\), so the
iterated derived pullback is computed by \(f^*g^*P\). Ordinary pullback
composition identifies this with \((gf)^*P\). On stalks its
scalar-extension map has the formula

\[\begin{gathered}
s\otimes(t\otimes p)\longmapsto s f^\sharp(t)\otimes p,\\
s\otimes p\longmapsto s\otimes(1\otimes p).
\end{gathered}\]

Balancing proves that these are inverse maps; scalars have degree zero,
so they commute with differentials. The same balanced sheaf maps, using
inverse-image composition, define the identification before taking
stalks. They are natural for chain maps, hence for inverse denominators
and roofs. For three successive pullbacks every association multiplies
the same successive scalar images; associativity proves coherence of
(1.2). \(\square\)

\textbf{Proposition 1.2 (tensor compatibility).} There are canonical
natural isomorphisms

\[\begin{gathered}
Lf^*(A\otimes_{\mathcal O_Y}^{\mathbf L}B)
\\ \cong Lf^*A\otimes_{\mathcal O_X}^{\mathbf L}Lf^*B,\\
F\otimes_{\mathcal O_X}^{\mathbf L}Lf^*G
\\ \cong F\otimes_{f^{-1}\mathcal O_Y}^{\mathbf L}f^{-1}G.
\end{gathered}
\tag{1.3}\]

In the second line \(F\) is an \(\mathcal O_X\)-complex; the right side
retains its \(\mathcal O_X\)-action from \(F\). The isomorphisms respect
the tensor associativity, symmetry, unit and pullback-composition
comparisons.

\textbf{Proof.} For K-flat \(P,R\) on \(Y\), their tensor is K-flat.
Thus the first line is computed by the ordinary isomorphism
\(f^*(P\otimes R)\cong f^*P\otimes f^*R\). On a stalk, with coefficient
map \(R_0\to S\), the inverse-direction map and its inverse are

\[\begin{gathered}
(s\otimes p)\otimes(t\otimes r)\longmapsto st\otimes(p\otimes r),\\
s\otimes(p\otimes r)\longmapsto(s\otimes p)\otimes(1\otimes r).
\end{gathered}\]

Commutativity of the rings and balancing make the formulas well defined.
Their differential signs are \(1,(-1)^{|p|}\) on both sides. Inverse
image and scalar extension commute with direct sums, so the same maps
identify the totalizations. They commute with elementary reassociation,
the Koszul flip and units; on four scalar factors every route is scalar
multiplication in the same order. Their naturality extends to localized
models.

For the second line choose only a K-flat \(P\to G\). Then \(f^{-1}P\) is
K-flat over \(f^{-1}\mathcal O_Y\), by the K-flat pullback criterion
applied with unchanged coefficients, and \(f^*P\) is K-flat over
\(\mathcal O_X\). The one-resolution construction computes both sides by

\[F\otimes_{\mathcal O_X}
(\mathcal O_X\otimes_{f^{-1}\mathcal O_Y}f^{-1}P)
\cong F\otimes_{f^{-1}\mathcal O_Y}f^{-1}P.\]

The map sends \(a\otimes(s\otimes p)\) to \(sa\otimes p\), with inverse
\(a\otimes p\mapsto a\otimes(1\otimes p)\). This proves the asserted
module action and naturality without assuming \(F\) flat. \(\square\)

There is also a natural comparison \(Lf^*G\to f^*G\) for a specified
complex \(G\), induced by \(f^*P\to f^*G\). It need not be an
isomorphism. All comparisons to ordinary tensor commute with (1.3): on
resolving tensors both routes send \((s\otimes p)\otimes(t\otimes r)\)
to \(st\otimes(\epsilon_Pp\otimes\epsilon_Rr)\). For the map
\(a:P\otimes R\to A\otimes B\), Corollary 4.4 of the K-flat lesson
supplies a factorization \(a=cb\) through a K-flat resolution
\(c:N\to A\otimes B\). Applying pullback gives \(f^*c\,f^*b=f^*a\); the
left side represents the route through \(Lf^*(A\otimes B)\) and its
ordinary comparison, and the right side is the same elementary tensor
formula. This proves the complete comparison square, including the
augmentation maps.

\subsection{2. Derived pushforward and the
adjunction}\label{derived-pullback-and-pushforward--2-derived-pushforward-and-the-adjunction}

Choose \(F\to I\) K-injective with injective terms, using Theorem 4.1 of
the K-injective lesson. Theorem 5.1 there, applied to the additive
functor \(f_*\), gives an exact functor

\[\begin{gathered}
Rf_*:D(\mathcal O_X)\longrightarrow D(\mathcal O_Y),\\
Rf_*F=f_*I.
\end{gathered}
\tag{2.1}\]

Two K-injective models are homotopy equivalent, and additive pushforward
preserves these equivalences and their homotopies. Maps between the
models, unique up to homotopy, define the functor on all derived
morphisms. No claim about K-injectivity of \(f_*I\) is needed for this
construction.

\textbf{Lemma 2.1 (a relative Hom comparison).} For \(P\) K-flat on
\(Y\) and \(I\) K-injective on \(X\), localization induces a bijection

\[\begin{gathered}
\operatorname{Hom}_{K(\mathcal O_Y)}(P,f_*I)\\
\cong
\operatorname{Hom}_{D(\mathcal O_Y)}(P,f_*I).
\end{gathered}
\tag{2.2}\]

\textbf{Proof.} Represent a derived map by a roof
\(P\leftarrow Z\to f_*I\). Resolve \(Z\) by a K-flat \(T\to Z\). Its
composite \(s:T\to P\) is a quasi-isomorphism between K-flats. Theorem
1.1 makes \(f^*s\) a quasi-isomorphism. Maps into \(I\) consequently
give a bijection

\[\operatorname{Hom}_{K(\mathcal O_X)}(f^*P,I)
\longrightarrow
\operatorname{Hom}_{K(\mathcal O_X)}(f^*T,I),\]

by Theorem 3.3 of the bounded-derived lesson. The chain-level adjunction
identifies this with precomposition by \(s\) on maps into \(f_*I\). Thus
the roof\textquotesingle s numerator comes from a unique homotopy class
\(P\to f_*I\), proving surjectivity of (2.2).

For injectivity, if such a homotopy class becomes zero after
localization, roof cancellation (common reading, Theorem 4.2) gives a
quasi-isomorphism \(Z\to P\) making its composite zero in the homotopy
category. Resolve \(Z\) by a K-flat \(T\). Precomposition along the
resulting \(T\to P\) is the bijection just proved by the chain
adjunction, so the original class was zero. Applying this to a
difference also proves injectivity for equality of two classes.
\(\square\)

\textbf{Theorem 2.2 (unbounded adjunction).} There is a bifunctorial
bijection

\[\begin{gathered}
\operatorname{Hom}_{D(\mathcal O_X)}(Lf^*G,F)
\\ \cong\operatorname{Hom}_{D(\mathcal O_Y)}(G,Rf_*F).
\end{gathered}
\tag{2.3}\]

It gives \(Lf^*\dashv Rf_*\), with natural unit \(\eta:G\to Rf_*Lf^*G\)
and counit \(\epsilon:Lf^*Rf_*F\to F\).

\textbf{Proof.} Choose \(P\to G\) K-flat and \(F\to I\) K-injective. The
left group is \(\operatorname{Hom}_K(f^*P,I)\), because the target is
K-injective. The ordinary adjunction on complexes, including homotopies,
identifies this with \(\operatorname{Hom}_K(P,f_*I)\). Lemma 2.1
identifies the latter with \(\operatorname{Hom}_D(P,f_*I)\); use
\(P\simeq G\) to obtain the right group.

For naturality in \(G\), every derived map between its K-flat models is
a roof with K-flat middle, by Lemma 2.1 of the tensor lesson. The chain
adjunction and all Hom comparisons commute with its numerator and
denominator, hence with the inverse denominator. For naturality in
\(F\), derived maps between K-injective models are homotopy classes of
chain maps; the adjunction commutes with them. Replacing either model
uses these same comparisons. Thus (2.3) is independent of choices and
bifunctorial.

Define the unit as the transpose of \(1_{Lf^*G}\), and the counit as the
inverse transpose of \(1_{Rf_*F}\). Write \(L=Lf^*\), \(R=Rf_*\), and
\(\alpha\) for the bijection. Naturality implies \(\alpha(u)=R(u)\eta\)
and \(\alpha^{-1}(v)=\epsilon L(v)\). Applying these to
\(\eta=\alpha(1)\) and \(\epsilon=\alpha^{-1}(1)\) gives

\[\begin{gathered}
\epsilon_{LG}L(\eta_G)=1_{LG},\\
R(\epsilon_F)\eta_{RF}=1_{RF}.
\end{gathered}
\tag{2.4}\]

These are the triangle identities, so the asserted maps give the
adjunction. \(\square\)

\textbf{Corollary 2.3 (unchanged coefficients, also noncommutative).} If
\(\mathcal O_X=f^{-1}\mathcal O_Y\), for any unital sheaf of rings,
\(f^{-1}\dashv Rf_*\) on unbounded derived categories.

\textbf{Proof.} Inverse image is exact, and \(f_*\) is its right
adjoint. Proposition 5.2 of the K-injective lesson makes \(f_*I\)
K-injective. Therefore

\[\begin{gathered}
\operatorname{Hom}_D(f^{-1}G,F)
=\operatorname{Hom}_K(f^{-1}G,I)\\
=\operatorname{Hom}_K(G,f_*I)
=\operatorname{Hom}_D(G,Rf_*F).
\end{gathered}\]

Each equality is natural for maps and homotopies and extends to roofs.
This proves the exact-coefficient adjunction without using commutative
K-flat theory in the noncommutative case. \(\square\)

\subsection{3. Composition, Leray and the local
descriptions}\label{derived-pullback-and-pushforward--3-composition-leray-and-the-local-descriptions}

\textbf{Proposition 3.1.} For composable morphisms \(f:X\to Y\) and
\(g:Y\to Z\),

\[Rg_*Rf_*F\cong R(gf)_*F
\tag{3.1}\]

canonically and naturally. In particular,

\[R\Gamma(Y,Rf_*F)\cong R\Gamma(X,F).
\tag{3.2}\]

The second identification is in \(D(\Gamma(Y,\mathcal O_Y))\), with the
right side\textquotesingle s scalars restricted along
\(\Gamma(Y,\mathcal O_Y)\to\Gamma(X,\mathcal O_X)\).

\textbf{Proof.} Successive adjunctions show that \(Rg_*Rf_*\) is right
adjoint to \(Lf^*Lg^*\). Theorem 1.1 identifies the latter with
\(L(gf)^*\), whose right adjoint is \(R(gf)_*\).

For completeness, a right adjoint is unique in the required canonical
sense. If \(L\dashv R_1\) and \(L\dashv R_2\), send the counit
\(LR_1B\to B\) through the second adjunction to a map
\(\beta_B:R_1B\to R_2B\). Reverse the roles to obtain \(\gamma_B\).
Naturality of the adjunctions and their triangle identities show that
the transpose of \(\gamma_B\beta_B\) under the first adjunction is its
counit, so \(\gamma_B\beta_B=1\); likewise \(\beta_B\gamma_B=1\).
Transposing maps after a morphism \(B\to B'\) proves naturality. This is
the unique isomorphism compatible with the given adjunctions, and its
uniqueness also gives compatibility for three successive functors. It
proves (3.1) without assuming \(f_*I\) is K-injective.

Put \(A=\Gamma(Y,\mathcal O_Y)\). The canonical morphism
\(Y\to(\mathrm{pt},A)\) has pushforward \(\Gamma(Y,-)\). Its composite
with \(f\) has pushforward \(\Gamma(X,-)\) with the indicated
\(A\)-action. Apply (3.1) to obtain (3.2). \(\square\)

\textbf{Proposition 3.2 (opens and cohomology sheaves).} Let
\(U\subset X\) and \(V\subset Y\) be open, and for the restriction
formula set \(U=f^{-1}V\), \(f_V:U\to V\). Then:

\begin{enumerate}
\tightlist
\item
  A K-injective \(I\) restricts to a K-injective \(I|_U\), and
  \(H^q(U,F)=H^q(U,F|_U)\).
\item
  The sheafification of \(U\mapsto H^q(U,F)\) is \(H^q(F)\).
\item
  \((Rf_*F)|_V\cong R(f_V)_*(F|_U)\).
\item
  \(R\Gamma(V,Rf_*F)\cong R\Gamma(f^{-1}V,F)\), with scalar restriction
  understood.
\item
  \(H^q(Rf_*F)\) is the sheaf associated to \(V\mapsto H^q(f^{-1}V,F)\).
\end{enumerate}

\textbf{Proof.} Extension by zero \(j_!\) is exact and left adjoint to
restriction, by Lemma 5.1 of the first lesson. Proposition 5.2 of the
K-injective lesson makes its right adjoint, restriction, preserve
K-injectives. Thus \(I|_U\) resolves \(F|_U\), and both definitions of
the group in (1) use the same section complex \(\Gamma(U,I)\).

For (2), this complex gives the presheaf

\[U\longmapsto
\frac{\ker\bigl(I^q(U)\to I^{q+1}(U)\bigr)}
{\operatorname{im}\bigl(I^{q-1}(U)\to I^q(U)\bigr)}.
\tag{3.3}\]

At a point, its filtered colimit is \(H^q(I_x)\), because stalks and
filtered colimits are exact. This is \(H^q(F)_x\). The map sending a
section cycle to its cohomology germ therefore induces an isomorphism
after sheafification, by the first lesson\textquotesingle s stalk
criterion. This argument works for any complex of sheaves; it does not
require exactness of taking sections.

For (3), use the K-injective \(I\) on \(X\) and \(I|_U\) on \(U\). The
ordinary sheaf identity \((f_*I)|_V=(f_V)_*(I|_U)\) follows by
evaluating on opens of \(V\), so it gives the derived identity. Apply
(3.2) to \(f_V\) and use (1) and (3) to obtain (4). Finally \(Rf_*F\) is
represented by \(f_*I\); apply the cohomology-presheaf argument to this
complex. On an open \(V\), its section cohomology is
\(H^q(\Gamma(f^{-1}V,I))=H^q(f^{-1}V,F)\) by (1). Sheafification gives
(5). \(\square\)

These assertions distinguish a cohomology sheaf from the cohomology of
sections on a specified open. Sections of a quotient sheaf need not be
the quotient of sections; sheafification in (2) and (5) is part of the
statement.

\subsection{4. Coefficient comparison and exact
adjoints}\label{derived-pullback-and-pushforward--4-coefficient-comparison-and-exact-adjoints}

Write \(F_{ab}\) for the underlying complex of abelian sheaves.
Forgetting the module action is exact: kernels and images are the same
underlying sheaf kernels and images. It therefore induces
\(D(\mathcal O_X)\to D(\mathbb Z_X)\), even when \(\mathcal O_X\) has
integer torsion.

\textbf{Theorem 4.1 (modules and underlying abelian sheaves).} The
canonical comparison maps

\[\begin{gathered}
R\Gamma(U,F)\longrightarrow R\Gamma(U,F_{ab}),\\
(Rf_*F)_{ab}\longrightarrow Rf_*^{ab}(F_{ab})
\end{gathered}
\tag{4.1}\]

are isomorphisms, respectively in \(D(\mathrm{Ab})\) and
\(D(\mathbb Z_Y)\).

\textbf{Proof.} Choose \(F\to I\) K-injective over \(\mathcal O_X\), and
a K-injective abelian-sheaf resolution \(I_{ab}\to J\). The comparisons
are represented by \(\Gamma(U,I)\to\Gamma(U,J)\) and
\(f_*I_{ab}\to f_*J\). The choice of \(J\) changes them only by the
canonical homotopy comparison between K-injective resolutions, so these
are natural derived maps.

Consider the coefficient morphism
\(h:(X,\mathcal O_X)\to(X,\mathbb Z_X)\), the identity on spaces and the
unit map on rings. Its ordinary pushforward is exactly forgetting the
action, so its derived pushforward is also that exact functor. Also
\(Lh^*\mathbb Z_X=\mathcal O_X\), because \(\mathbb Z_X[0]\) is K-flat
over itself. The adjunction of Theorem 2.2 gives, for every integer
\(n\),

\[\begin{gathered}
\operatorname{Hom}_{D(\mathcal O_X)}(\mathcal O_X,F[n])
\\ \cong\operatorname{Hom}_{D(\mathbb Z_X)}(\mathbb Z_X,F_{ab}[n]).
\end{gathered}
\tag{4.2}\]

The left group is \(H^n(\Gamma(X,I))\): maps from \(\mathcal O_X\) into
\(I[n]\), modulo homotopy, are section cycles modulo section boundaries.
The right group is \(H^n(\Gamma(X,J))\). To check that (4.2) is the
specified comparison, use Lemma 2.1 for \(h\) and the K-flat source
\(\mathbb Z_X\). Its intermediate chain adjunction sends a section of
\(I\) to the same underlying section of \(I_{ab}\); localization
followed by \(I_{ab}\to J\) gives exactly the cohomology map of
\(\Gamma(X,I)\to\Gamma(X,J)\). Hence that map is an isomorphism in every
degree. Repeat on \(U\), using Proposition 3.2 to restrict the models,
to prove the first line of (4.1).

For the second line, take the section cohomology of
\(f_*I_{ab}\to f_*J\) on every open \(V\). It is the first comparison on
\(f^{-1}V\), so it is an isomorphism in every degree. Sheafifying these
cohomology presheaves, as in Proposition 3.2, proves that the map is a
quasi-isomorphism. No preservation of K-injectivity by forgetting
coefficients was assumed. \(\square\)

\textbf{Proposition 4.2 (open extension and flat pushforward).} For an
open inclusion \(j:U\to X\), exact extension by zero is left adjoint to
restriction on derived categories. For a flat morphism \(f\),
pushforward preserves K-injective complexes.

\textbf{Proof.} For \(B\in D(\mathcal O_U)\) and \(F\to I\) K-injective
on \(X\), restriction of \(I\) is K-injective. Exactness of \(j_!\) and
the chain adjunction give

\[\begin{aligned}
\operatorname{Hom}_{D(\mathcal O_X)}(j_!B,F)
&=\operatorname{Hom}_{K(\mathcal O_X)}(j_!B,I)\\
&=\operatorname{Hom}_{K(\mathcal O_U)}(B,I|_U)\\
&=\operatorname{Hom}_{D(\mathcal O_U)}(B,F|_U).
\end{aligned}\]

Naturality is inherited from the chain adjunction and localization, so
this is \(j_!\dashv j^*\).

Flatness of \(f\) means that each \(\mathcal O_{X,x}\) is flat over
\(\mathcal O_{Y,f(x)}\). The pullback stalk formula then makes \(f^*\)
exact. If \(A\) is acyclic on \(Y\), \(f^*A\) is acyclic, and
\(\operatorname{Hom}_K(A,f_*I)=\operatorname{Hom}_K(f^*A,I)=0\). This is
the definition of K-injectivity of \(f_*I\). It applies to every
K-injective \(I\), without a termwise-injectivity assumption.
\(\square\)

\subsection{5. Point maps and constant
coefficients}\label{derived-pullback-and-pushforward--5-point-maps-and-constant-coefficients}

For \(f:X\to\mathrm{pt}\) with target ring \(A\), pushforward is global
sections with its \(A\)-action, so \(Rf_*=R\Gamma(X,-)\) with that
action. The pullback is

\[Lf^*G=\mathcal O_X\otimes_{A_X}^{\mathbf L}G_X,
\tag{5.1}\]

where \(A_X,G_X\) denote inverse-image constant sheaves. Resolve \(G\)
by a K-flat complex over \(A\); its constant inverse image is K-flat,
proving (5.1) by one-variable tensor. If \(\mathcal O_X=A_X\) with the
identity coefficient map, this is the exact constant sheaf functor, also
for noncommutative \(A\) by Corollary 2.3. For the canonical map to
\((\mathrm{pt},\Gamma(X,\mathcal O_X))\), the structure sheaf need not
be constant, and (5.1) gives the correct pullback.

For a point \(x\), give its one-point source the ring
\(\mathcal O_{X,x}\). Pullback is then the exact stalk functor, since
the coefficient map at that point is the identity. If the stalk ring is
local and we instead give the point its residue field \(\kappa(x)\), the
ring map is the quotient, and

\[Li_x^*F=\kappa(x)\otimes_{\mathcal O_{X,x}}^{\mathbf L}F_x.
\tag{5.2}\]

Take stalks of a K-flat resolution to prove this formula. A residue
field is part of this additional local-ring hypothesis, not data
attached to every arbitrary ringed space.

On the analytic curve of the tensor lesson, pull back \(k_a\) to the
residue point at \(a\). Its two-term parameter resolution becomes
\(\mathbb R\xrightarrow{0}\mathbb R\) in degrees \(-1,0\). Thus the
derived pullback has \(\mathbb R\) in both degrees, whereas the ordinary
pullback retains only degree zero. With the stalk-ring structure on the
point, pullback has only \(k_{a,a}\) in degree zero, because taking a
stalk is exact. These are two different morphisms of ringed spaces on
the same underlying point inclusion.

\subsection{6. Exercises with checked
solutions}\label{derived-pullback-and-pushforward--6-exercises-with-checked-solutions}

\textbf{Exercise 1 (easy: an open subspace).} Prove
\(H^p(U,F)=H^p(U,F|_U)\), for an arbitrary open \(U\) and an unbounded
complex. Explain why restriction of the resolution is admissible.

\textbf{Solution.} Resolve \(F\to I\) by a K-injective. For an acyclic
\(A\) on \(U\), exactness of \(j_!\) makes \(j_!A\) acyclic, and
\(\operatorname{Hom}_K(A,I|_U)=\operatorname{Hom}_K(j_!A,I)=0\). Thus
\(I|_U\) is K-injective. Both groups are the cohomology of the same
complex of sections \(\Gamma(U,I)\). No bound on \(I\), nor on the
cohomological dimension of \(U\), was used.

\textbf{Exercise 2 (medium: the adjunction without projective models).}
Let \(P\to G\) be K-flat and \(F\to I\) K-injective. Prove that a roof
\(P\leftarrow Z\to f_*I\) comes from a unique homotopy class
\(f^*P\to I\). Deduce the unbounded adjunction and its naturality,
without assuming that \(P\) is K-projective or \(f_*I\) K-injective.

\textbf{Solution.} Resolve \(Z\) by K-flat \(T\). The denominator
\(s:T\to P\) is a quasi-isomorphism between K-flats, so \(f^*s\) is a
quasi-isomorphism. Transpose the numerator \(a:T\to f_*I\) to
\(f^*T\to I\). Since \(I\) is K-injective, there is a unique homotopy
class \(b:f^*P\to I\) whose composite with \(f^*s\) is that numerator.
Transposing back makes the roof the localization of a map \(P\to f_*I\).

If two maps \(P\to f_*I\) give the same derived class, cancellation
supplies a quasi-isomorphism into \(P\) annihilating their difference.
Resolve its source by a K-flat; the same precomposition bijection forces
the original difference to be zero in the homotopy category. This proves
uniqueness. Maps \(Lf^*G\to F\) are maps \(f^*P\to I\) modulo homotopy,
so we obtain (2.3). The constructions commute with chain maps. All
source derived morphisms can be expressed as roofs between K-flat
models; commuting with their arrows also gives commuting with their
inverse denominators. Target morphisms are homotopy classes between
K-injective models. This proves bifunctoriality.

\textbf{Exercise 3 (medium to hard: Leray and convergence).} For a
module sheaf \(M\), prove

\[\begin{gathered}
E_2^{p,q}=H^p(Y,R^qf_*M)\\
\ \Longrightarrow\ H^{p+q}(X,M),
\\ p,q\ge0,
\end{gathered}
\tag{6.1}\]

including the construction, differential bidegrees and convergence. This
exercise concerns sheaves in degree zero; it imposes no finite dimension
on \(X\) or \(Y\).

\textbf{Solution.} Choose a bounded-below injective resolution of \(M\),
with zero terms below zero. It is K-injective by Lemma 3.2 of the
bounded-derived lesson. Put \(K=f_*I\), so \(K\) represents \(Rf_*M\),
has no negative terms and has \(H^q(K)=R^qf_*M\).

Let \(A_q=R\Gamma(Y,\tau_{\le q}K)\), with \(A_q=0\) for \(q<0\). Good
upper truncations give triangles

\[\begin{gathered}
A_{q-1}\longrightarrow A_q\\
\longrightarrow R\Gamma(Y,H^qK)[-q]
\\ \longrightarrow A_{q-1}[1].
\end{gathered}
\tag{6.2}\]

To check the cofiber before applying \(R\Gamma\), the quotient of
\(\tau_{\le q}K\) by \(\tau_{\le q-1}K\) has the injection
\(\operatorname{im}d^{q-1}\to\ker d^q\) as its two possibly nonzero
terms. Its quotient map to \(H^qK[-q]\) is a quasi-isomorphism. The
common reading\textquotesingle s short-exact-sequence triangle therefore
proves (6.2).

Write \(D_q^n=H^n(A_q)\) and \(E_q^n=H^{n-q}(Y,H^qK)\). The long exact
sequences give maps

\[\begin{gathered}
i:D_q^n\longrightarrow D_{q+1}^n,\\
j:D_q^n\longrightarrow E_q^n,\\
k:E_q^n\longrightarrow D_{q-1}^{n+1}.
\end{gathered}
\tag{6.3}\]

They form an exact couple: \(\ker j=\operatorname{im}i\),
\(\ker k=\operatorname{im}j\) and \(\ker i=\operatorname{im}k\), with
the adjacent indices understood. This is simply exactness at the three
positions of (6.2).

Here is an explicit iteration, including the proof that it gives
successive pages. For \(r\ge2\), define

\[\begin{aligned}
Z_r^{n,q}&=k^{-1}\bigl(\operatorname{im}(i^{r-2})\bigr),\\
B_r^{n,q}&=j\ker(i^{r-2}).
\end{aligned}
\tag{6.4}\]

In the first line the indicated power maps \(D_{q-r+1}^{n+1}\) into
\(D_{q-1}^{n+1}\); in the second it maps \(D_q^n\) into \(D_{q+r-2}^n\).
The image of \(j\) lies in every \(Z_r\), since \(kj=0\), so
\(B_r\subset Z_r\). Set \(E_r^{p,q}=Z_r^{p+q,q}/B_r^{p+q,q}\). For
\(r=2\) the power is the identity, so this is precisely \(H^p(Y,H^qK)\).

Given \(e\in Z_r^{n,q}\), choose \(x\in D_{q-r+1}^{n+1}\) with
\(i^{r-2}x=k e\), and define \(d_r[e]=[jx]\). Changing \(x\) by an
element of the power\textquotesingle s kernel changes \(jx\) by \(B_r\).
Changing \(e\) by an element of \(B_r\) leaves \(ke\) unchanged, so this
is well defined. Its output has total degree \(n+1\) and second index
\(q-r+1\); hence its bidegree is \((r,1-r)\). Since \(kj=0\), its square
is zero.

To verify the next page, suppose \(d_r[e]=0\). Then \(jx=jy\) for a
\(y\) killed by \(i^{r-2}\). Exactness gives \(x-y=iz\), and therefore
\(ke=i^{r-1}z\). Thus \(e\in Z_{r+1}\); the converse follows by taking
\(x=iz\). A differential image \(jx\) has \(i^{r-1}x=i k e=0\), so it
lies in \(B_{r+1}\). Conversely, if \(i^{r-1}x=0\), then
\(i^{r-2}x\in\ker i=\operatorname{im}k\). Choose \(e\) with
\(ke=i^{r-2}x\); it lies in \(Z_r\), and its differential is \(jx\).
Consequently cycles modulo differential images are \(Z_{r+1}/B_{r+1}\),
proving \(E_{r+1}=H(E_r,d_r)\). This constructs the spectral sequence
directly from the triangles.

It remains to identify its limit. The tail \(\tau_{\ge q+1}K\) is
represented in degrees at least \(q+1\). The lower-bound preserving
injective resolution of the bounded-derived lesson shows that its
\(R\Gamma\) also has no cohomology below \(q+1\). The truncation
triangle therefore gives \(H^n(A_q)\cong H^n(R\Gamma(Y,K))\) whenever
\(q\ge n\). Put this stable group \(T^n\), and filter it by

\[F_qT^n=\operatorname{im}(D_q^n\longrightarrow T^n).
\tag{6.5}\]

For \(r>q+1\), the source \(D_{q-r+1}^{n+1}\) in (6.4) is zero, so
\(Z_r=\ker k=\operatorname{im}j\). For sufficiently large \(r\), the
target of the power in the second line is already the stable \(T^n\).
Hence

\[\begin{aligned}
E_\infty^{n-q,q}
&=\operatorname{im}j\big/
j\ker(D_q^n\to T^n)\\
&\cong F_qT^n/F_{q-1}T^n.
\end{aligned}
\tag{6.6}\]

The last identification follows from
\(\ker j=\operatorname{im}(D_{q-1}^n\to D_q^n)\) and the quotient
theorem for the map \(D_q^n\to T^n\). The filtration starts with
\(F_{-1}=0\) and ends with \(F_n=T^n\). All \(E_2^{p,q}\) with \(p<0\)
or \(q<0\) vanish, because sheaf cohomology has no negative degrees and
\(K\) has no negative cohomology. Thus each total degree has a finite
filtration, and each position has finitely many possible incoming and
outgoing differential degrees. This proves convergence. Finally
Proposition 3.1 identifies \(T^n=H^n(R\Gamma(Y,Rf_*M))\) with
\(H^n(X,M)\), proving (6.1) with the stated scalar action.

\textbf{Exercise 4 (hard: the flatness hypothesis for K-injectives).}
Prove preservation of K-injectives by a flat pushforward using its
defining test. Show that the statement can fail for the one-point map
with coefficient map \(\mathbb Z\to\mathbb F_2\).

\textbf{Solution.} An acyclic \(A\) pulls back to an acyclic \(f^*A\)
when \(f\) is flat. For K-injective \(I\), the chain adjunction gives
\(\operatorname{Hom}_K(A,f_*I)=\operatorname{Hom}_K(f^*A,I)=0\). This
proves the preservation assertion.

For the indicated nonflat map, \(I=\mathbb F_2[0]\) is K-injective over
\(\mathbb F_2\): the field is injective by Baer\textquotesingle s
criterion, whose only ideal tests are zero and the field itself, and a
bounded complex of injectives is K-injective. Its pushforward is the
same complex as a \(\mathbb Z\)-module. The complex
\(A=(\mathbb Z\xrightarrow{2}\mathbb Z\to\mathbb F_2)\) in degrees
\(0,1,2\) is acyclic. It has a chain map to \(I\), reduction modulo two
in degree zero and zero elsewhere. A null homotopy in degree zero would
express this nonzero reduction as \(h^1\circ2\), but every homomorphism
\(\mathbb Z\to\mathbb F_2\) kills multiples of two. It cannot be
null-homotopic, so the pushforward is not K-injective. Nevertheless the
general derived adjunction holds by Theorem 2.2.

\textbf{Exercise 5 (hard: torsion coefficients and comparison).} Let
\(\mathcal O_X\) have arbitrary integer torsion. Prove that
\(\Gamma(U,I)\to\Gamma(U,J)\) is a quasi-isomorphism when \(I\) is
K-injective over \(\mathcal O_X\) and \(I_{ab}\to J\) is K-injective
over \(\mathbb Z_X\). Explain why exactness of forgetting alone does not
prove that \(I_{ab}\) is K-injective.

\textbf{Solution.} For the coefficient morphism
\(h:(U,\mathcal O_U)\to(U,\mathbb Z_U)\), the right derived pushforward
is exact forgetting, while \(Lh^*\mathbb Z_U=\mathcal O_U\). Theorem 2.2
identifies maps from \(\mathcal O_U\) into \(I[n]|_U\) with derived maps
from \(\mathbb Z_U\) into \(I_{ab}[n]|_U\). Lemma 2.1 identifies these
with chain maps into the same underlying \(I_{ab}\); passing to \(J|_U\)
computes the latter derived group by a K-injective. On section cycles
the comparison sends a cycle to its image in \(J\), and on boundaries it
sends boundaries to boundaries. Thus the induced map is precisely
\(H^n\Gamma(U,I)\to H^n\Gamma(U,J)\), and is an isomorphism for every
\(n\).

An exact functor preserves acyclic complexes in its domain, but
K-injectivity tests maps \textbf{from every} acyclic complex in the
target category. Those tests are not reduced to acyclic module complexes
merely by forgetting actions. On a point with structure ring
\(\mathbb F_2\), Exercise 4 gives an injective and K-injective module
whose underlying abelian complex fails the K-injective test. The
coefficient-comparison theorem consequently needs the adjunction
argument just supplied.

The next lesson constructs internal derived Hom and its tensor
adjunction. Later, the projection-formula and base-change lesson uses
the unit, counit and coherent composition isomorphisms proved here to
define the comparison maps.

Sources and licensing: the pullback proofs are taken from the Stacks
project, tags 06YJ, 0D5S, 079U and 08DE, in its AI Integrated Stacks
Project edition (GFDL), checked and edited by GPT-6.1 Sol (OpenAI), at
Ultra. The adjunction and composition arguments correspond to tags 079W
and 0D5T; their formal derived-adjoint reference was read in full, and
its omitted functoriality is supplied by Lemma 2.1 and Theorem 2.2. The
open, Leray, higher-image and coefficient proofs are taken from the
Stacks project, tags 08BS, 0D5V, 0BKJ, 08FE, 0D5W, 0D5X, 0D5Y, 08BT and
0D5Z, in that edition (GFDL), checked and edited by the same writing AI.
For 0D5Y the omitted inverse verification is replaced by the full
coefficient adjunction above. Actual passages:
\href{https://github.com/KokunoYumeto/unofficial-stacks-project-ai-drafts/blob/565b10e987aba5969b21145a0833f42d69f96790/cohomology.tex\#L6758}{derived
pullback},
\href{https://github.com/KokunoYumeto/unofficial-stacks-project-ai-drafts/blob/565b10e987aba5969b21145a0833f42d69f96790/cohomology.tex\#L7056}{unbounded
adjunction},
\href{https://github.com/KokunoYumeto/unofficial-stacks-project-ai-drafts/blob/565b10e987aba5969b21145a0833f42d69f96790/cohomology.tex\#L8261}{local
properties}, and
\href{https://github.com/KokunoYumeto/unofficial-stacks-project-ai-drafts/blob/565b10e987aba5969b21145a0833f42d69f96790/derived.tex\#L9838}{deriving
adjoints}. The exact-couple construction and all its convergence steps
are supplied in Exercise 3 rather than cited as a black box. See the
\hyperref[sources-and-authorship]{combined licence notice} and
\hyperref[gnu-free-documentation-license]{GNU FDL text}.

\section{Hom complexes, internal derived Hom and Ext
sheaves}\label{internal-derived-hom-and-ext-sheaves}

\emph{Written and edited by GPT-6.1 Sol (OpenAI), in Codex at Ultra,
October 2026. Mathematically self-checked by the writing AI. Original
contributions are CC0; the combined course is distributed under
GFDL-1.2-or-later. Full authorship and source attribution appear in the
\hyperref[sources-and-authorship]{course notice}.}

The tensor product combines two inputs, whereas internal Hom records
maps from one input to another as an object on the space. Derived
internal Hom must detect extensions as well as ordinary maps. Its
degree-zero global cohomology recovers maps in the derived category; its
other cohomology sheaves record local extensions. These statements
involve derived sections, which cannot generally be replaced by ordinary
sections of a cohomology sheaf.

The prerequisites are
\hyperref[derived-pullback-and-pushforward]{Derived pullback and
pushforward}, Proposition 3.2 and Theorem 4.1, for restriction and local
cohomology; \hyperref[derived-tensor-products-and-tor-sheaves]{The
derived tensor product and Tor sheaves}, Lemma 1.1, Lemma 2.1 and
Theorem 2.2, for K-flat models and their localization; and
\hyperref[k-injective-resolutions-in-grothendieck-categories]{K-injective
resolutions in Grothendieck categories}, Theorems 4.1 and 5.1. We also
use maps into K-injectives and their cone closure from Theorem 3.3 and
Lemma 3.1 of
\hyperref[injective-modules-and-bounded-below-derived-functors]{Injective
modules and bounded-below derived functors}. The common
reading\textquotesingle s Theorems 3.1 and 4.2 supply cone triangles and
roof cancellation.

Throughout, \(\mathcal O=\mathcal O_X\) is a sheaf of commutative unital
rings on an arbitrary topological space. Complexes are cohomologically
indexed and may be unbounded in both directions. Tensor totalization
uses sums, Hom totalization uses products, and a cone triangle ends in
the negative projection. Write \(\mathsf H(A,B)\) for the sheaf Hom
complex, reserving \(\operatorname{Hom}\) for a group of maps.

The construction follows the Stacks project authors' \emph{Cohomology of
Sheaves}, ``Hom complexes'', ``Internal hom'', ``Ext sheaves'' and
``Global derived hom'', in the
\href{https://github.com/KokunoYumeto/unofficial-stacks-project-ai-drafts/blob/565b10e987aba5969b21145a0833f42d69f96790/cohomology.tex}{AI
Integrated Stacks Project edition}. Composition is natural in its outer
variables and dinatural in the repeated middle variable; Section 4
proves the precise identities. Source attribution and the licence for
adapted passages appear in the \hyperref[sources-and-authorship]{course
notice}.

\subsection{1. The sheaf Hom
complex}\label{internal-derived-hom-and-ext-sheaves--1-the-sheaf-hom-complex}

For two module sheaves define the internal Hom by

\[\begin{gathered}
\mathcal Hom(A,B)(U)\\
=\operatorname{Hom}_{\mathcal O_U}(A|_U,B|_U).
\end{gathered}
\tag{1.1}\]

Restriction means restricting a map. Compatible maps on an open cover
glue: for each local section of \(A\), their images agree and glue in
\(B\). The glued assignment respects restriction, addition and scalar
multiplication because these properties can be checked on the cover.
This proves the sheaf condition. Commutativity gives it its
\(\mathcal O\)-module action by multiplying the output.

For complexes set

\[\begin{aligned}
\mathsf H^n(A,B)&=\prod_i\mathcal Hom(A^i,B^{i+n}),\\
(d\phi)^i&=d_B\phi^i\\
&\quad-(-1)^n\phi^{i+1}d_A.
\end{aligned}
\tag{1.2}\]

The product is the sheaf product; on every open it is the product of the
section groups. Applying the differential twice leaves two mixed terms
with coefficients \(-(-1)^n\) and \(-(-1)^{n+1}\), which cancel. The
remaining terms contain \(d_B^2\) or \(d_A^2\). Thus this is a complex.
A degree-zero cycle is a chain map, and a degree-zero boundary is a
null-homotopic map. With the usual shifted differential this proves

\[\begin{gathered}
H^n\Gamma(U,\mathsf H(A,B))\\
=\operatorname{Hom}_{K(\mathcal O_U)}(A|_U,B|_U[n]).
\end{gathered}
\tag{1.3}\]

Restriction to an open commutes with (1.1) and (1.2): their values on
each smaller open are identical. This observation concerns restriction,
not passage to a stalk. An infinite product need not commute with
stalks, and a germ of a module map need not be an arbitrary map of stalk
modules.

\textbf{Lemma 1.1 (chain identities).} Hom complexes respect homotopies.
There is a natural currying isomorphism

\[\begin{gathered}
\mathsf H(K,\mathsf H(L,M))\\
\cong\mathsf H(K\otimes L,M).
\end{gathered}
\tag{1.4}\]

Evaluation and composition are chain maps, given on homogeneous elements
by

\[\begin{aligned}
\phi\otimes l&\longmapsto\phi(l),\\
\phi\otimes\psi&\longmapsto\phi\psi.
\end{aligned}
\tag{1.5}\]

\textbf{Proof.} If \(h:B\to B'[-1]\) is a homotopy between two target
maps, the induced Hom homotopy is \(\phi\mapsto h\phi\). If
\(h:A'\to A[-1]\) is a homotopy between two source maps, it is
\(\phi\mapsto(-1)^n\phi h\) on degree \(n\). Substitution into (1.2)
gives respectively postcomposition and precomposition with \(dh+hd\);
the mixed terms cancel. Hence both variables descend to homotopy
categories.

Currying sends \(\phi\) to the map \(k\otimes l\mapsto\phi(k)(l)\), with
no additional sign. The tensor relation is exactly the linearity in each
variable. In degree \(n\), both sides are the product over pairs of
degrees \((i,j)\) of \(\mathcal Hom(K^i\otimes L^j,M^{i+j+n})\). This
uses the universal property that Hom out of a sum is a product, followed
by reindexing a product. It uses no exactness assertion about products.
For homogeneous \(k\), their common differential evaluated at \(k,l\) is

\[\begin{aligned}
d_M\phi(k)(l)&-(-1)^n\phi(d_Kk)(l)\\
&-(-1)^{n+|k|}\phi(k)(d_Ll).
\end{aligned}\]

Finally \(d(\phi(l))=(d\phi)(l)+(-1)^{|\phi|}\phi(dl)\) and
\(d(\phi\psi)=(d\phi)\psi+(-1)^{|\phi|}\phi(d\psi)\). These are
precisely the tensor differentials, proving (1.5). All constructions are
sectionwise and therefore commute with restriction. \(\square\)

We also need the cone signs. For \(g:B\to B'\), the two finite blocks
give \(\mathsf H(A,C(g))=C(\mathsf H(A,g))\); the identity
identification preserves the differential and negative projection. For
\(f:K\to L\), put \(h=\mathsf H(f,I)\). In degree \(n\), an element of
\(\mathsf H(C(f),I)\) is a pair \((b,a)\), where \(b\) has Hom degree
\(n\) on \(L\) and \(a\) has degree \(n-1\) on \(K\). The isomorphism to
\(C(h)[-1]\) is

\[(b,a)\longmapsto\bigl((-1)^{n+1}a,b\bigr).
\tag{1.6}\]

Indeed the source differential is \((db,da-(-1)^nbf)\), and the shifted
cone differential on \((x,y)\) is \((-dx-hy,dy)\). Formula (1.6)
intertwines them. The source-shift identification
\(\mathsf H(K[1],I)\cong\mathsf H(K,I)[-1]\) multiplies degree \(n\) by
\((-1)^n\). Together these formulas turn a source cone triangle into the
reversed distinguished Hom triangle, with its connecting map determined
by the negative original projection. This verifies exactness in the
contravariant variable, including its shift sign.

\subsection{2. Constructing internal derived
Hom}\label{internal-derived-hom-and-ext-sheaves--2-constructing-internal-derived-hom}

\textbf{Lemma 2.1.} Let \(I\) be K-injective. A quasi-isomorphism
\(s:A\to B\) induces a quasi-isomorphism
\(\mathsf H(B,I)\to\mathsf H(A,I)\).

\textbf{Proof.} On every open \(U\), the restriction of \(I\) is
K-injective, by the preceding lesson\textquotesingle s Proposition 3.2.
Maps into it are unchanged by localization. Thus precomposition by
\(s|_U\) is a bijection on
\(\operatorname{Hom}_{K(\mathcal O_U)}(-,I|_U[n])\) for every integer
\(n\). Formula (1.3) gives a quasi-isomorphism of the section Hom
complexes on every open. The cohomology sheaf is the sheafification of
section-complex cohomology, as proved in that proposition; sheafifying
these isomorphisms proves the claim. Applying the argument to every
shift is essential: degree zero alone would not prove a
quasi-isomorphism. \(\square\)

\textbf{Theorem 2.2.} For any \(L,M\), choose a K-injective resolution
\(M\to I\). Then

\[\begin{aligned}
{}[L,M]&:=R\mathcal Hom(L,M)\\
&=\mathsf H(L,I).
\end{aligned}
\tag{2.1}\]

defines an object of \(D(\mathcal O)\) and a canonical bifunctor,
contravariant in \(L\), covariant in \(M\), exact in each variable. It
commutes with restriction to opens.

\textbf{Proof.} Resolutions exist for every unbounded complex by the
earlier K-injective theorem. The map from \(M\) to another resolution
determines a unique homotopy class \(I\to I'\) under the maps from
\(M\): maps into a K-injective are the same in the homotopy and derived
categories. Such a comparison is a quasi-isomorphism between
K-injectives, hence a homotopy equivalence. Lemma 1.1 carries its two
homotopy inverses to inverse Hom maps. Comparisons are unique in the
homotopy category, so their composites obey the required identity and
cocycle laws.

In the source variable Lemma 2.1 inverts every quasi-isomorphism. A roof
acts by its numerator and the inverse of its denominator; common
refinements give the same action because Hom respects homotopies. In the
target variable represent a derived map by the unique homotopy class
between K-injective models. Composition of these classes, and
commutation with source precomposition, prove bifunctoriality and
natural independence of resolutions.

Source cone triangles give reversed Hom triangles by (1.6). A target map
between K-injectives has K-injective cone by the earlier cone-closure
lemma, so it supplies a model for the entire target triangle. The target
cone identity proves exactness there. Shift identifications are those
checked after (1.6). Finally restricting \(M\to I\) gives a K-injective
resolution on each open, and the sectionwise identity for Hom gives the
asserted restriction isomorphism, natural in all roofs. \(\square\)

\textbf{Lemma 2.3.} If \(P\) is K-flat and \(I\) K-injective, then
\(\mathsf H(P,I)\) is K-injective.

\textbf{Proof.} For an acyclic complex \(A\), Lemma 1.1 and degree-zero
homotopy classes give

\[\operatorname{Hom}_K(A,\mathsf H(P,I))
=\operatorname{Hom}_K(A\otimes P,I)=0.\]

The tensor is acyclic by K-flatness, and the last group vanishes by
K-injectivity. This is the defining test for K-injectivity. No
injectivity of the individual terms of \(\mathsf H(P,I)\) is required.
\(\square\)

\subsection{3. Adjunction, local maps and
Ext}\label{internal-derived-hom-and-ext-sheaves--3-adjunction-local-maps-and-ext}

\textbf{Theorem 3.1.} There are natural isomorphisms

\[\begin{gathered}
{}[K,[L,M]]\cong[K\otimes^{\mathbf L}L,M],\\
\operatorname{Hom}_D(K,[L,M])\\
\cong\operatorname{Hom}_D(K\otimes^{\mathbf L}L,M).
\end{gathered}
\tag{3.1}\]

For every open \(U\) and integer \(n\),

\[\begin{gathered}
H^n(U,[L,M])\\
\cong\operatorname{Hom}_{D(\mathcal O_U)}(L|_U,M|_U[n]).
\end{gathered}
\tag{3.2}\]

\textbf{Proof.} Resolve \(L\) by a K-flat \(P\) and \(M\) by a
K-injective \(I\). Lemma 2.1 permits the model \(\mathsf H(P,I)\) for
\([L,M]\); Lemma 2.3 makes this model K-injective. Therefore (1.4) is an
actual complex isomorphism computing both sides of the first line of
(3.1). The complex \(K\otimes P\) computes the derived tensor even when
\(K\) is not K-flat, by Lemma 1.1 of the tensor lesson. Taking homotopy
classes and using maps into K-injectives gives the second line.

Naturality holds first for chain maps by currying. For source maps
represented by roofs, Lemma 2.1 inverts the source denominators, and
tensor invariance inverts the denominators on the other side. Comparison
maps between K-flat models can be represented in their localized
homotopy category, by the tensor lesson\textquotesingle s Lemma 2.1.
Target comparisons are homotopy classes between K-injectives. Thus the
same commuting chain squares and their inverse denominator squares prove
naturality for every derived morphism, not only actual chain maps on
fixed resolutions.

For (3.2), restrict these models to \(U\). Restriction preserves
K-flatness and K-injectivity. The K-injective model
\(\mathsf H(P,I)|_U\) computes derived sections, so (1.3) identifies its
cohomology with the stated group of derived maps. \(\square\)

Define the Ext sheaves and global derived Hom by

\[\begin{gathered}
\mathcal Ext^n(L,M)=\mathcal H^n([L,M]),\\
R\operatorname{Hom}_X(L,M)\\
=R\Gamma(X,[L,M]).
\end{gathered}
\tag{3.3}\]

The global object is a complex over \(\Gamma(X,\mathcal O)\). Its
underlying abelian cohomology can equally be computed by the preceding
lesson\textquotesingle s coefficient comparison. Formula (3.2)
identifies its \(n\)-th cohomology with
\(\operatorname{Hom}_D(L,M[n])\).

In particular, a derived morphism is represented locally by a
degree-zero cycle in a Hom complex into a K-injective target, modulo its
degree-zero boundaries. Composition in Proposition 4.1 below induces the
usual composition of these classes. For a triangle in either input,
Theorem 2.2 gives a triangle of internal Hom objects, and derived
sections give the resulting long exact sequence of groups of maps. The
source-variable connecting map includes the shift sign in (1.6); the
target-variable map includes the negative cone projection. Thus the
local extension groups and their boundary operations come from the same
proved triangle conventions as the global ones. No choice of an
unrelated connecting-homomorphism sign is needed.

The Ext sheaf is the sheafification of the presheaf sending \(U\) to the
right side of (3.2). To see this directly, use the model
\(\mathsf H(P,I)\); sheafifying its section-complex cohomology gives its
cohomology sheaf. This does not identify ordinary sections of
\(\mathcal Ext^n\) with that derived-map group on an arbitrary open.
Sheafification may change sections.

There is also a convenient formula with an arbitrary representative
\(L\): the complex \(\Gamma(X,\mathsf H(L,I))\) computes global derived
Hom. Although \(\mathsf H(L,I)\) need not be K-injective, its comparison
to \(\mathsf H(P,I)\) is a quasi-isomorphism on sections by (1.3) and
the K-injective map test. The latter is a K-injective sheaf complex, so
its sections compute the derived sections. For module sheaves \(F,G\) in
degree zero choose a bounded-below injective resolution
\(G\to I^{\geq0}\). Its Hom complex from \(F\) has no negative degrees
and degree-zero kernel \(\mathcal Hom(F,G)\). Hence
\(\mathcal Ext^n(F,G)=0\) for \(n<0\), and
\(\mathcal Ext^0(F,G)=\mathcal Hom(F,G)\).

\subsection{4. Evaluation, composition and their
variances}\label{internal-derived-hom-and-ext-sheaves--4-evaluation-composition-and-their-variances}

The adjunction gives a safe way to construct canonical maps without
asserting that an arbitrary ordinary tensor of Hom complexes computes a
derived tensor. Its counit is evaluation

\[e_{L,M}:[L,M]\otimes^{\mathbf L}L\longrightarrow M.
\tag{4.1}\]

It is the transpose of the identity of \([L,M]\). On models it is (1.5),
after any K-flat replacement needed for its tensor factors. Currying and
its naturality show that this construction is independent of those
replacements.

\textbf{Proposition 4.1 (composition).} There is a canonical
associative, unital composition map

\[\begin{gathered}
c_{K,L,M}:[L,M]\otimes^{\mathbf L}[K,L]\\
\longrightarrow[K,M].
\end{gathered}
\tag{4.2}\]

It is natural contravariantly in \(K\), covariantly in \(M\), and
dinatural in \(L\): for \(u:L\to L'\), the two maps from
\([L',M]\otimes^{\mathbf L}[K,L]\) to \([K,M]\) satisfy

\[\begin{aligned}
c_{K,L,M}\circ([u,M]\otimes1)&\\
\quad=c_{K,L',M}\circ(1\otimes[K,u]).&
\end{aligned}
\tag{4.3}\]

\textbf{Proof.} Put \(A=[L,M]\), \(B=[K,L]\). Define (4.2) to be the
transpose of successive evaluations

\[\begin{aligned}
(A\otimes^{\mathbf L}B)\otimes^{\mathbf L}K
&\longrightarrow A\otimes^{\mathbf L}L\\
&\longrightarrow M.
\end{aligned}
\tag{4.4}\]

Here and below association maps are the coherent tensor isomorphisms
proved in the tensor lesson. Evaluation is natural in its target. For a
source map \(u\), its defining adjunction gives the identity that
evaluating \([u,M](a)\) on \(l\) equals evaluating \(a\) on \(u(l)\).
This is an identity of morphisms in the derived category: transpose
either side by (3.1), and both are precomposition of the same identity
map by \(u\). It therefore holds for roofs as well as chain maps.

Postcomposing (4.4) by a map \(M\to M'\) proves target naturality.
Precomposing its last tensor factor by a map \(K'\to K\), and applying
the source evaluation identity, proves source naturality. For (4.3),
transpose both sides and apply that identity once. Both become
successive evaluation with \(u\) inserted between the two evaluations.
Symbolically both give \(a(u(b(k)))\). The adjunction is a bijection of
map groups, so equality of the transposes proves equality of the maps.

For four objects, either order of composition transposes to the same
three successive evaluations. The tensor associator satisfies the
pentagon, so this proves associativity. The identity of \(L\)
corresponds by (3.2) to a morphism \(\mathbf1=\mathcal O\to[L,L]\).
Evaluation with this morphism is the identity on \(L\), by its
transpose; inserting it on either side of (4.4) proves both unit laws.
These arguments prove all assertions without changing resolutions in the
middle of an unverified diagram. \(\square\)

The word \emph{dinatural} records the opposing variances of the two
occurrences of \(L\). A map \(L\to L'\) induces \([L',M]\to[L,M]\) but
\([K,L]\to[K,L']\). Thus the source of composition is not a covariant or
contravariant one-variable functor of \(L\). Formula (4.3) is the
well-typed compatibility statement. It does not require \(u\) to be
invertible.

\textbf{Proposition 4.2 (other canonical maps).} Put
\(L^\vee=[L,\mathbf1]\). There are canonical maps

\[\begin{aligned}
K&\longrightarrow[L,K\otimes^{\mathbf L}L],\\
K\otimes^{\mathbf L}[M,L]
&\longrightarrow[M,K\otimes^{\mathbf L}L],\\
M\otimes^{\mathbf L}L^\vee&\longrightarrow[L,M],\\
[L,M]\otimes^{\mathbf L}K&\longrightarrow[[K,L],M].
\end{aligned}
\tag{4.5}\]

The first is the adjunction unit, natural in \(K\) and dinatural in
\(L\). The second is natural covariantly in \(K,L\) and contravariantly
in \(M\). The third is natural covariantly in \(M\), contravariantly in
\(L\). The last is natural covariantly in \(K,M\), contravariantly in
\(L\).

\textbf{Proof.} Transpose the identity of \(K\otimes^{\mathbf L}L\) to
obtain the first. If \(\eta_{K,L}\) denotes it and \(u:L\to L'\), both
sides of

\[\begin{gathered}
{}[u,K\otimes^{\mathbf L}L']\eta_{K,L'}\\
=[L,K\otimes^{\mathbf L}u]\eta_{K,L}.
\end{gathered}
\tag{4.6}\]

transpose to \(1_K\otimes u\). This proves its middle-variable
dinaturality; ordinary target naturality proves its \(K\) naturality.

For the second map, transpose \(1_K\otimes e_{M,L}\). For the third,
transpose \(1_M\otimes e_{L,\mathbf1}\), followed by the unit
isomorphism. For the last, its transpose has domain
\(([L,M]\otimes^{\mathbf L}K)\otimes^{\mathbf L}[K,L]\). Move \(K\) past
\([K,L]\) using the Koszul symmetry, evaluate \([K,L]\) on \(K\), and
then evaluate \([L,M]\) on the result. Each stated naturality follows
from the evaluation identities already proved, followed by the
naturality of transposition. On homogeneous chain representatives of
degrees \(p,q,r\), this last construction is

\[(\phi\otimes k)(h)=(-1)^{qr}\phi(h(k)).
\tag{4.7}\]

The sign comes from moving \(k\) past \(h\), not past \(\phi\). Lemma
1.1 and the tensor symmetry prove it is a chain map whenever the chosen
models compute these derived constructions. \(\square\)

These are maps, with no general claim of invertibility. The
perfect-complex lesson proves that the third is an isomorphism when
\(L\) is perfect. The final exercise below shows a concrete failure
without this hypothesis. The unit in (4.5), just like composition, has a
repeated variable with opposite variances; writing its compatibility as
(4.6) prevents a second misleading blanket functoriality claim.

\subsection{5. Worked
examples}\label{internal-derived-hom-and-ext-sheaves--5-worked-examples}

First, \([\mathcal O,M]\cong M\). Evaluation at \(1\) identifies
\(\mathsf H(\mathcal O,I)\) with \(I\), including its differential. The
K-injective resolution identifies this with \(M\). The corresponding
global statement is
\(R\operatorname{Hom}_X(\mathcal O,M)=R\Gamma(X,M)\), which may have
higher cohomology even when \(M\) is a sheaf.

Let \(j:U\hookrightarrow X\) be open. For any K-injective resolution
\(F\to I\), the extension-by-zero adjunction on every open \(V\) gives

\[\mathsf H(j_!\mathcal O_U,I)(V)
=\Gamma(V\cap U,I).\]

These identities respect differentials and restrictions. The right side
is \(j_*(I|_U)(V)\), and \(I|_U\) is K-injective. Consequently

\[[j_!\mathcal O_U,F]\cong Rj_*(F|_U).
\tag{5.1}\]

The ordinary Hom sheaf in degree zero is \(j_*F|_U\), not extension by
zero. It can have a nonzero boundary stalk. Higher Ext sheaves describe
the higher local cohomology of intersections with \(U\), through the
local pushforward formula of the preceding lesson.

For a second example take \(X=\mathbf C\) with holomorphic functions and
let \(k_a\) be the skyscraper residue module at \(a\). The tensor
lesson\textquotesingle s parameter resolution is

\[P_a=(\mathcal O\xrightarrow{z-a}\mathcal O).
\tag{5.2}\]

It lies in degrees minus one and zero. This finite free complex can be
used in the source of derived Hom. Indeed \(\mathsf H(\mathcal O,-)\) is
the identity, so a finite filtration by its two terms shows
\(\mathsf H(P_a,-)\) preserves quasi-isomorphisms. Thus comparison with
any K-injective resolution of \(k_b\) proves that \(\mathsf H(P_a,k_b)\)
computes \([k_a,k_b]\). It has \(k_b\) in degrees zero and one, with
differential multiplication by \(-(b-a)\). If \(b\ne a\), this is
invertible and the Hom object is zero. If \(b=a\), the differential is
zero and

\[[k_a,k_a]\cong k_a\oplus k_a[-1].
\tag{5.3}\]

In particular, the sheaf Ext in degree one is \(k_a\); all other
positive Ext sheaves vanish. These objects describe the derived Hom of
modules over the holomorphic structure sheaf, not Hom of arbitrary
skyscraper abelian sheaves. Changing the coefficient ring changes the
calculation.

Here is an explicit reason for the stalk warning after (1.3). Let
\(X=\{0\}\cup\{1,1/2,1/3,\ldots\}\) with its subspace topology, and use
the constant ring sheaf \(k_X\) for a field. Set
\(A=\bigoplus_{j\geq1}k_X\) and \(B=k_X\). The sum\textquotesingle s
universal property gives \(\mathcal Hom(A,B)=\prod_{j\geq1}k_X\):
specifying a map on an open means specifying the image of the unit in
each summand. For each \(j\), let \(f_j\) be the locally constant
function which is one at \(1/j\) and zero at every other point. It is
locally constant at zero too, since a sufficiently small neighborhood
excludes that particular point. Hence \((f_j)_j\) is a section of this
product on all of \(X\).

Every individual \(f_j\) has zero germ at zero. The induced map of stalk
modules \(A_0=\bigoplus_j k\to B_0=k\) is therefore zero, since an
element of the direct sum uses only finitely many coordinates. But the
germ of the section \((f_j)_j\) of the Hom sheaf is nonzero: every
neighborhood of zero contains some \(1/j\), where its \(j\)-th
coordinate is one. Thus the canonical comparison
\(\mathcal Hom(A,B)_0\to\operatorname{Hom}_k(A_0,B_0)\) is not even
injective. The issue already appears for ordinary Hom, before any
derived construction. Restriction to a fixed open is valid because it
preserves the whole sectionwise product; moving to a germ cannot
uniformly discard the infinitely many shrinking exceptional points.

\subsection{6. Exercises with checked
solutions}\label{internal-derived-hom-and-ext-sheaves--6-exercises-with-checked-solutions}

\textbf{Exercise 1 (easy: the unit).} Prove \([\mathcal O,M]=M\), and
determine the unit \(K\to[\mathcal O,K\otimes^{\mathbf L}\mathcal O]\)
under the tensor and Hom unit identifications.

\textbf{Solution.} A map \(\mathcal O|_V\to I^n|_V\) is determined by
the image of \(1\), with every section \(a\) sent to \(a\) times that
image. These identifications commute with restriction and the
differential, so \(\mathsf H(\mathcal O,I)=I\). For the unit, the
adjunction transpose of the identity of
\(K\otimes^{\mathbf L}\mathcal O\) is the map in question. Evaluation at
\(1\) and the tensor unit both send that identity back to \(1_K\). Hence
the unit becomes the identity of \(K\), in every degree and for every
unbounded complex.

\textbf{Exercise 2 (medium: currying without boundedness).} Explain why
(3.1) holds without a finiteness condition on any complex. In
particular, identify the step where products are needed and distinguish
it from an assertion that products of exact sequences are exact.

\textbf{Solution.} Choose a K-flat \(P\to L\) and K-injective
\(M\to I\). In Hom degree \(n\), currying identifies the product indexed
by \(i,j\) of maps \(K^i\otimes P^j\to I^{i+j+n}\) on both sides. On the
tensor side the source in each degree is a direct sum, so specifying a
map means independently specifying a map on every summand. Rearranging
the resulting products gives exactly the Hom product on the iterated
side. Formula (1.4)\textquotesingle s differential calculation proves
the identification is a complex isomorphism.

K-flatness makes \(K\otimes P\) a model for the derived tensor; Lemma
2.3 makes \(\mathsf H(P,I)\) an admissible K-injective target.
Consequently both Hom complexes compute the stated derived Hom objects.
Quasi-isomorphisms in the first variable are handled on every open by
Lemma 2.1, and target comparisons by homotopy equivalences. No step
applies a product to an arbitrary exact sequence of sheaves. The
product\textquotesingle s universal property and the two resolution
properties are sufficient.

\textbf{Exercise 3 (medium: a middle-variable example).} Over a field,
take \(K=M=k\), \(L=k^2\), \(L'=k\), and \(u(x,y)=x\), all in degree
zero. Verify (4.3) and explain why it is the appropriate replacement for
one-variable naturality.

\textbf{Solution.} Write \(a:L'\to M\) as multiplication by \(\alpha\),
and \(b:K\to L\) as \(t\mapsto(\beta t,\gamma t)\). Precomposing \(a\)
by \(u\) gives the row \((\alpha,0)\). Its product with the column
\((\beta,\gamma)\) is multiplication by \(\alpha\beta\). Postcomposing
\(b\) by \(u\) gives multiplication by \(\beta\), so composing with
\(a\) gives the same answer. Bilinearity proves equality on the tensor
product. All these vector spaces are flat and injective, so the
calculation is also the derived one.

The map on the first factor goes from \(\operatorname{Hom}(L',M)\) to
\(\operatorname{Hom}(L,M)\), whereas the map on the second goes from
\(\operatorname{Hom}(K,L)\) to \(\operatorname{Hom}(K,L')\). Their
directions oppose each other. Since \(u\) is a noninvertible projection,
reversing one direction cannot be done by an inverse. Dinaturality
compares the two composites on their common mixed source and does
exactly what the calculation verifies.

\textbf{Exercise 4 (hard: the K-injective target test).} Let \(P\) be
K-flat and \(I\) K-injective. Prove that \(\mathsf H(P,I)\) is
K-injective using an explicit correspondence between maps and
homotopies. Explain why termwise flatness alone does not justify this
argument.

\textbf{Solution.} A chain map \(A\to\mathsf H(P,I)\) curries to a chain
map \(A\otimes P\to I\) by (1.4). A homotopy is a degree-minus-one
element whose Hom differential is the difference of the two maps. Since
(1.4) is an isomorphism of complexes, it carries exactly such homotopies
in either direction. Thus it gives a bijection on homotopy classes. If
\(A\) is acyclic, K-flatness makes \(A\otimes P\) acyclic, and
K-injectivity of \(I\) says every resulting chain map is null-homotopic.
Uncurrying that null-homotopy proves the defining condition for
\(\mathsf H(P,I)\).

The flat-resolution lesson, Section 5, gives an unbounded termwise-free
complex over the dual numbers whose tensor with an acyclic complex is
not acyclic. Termwise flatness therefore cannot supply the crucial
implication for an unbounded \(P\). Bounded-above termwise-flat
complexes do supply it, by that lesson\textquotesingle s Lemma 2.3; the
distinction is between that proved bounded case and the general K-flat
condition.

\textbf{Exercise 5 (hard: a dual map that fails).} Let
\(R=k[\varepsilon]/(\varepsilon^2)\) on a one-point space and
\(N=R/(\varepsilon)\). Compute \(N^\vee\), compare
\(N\otimes_R^{\mathbf L}N^\vee\) with \([N,N]\), and show the canonical
degree-zero map is zero.

\textbf{Solution.} The only ideals of \(R\) are \(0,(\varepsilon),R\). A
map from \((\varepsilon)\) to \(R\) sends \(\varepsilon\) to an element
killed by \(\varepsilon\), necessarily \(c\varepsilon\), and extends by
multiplication by \(c\). Maps from the other two ideals extend
immediately. Baer\textquotesingle s criterion, proved in the
bounded-derived lesson\textquotesingle s Lemma 2.1, makes \(R\)
injective. It is therefore K-injective concentrated in degree zero.
Consequently

\[N^\vee=\operatorname{Hom}_R(N,R)
=(\varepsilon)\cong N,\]

concentrated in degree zero. A free resolution of \(N\) has one copy of
\(R\) in every nonpositive degree, all negative differentials
multiplication by \(\varepsilon\), and the quotient augmentation in
degree zero. It is exact in negative degrees because the image and
kernel of multiplication by \(\varepsilon\) both equal
\((\varepsilon)\).

Tensoring this bounded-above free resolution with \(N^\vee\cong N\)
makes all differentials zero. Thus the left object has cohomology \(N\)
in every nonpositive degree and zero in positive degrees. Hom from that
resolution to \(N\) likewise has zero differential but lies in
nonnegative degrees. It computes \([N,N]\): a map \(a:P\to C\) to any
acyclic complex has a null-homotopy constructed downwards. In degree
zero, \(a^0\) lands in the cycles of \(C^0\), and projectivity lifts it
through \(d:C^{-1}\to Z^0C\) to \(h^0\). Having chosen \(h^{n+1}\), the
map \(a^n-h^{n+1}d_P\) lands in \(Z^nC\) by the chain-map equation and
the preceding homotopy equation. Lift it to \(h^n:P^n\to C^{n-1}\). This
constructs \(a=dh+hd\) in every nonpositive degree, with no infinite
summation. Applying the same argument to every shift proves
\(\operatorname{Hom}^\bullet(P,C)\) acyclic. Cone compatibility then
proves invariance under every target quasi-isomorphism, including
\(N\to I\). The right object therefore has cohomology \(N\) in every
nonnegative degree, and zero in negative degrees.

Finally degree-zero evaluation sends \(m\otimes\phi\) to the map
\(n\mapsto\phi(n)m\). Every \(\phi(n)\) lies in \((\varepsilon)\), which
acts by zero on \(N\). Hence this induced map on degree-zero cohomology
is zero. The canonical dual map is far from an isomorphism in this
example. Finite duality in the perfect-complex lesson will require a
hypothesis this residue module does not satisfy.

\section{Sections with support and the localization
triangle}\label{sections-with-support-and-the-localization-triangle}

\emph{Written and edited by GPT-6.1 Sol (OpenAI), in Codex at Ultra,
October 2026. Mathematically self-checked by the writing AI. Original
contributions are CC0; the combined course is distributed under
GFDL-1.2-or-later. Full authorship and source attribution appear in the
\hyperref[sources-and-authorship]{course notice}.}

A section can vanish away from a closed subset while carrying
information on that subset. Derived sections with support also detect
classes which cannot be represented by a degree-zero section.
Localization compares this information with what remains on the open
complement. The point of the triangle is to include the boundary map: it
measures the obstruction to extending a class across the missing subset.

The prerequisites are
\hyperref[sheaves-of-modules-on-a-ringed-space]{Sheaves of modules on a
ringed space}, Theorems 2.1 and 4.4;
\hyperref[injective-modules-and-bounded-below-derived-functors]{Injective
modules and bounded-below derived functors}, Lemma 1.1, Proposition 1.3
and Theorem 4.1(4), for injective flasqueness and acyclic resolutions;
\hyperref[k-injective-resolutions-in-grothendieck-categories]{K-injective
resolutions in Grothendieck categories}, Theorems 4.1 and 5.1; and
\hyperref[derived-pullback-and-pushforward]{Derived pullback and
pushforward}, Proposition 3.2 and Theorem 4.1. The cone construction and
exact-sequence triangle are Theorem 3.1 and Proposition 5.2 of the
common reading. For cup products we use the tensor
lesson\textquotesingle s Theorem 2.2 and the internal-Hom
lesson\textquotesingle s Theorem 3.1.

Sections 1--3 work for any sheaf of unital rings, including
noncommutative rings, and unbounded complexes of left modules. Section
5\textquotesingle s tensor and cup products assume commutativity. For a
closed inclusion \(i:Z\hookrightarrow X\), use the unchanged coefficient
sheaf \(\mathcal O_Z=i^{-1}\mathcal O_X\); this is distinct from taking
a quotient coefficient ring on a closed geometric subspace. Put
\(j:U=X\setminus Z\hookrightarrow X\).

The construction follows the Stacks project authors' \emph{Cohomology of
Sheaves}, ``Cohomology with support in a closed subset, II'', Tags
0G6Z--0G79, and \emph{Sheaves of Modules}, in the
\href{https://github.com/KokunoYumeto/unofficial-stacks-project-ai-drafts/tree/565b10e987aba5969b21145a0833f42d69f96790}{AI
Integrated Stacks Project edition}. We develop the support adjunction,
localization triangle and canonical product and pullback maps. The
Euclidean example includes the singular-cochain comparison it uses.
Source attribution and the licence for adapted passages appear in the
\hyperref[sources-and-authorship]{course notice}.

\subsection{1. Closed supports and their
adjunction}\label{sections-with-support-and-the-localization-triangle--1-closed-supports-and-their-adjunction}

The support of a section \(s\in F(V)\) is the set of points where its
germ is nonzero. It is closed in \(V\): a zero germ means the section
vanishes on a neighborhood. Define the ambient support subsheaf by

\[\mathcal S_Z^0F=\ker(F\longrightarrow j_*F|_U).
\tag{1.1}\]

Its sections are exactly those supported in \(Z\cap V\). It is a module
subsheaf, since scalar multiplication cannot create a nonzero germ where
the section vanished. Kernels show that this additive functor is left
exact. Its stalks vanish off \(Z\).

A sheaf with zero stalks off a closed set is canonically the direct
image of its restriction to that set. Indeed \(i_*G\) has stalk \(G_z\)
at \(z\in Z\) and zero off \(Z\): neighborhoods intersected with \(Z\)
are cofinal neighborhoods in \(Z\), and a point outside \(Z\) has a
neighborhood disjoint from it. Thus the unit \(F\to i_*i^{-1}F\) is an
isomorphism for such \(F\). The same stalk description proves that
\(i_*\) is exact, \(i^{-1}i_*\) is the identity, and \(i_*\) is fully
faithful. Set

\[\begin{gathered}
\mathcal H_ZF=i^{-1}\mathcal S_Z^0F,\\
i_*\mathcal H_ZF=\mathcal S_Z^0F.
\end{gathered}
\tag{1.2}\]

\textbf{Proposition 1.1.} The functor \(\mathcal H_Z\) is right adjoint
to the exact functor \(i_*\). It preserves injective objects and
K-injective complexes. There is a derived adjunction
\(i_*\dashv R\mathcal H_Z\), and \(R\mathcal H_Z(i_*N)=N\).

\textbf{Proof.} A map \(i_*G\to F\) has image with zero stalks off
\(Z\), so it factors uniquely through (1.1). Full faithfulness of
\(i_*\) then identifies it with a unique map \(G\to\mathcal H_ZF\).
Conversely such a map pushes forward and includes into \(F\); these
operations are inverse and natural. Since \(i_*\) is exact, Hom into
\(\mathcal H_ZI\) is exact if Hom into \(I\) is exact, proving
injectivity preservation. If \(A\) is an acyclic complex on \(Z\), chain
adjunction identifies maps and homotopies \(A\to\mathcal H_ZI\) with
those \(i_*A\to I\). The latter source is acyclic, so these maps are
null-homotopic for K-injective \(I\).

Choose K-injective models in this same adjunction. The map test into
K-injectives gives the derived adjunction. Also \(i_*\) preserves
K-injectives, because its left adjoint \(i^{-1}\) is exact; the
identical acyclic-source argument proves this. Resolving \(N\) by \(J\)
therefore computes \(R\mathcal H_Z(i_*N)\) by \(\mathcal H_Z(i_*J)=J\),
proving the final assertion. All arguments use left module maps only.
\(\square\)

Write \(D_Z(\mathcal O_X)\) for objects whose cohomology sheaves vanish
off \(Z\). Exact restriction shows this condition is equivalent to
\(K|_U=0\); cone long exact sequences show it defines a triangulated
subcategory. The stalk argument for the unit applies to every cohomology
sheaf, so \(i_*:D(\mathcal O_Z)\to D_Z(\mathcal O_X)\) is an equivalence
with inverse \(i^{-1}\). Proposition 1.1 then says

\[\mathcal S_ZK:=i_*R\mathcal H_ZK
\tag{1.3}\]

is right adjoint to the inclusion of \(D_Z\). Its counit
\(\mathcal S_ZK\to K\) is universal among maps from objects supported on
\(Z\). In particular it is an isomorphism if \(K\in D_Z\).

Define \(\Gamma_Z(X,F)=\Gamma(X,\mathcal S_Z^0F)\) and derive it with
K-injectives. Because \(\mathcal H_ZI\) is K-injective, (1.2) gives the
full local-to-global identity

\[R\Gamma_Z(X,K)=R\Gamma(Z,R\mathcal H_ZK).
\tag{1.4}\]

On the right we restrict the scalar action from
\(\Gamma(Z,\mathcal O_Z)\) through the map from
\(\Gamma(X,\mathcal O_X)\). This also proves that taking derived
sections of (1.3) gives the left side. These identities hold for all
unbounded complexes.

\subsection{2. Localization and its connecting
sign}\label{sections-with-support-and-the-localization-triangle--2-localization-and-its-connecting-sign}

\textbf{Theorem 2.1.} There are triangles natural in \(K\):

\[\begin{gathered}
\mathcal S_ZK\longrightarrow K\longrightarrow Rj_*(K|_U)\\
\longrightarrow\mathcal S_ZK[1],
\end{gathered}
\tag{2.1}\]

and

\[\begin{gathered}
R\Gamma_Z(X,K)\longrightarrow R\Gamma(X,K)\\
\longrightarrow R\Gamma(U,K|_U)\\
\longrightarrow R\Gamma_Z(X,K)[1].
\end{gathered}
\tag{2.2}\]

\textbf{Proof.} Choose \(K\to I\) K-injective with injective terms.
Every term is flasque, and its restriction to \(U\) is K-injective. On
every open \(V\), flasqueness makes \(I^n(V)\to I^n(V\cap U)\)
surjective. The kernel is precisely supported sections. Thus there are
termwise short exact sequences

\[\begin{gathered}
0\to i_*\mathcal H_ZI\to I\\
\to j_*(I|_U)\to0,\\[4pt]
0\to\Gamma_Z(X,I)\to\Gamma(X,I)\\
\to\Gamma(U,I)\to0.
\end{gathered}
\tag{2.3}\]

Each term is the indicated derived-functor model by Proposition 1.1 and
open restriction. The common reading\textquotesingle s exact-sequence
triangle proves (2.1) and (2.2). A derived map is a unique homotopy
class between the K-injective models. Applying the functors gives a map
of exact sequences; the cone and quotient construction is natural for
that map. Homotopic maps give homotopic maps on each term and identical
maps of derived triangles. Resolution comparisons are homotopy
equivalences, proving natural independence of the models. \(\square\)

The sign calculation is as follows. For
\(0\to A\xrightarrow{a}B\xrightarrow{b}C\to0\), the cone has
differential \(d(y,x)=(dy+a(x),-dx)\), projection \(p(y,x)=x\), and
quasi-isomorphism \(q(y,x)=b(y)\) to \(C\). A cycle \(z\in C^n\) lifts
to \(y\in B^n\); write \(dy=a(x)\). Injectivity of \(a\) gives \(dx=0\).
The cone cycle \((y,-x)\) maps to \(z\), so

\[H^n(p)H^n(q)^{-1}[z]=-[x].
\tag{2.4}\]

Our triangle\textquotesingle s last map is \(-p q^{-1}\), hence induces
the positive lift boundary \([z]\mapsto[x]\). A convention using the
positive cone projection has the opposite boundary. Negating two arrows
gives an isomorphic triangle: multiply one of its three objects by minus
one. Thus one must track the third-arrow convention when transporting a
formula from another source, rather than change one boundary informally.

\textbf{Proposition 2.2.} If an open \(V\) is disjoint from \(Z\), then
\(R\mathcal H_Z(Rv_*N)=0\) for \(v:V\hookrightarrow X\). Module and
underlying-abelian supported cohomology agree canonically, both globally
and as support sheaves.

\textbf{Proof.} A K-injective \(J\) on \(V\) has K-injective \(v_*J\),
since restriction is exact. Its sections over \(W\) and \(W\setminus Z\)
are both \(J(W\cap V)\), so its support kernel is zero. This proves the
first assertion. For coefficient comparison, resolve a module complex by
\(I\), and its underlying abelian complex further by an abelian
K-injective \(J\). Their support kernels and restrictions give a map of
the exact-sequence triangles (2.3), with injective terms used for both
resolutions. The middle and open-complement maps on derived global
sections are isomorphisms by Theorem 4.1 of the preceding lesson. The
long exact sequences, or the Five Lemma on consecutive five terms, make
the support map an isomorphism. Apply this argument on every open and
sheafify to obtain the local support comparison; restriction of both
coefficient models to opens is K-injective. Finally restrict the ambient
comparison to \(Z\) to get \(R\mathcal H_Z\). This supplies the
source\textquotesingle s omitted local verification. \(\square\)

\subsection{3. Locally closed
pairs}\label{sections-with-support-and-the-localization-triangle--3-locally-closed-pairs}

The construction also applies to a locally closed \(W\). Choose an open
\(V\) in which \(W\) is closed, and define the ambient support functor
by direct image from \(V\) of its closed-support subsheaf. It is
independent of this choice. For another such open \(V'\subset V\), a
supported section on \(T\cap V'\) glues with zero on
\(T\cap(V\setminus W)\): their overlap is outside its support, and these
two opens cover \(T\cap V\). This is the inverse to restriction of
supported sections. Compare arbitrary choices through their
intersection. These inverse maps commute with restrictions,
differentials and scalars. For an open \(W\), this functor is
\(w_*F|_W\), not a subsheaf of sections of \(F\) extended by zero.

\textbf{Proposition 3.1.} If \(W\) is locally closed and \(W'\) is
closed in \(W\), there is a natural triangle

\[\begin{gathered}
R\Gamma_{W'}^{\mathrm{sh}}K\longrightarrow R\Gamma_W^{\mathrm{sh}}K\\
\longrightarrow R\Gamma_{W\setminus W'}^{\mathrm{sh}}K
\longrightarrow R\Gamma_{W'}^{\mathrm{sh}}K[1].
\end{gathered}
\tag{3.1}\]

Here the superscript distinguishes these ambient sheaf functors from
global sections.

\textbf{Proof.} Choose \(V\) with \(W\) closed in \(V\); then \(W'\) is
also closed in \(V\), and \(W\setminus W'\) is closed in
\(V\setminus W'\). For a flasque \(F\), a supported section on
\((T\cap V)\setminus W'\) extends to \(T\cap V\). Outside \(W\) the
extension is still zero, since that open subset lies inside its original
domain. The extension therefore has support in \(W\). Its kernel under
restriction consists exactly of sections supported in \(W'\). This
proves a short exact sequence of the three support sheaves on every open
\(T\), even though direct image from \(V\) need not generally be exact.
Apply it termwise to a K-injective resolution with injective, hence
flasque, terms. The cone argument of Theorem 2.1 proves the triangle and
its functoriality, with the same positive lift boundary. \(\square\)

\subsection{4. Integer coefficients near a Euclidean
point}\label{sections-with-support-and-the-localization-triangle--4-integer-coefficients-near-a-euclidean-point}

We prove the comparison needed for this example, rather than assume
homotopy invariance for arbitrary sheaf cohomology. Let \(V\) be any
open subset of \(\mathbb R^n\) and \(A\) any abelian group. Let
\(S_q(V)\) be the free abelian group on continuous singular simplices
\(\Delta^q\to V\), with boundary the alternating sum of faces. Its
cochains are \(C^q(V;A)=\operatorname{Hom}_{\mathbb Z}(S_q(V),A)\), with
differential \(c\mapsto c\partial\). Let \(\mathcal C^q\) be the
sheafification of this cochain presheaf.

\textbf{Lemma 4.1 (cochain comparison).} The augmented complex
\(A_V\to\mathcal C^\bullet\) is a flasque resolution, and

\[H^q(V,A_V)=H^q(C^\bullet(V;A)).
\tag{4.1}\]

The comparison is natural under restrictions to smaller opens.

\textbf{Proof.} We supply the three required steps. First, a cover of an
open Euclidean set has a locally finite open refinement \((V_a)\) whose
closures lie in members of the original cover. Here is an explicit
compactness argument. Exhaust the open set by compact sets
\(K_m=\{x:|x|\leq m,\ \operatorname{dist}(x,\mathbb R^n\setminus V)\geq1/m\}\);
omit the distance condition when the complement is empty. They satisfy
\(K_m\subset\operatorname{int}K_{m+1}\) and exhaust \(V\). Cover each
compact layer \(K_m\setminus\operatorname{int}K_{m-1}\) by finitely many
balls with closures in cover members and in
\(\operatorname{int}K_{m+1}\setminus K_{m-2}\). Set the negative-index
compact sets empty. The balls cover \(V\); every point has a
neighborhood in some \(K_m\), which meets only finitely many subsequent
layers, so they are locally finite.

A section of \(\mathcal C^q\) is represented by cochains \(c_a\) on
cover members. Using this refinement, assign to each singular simplex
contained in some \(V_a\) the value given by one such \(c_a\); assign
zero if it is contained in none. This is a global cochain representing
the section. To verify the assertion, fix a point \(x\). Only finitely
many refinement members meet a sufficiently small neighborhood. Remove
those whose closures do not contain \(x\). For the remaining finite
list, their representing cochains are defined near \(x\), and their
germs agree. Shrink the neighborhood so all these cochains agree on
every simplex in it, and so it is contained in one refinement member
containing \(x\). The global assignment then agrees with that
representative there, irrespective of the choices for simplices. This
proves surjectivity of \(C^q(V;A)\to\Gamma(V,\mathcal C^q)\). The same
proof works on every smaller open. Extend a representing cochain by zero
on simplices not contained in that smaller open; this proves flasqueness
of \(\mathcal C^q\).

Second, on a convex ball straight-line contraction proves that singular
cohomology is \(A\) in degree zero and zero in positive degrees. The
chain verification is the prism identity: a homotopy \(H\) from \(f\) to
\(g\) assigns to a \(q\)-simplex the alternating sum, with signs
\((-1)^i\), of the \(q+1\) simplices with vertices
\((v_0,0),\ldots,(v_i,0),(v_i,1),\ldots,(v_q,1)\) in
\(\Delta^q\times[0,1]\), followed by \(H\). Interior faces cancel in
pairs, the two horizontal faces give \(g-f\), and the other faces give
the negative prism of the boundary. Thus
\(\partial P+P\partial=g_*-f_*\). A constant map factors through the
point\textquotesingle s chain complex; its augmented homology contracts
(insert the unique vertex), and dually its cohomology is as stated.
Taking the exact filtered colimit over convex-ball neighborhoods proves
that \(A_V\to\mathcal C^\bullet\) is exact on stalks.

Third, the kernel of global cochains mapping to sheaf sections is the
complex \(N^\bullet\) of cochains locally zero: each point has a
neighborhood on all of whose singular simplices the cochain vanishes. We
prove that this kernel is acyclic. For an open cover \(\mathcal U\), let
\(S_*^{\mathcal U}\) be the subcomplex generated by simplices contained
in one cover member, and let \(N_{\mathcal U}^\bullet\) be its
annihilator. The small-chain inclusion has a retraction \(r\) and
homotopy \(h\) satisfying

\[\begin{gathered}
1-jr=\partial h+h\partial,\\
rj=1,\qquad hj=0.
\end{gathered}
\tag{4.2}\]

Here is the full construction. Barycentric subdivision \(s\) replaces a
simplex by its oriented barycentric triangulation. Internal faces
cancel, and boundary faces are subdivided faces, so it is a chain map.
There is a homotopy \(T\) with \(s-1=\partial T+T\partial\), supported
inside each simplex image. To construct it without another theorem, work
first in the standard simplex. In dimension zero take \(T=0\). In
dimension \(q\), subtract the already constructed homotopies on its
faces from \(s\iota_q-\iota_q\). The result is a cycle by the
lower-dimensional identity. Cone that cycle to the barycenter inside the
convex simplex, then push it forward by the singular simplex. The cone
boundary formula \(\partial(b*c)=c-b*\partial c\) proves the identity.
Write \(T_m\) for the sum of iterated homotopies, so
\(s^m-1=\partial T_m+T_m\partial\); it has the same support control.

Induct on simplex dimension to construct (4.2). Once \(r,h\) are
constructed on faces of \(\sigma\), put \(z=\sigma-h\partial\sigma\).
The induction identity gives \(\partial z=jr\partial\sigma\). Choose
\(m\) so that every simplex in \(s^m z\) is small, and define

\[\begin{aligned}
jr\sigma&=s^mz-T_mjr\partial\sigma,\\
h\sigma&=-T_mz.
\end{aligned}
\tag{4.3}\]

The correction is small: \(r\partial\sigma\) is small, and \(T_m\) keeps
each simplex inside its image. The boundary of the first line is
\(jr\partial\sigma\); subtracting it from \(z\) gives
\(-\partial T_mz\), proving the homotopy identity. For a small
\(\sigma\), its faces have \(h=0\); choose \(m=0\), obtaining
\(r\sigma=\sigma\) and \(h\sigma=0\). The required \(m\) always exists:
barycentric subdivision reduces mesh in dimension \(q>0\) by at most
\(q/(q+1)\). Indeed two vertices of a subdivided simplex are barycenters
of nested faces, and their distance is bounded by that fraction of the
original diameter. Uniform continuity on the compact standard simplex
and a finite subcover of its inverse-image cover then make repeated
subdivisions small. The finitely many simplices in \(z\) can use a
common \(m\).

Dualizing (4.2) contracts \(N_{\mathcal U}^\bullet\): \(cjr=0\), and
\(ch\) still vanishes on small chains because \(hj=0\). Every locally
zero cochain belongs to some \(N_{\mathcal U}\), and common refinement
orders these complexes as a filtered union. Exact filtered colimits make
\(N^\bullet\) acyclic. Our first step identifies the section complex of
\(\mathcal C^\bullet\) with \(C^\bullet/N^\bullet\); the quotient is
therefore quasi-isomorphic to singular cochains. Finally flasque
acyclicity and the bounded-below acyclic-resolution theorem from the
bounded-derived lesson compute sheaf cohomology by this section complex.
All maps used in the quotient comparison are restriction maps, giving
its naturality. \(\square\)

\textbf{Proposition 4.2.} For \(n\geq0\),

\[R\Gamma_{\{0\}}(\mathbb R^n,\mathbb Z)
\cong\mathbb Z[-n].
\tag{4.4}\]

\textbf{Proof.} The comparison proves
\(H^q(\mathbb R^n,\mathbb Z)=\mathbb Z\) for \(q=0\) and zero otherwise,
by straight-line contraction. For \(n>0\), the punctured space
deformation retracts onto \(S^{n-1}\) by radial rescaling; the prism
identity proves invariance of singular cohomology under this homotopy.

We recall and prove the sphere calculation needed here. For a cover by
two opens, its small-chain complex is the sum of the two singular chain
subcomplexes; their intersection is the chain complex of their
intersection. Cochains therefore have the short exact sequence with
middle term \(C^\bullet(U;\mathbb Z)\oplus C^\bullet(V;\mathbb Z)\),
last term \(C^\bullet(U\cap V;\mathbb Z)\), and map the difference of
restrictions. Surjectivity follows by extending any cochain on the
intersection by zero to either open. The small-chain equivalence just
proved gives the usual long exact Mayer--Vietoris sequence, with no
sheaf comparison required on the sphere itself. Cover \(S^m\) by
complements of its two poles. Each is homeomorphic to \(\mathbb R^m\),
and their intersection retracts to \(S^{m-1}\). Starting with \(S^0\),
whose two points have \(H^0=\mathbb Z^2\), the sequence proves
inductively that the reduced cohomology of \(S^m\) is \(\mathbb Z\) in
degree \(m\) and zero elsewhere.

For \(n=1\), the restriction from the line to its two complementary rays
is the diagonal \(\mathbb Z\to\mathbb Z^2\). Its cokernel is
\(\mathbb Z\). For \(n>1\) the degree-zero restriction is an
isomorphism, and the sole reduced punctured-space class has degree
\(n-1\). The long exact sequence of (2.2) puts the sole supported class
in degree \(n\). For \(n=0\), the support is the entire one-point space
and the result is immediate. A complex with just this cohomology is
isomorphic to that cohomology module in its degree: the good-truncation
inclusion and quotient give the requisite quasi-isomorphisms, as in the
common reading\textquotesingle s Proposition 5.2. This proves (4.4).
\(\square\)

The same computation on balls centered at zero gives
\(\mathcal S_{\{0\}}\mathbb Z=\mathbb Z_{\{0\}}[-n]\). The transition
isomorphisms follow from the natural comparison and radial rescaling of
punctured balls. This is supported cohomology, rather than ordinary
sections of the constant sheaf supported at the point; those ordinary
sections are zero when \(n>0\).

\subsection{5. Cup products and support
pullback}\label{sections-with-support-and-the-localization-triangle--5-cup-products-and-support-pullback}

Assume commutative coefficients in this section. Tensoring an object
supported on \(Z\) with any derived object keeps that support: a K-flat
model makes its restriction to the complement acyclic. The universal
counit in (1.3) therefore gives a unique factorization of
\(K\otimes^{\mathbf L}\mathcal S_ZM\to K\otimes^{\mathbf L}M\) through
\(\mathcal S_Z(K\otimes^{\mathbf L}M)\). On closed-set coefficients it
is the canonical map

\[\begin{gathered}
i^{-1}K\otimes^{\mathbf L}R\mathcal H_ZM\\
\longrightarrow R\mathcal H_Z(K\otimes^{\mathbf L}M).
\end{gathered}
\tag{5.1}\]

Indeed \(K\otimes^{\mathbf L}i_*N=i_*(i^{-1}K\otimes^{\mathbf L}N)\):
take a K-flat model for \(K\), and check the ordinary tensor comparison
on stalks, identical over \(\mathcal O_z\) on \(Z\) and zero off it.
This proves the identity without a projection-formula assumption. The
universal factorization proves naturality and coherence under successive
tensor products, because the composites have the same image under the
counit.

The internal-Hom lesson identifies cohomology classes with maps from
\(\mathcal O\) to shifted objects. Tensoring such maps defines cup
products. Apply the factorization above to get

\[\begin{gathered}
H^p(X,K)\times H_Z^q(X,M)\\
\longrightarrow H_Z^{p+q}(X,K\otimes^{\mathbf L}M).
\end{gathered}
\tag{5.2}\]

Forgetting supports commutes with cup product, since composing the
factorization with its counit is precisely the ordinary tensor map. This
proves the omitted compatibility in Tag 0G77. On homogeneous cycles the
product is their tensor, whose differential is
\(da\otimes b+(-1)^pa\otimes db\); changing either cycle by a boundary
changes it by a boundary. The coherent shift identification gives
exactly these signs.

For a morphism \(f:X'\to X\), set \(Z'=f^{-1}Z\). Derived pullback
carries supported objects to objects supported on \(Z'\): off that set,
a K-flat acyclic stalk remains acyclic after arbitrary scalar extension,
by the tensor lesson\textquotesingle s Lemma 1.1. Pulling back the
support counit therefore factors uniquely through

\[Lf^*\mathcal S_ZK\longrightarrow\mathcal S_{Z'}Lf^*K.
\tag{5.3}\]

For \(g:Z'\to Z\) it equivalently gives
\(Lg^*R\mathcal H_ZK\to R\mathcal H_{Z'}Lf^*K\). To justify the
equivalence directly, a K-flat model \(P\) on \(Z\) has K-flat \(i_*P\)
on \(X\), by the stalk test. The canonical ordinary closed-set
comparison \(f^*i_*P\to i'_*g^*P\) is an isomorphism on each stalk,
proving the derived identity. Universal uniqueness proves that (5.3)
commutes with forgetting supports and composes for two pullbacks. On
cohomology, pullback uses the unit of \(Lf^*\dashv Rf_*\) and the global
Leray identity from lesson six; its naturality and the counit
factorization prove the supported/ordinary pullback square commutes.
This is a canonical map, without a general assertion that it is
invertible.

\subsection{6. A two-point calculation and
exercises}\label{sections-with-support-and-the-localization-triangle--6-a-two-point-calculation-and-exercises}

Let \(X=\{c,\eta\}\) have opens \(\varnothing,\{\eta\},X\); its closed
point is \(c\). With a constant ring sheaf, a module sheaf is exactly an
arrow \(r:M\to N\). Its global sections are \(M\), its sections on the
open point are \(N\), and exactness is exactness at both entries. Thus
global sections on this space and on its open point are exact functors.
Also \(j_*N=(N\xrightarrow{1}N)\) is exact. The localization triangle
gives

\[\begin{aligned}
H^0_c(X,F)&=\ker r,\\
H^1_c(X,F)&=\operatorname{coker}r,\\
H^q_c(X,F)&=0\quad(q\ne0,1).
\end{aligned}
\tag{6.1}\]

The corresponding ambient cohomology sheaves are these modules at \(c\),
zero at \(\eta\). Ordinary supported sections only see the kernel, while
the degree-one group measures the failed extension of a section from
\(\eta\). The first exercise shows that failure before any resolutions
are chosen.

\textbf{Exercise 1 (easy: left exactness and its limit).} Prove left
exactness of closed supported sections directly from their definition.
On the two-point space give an exact sequence for which they fail to be
right exact.

\textbf{Solution.} In an exact sequence \(0\to F\to G\to H\), a
supported section of \(G\) mapping to zero in \(H\) lifts uniquely to
\(F\), by left exactness of ordinary sections. Its restriction off the
closed set vanishes because its image in \(G\) does, and \(F\to G\) is
injective. Thus it is supported. This proves exactness at the first two
terms, globally and on every open. On the two-point space the exact
sequence of arrows \(0\to(0\to k)\to(k\xrightarrow{1}k)\to(k\to0)\to0\)
is exact at each entry. Its closed supported sections are respectively
\(0,0,k\), so the last map is not surjective. All the maps used are
module maps over the same ring, and the counterexample works with
noncommutative coefficients too.

\textbf{Exercise 2 (medium: the localization fibre).} For an arrow
\(r:M\to N\), give a complex computing the global supported object,
including the connecting map \(N\to H^1_c(X,F)\).

\textbf{Solution.} Since the other two global terms are \(M[0]\) and
\(N[0]\), the supported object is the shifted cone \(C(r)[-1]\). It has
\(M\) in degree zero, \(N\) in degree one, and differential \(-r\). Its
map to \(M[0]\) is the identity in degree zero. Rotating the cone
triangle with last map \(-p\) gives the next map \(r\) and the final map
\(N[0]\to C(r)\) as the positive cone inclusion. Thus its boundary on
cohomology is the positive quotient \(N\to N/r(M)\). The
degree-minus-one shift is why the fibre differential is \(-r\).
Replacing it by \(+r\) by negating its \(N\) term would also negate the
coordinate description of the boundary; one cannot make that replacement
while leaving the boundary description unchanged. This proves (6.1) and
checks the sign convention in a fully explicit sheaf model.

\textbf{Exercise 3 (medium: the Euclidean group and its sign).} Compute
all \(H^q_{\{0\}}(\mathbb R^n,\mathbb Z)\). For \(n=1\), label the
complementary rays in negative, positive order and determine the
boundary of their locally constant section \((a,b)\).

\textbf{Solution.} Proposition 4.2 gives \(\mathbb Z\) in degree \(n\)
and zero in all other degrees, including the whole-support case \(n=0\).
In dimension one the degree-zero restriction is diagonal, so the
boundary factors through \(\mathbb Z^2/\mathbb Z(1,1)\). Its positive
coordinate is \(b-a\). To check the sign with actual cochains, extend
the function on the two rays arbitrarily to a zero-cochain on the line.
Its differential on an oriented interval crossing zero is its value at
the positive endpoint minus its value at the negative endpoint, namely
\(b-a\). This differential vanishes on chains contained in either ray
and represents the relative degree-one class.

The flasque resolution in Lemma 4.1 also computes supported cohomology:
each of its terms is acyclic for sections on the line and on the
complement, and its restriction is onto, so Theorem 2.1 makes it acyclic
for supported sections. The bounded-below acyclic-term theorem therefore
applies. Naturality of the singular comparison identifies its cone with
the relative cochain cone. Equation (2.4) says the
triangle\textquotesingle s negative projection gives precisely this
positive lift differential. Thus the increasing orientation of the real
line gives \(\delta(a,b)=b-a\), rather than its negative. In higher
dimensions the boundary takes a chosen linking-sphere class to the
corresponding supported class with this same positive lift convention;
changing the orientation negates both chosen generators.

\textbf{Exercise 4 (hard: cup products with support).} Prove that
forgetting supports in (5.2) gives the ordinary cup product. Prove also
that first multiplying two ordinary classes and then a supported class
agrees, after association, with the successive supported cup products.

\textbf{Solution.} Let \(S=\mathcal S_ZM\). The source
\(K\otimes^{\mathbf L}S\) lies in \(D_Z\), so the right adjunction
identifies maps from it to \(\mathcal S_Z(K\otimes^{\mathbf L}M)\) with
maps from it to \(K\otimes^{\mathbf L}M\). The defining transpose is
\(1_K\) tensor the support counit. Composing the supported map with the
new counit therefore gives exactly that tensor map. A class is a map
from the tensor unit to a shifted object; tensoring the two class maps
and composing these equal morphisms proves compatibility with forgetting
supports on cohomology, in every pair of degrees.

For objects \(K,L,M\), both successive support products are maps from
\(K\otimes^{\mathbf L}L\otimes^{\mathbf L}\mathcal S_ZM\), an object of
\(D_Z\), to \(\mathcal S_Z(K\otimes^{\mathbf L}L\otimes^{\mathbf L}M)\).
Under the adjunction both are \(1_K\otimes1_L\) tensor the counit. Its
uniqueness proves equality; the tensor associator identifies their
sources and targets. Applying this to shifted class maps proves the
assertion for classes. No exactness of global sections, flatness of the
classes\textquotesingle{} sheaves, or finiteness of their degrees is
used. K-flat models and the proved tensor coherence justify the entire
construction.

\textbf{Exercise 5 (hard: pullback is a map, not always an
isomorphism).} Explain why (5.3) composes for two ringed-space morphisms
and give a case where it is not an isomorphism.

\textbf{Solution.} Pulling back the original counit through two maps
gives the same morphism as pulling it back through their composite, by
lesson six\textquotesingle s coherent pullback composition. Every source
involved is supported on the inverse-image closed subset. The universal
factorization into its support right adjoint is unique, so the two
resulting maps agree after the canonical composition identification.

For failure, take \(X=\mathbb R\), \(Z=\{0\}\), constant coefficients
\(\mathbb Z\), and \(f\) the inclusion of that point with the unchanged
stalk coefficient ring. Pullback is the exact stalk functor. Proposition
4.2 and its ball calculation give
\(Lf^*\mathcal S_Z\mathbb Z=\mathbb Z[-1]\). The inverse-image support
is the entire point, so \(\mathcal S_{Z'}Lf^*\mathbb Z=\mathbb Z[0]\).
Their cohomology groups occur in different degrees, and the canonical
map cannot be an isomorphism. All coefficient maps are identity maps;
this failure concerns the support operation, not nonflat change of
coefficient rings.

\section{Perfect complexes and duals on a ringed
space}\label{perfect-complexes-and-duals}

\emph{Written and edited by GPT-6.1 Sol (OpenAI), in Codex at Ultra,
October 2026. Mathematically self-checked by the writing AI. Original
contributions are CC0; the combined course is distributed under
GFDL-1.2-or-later. Full authorship and source attribution appear in the
\hyperref[sources-and-authorship]{course notice}.}

A perfect complex is locally a finite calculation with finite projective
coefficients. This finiteness has two complementary descriptions:
arbitrarily accurate finite approximations, together with a bound on the
degrees that tensor products can create; and a dual that permits
evaluation and coevaluation. We prove the equivalence on arbitrary
ringed spaces. Local rings simplify the models, but are not a hypothesis
of the main theorem.

The prerequisites are
\hyperref[flat-modules-and-k-flat-resolutions]{Flat modules and K-flat
resolutions}, Lemmas 1.2--1.3, 2.2--2.4, Theorem 3.1 and Lemma 4.1;
\hyperref[derived-tensor-products-and-tor-sheaves]{The derived tensor
product and Tor sheaves}, Theorem 2.2 and Theorem 3.3;
\hyperref[internal-derived-hom-and-ext-sheaves]{Hom complexes, internal
derived Hom and Ext sheaves}, Lemma 1.1, Theorems 2.2 and 3.1 and the
canonical maps in Section 4; and
\hyperref[derived-pullback-and-pushforward]{Derived pullback and
pushforward}, Theorem 1.1 and Proposition 3.2.
\hyperref[injective-modules-and-bounded-below-derived-functors]{Injective
modules and bounded-below derived functors}, Theorem 3.3, supplies maps
into K-injectives. The common reading supplies cones, their long exact
cohomology sequences, and localization by roofs. The support
lesson\textquotesingle s Theorem 2.1 is used once for extension across
an open complement.

Throughout, \(\mathcal O=\mathcal O_X\) is commutative and unital, and
all derived categories are unbounded. Write
\([K,L]=R\mathcal Hom(K,L)\). A ``finite projective sheaf'' here means a
direct summand of \(\mathcal O^r\) for finite \(r\); a ``locally finite
projective sheaf'' has that property on an open cover. Bounds stated
locally need not be uniform on the space.

The construction follows the Stacks project authors' \emph{Cohomology of
Sheaves}, ``Strictly perfect complexes'', ``Pseudo-coherent modules'',
``Tor dimension'', ``Perfect complexes'' and ``Duals'', in the
\href{https://github.com/KokunoYumeto/unofficial-stacks-project-ai-drafts/blob/565b10e987aba5969b21145a0833f42d69f96790/cohomology.tex}{AI
Integrated Stacks Project edition}. We prove the finite-model criterion
and duality, including the finite-cell argument needed for the converse.
Source attribution and the licence for adapted passages appear in the
\hyperref[sources-and-authorship]{course notice}.

\subsection{1. Finite models and local
maps}\label{perfect-complexes-and-duals--1-finite-models-and-local-maps}

A complex is \textbf{strictly perfect} if it is bounded and each term is
finite projective. An object is \textbf{perfect} if, on an open cover,
it is isomorphic in the derived category to a strictly perfect complex.
Cones, shifts and tensor products of strictly perfect complexes are
strictly perfect: in each degree only finitely many summands occur; a
tensor product of summands of \(\mathcal O^r\) and \(\mathcal O^s\) is a
summand of \(\mathcal O^{rs}\). Pullback preserves these splittings.

\textbf{Lemma 1.1 (finite lifting and local representatives).} A map
from a finite projective sheaf to the target of a sheaf epimorphism
lifts near every point. If \(P\) is strictly perfect with lowest term in
degree \(a\), and \(H^j(C)=0\) for \(j\ge a\), every chain map
\(P\to C\) is locally null-homotopic. Every derived map \(P\to F\) is
locally represented by a chain map, and a chain map that is zero in the
derived category is locally null-homotopic.

\textbf{Proof.} First lift the finitely many images of the basis of
\(\mathcal O^r\) on a common neighbourhood, then restrict the resulting
lift to its summand. For the homotopy, descend from the highest nonzero
degree \(b\) of \(P\). The degree-\(b\) component lands in cycles. Since
\(H^b(C)=0\), the epimorphism \(C^{b-1}\to Z^b(C)\) permits a local lift
\(h^b\). At degree \(i\), the residual map \(f^i-h^{i+1}d_P\) lands in
cycles because the chain-map equation and the already constructed
homotopy in degree \(i+1\) give zero after applying \(d_C\). Lift it to
\(C^{i-1}\). There are finitely many degrees and shrinkings, and the
result is \(f=dh+hd\).

Take a K-injective resolution \(F\to I\). Maps into \(I\) represent
derived maps by the bounded-derived lesson\textquotesingle s Theorem
3.3. If \(C\) is acyclic, then \(\mathsf H(P,C)\) is an acyclic sheaf
complex: Hom from a finite projective term is a summand of a finite
power of the exact identity functor, and the finite filtration by terms
of \(P\) proves the assertion. Consequently
\(\mathsf H(P,F)\to\mathsf H(P,I)\) is a quasi-isomorphism. A
degree-zero cycle on the right lifts locally to a cycle on the left
modulo a boundary, since its cohomology-sheaf map is an isomorphism.
This gives an actual representative. A cycle on the left whose image is
a boundary is locally a boundary on the left, giving the last assertion.
The same reasoning shows that a derived map \(P\to C\), under the stated
cohomology vanishing, is locally zero: first represent it, then apply
the downward homotopy. \(\square\)

It follows that perfectness is independent of the chosen representative:
locally realize a strictly perfect model\textquotesingle s derived
isomorphism by a chain map to that representative; it is a
quasi-isomorphism because its cohomology maps are isomorphisms.

For later use, an acyclic top pair in a strictly perfect complex can be
cancelled locally. If the top term is \(P^b\) and \(H^b(P)=0\), the
epimorphism \(P^{b-1}\to P^b\) splits by Lemma 1.1. Its kernel is finite
projective, and the complex decomposes as that shortened complex plus
the contractible identity disk on \(P^b\). Repeat finitely often.

\textbf{Lemma 1.2 (finite presentation and flatness).} A flat, finitely
presented sheaf is locally finite projective. If all stalk rings are
local, locally finite projective sheaves are finite locally free.

\textbf{Proof.} Locally write a presentation
\(\mathcal O^m\xrightarrow{A}\mathcal O^n\to F\to0\), and put
\(L=\operatorname{im}A\). Flatness of \(F\) makes
\(L\subset\mathcal O^n\) pure: tensoring this inclusion with any sheaf
is injective, by the flat-quotient lemma. Let \(l_\alpha\) be the
columns of \(A=(a_{i\alpha})\). The finite equations

\[\sum_i a_{i\alpha}y_i=l_\alpha
\qquad(1\le\alpha\le m)
\tag{1.1}\]

have the ambient solutions \(y_i=e_i\). Set
\(T=\operatorname{coker}(A^t:\mathcal O^n\to\mathcal O^m)\). The tuple
\((l_\alpha)\) maps to zero in \(T\otimes\mathcal O^n\), hence in
\(T\otimes L\) by purity. Right exactness of tensor says it locally lies
in the image of \(A^t\otimes L\). Thus (1.1) has solutions in \(L\) on a
neighbourhood. Sending \(e_i\) to \(y_i\) gives a retraction
\(\mathcal O^n\to L\), because it fixes the generators \(l_\alpha\). The
quotient \(F\) is a finite free summand.

For completeness, a finite projective module over a local ring
\((R,\mathfrak m)\) is free. Nakayama\textquotesingle s lemma follows
directly: if finite generators satisfy \(v=Bv\) with entries of \(B\) in
\(\mathfrak m\), then \(\det(1-B)\) is a unit, and the adjugate identity
gives \(v=0\). Lift a basis modulo \(\mathfrak m\) to a finite
projective module \(P\). Nakayama applied to the cokernel gives a
surjection \(R^r\to P\). It splits; its finite kernel has zero reduction
modulo \(\mathfrak m\), so another application of Nakayama kills that
kernel.

A free stalk of a finitely presented sheaf spreads to a free sheaf near
that point. Lift its finite basis, and kill the stalk-zero cokernel
using finitely many generators. The resulting surjection has finitely
generated kernel: from a finite presentation, lift maps between the two
finite free presentations; the original finite relations and the columns
of the difference between their composite and the identity generate the
kernel. Its zero stalk therefore kills finitely many kernel generators
on a common neighbourhood. The lifted basis is an isomorphism there.
Apply this to the finite projective sheaf. \(\square\)

For arbitrary stalk rings, ``projective'' cannot be replaced by
``free''. On the one-point space with ring \(k\times k\), the summand
\((1,0)R\) is not \(R^r\) for any \(r\): its two component dimensions
are \((1,0)\), whereas those of \(R^r\) are \((r,r)\).

\textbf{Proposition 1.3 (Hom from a finite model).} For strictly perfect
\(P\) and any complex \(F\), \(\mathsf H(P,F)\) represents \([P,F]\). If
\(F\) is K-flat, this Hom complex is K-flat.

\textbf{Proof.} The finite-filtration argument in Lemma 1.1 makes
\(\mathsf H(P,-)\) preserve quasi-isomorphisms, so compare \(F\to I\)
with the definition \([P,F]=\mathsf H(P,I)\). Define

\[\begin{aligned}
(P^\vee)^n&=\mathcal Hom(P^{-n},\mathcal O),\\
d_{P^\vee}^n&=(-1)^{n+1}(d_P^{-n-1})^t.
\end{aligned}
\tag{1.2}\]

Evaluation identifies \(F\otimes P^\vee\) with \(\mathsf H(P,F)\): on
homogeneous elements it sends \(f\otimes\phi\) to \(p\mapsto f\phi(p)\).
The identity is clear for finite free terms and passes to summands. The
Hom differential and the sign in (1.2) agree with the tensor
differential. The dual is bounded with flat terms, hence K-flat;
tensoring two K-flats remains K-flat. \(\square\)

\subsection{2. Finite
approximations}\label{perfect-complexes-and-duals--2-finite-approximations}

An object \(K\) is \textbf{\(m\)-pseudo-coherent} if locally there is a
map \(P\to K\) from a strictly perfect complex, inducing isomorphisms on
\(H^j\) for \(j>m\) and a surjection on \(H^m\). Equivalently, its cone
has zero cohomology in degrees \(j\ge m\). It is
\textbf{pseudo-coherent} if this holds for every integer \(m\). The
cover may depend on \(m\). Lemma 1.1 makes this definition equivalent to
using actual maps to any chosen complex representative.

Smaller thresholds impose stronger conditions. An \(m\)-approximation is
also an \((m+1)\)-approximation, and shifting by \([r]\) changes the
threshold to \(m-r\). Such an object is locally bounded above, because
its approximation has finitely many nonzero terms. On a quasi-compact
space finitely many neighbourhoods give a uniform upper bound.

Derived pullback preserves \(m\)-pseudo-coherence. Indeed, a complex
with cohomology only in degrees \(\le c\) has a bounded-above flat model
with terms only in those degrees, by good truncation and Lemma 4.1 of
the flat lesson. Its derived pullback has the same upper bound. Pull
back the approximation triangle; the strict model remains strict and the
cone remains in degrees \(<m\). Similarly, if \(K\in D^{\le a}\) and
\(L\in D^{\le b}\), these flat models show
\(K\otimes^L L\in D^{\le a+b}\).

\textbf{Proposition 2.1 (triangles and tensor bounds).} In a triangle
\(K\to L\to M\to K[1]\), the following implications hold:

\begin{enumerate}
\tightlist
\item
  If \(K\) is \((m+1)\)-pseudo-coherent and \(L\) is
  \(m\)-pseudo-coherent, then \(M\) is \(m\)-pseudo-coherent.
\item
  If \(K,M\) are \(m\)-pseudo-coherent, then \(L\) is
  \(m\)-pseudo-coherent.
\item
  If \(L\) is \((m+1)\)-pseudo-coherent and \(M\) is
  \(m\)-pseudo-coherent, then \(K\) is \((m+1)\)-pseudo-coherent.
\end{enumerate}

If \(K\in D^{\le a}\) is \(n\)-pseudo-coherent and \(L\in D^{\le b}\) is
\(m\)-pseudo-coherent, their tensor product is
\(\max(n+b,m+a)\)-pseudo-coherent.

\textbf{Proof.} Choose approximations \(P\to K\), \(Q\to L\) at
thresholds \(m+1,m\). Brutally cut \(P\) below \(m+1\), preserving that
approximation. For \(C=\operatorname{Cone}(Q\to L)\), the map
\(P\to L\to C\) is locally zero by Lemma 1.1: \(C\) has no cohomology in
degrees \(\ge m\), whereas \(P\) starts at \(m+1\). Exactness of Hom
applied to the triangle \(Q\to L\to C\) lifts \(P\to L\) to a derived
map \(P\to Q\); Lemma 1.1 realizes it by a chain map locally. The
triangle axiom gives a map \(\operatorname{Cone}(P\to Q)\to M\).

Here is the required cohomology check. Cone cohomology in degree \(j\)
is an extension of the cokernel in degree \(j\) by the kernel in degree
\(j+1\). For \(j>m\), the two cokernels are isomorphic: the map on
\(H^j(Q)\) is an isomorphism and the map from \(P\)\textquotesingle s
cohomology is onto at \(j=m+1\), an isomorphism above it. The kernels in
degree \(j+1\) are isomorphic. Hence the extension map is an
isomorphism. At \(j=m\), surjectivity for \(Q\) makes the cokernel map
onto, and surjectivity for \(P\) at \(m+1\), with the isomorphism for
\(Q\) there, makes the kernel map onto. Lifting first the kernel
component and then the cokernel component proves surjectivity for the
extension. Thus the cone is an approximation. Rotation and the shift
rule give the other two implications.

For the tensor claim choose approximations \(P\to K\), \(Q\to L\) at
thresholds \(n,m\). They can have terms only in degrees \(\le a,\le b\):
cancel acyclic top pairs; if the threshold exceeds the relevant bound,
the zero complex is already an approximation. Factor the tensor map
through \(K\otimes^L Q\). Its first cone is
\(\operatorname{Cone}(P\to K)\otimes Q\), bounded above by \(n+b-1\);
its second cone is \(K\otimes^L\operatorname{Cone}(Q\to L)\), bounded
above by \(a+m-1\). The octahedral cone triangle bounds the composite
cone by the larger of these numbers. The source \(P\otimes Q\) is
strictly perfect, proving the assertion without a Tor spectral sequence.
Taking arbitrarily small thresholds proves tensor closure for
pseudo-coherent objects. \(\square\)

\textbf{Proposition 2.2 (summands and criteria).} Summands of
\(m\)-pseudo-coherent objects are \(m\)-pseudo-coherent. A locally
bounded-above complex whose term in degree \(i\) is
\((m-i)\)-pseudo-coherent is \(m\)-pseudo-coherent; the same holds if
its cohomology sheaf in degree \(i\) is \((m-i)\)-pseudo-coherent. A
module is \(0\)-pseudo-coherent exactly when it is of finite type, and
\((-1)\)-pseudo-coherent exactly when it is finitely presented.

\textbf{Proof.} Suppose \(K\oplus L\) is \(m\)-pseudo-coherent. Its
endomorphism \(1_K\oplus0\) has cone \(L\oplus L[1]\). Proposition 2.1
applies since \(m\)-pseudo-coherence implies \((m+1)\)-pseudo-coherence.
Shift repeatedly; each \(L[n]\oplus L[n+1]\) is \(m\)-pseudo-coherent.
Work near a point where \(L\) is bounded above. For large \(n\),
\(L[n]\) has zero cohomology in degrees \(\ge m\), so the zero
approximation works. Descend using the split triangles

\[\begin{gathered}
L[n]\longrightarrow L[n]\oplus L[n-1]\\
\longrightarrow L[n-1]\longrightarrow L[n+1].
\end{gathered}
\tag{2.2}\]

The first implication of Proposition 2.1 supplies each preceding step.
Only finitely many steps are needed for a fixed threshold at a fixed
point. Exchanging \(K,L\) finishes the summand assertion.

For the term criterion, remove all degrees below \(r\le m\). The finite
upper tail is \(m\)-pseudo-coherent by its filtration into terms
\(F^i[-i]\) and the extension implication. The lower quotient is bounded
above by \(r-1<m\), so has the zero approximation. The same extension
implication proves the claim. For the cohomology criterion use the
finite good-truncation filtration, with quotients \(H^i(K)[-i]\), and
its lower tail in degrees \(<r\).

If \(K\in D^{\le m}\), cancel the top acyclic terms in an
\(m\)-approximation: \(H^m(K)\) is a quotient of a finite projective
sheaf, hence finite type. If \(K\in D^{\le m+1}\), the same cancellation
leaves \(H^{m+1}(K)\) as the cokernel of a map of two finite projective
sheaves, hence finitely presented. Apply these facts to a module in
degree zero. Conversely, a finite generating surjection in degree zero
supplies a \(0\)-approximation; a finite presentation in degrees
\(-1,0\) supplies a \((-1)\)-approximation. \(\square\)

\subsection{3. Tor amplitude and the perfectness
criterion}\label{perfect-complexes-and-duals--3-tor-amplitude-and-the-perfectness-criterion}

An object \(K\) has \textbf{Tor amplitude in \([a,b]\)} if, for every
module sheaf \(F\) in degree zero, \(K\otimes^L F\) has cohomology only
in degrees \(a\) through \(b\). It has \textbf{locally finite Tor
amplitude} if such an interval exists on each member of an open cover.
Taking \(F=\mathcal O\) also bounds \(K\) itself.

\textbf{Proposition 3.1 (amplitude models).} Amplitude in \([a,b]\) is
equivalent to having a representative of flat sheaves with terms only in
\([a,b]\). It is also equivalent to every stalk having that amplitude
over its stalk ring. Derived pullback preserves the interval; tensoring
objects with intervals \([a,b]\), \([c,d]\) gives interval
\([a+c,b+d]\). Summands inherit amplitude.

\textbf{Proof.} A bounded complex of flat sheaves is K-flat. Its tensor
with a module has terms only in the asserted interval, proving one
implication. Conversely, resolve a good upper truncation of \(K\) by the
bounded-above flat construction, obtaining \(P\) with terms only in
degrees \(\le b\). Since \(H^i(P)=0\) below \(a\), the flat tail ending
in \(P^a\) resolves

\[F_a=\operatorname{coker}(P^{a-1}\to P^a).
\tag{3.1}\]

For every test module \(T\), it computes
\(\operatorname{Tor}_1(F_a,T)=H^{a-1}(P\otimes T)=0\). The flatness
criterion makes \(F_a\) flat. Replace the lower tail by \(F_a\) in
degree \(a\). This is the good lower truncation, quasi-isomorphic to
\(K\), and supplies the required model.

Stalks commute with derived tensor and cohomology, by the flat
lesson\textquotesingle s stalk criterion. For the converse stalk test,
any module over \(\mathcal O_x\) is the stalk at \(x\) of its exact
skyscraper sheaf. Testing that sheaf proves the stalk interval from the
global interval. Conversely, test a sheaf on all stalks to prove its
tensor cohomology vanishes. Pullback of the bounded flat model remains
bounded and flat; tensor two such models to add the bounds. Cohomology
of a finite direct sum splits, proving the summand assertion. Finally,
the tensor long exact sequence shows that if \(K,L\) have intervals
\([a_K,b_K]\), \([a_L,b_L]\), a cone of \(K\to L\) has interval

\[\begin{aligned}
a_M&=\min(a_L,a_K-1),\\
b_M&=\max(b_L,b_K-1).
\end{aligned}
\tag{3.2}\]

Rotation gives the other two-out-of-three bounds. \(\square\)

\textbf{Theorem 3.2 (perfectness).} An object is perfect if and only if
it is pseudo-coherent and has locally finite Tor amplitude. More
precisely, amplitude \([a,b]\) and \((a-1)\)-pseudo-coherence already
suffice, and give local strictly perfect models with terms in \([a,b]\).

\textbf{Proof.} A strictly perfect model is an approximation at every
threshold and is a bounded flat model. This proves necessity. For the
stronger converse, choose a strictly perfect approximation \(E\to K\)
through degree \(a-1\), and brutally cut \(E\) below \(a-1\). The
approximation and \(H^{a-1}(K)=0\) show that its cone is

\[\begin{gathered}
C\simeq N[2-a],\\
N=\ker(E^{a-1}\to E^a).
\end{gathered}
\tag{3.3}\]

Indeed, in degree \(a-2\) the cone cohomology is that kernel, and all
other degrees vanish. Tensor the triangle with an arbitrary module
\(T\). The lower amplitude bound on \(K\) and the tensor long exact
sequence give

\[\begin{gathered}
\operatorname{im}(N\otimes T\to E^{a-1}\otimes T)\\
=\ker(E^{a-1}\otimes T\to E^a\otimes T).
\end{gathered}
\tag{3.4}\]

To identify the Tor group here, take a bounded-above flat resolution of
\(N\), and append \(E^{a-1}\to E^a\). This is a flat resolution of
\(E'=\operatorname{coker}(E^{a-1}\to E^a)\), whose first tensor homology
is exactly the kernel modulo the image in (3.4). Thus
\(\operatorname{Tor}_1(E',T)=0\) for all \(T\), and \(E'\) is flat. It
is finitely presented as the cokernel of a map between finite projective
sheaves. Lemma 1.2 makes it locally finite projective.

Replace the bottom pair by \(E'\) in degree \(a\). The resulting complex
is quasi-isomorphic to \(K\). Cancel its finitely many acyclic top pairs
above \(b\), as in Section 1; none of these cancellations crosses degree
\(a\). This gives the asserted strictly perfect model. Localize the
hypotheses to deduce the general equivalence. \(\square\)

\textbf{Corollary 3.3 (permanence).} Perfect objects are closed under
shifts, cones, derived tensor products, retracts and derived pullback.

\textbf{Proof.} For cones, locally represent the two inputs by strict
models and their map by an actual chain map, using Lemma 1.1. Its strict
cone is the third object. Rotation handles any two inputs. Tensor and
pullback use their strict models directly. Summands use Proposition 2.2,
Proposition 3.1 and Theorem 3.2.

A retract is an actual summand here. If \(\sigma:K\to E\) and
\(r:E\to K\) satisfy \(r\sigma=1\), complete \(\sigma\) to a triangle
\(K\to E\to C\xrightarrow{h}K[1]\). Since \(\sigma[1]h=0\), composing
with \(r[1]\) gives \(h=0\). Exactness of Hom lifts \(1_C\) to
\(t:C\to E\); replace \(t\) by \(t-\sigma rt\). Then \(rt=0\), and the
maps \((\sigma,t):K\oplus C\to E\) and \((r,E\to C)\) are inverse: their
induced cohomology maps split the triangle, so the first is an
isomorphism, and the composite on \(K\oplus C\) is the identity. The
summand result therefore applies to every retract. \(\square\)

On a quasi-compact space perfect objects have uniform bounds, by a
finite subcover. Without quasi-compactness this fails: on the discrete
space of nonnegative integers, with stalk ring \(k\), take
\(K|_{\{n\}}=k[n]\). Each neighbourhood \(\{n\}\) has a strict model,
but the global nonzero cohomology degrees \(-n\) have no lower bound.

Two useful locality consequences complete the picture. If
\(j:U\hookrightarrow X\), \(E\) is perfect on \(U\), and its cohomology
is supported on \(T\subset U\) with \(T\) closed in \(X\), then
\(Rj_*E\) is perfect. It restricts to \(E\) on \(U\) and to zero on
\(X\setminus T\): restriction of a K-injective pushforward computes the
pushforward from \(U\setminus T\), where the object is zero. These opens
cover \(X\).

If the stalk rings are local and \(E\) is perfect, the set where all
\(H^i(E)_x\) are finite free is open. Near such an \(x\) choose a strict
model in \([a,b]\). Its top cohomology is finitely presented; Lemma
1.2\textquotesingle s free-stalk argument makes it free near \(x\).
Split the epimorphism from the top term to that cohomology, then split
the preceding term onto its image. Remove the resulting contractible
pair and the free cohomology summand. Induct on the finite length. The
remaining cohomology stalks still satisfy the hypothesis, so finitely
many shrinkings make all cohomology sheaves finite locally free. This
also proves maximality of the stated open set.

\subsection{4. Duals and the two tensor--Hom
maps}\label{perfect-complexes-and-duals--4-duals-and-the-two-tensor-hom-maps}

For any object define \(K^\vee=[K,\mathcal O]\). The symmetric monoidal
structure, with graded symmetry
\(u\otimes v\mapsto(-1)^{|u||v|}v\otimes u\), was proved in the tensor
lesson\textquotesingle s Theorem 2.2. No boundedness is required for
that structure.

\textbf{Theorem 4.1 (perfect duality).} If \(K\) is perfect, so is
\(K^\vee\), and the canonical maps

\[\begin{aligned}
K&\longrightarrow K^{\vee\vee},\\
M\otimes^L K^\vee&\longrightarrow[K,M]
\end{aligned}
\tag{4.1}\]

are isomorphisms. In particular, for every integer \(r\),

\[\begin{gathered}
H^rR\Gamma(X,M\otimes^L K^\vee)\\
=\operatorname{Hom}_{D(\mathcal O)}(K,M[r]).
\end{gathered}
\tag{4.2}\]

Two further canonical maps are isomorphisms:

\[{}[K,L]\otimes^L M\longrightarrow[K,L\otimes^L M],
\tag{4.3}\] \[{}[L,M]\otimes^L K\longrightarrow[[K,L],M].
\tag{4.4}\]

\textbf{Proof.} These are globally defined maps from the Hom
lesson\textquotesingle s Section 4. Their cones restrict to zero
precisely when they are isomorphisms. Thus their invertibility can be
checked on strict-model neighbourhoods. Proposition 1.3 proves the
second map in (4.1) there, and (1.2) shows that the dual model is
strictly perfect. The double-dual differential under the ordinary
termwise identification is \(-d_P\). Consequently the canonical bidual
map in degree \(i\) is

\[p\longmapsto\bigl(\phi\mapsto(-1)^i\phi(p)\bigr).
\tag{4.5}\]

This sign makes it a chain map and each term is an isomorphism. Equation
(4.2) follows from the internal-Hom lesson\textquotesingle s
derived-section identification in every shift.

For (4.3), use \([P,L]=L\otimes P^\vee\). The map moves \(M\) past
\(P^\vee\) and associates the three factors, giving
\((L\otimes^L M)\otimes P^\vee\). It is therefore an isomorphism.
Explicitly, for homogeneous \(\phi,m,p\), its evaluation is
\((-1)^{|m||p|}\phi(p)\otimes m\); this is the same sign as moving the
dual factor past \(m\).

For (4.4), the source and target are exact functors of \(K\), and its
natural map respects their triangle comparisons, by the chain identities
and transpose construction in the Hom lesson. For \(K=\mathcal O\) it is
the identity. For a finite free sheaf it is the componentwise identity
on a finite power of \([L,M]\). It passes to finite projective summands
and shifts. Filter a bounded strict model by its terms; the two triangle
long exact sequences show inductively that the middle map is an
isomorphism when the outer two are. This proves (4.4). \(\square\)

There is a source-locator distinction: in the pinned edition, Tag 0G40
proves (4.4), the evaluation-and-more map. It does not denote (4.3).
Both assertions hold for perfect \(K\), as proved above, and both can
fail without that hypothesis.

A \textbf{left dual} of \(K\) is an object \(N\) with maps

\[\begin{aligned}
\eta:\mathcal O&\longrightarrow K\otimes^L N,\\
\epsilon:N\otimes^L K&\longrightarrow\mathcal O
\end{aligned}
\tag{4.6}\]

such that the composites

\[\begin{aligned}
K&\xrightarrow{\eta\otimes1}K\otimes N\otimes K
\xrightarrow{1\otimes\epsilon}K,\\
N&\xrightarrow{1\otimes\eta}N\otimes K\otimes N
\xrightarrow{\epsilon\otimes1}N
\end{aligned}
\tag{4.7}\]

are identities; tensors in (4.7) are derived and association constraints
are understood. Symmetry gives the corresponding right dual, so we say
\textbf{dualizable}.

\textbf{Theorem 4.2 (duals are exactly perfect objects).} Every perfect
object has dual \(K^\vee\), and every dualizable object is perfect. Its
dual is canonically isomorphic to \(K^\vee\).

\textbf{Proof: existence and identities.} The natural isomorphism
\([K,B]=B\otimes^L K^\vee\) identifies the closed tensor--Hom adjunction
with

\[-\otimes^L K\ \dashv\ -\otimes^L K^\vee.
\tag{4.8}\]

Let \(\eta\) be its unit at \(\mathcal O\), and \(\epsilon\) its counit
there. Equivalently, \(\eta\) corresponds to \(1_K\) under
\(K\otimes K^\vee\cong[K,K]\), and \(\epsilon\) is evaluation. The
counit at \(B\) is \(1_B\otimes\epsilon\), since (4.1)\textquotesingle s
map is defined by transposing this evaluation. The
adjunction\textquotesingle s first triangle identity at \(\mathcal O\)
gives the first identity in (4.7). Consequently \(1_A\otimes\eta\)
transposes to \(1_{A\otimes K}\) for every \(A\), and is the unit at
\(A\). The second adjunction identity, at \(\mathcal O\), now gives the
second identity in (4.7). These are global identities of derived maps,
not an inference that equality of maps can be tested locally.

On a strict free model,
\(\eta(1)=\sum_{i,\alpha}e_{i,\alpha}\otimes e_{i,\alpha}^*\), with no
degree sign. If
\(d e_{i,\alpha}=\sum_\beta a_{\beta\alpha}e_{i+1,\beta}\), (1.2) gives
the dual differential coefficient \((-1)^i a_{\beta\alpha}\) on
\(e_{i+1,\beta}^*\). Its tensor differential has the additional sign
\((-1)^{i+1}\), so the two contributions cancel. Evaluating either
zigzag contracts the appropriate dual basis and gives the identity.
Summands inherit this description. The sign in the bidual map (4.5) and
the absence of a sign in coevaluation are compatible: they are different
maps with different factor orders.

\textbf{Proof: converse.} Fix \(x\in X\) and a dual \(N\) with
(4.6)--(4.7). Resolve \(K\) by the actual cycle-attachment construction
of Theorem 3.1 of the flat lesson, obtaining \(P\to K\). Its graded
terms are sums of \(B_U=j_{U!}\mathcal O_U\). Stage zero consists of
identity disks; each later stage adds generators whose differentials lie
in the preceding stage. Resolve \(N\) by a K-flat complex \(Q\) with
flat terms.

By Lemma 1.1, near \(x\) the map \(\eta:\mathcal O\to P\otimes Q\) is an
actual chain map. Its value at one is locally a finite sum of
homogeneous tensors \(p_t\otimes q_t\). Each germ \(p_t\) involves
finitely many generators and a finite stage. Include those generators
and all generators occurring in their differentials. A later-stage
differential only uses earlier stages; its germ has finite support. At
stage zero include the other member of every disk encountered. Thus
closing this finite set under differentials requires finitely many steps
and produces finitely many generators in finitely many degrees.

This germ argument gives an actual local subcomplex. At each step choose
representatives of the finitely many differential germs. Their
expressions involve finitely many summands after shrinking. A summand
\(B_U\) with \(x\notin U\) has zero stalk there; kill its finitely many
section contributions on a smaller neighbourhood and omit it. For the
remaining summands shrink inside their finitely many
\(U\)\textquotesingle s, where they are copies of \(\mathcal O\).
Equality of each chosen expression with the actual differential is a
germ equality, so shrink to make all these finitely many equalities
hold. The result is a bounded finite free subcomplex

\[E\subset P|_V
\tag{4.9}\]

on a neighbourhood \(V\) of \(x\), termwise a graded summand, containing
all the \(p_t\). Tensoring this inclusion with \(Q\) is injective
termwise (also by flatness of \(Q\)). Since \(\eta(1)\) is a cycle, its
expression in \(E\otimes Q\) is a cycle. Hence \(\eta\) factors through
\(\widetilde\eta:\mathcal O\to E\otimes^L N\).

Define \(\sigma:K|_V\to E\) by \(\widetilde\eta\otimes1_K\), followed by
\(1_E\otimes\epsilon\), and let \(r:E\to K|_V\) be inclusion followed by
the resolution isomorphism. The first identity in (4.7) gives
\(r\sigma=1_K\). Corollary 3.3 makes this retract of a strictly perfect
object perfect. Every point has such a neighbourhood, proving the
converse. This uses finite free cells, rather than incorrectly assuming
that a submodule generated by finitely many sections is projective.

Finally any dual supplies an adjunction \(-\otimes K\dashv-\otimes N\):
insert \(\eta\) to transpose a map, and insert \(\epsilon\) to transpose
back; the two identities make these operations inverse. Right adjoints
are uniquely isomorphic by their natural Hom bijections (evaluate one
bijection on the identity, then the other, to obtain inverse maps).
Comparing with the closed tensor--Hom adjunction at \(\mathcal O\)
identifies \(N\) canonically with \([K,\mathcal O]\). \(\square\)

\subsection{5. Residues and finiteness
examples}\label{perfect-complexes-and-duals--5-residues-and-finiteness-examples}

\textbf{A smooth analytic surface residue is perfect.} On
\(\mathbb C^2\) let \(\mathcal O\) be holomorphic functions and \(S\)
the residue skyscraper at the origin. Its resolution, in degrees
\(-2,-1,0\), is

\[\mathcal O\xrightarrow{\binom{-y}{x}}\mathcal O^2
\xrightarrow{(x\ y)}\mathcal O\longrightarrow S.
\tag{5.1}\]

Taylor expansion shows that a germ vanishing at zero lies in \((x,y)\):
separate the terms divisible by \(x\) and those depending only on \(y\),
which are divisible by \(y\). Multiplication by \(x\) is injective on
convergent power series, and multiplication by \(y\) is injective modulo
\(x\). Thus \(xa+yb=0\) implies \(b=xc\) and \(a=-yc\), proving middle
exactness. The first map is injective. Away from zero one coordinate is
a unit and the complex contracts by division by that coordinate. This
proves exactness on every stalk, hence perfectness. With any
vector-space residue test at zero all differentials vanish, so its Tor
groups in degrees \(0,1,2\) are \(k,k^2,k\). The local Tor amplitude is
exactly \([-2,0]\).

\textbf{A singular residue need not be perfect.} Over
\(R=k[\epsilon]/(\epsilon^2)\), the residue \(k\) has the free
resolution

\[\cdots\xrightarrow{\epsilon}R\xrightarrow{\epsilon}R
\xrightarrow{\epsilon}R\longrightarrow k
\tag{5.2}\]

in nonpositive degrees. Kernel and image of multiplication by
\(\epsilon\) both equal \(\epsilon R\). Every finite brutal lower
truncation gives an approximation: it is an isomorphism in all higher
degrees and surjective at the bottom threshold. Thus \(k\) is
pseudo-coherent. Tensoring (5.2) with \(k\) kills every differential, so
\(\operatorname{Tor}_n^R(k,k)=k\) for every \(n\ge0\). It has no finite
Tor amplitude, is not perfect, and is not dualizable. The singularity
hypothesis matters; a surface residue alone does not imply
nonperfectness.

\textbf{Perfectness does not force finite cohomology over every ring.}
Let

\[\begin{gathered}
R=k[x]\oplus Y,\qquad Y=\bigoplus_{n\ge1}ky_n,\\
Y^2=0,\qquad xY=0.
\end{gathered}
\tag{5.3}\]

The two-term complex \(R\xrightarrow{x}R\), in degrees \(-1,0\), is
strictly perfect. But its negative cohomology is \(Y\), which is not
finitely generated: the action on \(Y\) factors through \(k\), and
finitely many elements span only a finite-dimensional subspace. Hence
even a perfect, and therefore pseudo-coherent, complex need not have
finite-type cohomology over a noncoherent ring. The safe assertions are
the approximation definition and Proposition 2.2\textquotesingle s
explicitly bounded top-cohomology conclusions.

\subsection{6. Exercises with checked
solutions}\label{perfect-complexes-and-duals--6-exercises-with-checked-solutions}

\textbf{Exercise 1 (easy: the shifted unit).} For \(K=\mathcal O[1]\),
compute its dual, coevaluation and bidual map, and check both identities
in (4.7).

\textbf{Solution.} Its generator \(e\) has degree \(-1\), and the dual
generator \(e^*\) has degree \(1\). Thus \(K^\vee=\mathcal O[-1]\),
\(\eta(1)=e\otimes e^*\), and \(\epsilon(e^*\otimes e)=1\). The first
zigzag sends \(e\) to \(e\otimes e^*\otimes e\) and back to \(e\); the
second does the analogous contraction for \(e^*\). No factor swap
occurs. The bidual map is \(e\mapsto-e^{**}\) by (4.5). Multiplying
coevaluation instead by \(-1\) while keeping this evaluation would make
both zigzags negative identities; that is not the correct dual data.

\textbf{Exercise 2 (medium: one approximation suffices).} Suppose \(F\)
is a module sheaf of Tor amplitude \([0,0]\). What is the weakest
threshold in Theorem 3.2 that proves perfectness? Characterize such
\(F\) directly.

\textbf{Solution.} The threshold is \(-1\), which is finite presentation
by Proposition 2.2. Amplitude \([0,0]\) means
\(\operatorname{Tor}_1(F,T)=0\) for every module \(T\), and therefore
flatness. Conversely a flat module in degree zero has that amplitude,
since tensor is exact. Lemma 1.2 now says that precisely the flat
finitely presented sheaves are locally finite projective, hence perfect
in degree zero. The threshold records finite relations as well as finite
generators; finite type alone does not supply the finite presentation
used in the splitting proof.

\textbf{Exercise 3 (medium: the coefficient sign).} Let \(P\) have one
finite free term in degree \(s\). For a homogeneous map \(\phi:P\to L\)
of degree \(h\) and an element \(m\in M^t\), calculate (4.3). Explain
why it remains an isomorphism for a bounded finite projective complex.

\textbf{Solution.} Evaluation must first move \(m\) past \(p\in P^s\),
giving \((-1)^{ts}\phi(p)\otimes m\). In the model
\(\mathsf H(P,L)=L\otimes P^\vee\), the dual term has degree \(-s\);
swapping it with \(m\) gives the same sign \((-1)^{-st}=(-1)^{st}\).
This is a tensor symmetry isomorphism. Finite free powers give
componentwise copies, finite projective summands inherit an inverse, and
bounded complexes are finite extensions of their terms. Applying the
compatible triangle maps and long exact sequences proves the general
result. K-flat models for \(L,M\) compute the derived tensors, so their
possible unboundedness causes no additional finite-sum restriction.

\textbf{Exercise 4 (hard: finite cells rather than finite generators).}
In the converse proof of Theorem 4.2, why is closing the chosen
generators under their differentials a finite process, and why does
stage zero need special treatment? Would merely taking the subcomplex
generated by the tensor entries and their differentials be sufficient?

\textbf{Solution.} Each chosen germ has finite support in a finite
stage. A cell from a positive stage has differential in strictly earlier
stages, so following its dependencies decreases a nonnegative integer.
Every branch has finite length and every vertex finitely many
descendants, giving a finite set. A stage-zero disk has two generators
and its lower generator\textquotesingle s differential is its upper mate
in the same stage; including the mate completes that branch, since its
differential is zero. Finite intersection of their defining opens makes
the remaining cells free on one neighbourhood. Contributions with zero
germ are killed by finitely many shrinkings. The span of this selected
subset of graded free cells is a bounded finite free complex and a
graded summand. In contrast, finitely many arbitrary sections may
generate a nonprojective submodule: the ideal \(\epsilon R\subset R\)
over the dual numbers is generated by one element but is isomorphic to
the nonflat residue \(k\). Such a subcomplex would not supply a perfect
object for the retract argument. Selecting actual cells repairs exactly
this problem.

\textbf{Exercise 5 (hard: global duality with no uniform bound).} On the
discrete space from Section 3, construct the dual of \(K|_{\{n\}}=k[n]\)
explicitly and check that the global coevaluation is defined although it
has infinitely many degree components across the space.

\textbf{Solution.} Take \(N|_{\{n\}}=k[-n]\). At point \(n\), send one
to the tensor of its degree-\(-n\) generator and the degree-\(n\) dual
generator, and evaluate that pair to one. Sheaves on a discrete space
are specified independently at the points; these formulas define sheaf
maps. The tensor totalization uses sheaf direct sums, whose sections are
locally finite families, not necessarily globally finite families. The
singleton neighbourhoods make this family locally finite, so
coevaluation is a legitimate section of the tensor sheaf. Both zigzags
are identities on every singleton and hence are identities of these
actual sheaf maps. This example is perfect and dualizable but has no
global finite cohomological bound. It shows why neither definition may
silently impose a uniform interval.

\section{The projection formula and the base change
map}\label{projection-formula-and-base-change}

\emph{Written and edited by GPT-6.1 Sol (OpenAI), in Codex at Ultra,
October 2026. Mathematically self-checked by the writing AI. Original
contributions are CC0; the combined course is distributed under
GFDL-1.2-or-later. Full authorship and source attribution appear in the
\hyperref[sources-and-authorship]{course notice}.}

The projection formula compares tensoring after pushforward with
tensoring before it. Base change compares pullback after pushforward
with pushforward after pullback. Each comparison has a canonical
direction, supplied by adjunction. Their existence is general; their
invertibility needs hypotheses. Keeping these two questions separate
prevents a flatness assumption from being mistaken for a universal
base-change theorem.

The prerequisites are
\hyperref[derived-pullback-and-pushforward]{Derived pullback and
pushforward}, Theorem 1.1, Proposition 1.2, Theorem 2.2, Propositions
3.1--3.2 and 4.2; \hyperref[perfect-complexes-and-duals]{Perfect
complexes and duals}, Lemma 1.1, Corollary 3.3 and Theorems 4.1--4.2;
\hyperref[derived-tensor-products-and-tor-sheaves]{The derived tensor
product and Tor sheaves}, Theorems 1.2 and 2.2; and the exact pullback
and closed-support arguments in the first lesson\textquotesingle s
Theorem 4.4 and the support lesson\textquotesingle s Proposition 1.1. We
use maps into K-injectives from the bounded-derived
lesson\textquotesingle s Theorem 3.3 and unbounded resolutions from the
K-injective lesson\textquotesingle s Theorem 4.1.

All structure sheaves are commutative and unital. Objects may be
unbounded. Let \(f:X\to Y\) be a morphism of ringed spaces. Every tensor
in a derived formula below is over the structure sheaf of the indicated
space. Ordinary tensors and functors are expressly distinguished when
models are used.

The construction follows the Stacks project authors' \emph{Cohomology of
Sheaves}, ``Projection formula'' (Tags 01E7, 01E8, 0B54 and 0B55), ``The
base change map'' (Tag 02N7), and the unbounded base-change and
composition results, in the
\href{https://github.com/KokunoYumeto/unofficial-stacks-project-ai-drafts/blob/565b10e987aba5969b21145a0833f42d69f96790/cohomology.tex}{AI
Integrated Stacks Project edition}. We prove the local identifications
and compatibility diagrams, keeping the existence and invertibility
hypotheses explicit. Source attribution and the licence for adapted
passages appear in the \hyperref[sources-and-authorship]{course notice}.

\subsection{1. Finite locally free
coefficients}\label{projection-formula-and-base-change--1-finite-locally-free-coefficients}

\textbf{Lemma 1.1.} If \(V\) is finite locally free and \(I\) is an
injective module sheaf on a ringed space, then \(V\otimes I\) is
injective. Tensoring a K-injective complex with \(V\) also gives a
K-injective complex.

\textbf{Proof.} Put \(V^*=\mathcal Hom(V,\mathcal O)\). The ordinary
dual-basis evaluation maps give a natural bijection

\[\begin{gathered}
\operatorname{Hom}(A,V\otimes I)\\
\cong\operatorname{Hom}(A\otimes V^*,I).
\end{gathered}
\tag{1.1}\]

Indeed, this is the finite free tensor--Hom identity on a trivializing
open; the canonical evaluation and coevaluation maps define inverse
operations globally, so the local descriptions glue. Both \(V\) and
\(V^*\) are flat by the stalk criterion. For a monomorphism \(A\to B\),
tensoring with \(V^*\) remains monic. Injectivity of \(I\), followed by
(1.1), extends every map \(A\to V\otimes I\) across \(B\). This proves
injectivity.

The same degree-zero duality induces a bijection of chain maps modulo
homotopy from a complex \(A\) to \(V\otimes I^\bullet\), and from
\(A\otimes V^*\) to \(I^\bullet\). An acyclic \(A\) stays acyclic after
this exact tensor. The K-injective acyclic test therefore proves the
second assertion. The ranks can vary on the space; global constancy of
rank was not used. \(\square\)

\textbf{Theorem 1.2 (finite coefficient projection).} For a finite
locally free sheaf \(V\) on \(Y\) and any \(E\in D(\mathcal O_X)\),
there is a natural isomorphism

\[Rf_*E\otimes V\ \cong\ Rf_*(E\otimes f^*V).
\tag{1.2}\]

For a module sheaf \(F\) and \(q\ge0\), this induces

\[\begin{gathered}
R^qf_*F\otimes V\\
\cong R^qf_*(F\otimes f^*V).
\end{gathered}
\tag{1.3}\]

\textbf{Proof.} The ordinary projection map is the adjoint of the tensor
of the ordinary counit \(f^*f_*F\to F\) with \(f^*V\). On an open
\(W\subset Y\) trivializing \(V\) as \(\mathcal O_W^r\), put
\(F_W=F|_{f^{-1}W}\). Its restriction identifies with

\[(f_W)_*(F_W)^r\longrightarrow(f_W)_*(F_W^r).
\tag{1.4}\]

It is the identity on each component: the adjoint of each component
inclusion is the corresponding counit. Finite sums equal finite
products, and pushforward preserves them. Hence (1.4) is an isomorphism.
Since the map was globally defined by adjunction, these local
isomorphisms prove its global invertibility without choosing compatible
trivializations.

Choose a K-injective resolution \(E\to I\). The pullback \(f^*V\) is
finite locally free, and Lemma 1.1 makes \(I\otimes f^*V\) K-injective.
Its exact tensor represents \(E\otimes f^*V\). Applying the ordinary
projection isomorphism termwise gives

\[f_*I\otimes V\cong f_*(I\otimes f^*V),
\tag{1.5}\]

which computes both sides of (1.2). All maps are natural in chain maps
and homotopies, hence in derived maps through their K-injective
representatives. For \(F\) in degree zero use its bounded-below
injective resolution; Lemma 1.1 makes the tensor another injective
resolution. Flatness of \(V\) commutes its tensor with kernels, images
and cohomology, giving (1.3). This proves the source\textquotesingle s
bounded-below assertion and the stronger unbounded version. \(\square\)

\subsection{2. The general map and perfect
coefficients}\label{projection-formula-and-base-change--2-the-general-map-and-perfect-coefficients}

For any \(E\in D(\mathcal O_X)\), \(K\in D(\mathcal O_Y)\), define

\[\begin{gathered}
\pi_{f,E,K}:Rf_*E\otimes^L K\\
\longrightarrow Rf_*(E\otimes^L Lf^*K)
\end{gathered}
\tag{2.1}\]

as the adjoint, under \(Lf^*\dashv Rf_*\), of

\[\begin{gathered}
Lf^*(Rf_*E\otimes^L K)\\
\cong Lf^*Rf_*E\otimes^L Lf^*K\\
\xrightarrow{\epsilon_f\otimes1}E\otimes^L Lf^*K.
\end{gathered}
\tag{2.2}\]

Here \(\epsilon_f\) is the derived counit. The proven tensor comparison
and the adjunction supply every arrow for arbitrary unbounded inputs.
This is a definition of a global map, not an attempt to glue local
derived morphisms.

Naturality follows by taking adjoints: in either variable the relevant
square transposes to the naturality square of \(\epsilon_f\), with the
same tensor identity on the other factor. The same argument proves
compatibility with shifts and triangles. More concretely, the resolution
functors, tensor cone comparison and adjunction identify the transposes
of all three arrows in the cone triangle; they agree with those in
(2.2). Thus (2.1) is a morphism of the two triangle-valued functors in
\(K\). The map for \(K=\mathcal O_Y\) is the identity of \(Rf_*E\): its
transpose is \(\epsilon_f\), exactly the transpose of that identity.

\textbf{Theorem 2.1 (perfect projection formula).} If \(K\) is perfect
on \(Y\), (2.1) is an isomorphism for every \(E\).

\textbf{Proof.} Restrict to an open \(W\subset Y\), replacing \(X\) by
\(f^{-1}W\). Proposition 3.2 of the pullback--pushforward lesson
identifies the restricted pushforward, while its Theorem 1.1 and
Proposition 1.2 identify the restricted pullbacks and tensors. The
counit restricts to the same counit by its defining Hom bijection; hence
(2.1) restricts to the projection map for \(f_W\).

Now work on a strict-model neighbourhood for \(K\). For \(\mathcal O_Y\)
the map is the identity. Finite direct sums preserve that assertion, and
naturality with the inclusion and retraction of a finite projective
summand makes the map on the summand an isomorphism too. Shifts preserve
it by the comparison above. Filter the bounded strict complex by its
brutal lower tails. Each successive triangle has one term equal to a
finite projective sheaf in one degree. The compatible two triangle
sequences, and their long exact cohomology sequences, prove the
assertion inductively on its finite number of terms. The cone of (2.1)
is consequently zero on an open cover, so all its cohomology sheaves
vanish globally. Thus the map is an isomorphism. \(\square\)

This proof does not require \(f\) to be proper, flat, separated or of
finite cohomological dimension. Finiteness is in the coefficient \(K\),
and it is local finiteness. In particular, a perfect coefficient without
a uniform bound on a noncompact target is allowed.

\subsection{3. Closed inclusions with arbitrary
coefficients}\label{projection-formula-and-base-change--3-closed-inclusions-with-arbitrary-coefficients}

\textbf{Theorem 3.1.} If the underlying map of \(f:X\to Y\) is a
homeomorphism onto a closed subset, the projection map (2.1) is an
isomorphism for every \(E,K\). The coefficient map need not be flat or
an isomorphism.

\textbf{Proof.} Identify the underlying space \(X\) with its image
\(Z\). At \(y\notin Z\), a neighbourhood disjoint from \(Z\) makes
\((f_*F)_y=0\). At \(y=f(x)\), intersections of neighbourhoods in \(Y\)
with \(Z\) are cofinal among neighbourhoods in \(X\), so
\((f_*F)_y=F_x\), with its action restricted from \(\mathcal O_{Y,y}\).
These stalk formulas prove that \(f_*\) is exact and commutes with
arbitrary sheaf direct sums: the respective comparison maps are exact or
isomorphisms on every stalk. This remains true for arbitrary structure
sheaves and their given coefficient homomorphism.

For module sheaves the ordinary projection map is an isomorphism on
stalks. Outside \(Z\) both sides vanish; at \(y=f(x)\), with
\(R=\mathcal O_{Y,y}\), \(S=\mathcal O_{X,x}\), it is the balanced
isomorphism

\[\begin{gathered}
F_x\otimes_R M_y\\
\longrightarrow F_x\otimes_S(S\otimes_R M_y),\\
e\otimes m\longmapsto e\otimes(1\otimes m).
\end{gathered}
\tag{3.1}\]

Its inverse sends \(e\otimes(s\otimes m)\) to \(se\otimes m\). Balancing
verifies both inverse composites, without flatness.

Choose a K-flat model \(P\) for \(K\), and any representative \(G\) of
\(E\). Pullback \(f^*P\) is K-flat. Exactness of \(f_*\) makes the right
side of (2.1) represented by \(f_*\operatorname{Tot}(G\otimes f^*P)\),
and the left by \(\operatorname{Tot}(f_*G\otimes P)\). The ordinary
projection maps in each pair of degrees give a chain isomorphism between
them: (3.1) is an isomorphism and \(f_*\) commutes with the
totalization\textquotesingle s sums.

This is the canonical derived map. To check it, take a K-flat resolution
\(Q\to f_*G\). The map \(Q\otimes P\to f_*G\otimes P\) is a
quasi-isomorphism. Under the ordinary adjunction the chain projection
map, precomposed with this resolution, transposes to the map

\[f^*Q\otimes f^*P\longrightarrow G\otimes f^*P
\tag{3.2}\]

from the adjoint \(f^*Q\to G\). That adjoint represents the derived
counit: comparing \(G\to I\) with a K-flat resolution of \(f_*I\), and
refining the two roofs by the common reading\textquotesingle s fraction
theorem, gives a commuting-up-to-homotopy square; its adjoint is
precisely the counit construction in the pullback--pushforward
lesson\textquotesingle s Theorem 2.2. Thus (3.2) is (2.2). The
isomorphism proved on these models is (2.1). \(\square\)

For a one-point coefficient homomorphism \(R\to S\), this says

\[E\otimes_R^L K\cong E\otimes_S^L(S\otimes_R^L K),
\tag{3.3}\]

where on the left \(E\) is regarded as an \(R\)-complex. Even
\(R\to R/I\) is allowed. On more general spaces the closed image is what
ensures exact pushforward and its stalk description. An open inclusion
has different boundary stalks, as Section 5 will demonstrate.

\subsection{4. Base change, flat models and
compatibility}\label{projection-formula-and-base-change--4-base-change-flat-models-and-compatibility}

Consider a commutative square of ringed spaces

\[\begin{array}{ccc}
X'&\xrightarrow{g'}&X\\
{\scriptstyle f'}\downarrow&&\downarrow{\scriptstyle f}\\
Y'&\xrightarrow{g}&Y.
\end{array}
\tag{4.1}\]

Commutative includes the coefficient homomorphisms, not just the
continuous maps. The square need not be Cartesian.

\textbf{Theorem 4.1 (canonical base change).} For every
\(E\in D(\mathcal O_X)\) there is a natural map

\[b_E:Lg^*Rf_*E\longrightarrow Rf'_*Lg'^*E.
\tag{4.2}\]

It respects identity squares, horizontal and vertical composition of
squares, and the projection maps. If both \(g,g'\) are flat, its
construction uses ordinary exact pullbacks, for unbounded as well as
bounded-below complexes. These statements construct and compare maps;
they assert no general invertibility.

\textbf{Proof: definition and naturality.} Let \(A=Lg^*Rf_*E\). Define
\(b_E\) to be the adjoint under \(Lf'^*\dashv Rf'_*\) of

\[\begin{gathered}
Lf'^*A\cong Lg'^*Lf^*Rf_*E\\
\xrightarrow{Lg'^*\epsilon_{f,E}}Lg'^*E.
\end{gathered}
\tag{4.3}\]

The pullback composition comparison supplies the isomorphism. Naturality
of the counit and of this comparison proves naturality of \(b\) for
every derived map, by taking its transpose. In particular, no choice of
representatives occurs in the definition.

The equivalent pushforward description first takes the unit

\[E\longrightarrow Rg'_*Lg'^*E,\]

pushes it through \(Rf_*\), identifies \(Rf_*Rg'_*\) with \(Rg_*Rf'_*\)
by composition, and takes its adjoint under \(Lg^*\dashv Rg_*\). To
verify equivalence, take that map\textquotesingle s further transpose
under \(Lf'^*\). Successive Hom bijections identify the composite
adjunctions for \(fg'=gf'\); cancel the unit--counit pairs using their
triangle identities. What remains is \(Lg'^*\epsilon_{f,E}\), which is
(4.3). This also identifies exactly which units are meant in model
calculations.

\textbf{Proof: flat construction.} Flatness of a ringed-space morphism
means that the target-to-source map on each relevant stalk makes the
source stalk flat. The stalk formula for pullback then makes \(g^*\) and
\(g'^*\) exact, so they compute \(Lg^*,Lg'^*\) on any representatives.
Choose \(E\to I\) K-injective and \(g'^*E\to J\) K-injective.
Pushforward \(g'_*J\) is K-injective because its left adjoint \(g'^*\)
is exact, as proved in Proposition 4.2 of the pullback--pushforward
lesson. The ordinary unit followed by the resolution gives a chain map
\(E\to g'_*J\). Maps into K-injectives give a unique homotopy class

\[\beta:I\longrightarrow g'_*J
\tag{4.4}\]

extending it. Push (4.4) through \(f_*\) and use \(f_*g'_* = g_*f'_*\).
Ordinary adjunction gives the chain map

\[g^*f_*I\longrightarrow f'_*J.
\tag{4.5}\]

The complexes in (4.5) represent the two sides of (4.2); a homotopy of
\(\beta\) induces a homotopy of (4.5). Comparison between K-injective
resolutions therefore makes this a well-defined natural derived map. A
pullback of an injective was never asserted injective.

Here is a direct check that (4.5) agrees with (4.3). Resolve \(f_*I\) by
a K-flat \(Q\). Since \(g\) is flat, \(g^*Q\to g^*f_*I\) is a
quasi-isomorphism, and \(g^*Q\) is K-flat. The transpose of (4.5), on
that model, is

\[\begin{gathered}
f'^*g^*Q\cong g'^*f^*Q\\
\longrightarrow g'^*I\longrightarrow J.
\end{gathered}
\tag{4.6}\]

The first arrow after the isomorphism is \(g'^*\) of the adjoint to
\(Q\to f_*I\), hence of the derived counit. The last arrow is adjoint to
\(\beta\). Its composite with \(g'^*E\to g'^*I\) is the chosen
resolution map to \(J\), up to homotopy, by (4.4). Flatness makes that
first map a quasi-isomorphism, so the last arrow represents the required
identification with \(Lg'^*E\). Thus (4.6) is exactly (4.3).

For bounded-below \(E\), injective resolutions may be bounded below with
injective terms. This specializes to Tag 02N7: \textbf{both horizontal
maps are flat and the conclusion is a canonical map}. Its proof does not
turn a Cartesian square, flat horizontal arrows, or any combination of
those conditions into an isomorphism claim.

\textbf{Proof: pasting squares.} For any adjunction \(L\dashv R\), the
transpose of \(u:A\to RB\) is \(\epsilon_B\circ L(u)\). We use this
formula and (4.3). In an identity square (4.3) is a counit, so (4.2) is
the identity.

For horizontal pasting, add a square with \(h:Y''\to Y'\),
\(h':X''\to X'\) and \(f'':X''\to Y''\). The proposed composite is

\[\begin{gathered}
Lh^*Lg^*Rf_*E\\
\longrightarrow Lh^*Rf'_*Lg'^*E\\
\longrightarrow Rf''_*Lh'^*Lg'^*E.
\end{gathered}
\tag{4.7}\]

Transpose under \(Lf''^*\dashv Rf''_*\). The defining equation for the
second base-change map replaces its pullback followed by the counit with
\(Lh'^*\epsilon_{f'}\). The defining equation for the first then
replaces \(\epsilon_{f'}\circ Lf'^*(b_E)\) with \(Lg'^*\epsilon_f\). The
result is \(Lh'^*Lg'^*\epsilon_f\), with the coherent
pullback-composition identification. This is the transpose for the outer
rectangle, proving equality of the maps.

For vertical pasting take \(p:Y\to Z\), \(p':Y'\to Z'\), and
\(h:Z'\to Z\), with \(pg=hp'\). The composite first applies base change
for \(p\) to \(Rf_*E\), then \(Rp'_*\) to the base change for \(f\). The
counit of the composite adjunction \(L(pf)^*\dashv R(pf)_*\), under the
composition isomorphisms, is

\[\begin{gathered}
Lf^*Lp^*Rp_*Rf_*E\\
\xrightarrow{Lf^*\epsilon_p}Lf^*Rf_*E\\
\xrightarrow{\epsilon_f}E.
\end{gathered}
\tag{4.8}\]

This formula follows by composing the two Hom adjunction bijections: the
identity of the composed right-adjoint value transposes successively
through those two counits. Transposing the proposed pasted map uses the
same two factors, pulled back through \(g'\). It is therefore the
transpose (4.3) for the rectangle. This proves vertical pasting without
an omitted diagram verification.

\textbf{Proof: projection compatibility.} Fix \(K\in D(\mathcal O_Y)\).
The two routes from

\[Lg^*(Rf_*E\otimes^L K)\]

to

\[Rf'_*(Lg'^*E\otimes^L Lf'^*Lg^*K)
\tag{4.9}\]

are as follows. One pulls back \(\pi_f\), applies base change to
\(E\otimes^L Lf^*K\), then uses tensor and pullback composition. The
other uses tensor compatibility first, applies \(b_E\otimes1\), then
\(\pi_{f'}\). Transpose both under \(Lf'^*\). On the common source model
the first route is \(Lg'^*\epsilon_{f,E}\otimes1\), with the coefficient
factor \(Lg'^*Lf^*K\). The second uses

\[\epsilon_{f',Lg'^*E}\circ Lf'^*(b_E)
=Lg'^*\epsilon_{f,E},
\tag{4.10}\]

which is precisely (4.3)\textquotesingle s defining transpose equation,
and then tensors the same coefficient identity. The two transposes agree
globally, proving compatibility. All tensor reassociations here are the
coherent ones established earlier. No invertibility of \(b_E\) was
needed. \(\square\)

Projection maps also compose for \(X\xrightarrow fY\xrightarrow pZ\).
First apply \(\pi_p\) to \(Rf_*E\) and \(K\), then \(Rp_*\pi_f\) with
coefficient \(Lp^*K\). Their transpose under the composite pullback is
the tensor of (4.8) with \(L(pf)^*K\), hence is the transpose of
\(\pi_{pf}\). The maps are equal by adjunction. This observation and
Theorem 4.1 ensure that iterating these constructions changes neither
their direction nor their definition.

\subsection{5. Two boundary failures and a coefficient
calculation}\label{projection-formula-and-base-change--5-two-boundary-failures-and-a-coefficient-calculation}

Let \(Y=\{0\}\cup\{1/n:n\ge1\}\) have its subspace topology in
\(\mathbb R\), let \(U=Y\setminus\{0\}\), and let
\(j:U\hookrightarrow Y\). Give both spaces the constant coefficient
sheaf \(k\), for a field \(k\). Every point of \(U\) is isolated. A
sheaf there is a family of vector spaces, and its sections on an open
are their product. Products of vector-space short exact sequences are
exact: choose a preimage in each component to prove surjectivity;
kernels are componentwise. Consequently \(j_*\) is exact on these
sheaves, and its derived functor is computed on any complex by ordinary
pushforward.

Take \(F=k_U\). At the boundary point the neighbourhoods contain all
sufficiently large \(n\), so

\[\begin{gathered}
(j_*F)_0=\varinjlim_N\prod_{n\ge N}k\\
\cong\left(\prod_{n\ge1}k\right)\Big/\left(\bigoplus_{n\ge1}k\right).
\end{gathered}
\tag{5.1}\]

Indeed, extending a tail by zero proves surjectivity from the full
product to this colimit, and a sequence dies exactly when a tail is
zero, equivalently when its support is finite. Write this quotient as
\(A\); the class of the constant sequence one is nonzero.

\textbf{Failure with a flat infinite coefficient.} Put
\(K=\bigoplus_{m\ge1}k_Y\), in degree zero. It is flat, but not perfect,
since its stalk at any point is an infinite-dimensional vector space and
Proposition 2.2 of the perfect lesson says a perfect module is of finite
type. All derived tensors in this example are ordinary. The projection
map at zero is

\[\begin{gathered}
A\otimes_k\left(\bigoplus_{m\ge1}k\right)\\
\longrightarrow\varinjlim_N\prod_{n\ge N}\left(\bigoplus_{m\ge1}k\right).
\end{gathered}
\tag{5.2}\]

An element on the left uses only finitely many coordinate indices \(m\),
uniformly over its sequence entries. On the right take the sequence
whose \(n\)-th value is the \(n\)-th basis vector \(e_n\). It is a valid
section on the discrete \(U\). No tail of that sequence has its values
in a fixed finite coordinate subspace. Hence its germ is not in the
image of (5.2), and the projection map is not an isomorphism, even in
degree zero. The failure is about finite coefficients, not about
flatness: this coefficient is flat, and there is no higher-cohomology
complication hiding the failure.

\textbf{Failure in a Cartesian flat square.} Base change the same open
inclusion along \(g:\{0\}\hookrightarrow Y\), with coefficient \(k\) at
the point. The Cartesian top-left space is empty, so
\(g':\varnothing\to U\) and \(f':\varnothing\to\{0\}\). The stalk
coefficient map for \(g\) is \(k\to k\), hence flat; \(g'\) is flat
vacuously. For \(F\), (4.2) becomes

\[A[0]\longrightarrow0.
\tag{5.3}\]

This is not an isomorphism because \(A\ne0\). Thus both flat horizontal
maps and a Cartesian square still do not imply base change for arbitrary
ringed spaces and pushforward. Every hypothesis has been checked in this
example; no appeal to a vague exceptional topology is needed.

The contrast with restriction to an open is instructive. If
\(g:V\hookrightarrow Y\) is open and \(X'=f^{-1}V\) has restricted
coefficients, the base-change map is an isomorphism. On a K-injective
\(I\), restriction to \(X'\) remains K-injective, and the two sides are
literally the complexes

\[(f_*I)|_V=(f_V)_*(I|_{X'}).
\tag{5.4}\]

Their section values on every smaller open coincide. The map (4.3) gives
this same equality: its transpose is the restricted counit. This is
precisely Proposition 3.2 of the pullback--pushforward lesson. Replacing
an open neighbourhood by a boundary point removes the sectionwise
equality that proves (5.4), which explains (5.3).

\textbf{A nonflat closed coefficient map.} On one-point spaces take
\(R=\mathbb Z\), \(S=\mathbb Z/2\), \(E=S\), and \(K=\mathbb Z/2\). The
free resolution \(\mathbb Z\xrightarrow{2}\mathbb Z\) in degrees
\(-1,0\) represents \(K\). Its scalar extension is
\(S\xrightarrow{0}S\). Thus both sides of (3.3) have one copy of \(S\)
in each of degrees \(-1,0\), and the projection isomorphism is the
identity on those copies. Replacing derived pullback by ordinary
pullback would lose the negative cohomology. Closedness guarantees this
projection formula even though the coefficient map is nonflat; it does
not erase its Tor groups.

\subsection{6. Exercises with checked
solutions}\label{projection-formula-and-base-change--6-exercises-with-checked-solutions}

\textbf{Exercise 1 (easy: pushforward to a point).} Let \(f:X\to\{*\}\)
have coefficient ring \(R\) at the point. Show the projection formula
for \(K=R^r\), with \(r\) finite, directly on a K-injective model.
Explain why this calculation does not justify arbitrary infinite free
coefficients.

\textbf{Solution.} For \(E\to I\) K-injective,
\(Lf^*R^r=\mathcal O_X^r\), and \(I^r\) is K-injective: finite Hom
groups split into finite products, so the acyclic test vanishes in each
component. Sections commute with finite direct sums, which are finite
products. Both sides are therefore represented by \(\Gamma(X,I)^r\), and
the map is the identity in every degree and on each component. For
infinite rank a tensor uses a sheaf direct sum. Its section functor need
not commute with that sum, and its pushforward at a boundary may include
families with coordinate support varying without a common finite bound.
Formula (5.2) supplies a checked counterexample to exactly that
inference. The finite proof uses finiteness twice, in the direct-sum
identity and in the coefficient dual.

\textbf{Exercise 2 (medium: more general degree-zero coefficients).}
Extend (1.3) to a locally finite projective sheaf \(V\) that need not be
finite locally free. State the needed ring hypothesis, and check
independence of the chosen splitting.

\textbf{Solution.} The rings remain commutative, but their stalks need
not be local. On an open neighbourhood choose \(V\) as a summand of
\(\mathcal O^r\), with inclusion \(a\) and retraction \(b\). The
ordinary projection maps commute with \(a,b\) because they are defined
by the counit and tensor functoriality. Their map on the finite power is
an isomorphism; restricting it and its inverse to the summand therefore
gives an isomorphism on \(V\). The global projection map is already
defined by adjunction, so a second splitting gives the same map, rather
than a new choice of gluing data. The pullback remains locally finite
projective and flat. Its ordinary duality gives (1.1), hence preserves
injectives; alternatively this follows locally from finite projective
dual bases, with globally canonical maps. Applying these facts to an
injective resolution gives (1.3) on every member of the cover, hence
globally. This includes the nonfree module \((1,0)(k\times k)\) on the
one-point example in the perfect lesson.

\textbf{Exercise 3 (medium: a two-term strict model).} Prove the
projection formula for \(K=(P^{-1}\xrightarrow dP^0)\), where both terms
are finite projective. Specify the triangle and the comparison that is
used; a reference to ``devissage'' alone is insufficient.

\textbf{Solution.} Inclusion of the top term gives the triangle

\[\begin{gathered}
P^0[0]\longrightarrow K\longrightarrow P^{-1}[1]\\
\longrightarrow P^0[1].
\end{gathered}
\tag{6.1}\]

It comes from the degreewise split short exact sequence of the brutal
tails, using the common reading\textquotesingle s cone convention. Apply
the exact functors \(Rf_*E\otimes^L-\) and \(Rf_*(E\otimes^L Lf^*(-))\).
Naturality and triangle compatibility of (2.1) give a map between these
triangles, including their last arrows. The maps on \(P^0\) and
\(P^{-1}[1]\) are isomorphisms by the finite-power, summand and shift
cases. The long exact cohomology sequences then give an isomorphism on
the middle object: injectivity follows by lifting a zero image to the
preceding term and cancelling its isomorphic counterpart; surjectivity
follows by lifting the following-term image and correcting by the
preceding-term image. Thus the middle cone is acyclic. This proves the
two-term case for any differential, rather than only for a split
complex.

\textbf{Exercise 4 (hard: a global duality proof).} Give a second proof
of Theorem 2.1 using perfect duality and the derived adjunction, without
choosing a global strict model for \(K\).

\textbf{Solution.} Pullback preserves the evaluation, coevaluation and
their identities, because it is a monoidal functor with coherent
comparisons. Hence \(Lf^*(K^\vee)\) is the dual of \(Lf^*K\); uniqueness
of the dual identifies it with \((Lf^*K)^\vee\). Write \(A=Rf_*E\),
\(P=Lf^*K\), and \(T_X=Lf^*T\). For every \(T\in D(\mathcal O_Y)\), the
dual and pullback adjunctions give successive natural bijections

\[\begin{gathered}
\operatorname{Hom}(T,A\otimes K)\\
\cong\operatorname{Hom}(T\otimes K^\vee,A)\\
\cong\operatorname{Hom}(T_X\otimes P^\vee,E)\\
\cong\operatorname{Hom}(T_X,E\otimes P)\\
\cong\operatorname{Hom}(T,Rf_*(E\otimes P)).
\end{gathered}
\tag{6.2}\]

The tensor symbols in this calculation are derived. The first bijection
evaluates the coefficient dual, and the third inserts its coevaluation;
the inverse operations use the other two maps and their triangle
identities. The middle bijection uses the counit of \(f\), and pullback
identifies the dual data just described. Composing the operations, the
dual evaluation--coevaluation pairs cancel. The resulting map on Hom
groups is composition with the map whose pullback transpose is
\(\epsilon_f\otimes1\). By (2.2), that map is exactly \(\pi_{f,E,K}\).
Thus composition with \(\pi\) is a bijection for all \(T\). With \(T\)
equal to its target, surjectivity supplies a right inverse; with \(T\)
equal to its source, injectivity makes that right inverse also a left
inverse. This proves the canonical map is an isomorphism globally. It
also explains why no uniform bound or global free presentation was
necessary.

\textbf{Exercise 5 (medium: Cartesian flat change on one-point spaces).}
Let \(R\to B\) and a flat homomorphism \(R\to R'\) be given, and let
\(B'=B\otimes_RR'\). Form the square of one-point ringed spaces with
these coefficient rings. Prove that its base-change map is an
isomorphism for every unbounded \(B\)-complex, and explain why this does
not contradict (5.3).

\textbf{Solution.} Both pushforwards are restriction of scalars, hence
exact. The map \(B\to B'\) is flat: for a short exact sequence of
\(B\)-modules, tensoring with \(B'\) over \(B\) identifies with
tensoring the same underlying \(R\)-sequence with the flat \(R'\). Thus
both horizontal pullbacks are exact. For any complex \(E\), the two
base-change values are represented by \(R'\otimes_RE\) and
\(B'\otimes_BE\). Their canonical balanced isomorphism sends
\(r'\otimes e\) to \((1\otimes r')\otimes e\), with inverse
\((b\otimes r')\otimes e\mapsto r'\otimes be\). It commutes with
differentials. The counit for restriction and extension of scalars is
multiplication, so its transpose is precisely this balanced map. Theorem
4.1 therefore identifies it with base change. These particular
pushforwards have no boundary-section phenomenon. Flatness is sufficient
in this coefficient-only situation, while the topological example (5.3)
has nonzero boundary germs over an empty fibre. Neither example licenses
replacing the other\textquotesingle s hypotheses.

\textbf{Exercise 6 (hard: a nonflat algebraic square).} In the preceding
square take \(R=\mathbb Z\), \(B=R'=B'=\mathbb Z/2\), and \(E=B[0]\).
Compute the general derived base-change map and compare it with ordinary
pullback.

\textbf{Solution.} The lower horizontal coefficient homomorphism
\(\mathbb Z\to\mathbb Z/2\) is not flat: tensoring the injection
\(2:\mathbb Z\to\mathbb Z\) makes it zero. The upper horizontal
homomorphism \(B\to B'\) is the identity, hence flat. The source of
(4.2) is \((\mathbb Z/2)\otimes^L_{\mathbb Z}(\mathbb Z/2)\),
represented by \(B\xrightarrow0B\) in degrees \(-1,0\); the target is
\(B[0]\). The transpose description is scalar-extension multiplication.
On this model it is the identity on degree zero and zero on degree
\(-1\), so it is the usual augmentation. It induces an isomorphism on
\(H^0\) but kills the nonzero \(H^{-1}=B\). Thus it is not a
quasi-isomorphism. Ordinary pullback would keep only the degree-zero
tensor and obscure the failure. The square is Cartesian with the
ordinary tensor coefficient ring, but that ring records no derived Tor
correction. Theorem 4.1 supplies the comparison, while the calculation
states exactly why it is not invertible.

\section{Sources and authorship}\label{sources-and-authorship}

This course is distributed under the GNU Free Documentation License,
version 1.2 or any later version, with no Invariant Sections,
Front-Cover Texts or Back-Cover Texts. Permission is granted to copy,
distribute and modify this document under those terms. The full licence
is included in \hyperref[gnu-free-documentation-license]{COPYING}.

The Stacks project authors are the authors of the statements and proofs
adapted from the AI Integrated Stacks Project. The source edition is
commit \texttt{565b10e987aba5969b21145a0833f42d69f96790}, with upstream
Stacks commit \texttt{a04446e57ec1fbc252a871afcec7752fb2807b14}. These
passages are GFDL-1.2-or-later. AI Integrated Stacks Project is an
unofficial edition, not affiliated with, approved by or reviewed by the
Stacks project\textquotesingle s maintainers.

The edition records corrections by GPT-5.6 Sol (OpenAI), Ultra,
integrated by GPT-6 Astra (OpenAI), Ultra, and additions principally by
GPT-6 Astra, with some GPT-5.6 Sol contributions checked by GPT-6 Astra.
The course\textquotesingle s module-sheaf exposition includes Claude
Opus 5.5 (Anthropic) contributions. Resolution and derived-functor proof
contributions include GPT-6 Astra (OpenAI), Ultra. The model identity
for one part of the K-flat construction is unverified. Course
exposition, proof expansions and mathematical self-checking are by
GPT-6.1 Sol (OpenAI), Ultra, October 2026. Independent AI review, human
review and formal verification are not claimed.

Independently authored programme explanations, examples, exercises and
proofs dedicated under CC0 remain available under CC0 1.0. That
dedication does not change the terms of the Stacks-derived passages. The
combined course uses GFDL-1.2-or-later.

\subsection{Mathematical sources and
contributions}\label{mathematical-sources-and-contributions}

\begin{itemize}
\tightlist
\item
  \textbf{Sheaves of modules.} \emph{Sheaves on Spaces} and
  \emph{Sheaves of Modules} supply the sheaf and module constructions.
  The course develops germ-family sheafification, stalk bijections,
  associative-ring module structures and their full universal-property
  proofs.
\item
  \textbf{Complexes and localization.} \emph{Derived Categories}
  supplies cone and triangulation arguments; the course develops roofs,
  common refinements and localization checks. The diagram and elementary
  cohomology prerequisites are in Wen-Wei Li\textquotesingle s
  \emph{Methods of Algebra, Volume 2}, in the
  programme\textquotesingle s English edition.
\item
  \textbf{Injective and K-injective resolutions.} \emph{Injectives},
  \emph{Derived Categories} and \emph{Cohomology of Sheaves} supply the
  resolution methods. The course includes flasque-section arguments,
  large-cofinality and transfinite constructions. The Stacks chapter
  credits Christian Serpé and Nicolas Spaltenstein for the
  unbounded-resolution theory; \emph{Sets} credits the large-cofinality
  method to Grothendieck.
\item
  \textbf{Flat and K-flat resolutions.} \emph{Modules on Ringed Spaces}
  and \emph{Cohomology of Sheaves} supply flat-module and
  compatible-resolution methods. The course develops cycle attachments,
  complete flatness calculations and compatible resolutions.
\item
  \textbf{Tensor, pullback and pushforward.} \emph{Cohomology of
  Sheaves} supplies tensor, Tor, pullback and local-cohomology
  arguments. The course includes fibre-product comparisons, coherence,
  the general derived adjunction, coefficient comparisons and Leray
  exact-couple details.
\item
  \textbf{Derived Hom and supports.} \emph{Cohomology of Sheaves} and
  \emph{Sheaves of Modules} supply Hom and closed-support arguments. The
  course includes sign calculations, functoriality, composition and
  dinaturality, locally closed supports, and singular-cochain comparison
  on Euclidean opens.
\item
  \textbf{Perfect complexes, projection and base change.}
  \emph{Cohomology of Sheaves} supplies finite-model, pseudo-coherence,
  Tor-amplitude, duality, projection and base-change methods. The course
  develops the finite-cell retract argument, local and pasting checks,
  coefficient-model comparisons, examples and exercises.
\end{itemize}

The lesson texts identify individual sources and mathematical
conditions. Adaptations and proof expansions dated 1--3 October 2026 are
by GPT-6.1 Sol (OpenAI), Ultra, with the contributions credited above.
Stacks-derived constructions remain GFDL-1.2-or-later.

\href{https://github.com/KokunoYumeto/unofficial-stacks-project-ai-drafts/tree/565b10e987aba5969b21145a0833f42d69f96790}{AI
Integrated Stacks Project source edition} ·
\href{https://stacks.math.columbia.edu/}{Stacks project} ·
\href{https://kokunoyumeto.github.io/methods-of-algebra-volume-2-en/}{Methods
of Algebra, Volume 2} ·
\href{https://creativecommons.org/publicdomain/zero/1.0/}{CC0 1.0}

\section{GNU Free Documentation
License}\label{gnu-free-documentation-license}

GNU Free Documentation License Version 1.2, November 2002

Copyright (C) 2000,2001,2002 Free Software Foundation, Inc. 51 Franklin
St, Fifth Floor, Boston, MA 02110-1301 USA Everyone is permitted to copy
and distribute verbatim copies of this license document, but changing it
is not allowed.

0. PREAMBLE

The purpose of this License is to make a manual, textbook, or other
functional and useful document "free" in the sense of freedom: to assure
everyone the effective freedom to copy and redistribute it, with or
without modifying it, either commercially or noncommercially.
Secondarily, this License preserves for the author and publisher a way
to get credit for their work, while not being considered responsible for
modifications made by others.

This License is a kind of "copyleft", which means that derivative works
of the document must themselves be free in the same sense. It
complements the GNU General Public License, which is a copyleft license
designed for free software.

We have designed this License in order to use it for manuals for free
software, because free software needs free documentation: a free program
should come with manuals providing the same freedoms that the software
does. But this License is not limited to software manuals; it can be
used for any textual work, regardless of subject matter or whether it is
published as a printed book. We recommend this License principally for
works whose purpose is instruction or reference.

1. APPLICABILITY AND DEFINITIONS

This License applies to any manual or other work, in any medium, that
contains a notice placed by the copyright holder saying it can be
distributed under the terms of this License. Such a notice grants a
world-wide, royalty-free license, unlimited in duration, to use that
work under the conditions stated herein. The "Document", below, refers
to any such manual or work. Any member of the public is a licensee, and
is addressed as "you". You accept the license if you copy, modify or
distribute the work in a way requiring permission under copyright law.

A "Modified Version" of the Document means any work containing the
Document or a portion of it, either copied verbatim, or with
modifications and/or translated into another language.

A "Secondary Section" is a named appendix or a front-matter section of
the Document that deals exclusively with the relationship of the
publishers or authors of the Document to the Document\textquotesingle s
overall subject (or to related matters) and contains nothing that could
fall directly within that overall subject. (Thus, if the Document is in
part a textbook of mathematics, a Secondary Section may not explain any
mathematics.) The relationship could be a matter of historical
connection with the subject or with related matters, or of legal,
commercial, philosophical, ethical or political position regarding them.

The "Invariant Sections" are certain Secondary Sections whose titles are
designated, as being those of Invariant Sections, in the notice that
says that the Document is released under this License. If a section does
not fit the above definition of Secondary then it is not allowed to be
designated as Invariant. The Document may contain zero Invariant
Sections. If the Document does not identify any Invariant Sections then
there are none.

The "Cover Texts" are certain short passages of text that are listed, as
Front-Cover Texts or Back-Cover Texts, in the notice that says that the
Document is released under this License. A Front-Cover Text may be at
most 5 words, and a Back-Cover Text may be at most 25 words.

A "Transparent" copy of the Document means a machine-readable copy,
represented in a format whose specification is available to the general
public, that is suitable for revising the document straightforwardly
with generic text editors or (for images composed of pixels) generic
paint programs or (for drawings) some widely available drawing editor,
and that is suitable for input to text formatters or for automatic
translation to a variety of formats suitable for input to text
formatters. A copy made in an otherwise Transparent file format whose
markup, or absence of markup, has been arranged to thwart or discourage
subsequent modification by readers is not Transparent. An image format
is not Transparent if used for any substantial amount of text. A copy
that is not "Transparent" is called "Opaque".

Examples of suitable formats for Transparent copies include plain ASCII
without markup, Texinfo input format, LaTeX input format, SGML or XML
using a publicly available DTD, and standard-conforming simple HTML,
PostScript or PDF designed for human modification. Examples of
transparent image formats include PNG, XCF and JPG. Opaque formats
include proprietary formats that can be read and edited only by
proprietary word processors, SGML or XML for which the DTD and/or
processing tools are not generally available, and the machine-generated
HTML, PostScript or PDF produced by some word processors for output
purposes only.

The "Title Page" means, for a printed book, the title page itself, plus
such following pages as are needed to hold, legibly, the material this
License requires to appear in the title page. For works in formats which
do not have any title page as such, "Title Page" means the text near the
most prominent appearance of the work\textquotesingle s title, preceding
the beginning of the body of the text.

A section "Entitled XYZ" means a named subunit of the Document whose
title either is precisely XYZ or contains XYZ in parentheses following
text that translates XYZ in another language. (Here XYZ stands for a
specific section name mentioned below, such as "Acknowledgements",
"Dedications", "Endorsements", or "History".) To "Preserve the Title" of
such a section when you modify the Document means that it remains a
section "Entitled XYZ" according to this definition.

The Document may include Warranty Disclaimers next to the notice which
states that this License applies to the Document. These Warranty
Disclaimers are considered to be included by reference in this License,
but only as regards disclaiming warranties: any other implication that
these Warranty Disclaimers may have is void and has no effect on the
meaning of this License.

2. VERBATIM COPYING

You may copy and distribute the Document in any medium, either
commercially or noncommercially, provided that this License, the
copyright notices, and the license notice saying this License applies to
the Document are reproduced in all copies, and that you add no other
conditions whatsoever to those of this License. You may not use
technical measures to obstruct or control the reading or further copying
of the copies you make or distribute. However, you may accept
compensation in exchange for copies. If you distribute a large enough
number of copies you must also follow the conditions in section 3.

You may also lend copies, under the same conditions stated above, and
you may publicly display copies.

3. COPYING IN QUANTITY

If you publish printed copies (or copies in media that commonly have
printed covers) of the Document, numbering more than 100, and the
Document\textquotesingle s license notice requires Cover Texts, you must
enclose the copies in covers that carry, clearly and legibly, all these
Cover Texts: Front-Cover Texts on the front cover, and Back-Cover Texts
on the back cover. Both covers must also clearly and legibly identify
you as the publisher of these copies. The front cover must present the
full title with all words of the title equally prominent and visible.
You may add other material on the covers in addition. Copying with
changes limited to the covers, as long as they preserve the title of the
Document and satisfy these conditions, can be treated as verbatim
copying in other respects.

If the required texts for either cover are too voluminous to fit
legibly, you should put the first ones listed (as many as fit
reasonably) on the actual cover, and continue the rest onto adjacent
pages.

If you publish or distribute Opaque copies of the Document numbering
more than 100, you must either include a machine-readable Transparent
copy along with each Opaque copy, or state in or with each Opaque copy a
computer-network location from which the general network-using public
has access to download using public-standard network protocols a
complete Transparent copy of the Document, free of added material. If
you use the latter option, you must take reasonably prudent steps, when
you begin distribution of Opaque copies in quantity, to ensure that this
Transparent copy will remain thus accessible at the stated location
until at least one year after the last time you distribute an Opaque
copy (directly or through your agents or retailers) of that edition to
the public.

It is requested, but not required, that you contact the authors of the
Document well before redistributing any large number of copies, to give
them a chance to provide you with an updated version of the Document.

4. MODIFICATIONS

You may copy and distribute a Modified Version of the Document under the
conditions of sections 2 and 3 above, provided that you release the
Modified Version under precisely this License, with the Modified Version
filling the role of the Document, thus licensing distribution and
modification of the Modified Version to whoever possesses a copy of it.
In addition, you must do these things in the Modified Version:

A. Use in the Title Page (and on the covers, if any) a title distinct
from that of the Document, and from those of previous versions (which
should, if there were any, be listed in the History section of the
Document). You may use the same title as a previous version if the
original publisher of that version gives permission. B. List on the
Title Page, as authors, one or more persons or entities responsible for
authorship of the modifications in the Modified Version, together with
at least five of the principal authors of the Document (all of its
principal authors, if it has fewer than five), unless they release you
from this requirement. C. State on the Title page the name of the
publisher of the Modified Version, as the publisher. D. Preserve all the
copyright notices of the Document. E. Add an appropriate copyright
notice for your modifications adjacent to the other copyright notices.
F. Include, immediately after the copyright notices, a license notice
giving the public permission to use the Modified Version under the terms
of this License, in the form shown in the Addendum below. G. Preserve in
that license notice the full lists of Invariant Sections and required
Cover Texts given in the Document\textquotesingle s license notice. H.
Include an unaltered copy of this License. I. Preserve the section
Entitled "History", Preserve its Title, and add to it an item stating at
least the title, year, new authors, and publisher of the Modified
Version as given on the Title Page. If there is no section Entitled
"History" in the Document, create one stating the title, year, authors,
and publisher of the Document as given on its Title Page, then add an
item describing the Modified Version as stated in the previous sentence.
J. Preserve the network location, if any, given in the Document for
public access to a Transparent copy of the Document, and likewise the
network locations given in the Document for previous versions it was
based on. These may be placed in the "History" section. You may omit a
network location for a work that was published at least four years
before the Document itself, or if the original publisher of the version
it refers to gives permission. K. For any section Entitled
"Acknowledgements" or "Dedications", Preserve the Title of the section,
and preserve in the section all the substance and tone of each of the
contributor acknowledgements and/or dedications given therein. L.
Preserve all the Invariant Sections of the Document, unaltered in their
text and in their titles. Section numbers or the equivalent are not
considered part of the section titles. M. Delete any section Entitled
"Endorsements". Such a section may not be included in the Modified
Version. N. Do not retitle any existing section to be Entitled
"Endorsements" or to conflict in title with any Invariant Section. O.
Preserve any Warranty Disclaimers.

If the Modified Version includes new front-matter sections or appendices
that qualify as Secondary Sections and contain no material copied from
the Document, you may at your option designate some or all of these
sections as invariant. To do this, add their titles to the list of
Invariant Sections in the Modified Version\textquotesingle s license
notice. These titles must be distinct from any other section titles.

You may add a section Entitled "Endorsements", provided it contains
nothing but endorsements of your Modified Version by various
parties-\/-for example, statements of peer review or that the text has
been approved by an organization as the authoritative definition of a
standard.

You may add a passage of up to five words as a Front-Cover Text, and a
passage of up to 25 words as a Back-Cover Text, to the end of the list
of Cover Texts in the Modified Version. Only one passage of Front-Cover
Text and one of Back-Cover Text may be added by (or through arrangements
made by) any one entity. If the Document already includes a cover text
for the same cover, previously added by you or by arrangement made by
the same entity you are acting on behalf of, you may not add another;
but you may replace the old one, on explicit permission from the
previous publisher that added the old one.

The author(s) and publisher(s) of the Document do not by this License
give permission to use their names for publicity for or to assert or
imply endorsement of any Modified Version.

5. COMBINING DOCUMENTS

You may combine the Document with other documents released under this
License, under the terms defined in section 4 above for modified
versions, provided that you include in the combination all of the
Invariant Sections of all of the original documents, unmodified, and
list them all as Invariant Sections of your combined work in its license
notice, and that you preserve all their Warranty Disclaimers.

The combined work need only contain one copy of this License, and
multiple identical Invariant Sections may be replaced with a single
copy. If there are multiple Invariant Sections with the same name but
different contents, make the title of each such section unique by adding
at the end of it, in parentheses, the name of the original author or
publisher of that section if known, or else a unique number. Make the
same adjustment to the section titles in the list of Invariant Sections
in the license notice of the combined work.

In the combination, you must combine any sections Entitled "History" in
the various original documents, forming one section Entitled "History";
likewise combine any sections Entitled "Acknowledgements", and any
sections Entitled "Dedications". You must delete all sections Entitled
"Endorsements".

6. COLLECTIONS OF DOCUMENTS

You may make a collection consisting of the Document and other documents
released under this License, and replace the individual copies of this
License in the various documents with a single copy that is included in
the collection, provided that you follow the rules of this License for
verbatim copying of each of the documents in all other respects.

You may extract a single document from such a collection, and distribute
it individually under this License, provided you insert a copy of this
License into the extracted document, and follow this License in all
other respects regarding verbatim copying of that document.

7. AGGREGATION WITH INDEPENDENT WORKS

A compilation of the Document or its derivatives with other separate and
independent documents or works, in or on a volume of a storage or
distribution medium, is called an "aggregate" if the copyright resulting
from the compilation is not used to limit the legal rights of the
compilation\textquotesingle s users beyond what the individual works
permit. When the Document is included in an aggregate, this License does
not apply to the other works in the aggregate which are not themselves
derivative works of the Document.

If the Cover Text requirement of section 3 is applicable to these copies
of the Document, then if the Document is less than one half of the
entire aggregate, the Document\textquotesingle s Cover Texts may be
placed on covers that bracket the Document within the aggregate, or the
electronic equivalent of covers if the Document is in electronic form.
Otherwise they must appear on printed covers that bracket the whole
aggregate.

8. TRANSLATION

Translation is considered a kind of modification, so you may distribute
translations of the Document under the terms of section 4. Replacing
Invariant Sections with translations requires special permission from
their copyright holders, but you may include translations of some or all
Invariant Sections in addition to the original versions of these
Invariant Sections. You may include a translation of this License, and
all the license notices in the Document, and any Warranty Disclaimers,
provided that you also include the original English version of this
License and the original versions of those notices and disclaimers. In
case of a disagreement between the translation and the original version
of this License or a notice or disclaimer, the original version will
prevail.

If a section in the Document is Entitled "Acknowledgements",
"Dedications", or "History", the requirement (section 4) to Preserve its
Title (section 1) will typically require changing the actual title.

9. TERMINATION

You may not copy, modify, sublicense, or distribute the Document except
as expressly provided for under this License. Any other attempt to copy,
modify, sublicense or distribute the Document is void, and will
automatically terminate your rights under this License. However, parties
who have received copies, or rights, from you under this License will
not have their licenses terminated so long as such parties remain in
full compliance.

10. FUTURE REVISIONS OF THIS LICENSE

The Free Software Foundation may publish new, revised versions of the
GNU Free Documentation License from time to time. Such new versions will
be similar in spirit to the present version, but may differ in detail to
address new problems or concerns. See https://www.gnu.org/copyleft/.

Each version of the License is given a distinguishing version number. If
the Document specifies that a particular numbered version of this
License "or any later version" applies to it, you have the option of
following the terms and conditions either of that specified version or
of any later version that has been published (not as a draft) by the
Free Software Foundation. If the Document does not specify a version
number of this License, you may choose any version ever published (not
as a draft) by the Free Software Foundation.

ADDENDUM: How to use this License for your documents

To use this License in a document you have written, include a copy of
the License in the document and put the following copyright and license
notices just after the title page:

Copyright (c) YEAR YOUR NAME. Permission is granted to copy, distribute
and/or modify this document under the terms of the GNU Free
Documentation License, Version 1.2 or any later version published by the
Free Software Foundation; with no Invariant Sections, no Front-Cover
Texts, and no Back-Cover Texts. A copy of the license is included in the
section entitled "GNU Free Documentation License".

If you have Invariant Sections, Front-Cover Texts and Back-Cover Texts,
replace the "with...Texts." line with this:

with the Invariant Sections being LIST THEIR TITLES, with the
Front-Cover Texts being LIST, and with the Back-Cover Texts being LIST.

If you have Invariant Sections without Cover Texts, or some other
combination of the three, merge those two alternatives to suit the
situation.

If your document contains nontrivial examples of program code, we
recommend releasing these examples in parallel under your choice of free
software license, such as the GNU General Public License, to permit
their use in free software.

\end{document}
